\documentclass[12pt,a4paper]{article}

% -------------------------------------------------
% PACOTES E CONFIGURAÇÕES
% -------------------------------------------------
\usepackage[utf8]{inputenc}
\usepackage[T1]{fontenc}

% Suporte a idioma com fallback automático para compilação em qualquer ambiente TeX Live
\IfFileExists{brazilian.ldf}{%
  \usepackage[brazilian]{babel}%
}{%
  \usepackage[english]{babel}%
  \addto\captionsenglish{%
    \renewcommand{\contentsname}{Sumário}%
    \renewcommand{\listfigurename}{Lista de Figuras}%
    \renewcommand{\listtablename}{Lista de Tabelas}%
    \renewcommand{\figurename}{Figura}%
    \renewcommand{\tablename}{Tabela}%
  }%
}

\usepackage{lmodern}
\usepackage{geometry}
\geometry{top=2.8cm,bottom=2.5cm,left=3cm,right=2.5cm}

% Microtype: melhora tipográfica (carregado se disponível)
\IfFileExists{microtype.sty}{\usepackage{microtype}}{}

\usepackage{setspace}
\onehalfspacing
\usepackage{graphicx}
\usepackage{float}
\usepackage{booktabs}
\usepackage{longtable}
\usepackage{array}
\usepackage{tabularx}

% Enumitem com fallback elegante
\IfFileExists{enumitem.sty}{%
  \usepackage{enumitem}%
  \setlist[enumerate,1]{label=\textbf{\arabic*.}, leftmargin=*}%
}{%
  \renewcommand{\labelenumi}{\textbf{\arabic{enumi}.}}%
}

\usepackage{xcolor}
\usepackage{hyperref}
\usepackage{bookmark}
\usepackage{listings}
\usepackage{amsmath}
\usepackage{amssymb}
\usepackage{fancyhdr}

% Lastpage com fallback seguro
\IfFileExists{lastpage.sty}{%
  \usepackage{lastpage}%
}{%
  \AtEndDocument{\label{LastPage}}%
}

\usepackage{etoolbox}
\renewcommand{\_}{\textunderscore\allowbreak}

% -------------------------------------------------
% CORES DO PROJETO LCQUI
% -------------------------------------------------
\definecolor{entityHead}{RGB}{30,64,110}
\definecolor{entityBg}{RGB}{235,242,250}
\definecolor{entityFrame}{RGB}{30,64,110}

\definecolor{noteHead}{RGB}{176,98,10}
\definecolor{noteBg}{RGB}{255,244,229}
\definecolor{noteFrame}{RGB}{176,98,10}

\definecolor{ruleHead}{RGB}{20,110,80}
\definecolor{ruleBg}{RGB}{230,247,240}
\definecolor{ruleFrame}{RGB}{20,110,80}

\definecolor{pkColor}{RGB}{176,40,40}
\definecolor{fkColor}{RGB}{30,90,160}
\definecolor{typeColor}{RGB}{60,60,60}
\definecolor{typeBg}{RGB}{225,225,232}

% -------------------------------------------------
% HYPERSETUP
% -------------------------------------------------
\hypersetup{
    colorlinks=true,
    linkcolor=blue!50!black,
    urlcolor=blue!60!black,
    citecolor=green!40!black,
    pdftitle={Projeto de Desenvolvimento do Site LCQUI},
    pdfauthor={Zadoque Carneiro}
}

\pagestyle{fancy}
\fancyhf{}
\lhead{Projeto LCQUI}
\rhead{Planejamento Técnico}
\cfoot{\thepage\ de \pageref{LastPage}}
\setlength{\headheight}{15pt}

% -------------------------------------------------
% DEFINIÇÃO DE LINGUAGENS (MERMAID E TYPESCRIPT)
% -------------------------------------------------
\lstdefinelanguage{Mermaid}{%
  morekeywords={flowchart,graph,sequenceDiagram,classDiagram,erDiagram,stateDiagram,participant,subgraph,end,style,click,Note,rect,alt,else,loop,par,and,autonumber},%
  sensitive=false%
}

\lstdefinelanguage{TypeScript}{
  keywords={typeof, new, true, false, catch, function, return, null, switch, var, if, in, while, do, else, case, break, const, let, export, import, from, interface, type, class, extends, implements, async, await, string, number, boolean, any, void, as},
  keywordstyle=\color{blue}\bfseries,
  ndkeywords={admin, onCall, PDFDocument, Date, Promise, reduce},
  ndkeywordstyle=\color{purple}\bfseries,
  identifierstyle=\color{black},
  sensitive=true,
  comment=[l]{//},
  morecomment=[s]{/*}{*/},
  commentstyle=\color{gray}\ttfamily\itshape,
  stringstyle=\color{green!50!black}\ttfamily,
  morestring=[b]',
  morestring=[b]",
  morestring=[b]`
}

\lstset{
    basicstyle=\ttfamily\small,
    breaklines=true,
    frame=single,
    columns=fullflexible,
    keepspaces=true,
    showstringspaces=false,
    literate=
        {á}{{\'a}}1 {à}{{\`a}}1 {ã}{{\~a}}1 {â}{{\^a}}1
        {é}{{\'e}}1 {ê}{{\^e}}1
        {í}{{\'i}}1
        {ó}{{\'o}}1 {õ}{{\~o}}1 {ô}{{\^o}}1
        {ú}{{\'u}}1 {ü}{{\"u}}1
        {ç}{{\c c}}1
        {Á}{{\'A}}1 {À}{{\`A}}1 {Ã}{{\~A}}1 {Â}{{\^A}}1
        {É}{{\'E}}1 {Ê}{{\^E}}1
        {Í}{{\'I}}1
        {Ó}{{\'O}}1 {Õ}{{\~O}}1 {Ô}{{\^O}}1
        {Ú}{{\'U}}1 {Ü}{{\"U}}1
        {Ç}{{\c C}}1
}

\newcommand{\placeholderfigure}[3]{%
\begin{figure}[H]
    \centering
    \fbox{\begin{minipage}[c][#2][c]{0.93\textwidth}
        \centering
        \textbf{Espaço reservado para diagrama/imagem}\\[0.5em]
        \textit{#1.png}\\[0.5em]
        \small #3
    \end{minipage}}
    \caption{#3}
\end{figure}}

% -------------------------------------------------
% CAIXAS DE ENTIDADE, OBSERVAÇÃO E REGRAS
% Suporta tcolorbox quando disponível, ou renderização nativa ultra-rápida
% -------------------------------------------------
\newif\ifusetcolorbox
\IfFileExists{tcolorbox.sty}{\usetcolorboxtrue}{\usetcolorboxfalse}

\ifusetcolorbox
  \usepackage{tcolorbox}
  \tcbuselibrary{skins,breakable}
  \newtcolorbox{entidade}[1]{%
    enhanced, breakable, colback=entityBg, colframe=entityFrame,
    colbacktitle=entityHead, coltitle=white, fonttitle=\bfseries\large,
    title={\ttfamily #1}, boxrule=0.9pt, arc=2mm,
    left=5mm, right=5mm, top=2mm, bottom=2mm, boxsep=1mm,
    attach boxed title to top left={xshift=3mm,yshift*=-1mm},
    boxed title style={boxrule=0pt,arc=1.5mm},
    before skip=10pt, after skip=14pt
  }
  \newcommand{\tipo}[1]{%
    \tcbox[nobeforeafter,boxrule=0pt,boxsep=1pt,left=3pt,right=3pt,top=1.5pt,bottom=1.5pt,
           colback=typeBg,colframe=typeBg,arc=1.5mm,
           fontupper=\scriptsize\ttfamily\color{typeColor}]{#1}%
  }
  \newcommand{\chave}[1]{%
    \ifstrequal{#1}{PK}{\tcbox[nobeforeafter,boxrule=0pt,boxsep=1pt,left=3pt,right=3pt,top=1.5pt,bottom=1.5pt,colback=pkColor,colframe=pkColor,arc=1.5mm,fontupper=\scriptsize\bfseries\color{white}]{PK}}{%
    \ifstrequal{#1}{FK}{\tcbox[nobeforeafter,boxrule=0pt,boxsep=1pt,left=3pt,right=3pt,top=1.5pt,bottom=1.5pt,colback=fkColor,colframe=fkColor,arc=1.5mm,fontupper=\scriptsize\bfseries\color{white}]{FK}}{%
    \tcbox[nobeforeafter,boxrule=0pt,boxsep=1pt,left=3pt,right=3pt,top=1.5pt,bottom=1.5pt,colback=gray!25,colframe=gray!25,arc=1.5mm,fontupper=\scriptsize\itshape\color{black!70}]{#1}%
    }}%
  }
  \newtcolorbox{explicacao}[1][Observação]{%
    enhanced, breakable, colback=noteBg, colframe=noteFrame,
    colbacktitle=noteHead, coltitle=white, fonttitle=\bfseries\small,
    title={\raisebox{-0.1em}{\textbf{i}}\ \ #1}, boxrule=0.8pt, arc=1.5mm,
    left=4mm, right=4mm, top=1.5mm, bottom=1.5mm, before skip=8pt, after skip=10pt
  }
  \newtcolorbox{regra}[1][Regra de negócio]{%
    enhanced, breakable, colback=ruleBg, colframe=ruleFrame,
    colbacktitle=ruleHead, coltitle=white, fonttitle=\bfseries\small,
    title={#1}, boxrule=0.8pt, arc=1.5mm,
    left=4mm, right=4mm, top=1.5mm, bottom=1.5mm, before skip=8pt, after skip=10pt
  }
\else
  % Renderização nativa limpa, rápida e estética (zero dependência de pacotes extras)
  \newcommand{\tipo}[1]{\colorbox{typeBg}{\scriptsize\ttfamily\color{typeColor}\strut\,#1\,}}
  \newcommand{\chave}[1]{%
    \ifstrequal{#1}{PK}{\colorbox{pkColor}{\scriptsize\bfseries\color{white}\strut\,PK\,}}{%
    \ifstrequal{#1}{FK}{\colorbox{fkColor}{\scriptsize\bfseries\color{white}\strut\,FK\,}}{%
    \colorbox{gray!25}{\scriptsize\itshape\color{black!70}\strut\,#1\,}%
    }}%
  }
  \newsavebox{\entidadeBox}
  \newenvironment{entidade}[1]{%
    \par\vspace{8pt}\noindent%
    \def\entidadeTitle{#1}%
    \begin{lrbox}{\entidadeBox}%
    \begin{minipage}{\dimexpr\textwidth-2\fboxrule-10mm\relax}%
    \vspace{2mm}%
  }{%
    \vspace{2mm}%
    \end{minipage}%
    \end{lrbox}%
    \setlength{\fboxrule}{0.9pt}%
    \setlength{\fboxsep}{0pt}%
    \fcolorbox{entityFrame}{entityBg}{%
      \begin{minipage}{\dimexpr\textwidth-2\fboxrule\relax}%
        \colorbox{entityHead}{%
          \begin{minipage}{\dimexpr\textwidth-2\fboxrule-8mm\relax}%
            \vspace{2mm}%
            \hspace{4mm}{\color{white}\bfseries\large\ttfamily \entidadeTitle}%
            \vspace{2mm}%
          \end{minipage}%
        }\par
        \vspace{2mm}%
        \hspace{5mm}\usebox{\entidadeBox}%
        \vspace{2mm}%
      \end{minipage}%
    }\par\vspace{10pt}%
  }
  \newsavebox{\explicacaoBox}
  \newenvironment{explicacao}[1][Observação]{%
    \par\vspace{8pt}\noindent%
    \def\explicacaoTitle{#1}%
    \begin{lrbox}{\explicacaoBox}%
    \begin{minipage}{\dimexpr\textwidth-2\fboxrule-8mm\relax}%
    \small\vspace{2mm}%
  }{%
    \vspace{2mm}%
    \end{minipage}%
    \end{lrbox}%
    \setlength{\fboxrule}{0.8pt}%
    \setlength{\fboxsep}{0pt}%
    \fcolorbox{noteFrame}{noteBg}{%
      \begin{minipage}{\dimexpr\textwidth-2\fboxrule\relax}%
        \colorbox{noteHead}{%
          \begin{minipage}{\dimexpr\textwidth-2\fboxrule-6mm\relax}%
            \vspace{1.5mm}%
            \hspace{3mm}{\color{white}\bfseries\small\raisebox{-0.1em}{\textbf{i}}\ \ \explicacaoTitle}%
            \vspace{1.5mm}%
          \end{minipage}%
        }\par
        \vspace{2mm}%
        \hspace{4mm}\usebox{\explicacaoBox}%
        \vspace{2mm}%
      \end{minipage}%
    }\par\vspace{10pt}%
  }
  \newsavebox{\regraBox}
  \newenvironment{regra}[1][Regra de negócio]{%
    \par\vspace{8pt}\noindent%
    \def\regraTitle{#1}%
    \begin{lrbox}{\regraBox}%
    \begin{minipage}{\dimexpr\textwidth-2\fboxrule-8mm\relax}%
    \small\vspace{2mm}%
  }{%
    \vspace{2mm}%
    \end{minipage}%
    \end{lrbox}%
    \setlength{\fboxrule}{0.8pt}%
    \setlength{\fboxsep}{0pt}%
    \fcolorbox{ruleFrame}{ruleBg}{%
      \begin{minipage}{\dimexpr\textwidth-2\fboxrule\relax}%
        \colorbox{ruleHead}{%
          \begin{minipage}{\dimexpr\textwidth-2\fboxrule-6mm\relax}%
            \vspace{1.5mm}%
            \hspace{3mm}{\color{white}\bfseries\small \regraTitle}%
            \vspace{1.5mm}%
          \end{minipage}%
        }\par
        \vspace{2mm}%
        \hspace{4mm}\usebox{\regraBox}%
        \vspace{2mm}%
      \end{minipage}%
    }\par\vspace{10pt}%
  }
\fi

% Linha de campo: nome | tipo postgres | chave/observação (opcional)
\newcommand{\campo}[3]{%
  \noindent\textit{\textbf{#1}}\hfill\tipo{#2}%
  \ifstrempty{#3}{\ \chave{NOT NULL}}{\ \chave{#3}}%
  \par\vspace{2pt}%
}

\newcommand{\subcampo}[1]{%
  \noindent\hspace{4mm}{\small\textit{#1}}\par\vspace{2pt}%
}

\title{\textbf{Projeto de Desenvolvimento do Site LCQUI}\\[0.5em]\large Planejamento Funcional, Arquitetural, de Dados e de Implementação}
\author{Zadoque Carneiro}
\date{\today}

\begin{document}

\maketitle
\tableofcontents
\newpage

% -------------------------------------------------
% SEÇÕES MODULARES DO DOCUMENTO
% -------------------------------------------------
\section{Introdução}

O LCQUI é o Laboratório de Ciências Químicas da Universidade Estadual do Norte Fluminense Darcy Ribeiro (UENF), vinculado ao Centro de Ciência e Tecnologia (CCT). Ele funciona como uma unidade de ensino, pesquisa e suporte técnico voltada para a área de química.

\section{Como vai ser o Sistema}

Criar um sistema online que contenha as informações dos bens patrimoniais, reagentes e materiais do almoxarifado do setor de Graduação do LCQUI, bem como criar uma plataforma geral de acesso aos roteiros das aulas experimentais desenvolvidas nos laboratórios desse setor. A intenção é, dessa forma, auxiliar os professores do LCQUI na consulta dos itens existentes, registro de novos itens com maior agilidade, colaborando na identificação de itens inservíveis, em escassez ou que necessitam de reposição ou reparos, auxiliando assim na gestão da chefia do LCQUI. Em relação aos roteiros experimentais, os professores poderão utilizar a plataforma para lançamento ou troca de arquivos e divulgação junto aos alunos.

%%-----------------------------%%
%%        STAKEHOLDERS         %%
%%-----------------------------%%

\section{Stakeholders}

\begin{explicacao}[Importante]
Um mesmo usuário pode possuir múltiplos papéis, observadas as combinações permitidas na Matriz de Papéis e Permissões da seção de Regras de Negócio. O papel \textit{Chefe Geral} é mutuamente exclusivo com os demais papéis para uma mesma conta. O papel \textit{Bolsista} possui como pré-condição existencial e estrutural o papel de \textit{Aluno} ($Bolsista \implies Aluno$).
\end{explicacao}

%%--------------------------------------------%%
%%       STAKEHOLDERS - CHEFE GERAL           %%
%%--------------------------------------------%%
\subsection{Chefe Geral}
Esse é o maior cargo do sistema: Esse cargo gerencia os outros cargos.
\begin{itemize}
    \item Poder cadastrar outros usuários e definir os papéis deles no sistema, de acordo com as regras descritas em \textit{Regras de Negócio} e na \textit{Descrição das telas}.
    \item Pode cadastrar um Almoxarifado. 
    \item Ao cadastrar um Gestor de Almoxarifado, pode colocá-lo responsável por mais de um almoxarifado. 
    \item Ao cadastrar um Gestor de Almoxarifado, pode vinculá-lo imediatamente para gerenciar um ou mais almoxarifados, ou apenas cadastrar o usuário (deixando a vinculação para um momento futuro, resultando em uma dashboard sem almoxarifados ativos inicialmente).
    \item Atribuir o papel de Bolsista a um usuário que já possua o papel Aluno.
\end{itemize}

%%-----------------------------------------------%%
%%     STAKEHOLDERS - GESTOR DE ALMOXARIFADO     %%
%%-----------------------------------------------%%
\subsection{Gestor de Almoxarifado}
Pode haver mais de um gestor, ele(s) pode(m):
\begin{itemize}
    \item Gerar um pdf de etiquetas de frasco 
    \item Adicionar/alterar a quantidade de um reagente no estoque.
    \item Pesquisar por um reagente no estoque.
    \item Cadastrar uma nova especificação de reagente, informando os dados obrigatórios definidos pelo tipo de substância e pelo estado físico.
    \item Marcar um reagente como removido (não exclui do banco, apenas muda o status).
    \item Alterar o status de um reagente: marcar empréstimo, devolução (através de pesagem), descarte, vazio ou quebrado.
    \item Autorizar expressamente o uso de frascos vencidos para fins didáticos, demonstrações e pesquisa excepcional conforme Q04.
    \item Gerar relatório mensal (ou de um período específico) para um almoxarifado ou para todos.
    \item Editar/adicionar propriedades de um reagente.
\end{itemize}

%%-----------------------------------------------%%
%%  STAKEHOLDERS - GESTOR DE BENS PATRIMONIAIS   %%
%%-----------------------------------------------%%
\subsection{Gestor de Bens Patrimoniais}
Pode haver mais de um gestor de bem patrimonial, ele(s) pode(m):
\begin{itemize}
    \item Visualizar as requisições dos professores para adicionar/editar um bem patrimonial.
    \item Quando a requisição do professor é para colocar um bem como inservível e o gestor aprova, o gestor deve dar baixa do bem patrimonial através do processo da UENF (abrir processo no SEI etc.). Após esse processo, ele marca no sistema que foi dada baixa.
    \item Adicionar um bem patrimonial e criar resumos de catálogo.
    \item Editar um bem patrimonial e suas propriedades.
    \item Gerar e baixar relatório de bens patrimoniais mensal (ou de período específico).
    \item Visualizar os bens patrimoniais existentes.
    \item Pesquisar bens patrimoniais por número de patrimônio ou nome.
    \item Filtrar bens patrimoniais por local (prédio, andar e sala), estado de conservação ou status.
    \item Imprimir Relatório mensal (ou período personalizado), ou sobre um Bem patrimonial específico podendo filtrar por prédio, andar, sala, Nome do responsável, etc.
\end{itemize}

%%-----------------------------------%%
%%        STAKEHOLDERS - PROFESSOR   %%
%%-----------------------------------%%
\subsection{Professor}
Com certeza haverá mais de um, eles podem:
\begin{itemize}
    \item Criar uma turma (especificando nome e capacidade máxima de alunos).
    \item Copiar o código de uma turma dele.
    \item Criar um post em uma turma (ao criar um post, o professor pode selecionar um roteiro que ele já tenha feito o upload, um roteiro que outro professor compartilhou com ele, ou ainda fazer upload de um roteiro do próprio computador -- isso abre um modal para cadastrar o roteiro).
    \item Adicionar usuários apenas como alunos por e-mail (mandando o link para eles trocarem a senha).
    \item Pesquisar por bens patrimoniais.
    \item Abrir uma requisição de adição ou edição de bem patrimonial (analisada por um gestor de bens patrimoniais).
    \item Adicionar novos Locais (Prédio, Andar e Sala) no sistema de forma autônoma (por exemplo, ao fazer uma requisição, selecionar a opção "Novo Local" e registrá-lo).
    \item Solicitar/receber empréstimo de reagentes do almoxarifado; o registro operacional da retirada é efetuado pelo Gestor de Almoxarifado.
    \item Adicionar um roteiro de experimento (formato PDF).
    \item Vincular um roteiro de experimento a uma turma.
    \item Compartilhar um roteiro de experimento com outros professores.
    \item Acessar a aba de roteiros compartilhados (dividida em: ``Compartilhados comigo'' e ``Eu compartilhei'').
    \item Baixar um roteiro (seja adicionado por ele mesmo ou compartilhado com ele por outro professor).
    \item Vincular um roteiro de experimento compartilhado com ele por outro professor a uma turma dele.
    \item Arquivar e desarquivar turmas sob sua responsabilidade.
\end{itemize}

%%-----------------------------------%%
%%        STAKEHOLDERS - ALUNO       %%
%%-----------------------------------%%
\subsection{Aluno}
Esses com certeza serão maioria. Eles podem:
\begin{itemize}
    \item Entrar em uma turma usando o código disponibilizado pelo professor (respeitado o limite de capacidade da turma).
    \item Visualizar suas turmas e o professor respectivo de cada turma.
    \item Baixar qualquer roteiro de qualquer turma da qual façam parte.
    \item Visualizar os colegas que estão na mesma turma.
    \item Comentar em um post (o comentário é visível tanto para o professor quanto para os colegas).
\end{itemize}

\subsection{Bolsista}
Seja de graduação ou pós-graduação, este papel identifica um usuário que possui o papel de Aluno e uma autorização adicional para retirada de reagentes ($Bolsista \implies Aluno$). Ele é mutuamente compatível com Aluno, mas possui restrições próprias para acumulação com papéis de gestor.
\begin{itemize}
    \item Ele pode ver os reagentes e seus estoques assim como os professores.
    \item O gestor de reagentes pode registrar retiradas e devoluções para esses usuários.
    \item Bolsista e Gestor de Almoxarifado são mutuamente exclusivos (SoD). Aluno sem Bolsista pode exercer Gestor de Almoxarifado, conforme DP-C01.
\end{itemize}

\subsection{Matriz de Papéis e Permissões}
Esta tabela é um resumo de stakeholders e UX. A fonte normativa única da
autorização, inclusive estado ativo, escopo, ownership, decisão transacional,
revogação e resultado de negação, é a matriz M9 da Seção~7. Papéis exibidos ou
informados pelo cliente nunca determinam permissão.


\small
\setlength{\tabcolsep}{4pt}
\small
\setlength{\tabcolsep}{3pt}
\begin{longtable}{>{\raggedright\arraybackslash}p{0.24\textwidth}*{6}{>{\centering\arraybackslash}p{0.105\textwidth}}}
\toprule
\textbf{Ação} & \textbf{Chefe} & \textbf{Gestor Almox.} & \textbf{Gestor Patr.} & \textbf{Professor} & \textbf{Aluno} & \textbf{Bolsista}\\
\midrule
\endfirsthead
\toprule
\textbf{Ação} & \textbf{Chefe} & \textbf{Gestor Almox.} & \textbf{Gestor Patr.} & \textbf{Professor} & \textbf{Aluno} & \textbf{Bolsista} (cont.)\\
\midrule
\endhead
\bottomrule
\endfoot
Gerenciar usuários e papéis & \checkmark & -- & -- & -- & -- & --\\
Gerenciar almoxarifados & \checkmark & -- & -- & -- & -- & --\\
Consultar estoque de reagentes & \checkmark & \checkmark & -- & \checkmark & \checkmark & \checkmark\\
Cadastrar/editar reagente e frasco & \checkmark & \checkmark & -- & -- & -- & --\\
Registrar retirada/devolução (operação de balcão) & \checkmark & \checkmark & -- & retirante & -- & retirante\\
Gerenciar bens patrimoniais & \checkmark & -- & \checkmark & leitura & -- & --\\
Abrir requisição de patrimônio & \checkmark & -- & -- & \checkmark & -- & --\\
Responder requisição de patrimônio & \checkmark & -- & \checkmark & -- & -- & --\\
Gerenciar turmas próprias & \checkmark & -- & -- & \checkmark & -- & --\\
Ingressar em turma & -- & -- & -- & -- & \checkmark & \checkmark\\
Criar/comentar posts permitidos & -- & -- & -- & \checkmark & \checkmark & \checkmark\\
Gerenciar roteiros próprios & \checkmark & -- & -- & \checkmark & -- & --\\
Compartilhar roteiros & \checkmark & -- & -- & \checkmark & -- & --\\
Cadastrar/editar matérias & \checkmark & -- & -- & \checkmark & -- & --\\
Gerar etiquetas de frasco & \checkmark & \checkmark & -- & -- & -- & --\\
\bottomrule
\end{longtable}


\begin{explicacao}[Regra de autorização]
A Cloud Function responsável por cada mutação deve resolver os papéis a partir dos registros persistidos e, quando aplicável, validar também o escopo do recurso. Por exemplo, um Gestor de Almoxarifado só atua sobre almoxarifados aos quais está vinculado. O frontend pode ocultar ou desabilitar ações, mas não constitui autorização.

A linha ``Gerenciar turmas próprias'' pertence ao professor dono. O Chefe Geral não possui turmas próprias (RN-ROLE-01) e não recebe ownership de turma alheia: sua atuação acadêmica é a intervenção excepcional de moderação Q13, com justificativa, auditoria e a marcação institucional nas publicações. O contrato de Turma, matrícula e convites é a Seção~\ref{sec:regras-turmas-m11}.
\end{explicacao}

\paragraph{Operador e retirante (AUDITORIA-7, C002).}
A autorização de Chefe Geral na matriz refere-se a operar o balcão, identificado por \texttt{id\_gestor\_retirada}. O destinatário \texttt{id\_usuario\_retirou} deve ser Professor ou Bolsista ativo. Chefe Geral não pode ser retirante, pois seu papel é exclusivo (RN-ROLE-01).

%%-----------------------------------%%
%%   MODELAGEM ENTIDADES SQL - 3FN   %%
%%-----------------------------------%%
\section{Modelagem Entidades SQL - 3FN}\label{sec:modelagem}

\begin{explicacao}[Como ler as caixas de entidade]
Cada caixa a seguir representa uma \textbf{tabela do PostgreSQL} na Terceira Forma Normal (3FN). Todos os atributos são \textbf{NOT NULL} por padrão, salvo quando explicitamente indicado \chave{NULL}. O nome do campo aparece à esquerda e o \textbf{tipo de dado} sugerido para o PostgreSQL aparece à direita, em cinza. Quando aplicável, um selo adicional indica \chave{PK} (chave primária), \chave{FK} (chave estrangeira) ou uma restrição como \chave{NULL} (aceita valor nulo) ou \chave{UNIQUE}.
\end{explicacao}

%%-----------------------------------%%
%%        ENTIDADES - USUÁRIO        %%
%%-----------------------------------%%
\subsection{Usuario}
\begin{entidade}{Usuario}
\campo{id}{SERIAL}{PK}
\campo{nome}{VARCHAR(150)}{NOT NULL}
\campo{email}{VARCHAR(150)}{UNIQUE, NOT NULL}
\campo{profile\_url\_photo}{TEXT}{NULL}
\end{entidade}

\subsection{Chefe Geral}
\begin{entidade}{Chefe\_Geral}
\campo{id\_usuario}{INTEGER}{PK, FK}
\end{entidade}

\subsection{Gestor de Almoxarifado}
\begin{entidade}{Gestor\_Almoxarifado}
\campo{id\_usuario}{INTEGER}{PK, FK}
\end{entidade}

\subsection{Tabela Gestor de Almoxarifado \texorpdfstring{\texorpdfstring{$\times$}{x}}{x} Almoxarifado (N para N)}
\begin{entidade}{Gestor\_Almoxarifado\_x\_Almoxarifado}
\campo{id\_gestor\_almoxarifado}{INTEGER}{PK, FK}
\campo{id\_almoxarifado}{INTEGER}{PK, FK}
\end{entidade}
\begin{regra}
A chave primária composta \textit{(id\_gestor\_almoxarifado, id\_almoxarifado)} impede vincular o mesmo gestor duas vezes ao mesmo almoxarifado.
\end{regra}

\subsection{Gestor de bens patrimoniais}
\begin{entidade}{Gestor\_Bens\_Patrimoniais}
\campo{id\_usuario}{INTEGER}{PK, FK}
\end{entidade}
%%-----------------------------------%%
%%        BENS PATRIMONIAIS          %%
%%-----------------------------------%%
\subsection{Bem patrimonial}
\begin{entidade}{Bem\_Patrimonial}
\campo{id}{SERIAL}{PK}
\campo{id\_resumo\_bem\_patrimonial}{INTEGER}{FK, NOT NULL}
\campo{numero\_patrimonio}{VARCHAR(30)}{UNIQUE, NOT NULL}
\campo{estado\_conservacao}{ENUM}{BOM, REGULAR, RUIM, NOT NULL}
\campo{id\_local}{INTEGER}{FK, NOT NULL}
\campo{photo\_url}{TEXT}{NOT NULL}
\campo{documento\_dado\_baixa\_pdf\_url}{TEXT}{NULL}
\campo{nome\_responsavel\_sei}{VARCHAR(150)}{NOT NULL}
\campo{status}{ENUM}{Ativo, Inservivel, Ja\_dado\_baixa, NOT NULL}
\campo{descricao\_complementar}{VARCHAR(200)}{NULL}
\campo{versao}{INTEGER}{DEFAULT 1, NOT NULL}
\end{entidade}

\begin{explicacao}
\textit{id\_resumo\_bem\_patrimonial}, \textit{numero\_patrimonio}, \textit{estado\_conservacao}, \textit{id\_local}, \textit{photo\_url}, \textit{nome\_responsavel\_sei} e \textit{status} são estritamente \textbf{NOT NULL}. Apenas \textit{documento\_dado\_baixa\_pdf\_url} admite NULL (preenchido exclusivamente quando o bem atinge o status \texttt{Ja\_dado\_baixa}) e \textit{descricao\_complementar} (opcional para notas individuais daquela unidade física).
\end{explicacao}

\subsection{Resumo bem patrimonial}
\begin{entidade}{Resumo\_Bem\_Patrimonial}
\campo{id}{SERIAL}{PK}
\campo{nome}{VARCHAR(150)}{NOT NULL}
\campo{descricao}{TEXT}{NOT NULL}
\end{entidade}

\begin{explicacao}[Regra de Identidade de Bens Patrimoniais]
\textit{Resumo\_Bem\_Patrimonial} representa o modelo catalográfico do equipamento (ex.: ``Microscópio Óptico Binocular Coleman''). Vários bens físicos compartilham o mesmo resumo. A alteração do nome do resumo é restrita ao Gestor de Bens e reflete em todos os itens do modelo. Caso um professor solicite alteração de nome para um bem individual, a operação representa reclassificação para outro resumo existente ou criação de um novo modelo, preservando a integridade dos demais equipamentos. Expressamente, resumos de catálogo não possuem vínculo com arquivos fotográficos, pertencendo a imagem unicamente à unidade física cadastrada em \textit{Bem\_Patrimonial}.
\end{explicacao}
\newpage
\subsection{Historico bem patrimonial}
\begin{entidade}{Historico\_Bem\_Patrimonial}
\campo{id}{SERIAL}{PK}
\campo{id\_bem\_patrimonial}{INTEGER}{FK, NOT NULL}
\campo{timestamp}{TIMESTAMP}{NOT NULL}
\campo{tipo}{ENUM}{cadastro, edicao, baixa, NOT NULL}
\campo{id\_usuario}{INTEGER}{FK, NOT NULL}
\end{entidade}

\subsection{Alteracao bem patrimonial}
\begin{entidade}{Alteracao\_Bem\_Patrimonial}
\campo{id}{SERIAL}{PK}
\campo{id\_historico}{INTEGER}{FK, NOT NULL}
\campo{campo}{VARCHAR(60)}{NOT NULL}
\campo{valor\_anterior}{TEXT}{NOT NULL}
\campo{valor\_novo}{TEXT}{NOT NULL}
\end{entidade}
\newpage
%%-----------------------------------%%
%%        REQUISIÇÕES DE BEM         %%
%%-----------------------------------%%
\subsection{Requisicao Professor edicao Bem patrimonial}
\begin{entidade}{Requisicao\_Edicao\_Bem\_Patrimonial}
\campo{id}{SERIAL}{PK}
\campo{id\_bem\_patrimonial}{INTEGER}{FK, NOT NULL}
\campo{status}{ENUM}{pendente, aprovada, rejeitada, NOT NULL}
\campo{feita\_em}{TIMESTAMP}{NOT NULL}
\campo{respondida\_em}{TIMESTAMP}{NULL}
\campo{id\_usuario\_solicitante}{INTEGER}{FK, NOT NULL}
\campo{id\_usuario\_respondente}{INTEGER}{FK, NULL}
\campo{novo\_nome}{VARCHAR(150)}{NULL}
\campo{novo\_status}{ENUM}{Ativo, Inservivel, NULL}
\campo{novo\_estado\_conservacao}{ENUM}{BOM, REGULAR, RUIM, NULL}
\campo{novo\_id\_local}{INTEGER}{FK, NULL}
\campo{nova\_photo\_url}{TEXT}{NULL}
\campo{motivo}{TEXT}{NOT NULL}
\campo{justificativa\_resposta}{TEXT}{NULL}
\campo{novo\_id\_resumo\_bem\_patrimonial}{INTEGER}{FK, NULL}
\campo{versao\_bem\_origem}{INTEGER}{NOT NULL}
\end{entidade}

\begin{regra}
RF10: \textit{UNIQUE(id\_bem\_patrimonial) WHERE status = 'pendente'} -- impede mais de uma requisição de edição pendente simultânea para o mesmo bem patrimonial. Um índice único parcial resolve isso no PostgreSQL; no Firestore, a checagem precisa ser feita na própria \textit{Cloud Function} de criação, dentro de uma transação (ver seção \textit{Tecnologia e Relatórios}).
\end{regra}
\newpage
\subsection{Requisicao Professor adicao Bem patrimonial}
\begin{entidade}{Requisicao\_Adicao\_Bem\_Patrimonial}
\campo{id}{SERIAL}{PK}
\campo{numero\_patrimonio\_proposto}{VARCHAR(30)}{NOT NULL}
\campo{status}{ENUM}{pendente, aprovada, rejeitada, NOT NULL}
\campo{feita\_em}{TIMESTAMP}{NOT NULL}
\campo{id\_bem\_patrimonial\_se\_aprovado}{INTEGER}{FK, NULL}
\campo{respondida\_em}{TIMESTAMP}{NULL}
\campo{id\_usuario\_solicitante}{INTEGER}{FK, NOT NULL}
\campo{id\_usuario\_respondente}{INTEGER}{FK, NULL}
\campo{estado\_conservacao\_proposto}{ENUM}{BOM, REGULAR, RUIM, NOT NULL}
\campo{photo\_url\_proposta}{TEXT}{NOT NULL}
\campo{nome\_responsavel\_proposto}{VARCHAR(150)}{NOT NULL}
\campo{id\_local}{INTEGER}{FK, NOT NULL}
\campo{id\_resumo\_bem\_patrimonial}{INTEGER}{FK, NULL}
\campo{nome\_resumo\_proposto}{VARCHAR(150)}{NULL}
\campo{descricao\_resumo\_proposta}{TEXT}{NULL}
\campo{motivo}{TEXT}{NOT NULL}
\campo{justificativa\_resposta}{TEXT}{NULL}
\end{entidade}

\begin{explicacao}
Se o local desejado não existir no banco, o professor deve primeiro criar o \textit{Local} no sistema, para então vincular a respectiva FK (\textit{id\_local}) nesta requisição.
\end{explicacao}
\begin{regra}
RF10b: \textit{UNIQUE(numero\_patrimonio\_proposto) WHERE status = 'pendente'} --
impede mais de uma requisição de adição pendente simultânea para o mesmo
número de patrimônio. No PostgreSQL, isso é implementado como índice
único parcial; no Firestore, a verificação é reimplementada na Cloud
Function de criação da requisição, dentro de uma transação.
\end{regra}
\subsection{Lock de Requisição de Patrimônio}
\begin{entidade}{Locks\_Requisicao\_Patrimonio}
\campo{id (docId determinístico)}{VARCHAR(100)}{PK}
\campo{criado\_em}{TIMESTAMP}{NOT NULL}
\end{entidade}

\begin{regra}[Lock Determinístico para Unicidade de Requisição Pendente]
O Firestore não adquire trava pessimista sobre \textbf{ausência} de documento em uma query transacional. Por isso, a garantia de que não existam duas requisições pendentes para o mesmo bem ou plaqueta \textbf{não pode} ser implementada pela consulta \texttt{WHERE status = 'pendente'} dentro de uma transação -- dois clientes simultâneos lerão conjunto vazio, ambos passarão na validação e ambos gravarão, violando RF10/RF10b.

A solução é o padrão \textit{deterministic lock}: a Cloud Function tenta criar um documento cujo \textbf{docId é derivado deterministicamente} do recurso disputado. O Firestore garante que apenas uma gravação com o mesmo docId vence quando executadas concorrentemente dentro de uma transação.

\begin{itemize}
  \item Para requisição de \textbf{edição}: docId = \texttt{bem\_edicao\_\{id\_bem\_patrimonial\}}
  \item Para requisição de \textbf{adição}: docId = \texttt{bem\_adicao\_\{numero\_patrimonio\_proposto\}}
\end{itemize}

O lock é removido na transação de aprovação/rejeição ordinária. O lock \texttt{Locks\_Requisicao\_Patrimonio/bem\_edicao\_\{id\}} DEVE ser removido pela transação também em caso de rejeição por divergência de versão do bem patrimonial, solucionando deadlocks permanentes. Lock e requisição são criados atomicamente; não há commit parcial dessa transação. Limpeza verifica o vínculo e jamais remove um lock apenas porque a requisição pendente tem mais de 24 horas.
\end{regra}

\subsection{Impressão de Etiquetas de Frasco}

\begin{entidade}{Impressao\_Etiqueta\_Frasco}
\campo{id}{SERIAL}{PK}
\campo{gerado\_em}{TIMESTAMP}{NOT NULL}
\campo{gerado\_por}{INTEGER}{FK, NOT NULL}
\campo{codigo\_inicial}{INTEGER}{NOT NULL}
\campo{codigo\_final}{INTEGER}{NOT NULL}
\end{entidade}

\begin{explicacao}
Esta entidade audita de forma isolada as sessões de geração em lote de etiquetas virgens contínuas. \textit{codigo\_inicial} e \textit{codigo\_final} referem-se aos números sequenciais gerados (ex.: \texttt{LCQUI-1} a \texttt{LCQUI-50}). Não há vínculo de chave estrangeira com a tabela \textit{Frasco\_Reagente}, dado que a impressão em branco é puramente utilitária e não cria reservas de ID. A transação e consolidação final do código só ocorrem no banco no momento do cadastro físico.
\end{explicacao}

%%-----------------------------------%%
%%        ALMOXARIFADO / REAGENTES   %%
%%-----------------------------------%%
\subsection{Almoxarifado}
\begin{entidade}{Almoxarifado}
\campo{id}{SERIAL}{PK}
\campo{id\_local}{INTEGER}{FK, NOT NULL}
\campo{nome}{VARCHAR(100)}{NOT NULL}
\campo{descricao}{VARCHAR(500)}{NOT NULL}
\campo{ativo}{BOOLEAN}{DEFAULT TRUE, NOT NULL}
\end{entidade}

\begin{explicacao}
O campo \textit{ativo} indica se o almoxarifado está operacional. Quando \texttt{ativo = false}, o almoxarifado está desativado e as regras de bloqueio da Seção~\ref{sec:regras-almoxarifado-inativo} se aplicam. A desativação é reversível e não exclui registros históricos.
\end{explicacao}

\subsection{Substância Química}
\begin{entidade}{Substancia\_Quimica}
\campo{id}{SERIAL}{PK}
\campo{nome}{VARCHAR(150)}{NOT NULL}
\campo{cas\_number}{VARCHAR(15)}{UNIQUE, NULL}
\campo{formula\_quimica}{VARCHAR(50)}{NULL}
\campo{ativo}{BOOLEAN}{DEFAULT TRUE, NOT NULL}
\end{entidade}

\begin{explicacao}
A entidade \textit{Substancia\_Quimica} representa uma substância quimicamente identificável, independentemente de fabricante, concentração, pureza, lote ou frasco.

O campo \textit{cas\_number} identifica a substância e permanece \textbf{UNIQUE} (permitindo múltiplos NULL, já que algumas substâncias podem não ter CAS cadastrado). Uma mesma substância pode participar da composição de diversas especificações de reagentes diferentes.
\end{explicacao}
\newpage
\subsection{Resumo do Reagente / Solvente}
\begin{entidade}{Resumo\_Reagente}
\campo{id}{SERIAL}{PK}
\campo{nome}{VARCHAR(150)}{NOT NULL}
\campo{tipo\_substancia}{ENUM}{PURA, MISTURA, NOT NULL}
\campo{natureza\_quimica}{ENUM}{ORGANICO, INORGANICO, ELEMENTO, HIBRIDO, NOT NULL}
\campo{requer\_pesagem\_frequente}{BOOLEAN}{DEFAULT FALSE, NOT NULL}
\campo{frequencia\_pesagem\_dias}{INTEGER}{NULL}
\campo{estado\_fisico}{ENUM}{SOLIDO, LIQUIDO, NOT NULL}
\campo{eh\_higroscopico}{BOOLEAN}{DEFAULT FALSE, NOT NULL}
\end{entidade}
\newpage
\begin{explicacao}[Regra Fechada de Identidade do Resumo de Reagente]
O \textit{Resumo\_Reagente} representa o item químico agrupador apresentado no catálogo do sistema, permitindo agrupar diferentes especificações comerciais ou formulações do mesmo reagente.

Dois reagentes pertencem ao \textbf{mesmo} \textit{Resumo\_Reagente} quando compartilham a mesma identidade química nominal de catálogo (ex.: ``Ácido Clorídrico'', ``Etanol'', ``Hidróxido de Sódio'').
\begin{itemize}
    \item Variações de pureza (P.A., Técnico, Grau HPLC) $\implies$ \textbf{mesmo resumo}, especificações distintas.
    \item Variações de concentração (HCl 37\%, HCl 10\%, Etanol 99.5\%, Etanol 70\%) $\implies$ \textbf{mesmo resumo}, especificações distintas.
    \item Fabricante, código de catálogo e densidade $\implies$ atributos de especificação, \textbf{mesmo resumo}.
\end{itemize}
Devem ser cadastrados como \textbf{resumos diferentes}:
\begin{itemize}
    \item Espécies químicas distintas com fórmulas ou números CAS bases diferentes (ex.: Sulfato de Cobre II vs Sulfato de Ferro II);
    \item Soluções formuladas consagradas com finalidade analítica própria (ex.: ``Reagente de Benedict'', ``Solução de Fehling A'', ``Solução de Lugol'') $\implies$ cada uma constitui um \textit{Resumo\_Reagente} próprio do tipo \texttt{MISTURA}.
\end{itemize}

Os campos \textit{medida\_total\_usado} e \textit{unidade\_de\_medida} foram removidos desta entidade: ambos são métricas agregadas derivadas de \textit{Frasco\_Reagente} e \textit{Emprestimo\_Reagente}, e devem existir apenas na view materializada descrita na seção \textit{Métricas Agregadas}, nunca como coluna editável de uma tabela cadastral.

A unidade operacional é derivada de Resumo\_Reagente.estado\_fisico: SOLIDO usa g e LIQUIDO usa mL. Não persistir cópia redundante da unidade no modelo 3FN.

DP-A01: estado\_fisico e eh\_higroscopico pertencem ao Resumo\_Reagente. Estados físicos distintos nas condições de catálogo (25 graus Celsius, 1 atm) exigem resumos distintos; solução aquosa e sólido puro não são o mesmo resumo. Densidade e concentração permanecem na especificação comercial.

\textit{natureza\_quimica} classifica o item do catálogo como um todo (não uma \textit{Substancia\_Quimica} individual): \texttt{ELEMENTO} para substâncias simples (não se classificam como orgânicas nem inorgânicas), \texttt{ORGANICO}/\texttt{INORGANICO} para os demais casos, e \texttt{HIBRIDO} para fluidos complexos que não se decompõem em uma lista simples de substâncias em \textit{Composicao\_Reagente} (ex.: Soro Fetal Bovino). Como \texttt{HIBRIDO} não é derivável a partir das substâncias componentes, o campo é curado no nível do resumo, e não em \textit{Substancia\_Quimica}, evitando duas fontes de verdade diferentes para itens PUROS e para MISTURA/HÍBRIDO. O rótulo exibido na interface (ex.: ``Híbrido'', ``Complexo'' ou ``Biológico'') é independente do valor salvo no banco e pode ser alterado livremente.

O campo \textit{letra\_inicial}, presente em versões anteriores deste documento, foi removido do modelo 3FN: por ser derivável do próprio \textit{nome}, sua persistência aqui violaria a forma normal. Ele reaparece como uma denormalização proposital exclusiva da implementação em Firestore -- ver seção \textit{Notas de Mapeamento para Firestore}.
\end{explicacao}
\newpage
\subsection{Composição do Reagente}
\begin{entidade}{Composicao\_Reagente}
\campo{id}{SERIAL}{PK}
\campo{id\_especificacao\_reagente}{INTEGER}{FK, NOT NULL}
\campo{id\_substancia\_quimica}{INTEGER}{FK, NOT NULL}
\campo{valor\_min}{NUMERIC(10,5)}{NULL}
\campo{valor\_max}{NUMERIC(10,5)}{NULL}
\campo{tipo\_concentracao}{ENUM}{M\_M,V\_V,M\_V,MOL\_L,MOL\_KG,PPM,PPB}
\campo{notacao\_original\_fabricante}{VARCHAR(50)}{NULL}
\end{entidade}

\begin{explicacao}
A entidade \textit{Composicao\_Reagente} representa a relação entre uma especificação de reagente cadastrada e as substâncias químicas que a compõem.

A relação é do tipo N:N:

\begin{center}
\textit{Especificacao\_Reagente} $\leftrightarrow$ \textit{Substancia\_Quimica}
\end{center}

A composição foi movida do nível de \textit{Resumo\_Reagente} para o nível de \textit{Especificacao\_Reagente}, pois a concentração (37\%, 10\%, 1 mol/L etc.) é uma característica que varia entre especificações de um mesmo resumo, não uma propriedade do resumo em si. Vincular a composição ao resumo impediria registrar corretamente, por exemplo, HCl 37\% e HCl 10\% sob o mesmo \textit{Resumo\_Reagente} com valores de composição diferentes.

Uma substância pura normalmente terá apenas uma linha associada, enquanto uma mistura poderá possuir diversas substâncias.

Q03/DP-A03: concentração pontual usa valor\_min = valor\_max ou valor\_max NULL; faixa usa limites numéricos ordenados. CHECK (valor\_max IS NULL OR valor\_min <= valor\_max). O servidor rejeita faixa com máximo informado e mínimo ausente; preserva a notação original do fabricante separadamente, sem usá-la como autoridade numérica.
\end{explicacao}

\begin{regra}
\textit{UNIQUE(id\_especificacao\_reagente, id\_substancia\_quimica)} -- impede duas linhas de composição para a mesma substância dentro da mesma especificação.
\end{regra}

\subsection{Especificação do Reagente / Solvente}
\begin{entidade}{Especificacao\_Reagente}
\campo{id}{SERIAL}{PK}
\campo{id\_resumo\_reagente}{INTEGER}{FK, NOT NULL}
\campo{descricao}{VARCHAR(150)}{NOT NULL}
\campo{fabricante}{VARCHAR(150)}{NULL}
\campo{codigo\_produto\_fabricante}{VARCHAR(100)}{NULL}
\campo{grau\_pureza}{VARCHAR(30)}{NULL}
\campo{densidade}{NUMERIC(8,4)}{NULL}
\campo{classe\_inflamabilidade}{ENUM}{NAO\_INFLAMAVEL,CLASSE\_1,CLASSE\_2,CLASSE\_3, NOT NULL}
\campo{eh\_controlado\_pf}{BOOLEAN}{DEFAULT FALSE, NOT NULL}
\campo{eh\_controlado\_eb}{BOOLEAN}{DEFAULT FALSE, NOT NULL}
\campo{link\_fds\_fispq}{VARCHAR(255)}{NULL}
\end{entidade}

\begin{explicacao}
A especificação referencia o resumo, fonte do estado físico; a unidade é derivada desse estado. Frasco e empréstimo chegam ao resumo pela cadeia de FKs. Snapshots de consulta e de cálculo são documentados separadamente no Firestore.

\textit{GASOSO} foi removido de \textit{estado\_fisico} nesta versão: nenhum reagente gasoso está listado para o almoxarifado atual, e reagentes gasosos precisam de campos e método de medição próprios (pressão, não peso em balança) em vez de reaproveitar os campos já pensados para sólido/líquido. O estudo detalhado para reintrodução está na seção \textit{Implementações em Estudo para Versões Futuras}.
\end{explicacao}

\subsection{Estoque Mínimo por Almoxarifado}
\begin{entidade}{Estoque\_Minimo\_Almoxarifado}
\campo{id\_especificacao\_reagente}{INTEGER}{FK, PK, NOT NULL}
\campo{id\_almoxarifado}{INTEGER}{FK, PK, NOT NULL}
\campo{qtd\_limiar\_escassez}{INTEGER}{NOT NULL}
\campo{ativo}{BOOLEAN}{DEFAULT TRUE, NOT NULL}
\campo{notificacao\_ativa}{BOOLEAN}{DEFAULT TRUE, NOT NULL}
\end{entidade}

\begin{explicacao}
A entidade \textit{Estoque\_Minimo\_Almoxarifado} estabelece o limiar de escassez (em quantidade de frascos) de uma \textit{Especificacao\_Reagente} para um determinado \textit{Almoxarifado}. O par Especificação × Almoxarifado é a identidade unívoca da configuração. No Firestore, esta entidade é implementada como a subcoleção \texttt{Almoxarifado/\{id\}/Estoques\_Configurados/\{idResumo\_idEspec\}}.
O campo \textit{ativo} indica que o par continua configurado ou intencionalmente estocado. 
O campo \textit{notificacao\_ativa} indica que a emissão automática do alerta está habilitada (útil para silenciar alertas sem excluir a intenção de estocar a especificação).
\end{explicacao}

\subsection{Lote}
\begin{entidade}{Lote}
\campo{id}{SERIAL}{PK}
\campo{id\_especificacao\_reagente}{INTEGER}{FK, NOT NULL}
\campo{data\_aquisicao}{TIMESTAMP}{NOT NULL}
\campo{qtd\_frascos\_comprados}{INTEGER}{DEFAULT 0, NOT NULL}
\campo{nome\_fornecedor}{VARCHAR(80)}{NOT NULL}
\campo{numero\_lote}{VARCHAR(100)}{NOT NULL}
\campo{nota\_fiscal}{VARCHAR(100)}{NOT NULL}
\campo{data\_fabricacao}{DATE}{NOT NULL}
\campo{data\_validade}{DATE}{NOT NULL}
\end{entidade}
\begin{regra}
\textit{UNIQUE(id\_especificacao\_reagente, nome\_fornecedor, numero\_lote)}.
\end{regra}
\begin{explicacao}
\textit{qtd\_frascos\_comprados} é dado de entrada (o que consta na nota fiscal) e permanece nesta tabela. Já \textit{qtd\_frascos\_cadastrados} -- quantos frascos deste lote já foram individualmente registrados em \textit{Frasco\_Reagente} -- é uma contagem agregada (\texttt{COUNT(*) FROM Frasco\_Reagente WHERE id\_lote = X}) e, por isso, \textbf{não} é coluna desta tabela cadastral: ela vive exclusivamente em \textit{Lote\_Materializado} (ver seção \textit{Materialized (Views)}), mantida por mecanismo transacional a cada frasco cadastrado ou removido.
\end{explicacao}
\subsection{Frasco de Reagente / Solvente}
\begin{entidade}{Frasco\_Reagente}
\campo{id}{SERIAL}{PK}
\campo{id\_almoxarifado}{INTEGER}{FK, NOT NULL}
\campo{id\_lote}{INTEGER}{FK, NULL}
\campo{id\_especificacao\_reagente}{INTEGER}{FK, NULL}
\campo{detalhe\_local\_armazenamento}{TEXT}{NULL}
\campo{codigo\_frasco}{TEXT}{UNIQUE, NOT NULL}
\campo{data\_ultima\_pesagem}{TIMESTAMP}{NULL}
\campo{conteudo\_nominal}{NUMERIC(10,3)}{NULL}
\campo{peso\_no\_cadastrado}{NUMERIC(10,3)}{NOT NULL}
\campo{peso\_atual}{NUMERIC(10,3)}{NOT NULL}
\campo{peso\_frasco\_vazio}{NUMERIC(10,3)}{NULL}
\campo{medida\_usada}{NUMERIC(10,3)}{DEFAULT 0, NOT NULL}
\campo{data\_abertura}{DATE}{NULL}
\campo{validade\_fechado}{DATE}{NULL}
\campo{validade\_efetiva}{DATE}{NULL}
\campo{validade\_desconhecida}{BOOLEAN}{DEFAULT FALSE, NOT NULL}
\campo{validade\_apos\_aberto\_dias}{INTEGER}{NULL}
\campo{estado\_fisico\_frasco}{ENUM}{FECHADO,ABERTO,VAZIO,QUEBRADO,DESCARTADO,EXTRAVIADO, NOT NULL}
\campo{disponibilidade}{ENUM}{DISPONIVEL, EMPRESTADO, INDISPONIVEL, NOT NULL}
\campo{vencido}{BOOLEAN}{DEFAULT FALSE, NOT NULL}
\campo{em\_quarentena}{BOOLEAN}{DEFAULT FALSE, NOT NULL}
\campo{uso\_vencido\_autorizado}{BOOLEAN}{DEFAULT FALSE, NOT NULL}
\campo{detalhe\_status}{TEXT}{NULL}
\campo{cadastrado\_em}{TIMESTAMP}{NOT NULL}
\campo{cadastrado\_por}{INTEGER}{FK, NOT NULL}
\campo{abertura\_historica\_desconhecida}{BOOLEAN}{DEFAULT FALSE, NOT NULL}
\end{entidade}

\begin{regra}[Constraint de Identidade Química do Frasco]
\textit{CHECK (id\_lote IS NOT NULL OR id\_especificacao\_reagente IS NOT NULL)} -- pelo menos um dos dois deve existir. Quando \textit{id\_lote} é informado, \textit{id\_especificacao\_reagente} fica NULL (a especificação é derivável via \textit{Lote.id\_especificacao\_reagente}). Quando \textit{id\_lote} é NULL, \textit{id\_especificacao\_reagente} é obrigatório para manter o caminho até a densidade, unidade de medida e resumo do reagente.
\end{regra}

\begin{explicacao}
\textit{conteudo\_nominal} (antigo \textit{capacidade\_nominal}) representa o conteúdo declarado no cadastro: para líquidos, é o volume nominal em mL; para sólidos, é a massa nominal em g. Se NULL, o valor original do rótulo é desconhecido ou ilegível (não usar zero para representar desconhecido).

\textit{peso\_no\_cadastrado} é o peso total (frasco + conteúdo) registrado no momento do cadastro. Este campo é \textbf{imutável} após o cadastro -- nunca é sobrescrito. \textit{peso\_atual} é o peso total corrente, inicializado com o valor de \textit{peso\_no\_cadastrado} e atualizado a cada devolução de empréstimo ou pesagem de rotina.

Os campos \textit{peso\_no\_cadastrado}, \textit{peso\_atual}, \textit{peso\_frasco\_vazio} e \textit{medida\_usada} são sempre leituras de balança em gramas, independentemente de o conteúdo ser sólido ou líquido, já que a operação real do laboratório consiste em pesar o frasco; o volume, quando aplicável, é sempre derivado por conversão via densidade.

O campo \textit{unidade\_de\_medida}, antes duplicado nesta entidade, foi removido: a unidade do produto (ml ou g) é obtida via \textit{Frasco\_Reagente} $\rightarrow$ \textit{Lote} $\rightarrow$ \textit{Especificacao\_Reagente} (ou diretamente via \textit{id\_especificacao\_reagente} quando o lote é NULL), chegando ao Resumo\_Reagente, cujo estado determina a unidade.

O antigo campo único \textit{status} (com valores combinados como \texttt{VENCIDO\_DISPONIVEL}, \texttt{VENCIDO\_EMPRESTADO}) foi decomposto em dimensões ortogonais -- \textit{estado\_fisico\_frasco}, \textit{disponibilidade}, \textit{vencido} e \textit{em\_quarentena} -- evitando o crescimento combinatorial do enum e simplificando consultas e regras de notificação.

O campo \textit{id\_usuario\_em\_uso} foi removido: quem está com o frasco no momento é obtido consultando \textit{Emprestimo\_Reagente} pelo empréstimo ativo (\textit{status} \texttt{EM\_USO} ou \texttt{ATRASADO}) daquele \textit{id\_frasco\_reagente}.

\textit{validade\_fechado} é mantido de forma independente de \textit{Lote.data\_validade} porque \textit{id\_lote} é opcional: um frasco pode ser cadastrado sem lote associado, e mesmo quando o lote existe, o formulário apenas pré-preenche este campo a partir dele -- o valor permanece editável e não é uma cópia somente-leitura.

O campo \textit{uso\_vencido\_autorizado} fecha a distinção de negócio entre um frasco vencido que permanece apto ao uso didático e um frasco vencido que aguarda descarte.
\end{explicacao}

\subsection{Histórico do Frasco de Reagente / Solvente}
\begin{entidade}{Historico\_Frasco\_Reagente}
\campo{id}{SERIAL}{PK}
\campo{id\_frasco\_reagente}{INTEGER}{FK, NOT NULL}
\campo{id\_almoxarifado}{INTEGER}{FK, NOT NULL}
\campo{id\_gestor}{INTEGER}{FK, NOT NULL}
\subcampo{Para ações automáticas do servidor (jobs de vencimento, atraso), no 3FN usa-se o id\_usuario = 0 (\'SISTEMA LCQUI\'); no Firestore, usa-se a string sentinela \_\_SISTEMA\_\_.}
\campo{tipo}{ENUM}{
  \shortstack[l]{%
    CADASTRO, SAIU, ENTROU, FICOU\_VAZIO, QUEBROU, FOI\_DESCARTADO,\\
    VENCEU, INVENTARIO\_ROTINA, EVAPORACAO, AJUSTE, CONCLUSAO,\\
    ENTROU\_EM\_QUARENTENA, LIBERADO\_QUARENTENA, PENDENTE\_DE\_DESCARTE,\\
    USO\_VENCIDO\_AUTORIZADO, EXTRAVIO, REENCONTRO%
  }
}{NOT NULL}
\campo{campo\_ajustado}{VARCHAR(30)}{NULL}
\campo{id\_emprestimo\_reagente}{INTEGER}{FK, NULL}
\campo{peso\_anterior}{NUMERIC(10,3)}{NULL}
\campo{peso\_novo}{NUMERIC(10,3)}{NULL}
\campo{medida\_ajustada}{NUMERIC(10,3)}{NULL}
\campo{unidade\_medida\_ajustada}{ENUM}{ml,g, NULL}
\campo{motivo}{TEXT}{NULL}
\campo{timestamp}{TIMESTAMP}{NOT NULL}
\end{entidade}

\begin{explicacao}
\textit{peso\_anterior} e \textit{peso\_novo} são sempre leituras de balança em gramas. Já \textit{unidade\_medida\_ajustada} qualifica exclusivamente \textit{medida\_ajustada}, indicando se o ajuste concluído é reportado em massa ou volume. Podem ser NULL caso o evento histórico não seja de natureza quantitativa. O campo \textit{motivo} é genérico e aceita NULL, porém é obrigatório para eventos que exigem justificativa (ex.: EXTRAVIO, REENCONTRO e ajustes manuais). No modelo relacional 3FN, reserva-se a linha canônica \texttt{Usuario.id = 0} com nome \'SISTEMA LCQUI\' para satisfazer a \textit{foreign key} não nula em rotinas automáticas de servidor; no Firestore, o documento recebe a string sentinela \texttt{\_\_SISTEMA\_\_}.
\end{explicacao}

\subsection{Empréstimo de Reagente / Solvente}
\begin{entidade}{Emprestimo\_Reagente}
\campo{id}{SERIAL}{PK}
\campo{id\_frasco\_reagente}{INTEGER}{FK, NOT NULL}
\campo{id\_usuario\_retirou}{INTEGER}{FK, NOT NULL}
\campo{id\_gestor\_retirada}{INTEGER}{FK, NOT NULL}
\campo{status}{ENUM}{EM\_USO,ATRASADO,DEVOLVIDO,DEVOLVIDO\_COM\_ATRASO,ENCERRADO\_EXTRAORDINARIO, NOT NULL}
\campo{data\_retirada}{TIMESTAMP}{NOT NULL}
\campo{data\_devolucao\_prevista}{DATE}{NOT NULL}
\campo{data\_devolucao\_efetuada}{TIMESTAMP}{NULL}
\campo{id\_local\_usado}{INTEGER}{FK, NOT NULL}
\campo{id\_usuario\_devolveu}{INTEGER}{FK, NULL}
\campo{id\_gestor\_devolucao}{INTEGER}{FK, NULL}
\campo{peso\_saida}{NUMERIC(10,3)}{NOT NULL}
\campo{peso\_retorno}{NUMERIC(10,3)}{NULL}
\campo{medida\_utilizada}{NUMERIC(10,3)}{NULL}
\campo{unidade\_medida\_utilizada}{ENUM}{ml,g, NOT NULL}
\campo{peso\_perda\_evaporacao}{NUMERIC(10,3)}{DEFAULT 0, NOT NULL}
\campo{uso\_vencido\_aceito}{BOOLEAN}{DEFAULT FALSE, NOT NULL}
\campo{finalidade\_uso}{ENUM}{NOT NULL}
\subcampo{AULA\_PRATICA, DEMONSTRACAO, PESQUISA\_TCC\_POS, ESTUDO\_DEGRADACAO\_RESIDUOS}
\campo{auto\_atendimento}{BOOLEAN}{DEFAULT FALSE, NOT NULL}
\campo{justificativa\_metodologica}{TEXT}{NULL}
\campo{tcr\_versao}{VARCHAR(100)}{NULL}
\campo{tcr\_aceito\_em}{TIMESTAMP}{NULL}
\campo{tcr\_aceito\_por}{INTEGER}{FK, NULL}
\campo{tcr\_auditoria\_id}{INTEGER}{FK, NULL}
\campo{tipo\_encerramento\_excepcional}{ENUM}{QUEBRA\_ACIDENTAL, EXTRAVIO\_SINISTRO, NULL}
\campo{motivo\_encerramento\_excepcional}{TEXT}{NULL}
\campo{id\_gestor\_encerramento}{INTEGER}{FK, NULL}
\campo{massa\_perda\_estimada\_g}{NUMERIC(10,3)}{NULL}
\end{entidade}

\begin{explicacao}
\textbf{Semântica dos atores da retirada:} \textit{id\_usuario\_retirou} identifica o \textbf{portador físico} -- o Professor ou Bolsista que receberá e usará o reagente. \textit{id\_gestor\_retirada} identifica o \textbf{operador do almoxarifado} -- o Gestor de Almoxarifado (ou Chefe Geral) que registrou a operação no sistema. No autoatendimento, os dois UIDs coincidem e a flag \textit{auto\_atendimento} fica \texttt{true}. Esta separação é essencial para rastreabilidade de auditoria: o relatório de atividade do gestor parte de \textit{id\_gestor\_retirada}, enquanto o painel ``Meus Reagentes'' do professor parte de \textit{id\_usuario\_retirou}.

\textbf{Semântica dos atores da devolução:} \textit{id\_usuario\_devolveu} identifica o portador físico que entrega o frasco ao balcão, enquanto \textit{id\_gestor\_devolucao} identifica o Gestor de Almoxarifado que operou e registrou a devolução na balança.

\textit{peso\_saida} e \textit{peso\_retorno} são sempre leituras de balança em gramas. \textit{unidade\_medida\_utilizada} qualifica exclusivamente \textit{medida\_utilizada} -- o consumo calculado, reportado em massa (g) para sólidos ou em volume (ml) para líquidos, via $V = \Delta m / \rho$.

O campo \textit{data\_devolucao\_prevista} em nível NoSQL físico armazena uma string civil no padrão YYYY-MM-DD (fuso America/Sao\_Paulo), e o vencimento legal encerra-se estritamente às 23:59:59.999 do referido dia civil.
\end{explicacao}



\subsection{Métricas Agregadas do Reagente}
\begin{explicacao}
As seguintes informações devem ser tratadas preferencialmente como \textbf{views materializadas, consultas agregadas ou campos derivados mantidos por mecanismo transacional}:
\begin{itemize}
\item \textit{medida\_total\_usado} e \textit{unidade\_de\_medida} (agregados por \textit{Resumo\_Reagente});
\item \textit{qtd\_frascos\_fechados\_disponiveis};
\item \textit{qtd\_frascos\_abertos\_disponiveis};
\item \textit{qtd\_frascos\_em\_uso};
\item \textit{qtd\_frascos\_vazios};
\item \textit{qtd\_frascos\_descartados};
\item \textit{qtd\_frascos\_vencidos};
\item \textit{qtd\_frascos\_quebrados}.
\end{itemize}
\end{explicacao}

%%-----------------------------------%%
%%        TURMAS / ROTEIROS          %%
%%-----------------------------------%%
\subsection{Turma}
\begin{entidade}{Turma}
\campo{id}{SERIAL}{PK}
\campo{id\_materia}{INTEGER}{FK, NOT NULL}
\campo{id\_professor}{INTEGER}{FK, NOT NULL}
\campo{status}{ENUM}{Ativo, Arquivada, NOT NULL}
\campo{nome\_turma}{VARCHAR(100)}{NOT NULL}
\campo{ano}{INTEGER}{NOT NULL}
\campo{semestre}{SMALLINT}{CHECK (semestre IN (1,2)), NOT NULL}
\campo{capacidade}{INTEGER}{NOT NULL}
\campo{codigo\_turma}{VARCHAR(20)}{UNIQUE, NOT NULL}
\campo{data\_criacao}{TIMESTAMP}{NOT NULL}
\end{entidade}

\begin{explicacao}
Os campos \textit{ano} e \textit{semestre} representam o per\'{i}odo letivo da turma (ex.: 2026/1). Ambos s\~{a}o exibidos na listagem de turmas arquivadas do professor -- eliminando a necessidade de parsing do \textit{nome\_turma} ou infer\^{e}ncia a partir de \textit{data\_criacao}. O modal de cria\c{c}\~{a}o de turma solicita esses dois campos al\'{e}m dos demais.
\end{explicacao}

\begin{regra}[RN-TUR-01: Controle de Capacidade da Turma]
O ingresso de aluno por código de turma exige obrigatoriamente que a quantidade de alunos matriculados ativos seja estritamente menor que a capacidade da turma:
\[
\text{COUNT}(\text{Aluno\_x\_Turma WHERE id\_turma} = X) < \text{Turma.capacidade}
\]
Caso a capacidade seja atingida, o sistema rejeita o ingresso por código. Apenas o professor responsável pode ultrapassar a capacidade mediante envio de convite nominal por e-mail, registrando justificativa em auditoria.
\end{regra}

\subsection{Historico de Alunos na Turma}
\begin{entidade}{Historico\_Alunos\_Turma}
\campo{id}{SERIAL}{PK}
\campo{id\_turma}{INTEGER}{FK, NOT NULL}
\campo{tipo}{ENUM}{inclusao\_aluno, exclusao\_aluno, NOT NULL}
\campo{id\_aluno}{INTEGER}{FK, NOT NULL}
\campo{modo\_ingresso}{ENUM}{CODIGO, CONVITE, NULL}
\campo{justificativa}{TEXT}{NULL}
\campo{timestamp}{TIMESTAMP}{NOT NULL}
\campo{removido\_por}{INTEGER}{FK, NULL}
\subcampo{Obrigatório quando tipo = 'exclusao\_aluno'}
\end{entidade}

\subsection{Historico de Posts na Turma}
\begin{entidade}{Historico\_Posts\_Turma}
\campo{id}{SERIAL}{PK}
\campo{id\_post}{INTEGER}{FK, NOT NULL}
\campo{editado\_por}{INTEGER}{FK, NOT NULL}
\campo{novo\_titulo}{VARCHAR(150)}{NULL}
\campo{novo\_id\_roteiro}{INTEGER}{FK, NULL}
\campo{nova\_descricao}{TEXT}{NULL}
\campo{antigo\_titulo}{VARCHAR(150)}{NULL}
\campo{antigo\_id\_roteiro}{INTEGER}{FK, NULL}
\campo{antiga\_descricao}{TEXT}{NULL}
\campo{timestamp}{TIMESTAMP}{NOT NULL}
\end{entidade}

\begin{explicacao}
O campo \textit{editado\_por} registra o UID do usuário que realizou a edição. Normalmente é o Professor dono da turma, mas o Chefe Geral também pode editar ou remover posts em caráter excepcional de moderação institucional (Q13). O \textit{id\_turma} não é armazenado nesta entidade 3FN porque pode ser obtido por \textit{Historico\_Posts\_Turma.id\_post} $\rightarrow$ \textit{Post.id\_turma}; se a consulta direta por turma for necessária no Firestore, esse vínculo poderá ser denormalizado na coleção de histórico.
\end{explicacao}

\subsection{Historico de Comentario em um Post}
\begin{entidade}{Historico\_Comentario}
\campo{id}{SERIAL}{PK}
\campo{id\_comentario}{INTEGER}{FK, NOT NULL}
\campo{novo\_texto}{TEXT}{NOT NULL}
\campo{texto\_antigo}{TEXT}{NOT NULL}
\campo{editado\_em}{TIMESTAMP}{NOT NULL}
\campo{editado\_por}{INTEGER}{FK, NOT NULL}
\end{entidade}

\subsection{Roteiro de Experimento}
\begin{entidade}{Roteiro\_Experimento}
\campo{id}{SERIAL}{PK}
\campo{id\_professor\_upload}{INTEGER}{FK, NOT NULL}
\campo{file\_url}{TEXT}{NOT NULL}
\campo{nome}{VARCHAR(150)}{NOT NULL}
\campo{descricao}{TEXT}{NOT NULL}
\end{entidade}

\subsection{Roteiro Professor Compartilhado}
\begin{entidade}{Roteiro\_Professor\_Compartilhado}
\campo{id\_roteiro\_experimento}{INTEGER}{PK, FK}
\campo{id\_professor}{INTEGER}{PK, FK}
\end{entidade}

\subsection{Professor}
\begin{entidade}{Professor}
\campo{id\_usuario}{INTEGER}{PK, FK}
\campo{centro}{ENUM}{CCT, CCTA, CBB, CCH, NOT NULL}
\campo{laboratorio}{VARCHAR(100)}{NOT NULL}
\end{entidade}

\subsection{Materia}
\begin{entidade}{Materia}
\campo{id}{SERIAL}{PK}
\campo{nome}{VARCHAR(100)}{NOT NULL}
\campo{codigo\_materia}{VARCHAR(10)}{UNIQUE, NOT NULL}
\end{entidade}

\subsection{Tabela Professor \texorpdfstring{$\times$}{x} Materia (N para N)}
\begin{entidade}{Professor\_x\_Materia}
\campo{id\_professor}{INTEGER}{PK, FK}
\campo{id\_materia}{INTEGER}{PK, FK}
\end{entidade}
\begin{regra}
A chave primária composta \textit{(id\_professor, id\_materia)} impede vincular o mesmo professor duas vezes à mesma matéria.
\end{regra}

\subsection{Aluno}
\begin{entidade}{Aluno}
\campo{id\_usuario}{INTEGER}{PK, FK}
\campo{numero\_matricula}{VARCHAR(20)}{UNIQUE, NOT NULL}
\end{entidade}

\subsection{Bolsista}
\begin{entidade}{Bolsista}
\campo{id\_usuario}{INTEGER}{PK, FK}
\end{entidade}

\subsection{Convite para Aluno}
\begin{entidade}{Convite\_Aluno}
\campo{id}{SERIAL}{PK}
\campo{id\_turma}{INTEGER}{FK, NULL}
\campo{email}{VARCHAR(150)}{NOT NULL}
\campo{convidado\_em}{TIMESTAMP}{NOT NULL}
\campo{status}{ENUM}{pendente, aceitado, expirado, NOT NULL}
\campo{expira\_em}{TIMESTAMP}{NOT NULL}
\campo{convidado\_por}{INTEGER}{FK, NOT NULL}
\campo{numero\_matricula}{VARCHAR(20)}{NULL}
\campo{token\_hash}{VARCHAR(100)}{NOT NULL}
\campo{ultimo\_reenvio\_por}{INTEGER}{FK, NULL}
\campo{exceder\_capacidade}{BOOLEAN}{DEFAULT FALSE, NOT NULL}
\campo{justificativa\_excecao}{TEXT}{NULL}
\campo{aceitado\_por}{INTEGER}{FK, NULL}
\campo{aceitado\_em}{TIMESTAMP}{NULL}
\end{entidade}

\begin{regra}
A combinação \textit{email + id\_turma} não pode possuir mais de um convite \textit{pendente}. Um mesmo e-mail pode possuir convites para turmas diferentes e um aluno pode participar de várias turmas; por isso, o e-mail não é globalmente único nesta entidade.
\end{regra}

\begin{explicacao}
\textit{numero\_matricula} aqui é opcional e \textbf{não é UNIQUE}, pois serve apenas para pré-declarar a matrícula quando o professor a conhece no momento do convite. Ao aceitar o convite, o sistema cria as linhas \textit{Usuario} e \textit{Aluno} correspondentes, copiando \textit{email} e \textit{numero\_matricula} para a nova linha de \textit{Aluno}, onde a restrição \textbf{UNIQUE} é rigorosamente aplicada.
\end{explicacao}

\subsection{Notificação (Unificada)}
\begin{entidade}{Notificacao}
\campo{id}{SERIAL}{PK}
\campo{id\_destinatario}{INTEGER}{FK, NOT NULL}
\campo{papel\_destinatario}{ENUM}{
  \shortstack[l]{%
    Aluno, Professor, Bolsista, Gestor\_Bens\_Patrimoniais,\\
    Gestor\_Almoxarifado, Chefe\_Geral%
  }
}{NOT NULL}
\campo{tipo}{ENUM}{
  \shortstack[l]{%
    ADICIONADO, POST, COMENTARIO, REMOVIDO,\\
    TURMA\_ARQUIVADA, TURMA\_DESARQUIVADA,\\
    REQUISICAO\_BEM, DATA\_DEVOLUCAO\_REAGENTE,\\
    ROTEIRO\_COMPARTILHADO,\\
    REQUISICAO\_EDICAO\_BEM, REQUISICAO\_ADICAO\_BEM,\\
    BEM\_INSERVIVEL,\\
    ENTREGA\_ATRASADA, FRASCOS\_VAZIOS, FRASCOS\_QUEBRADOS,\\
    FRASCOS\_VENCIDOS, FRASCOS\_A\_SEREM\_PESADOS,\\
    FRASCOS\_EM\_QUARENTENA, ESCASSEZ\_ESTOQUE,\\
    AUTO\_ATENDIMENTO\_RETIRADA%
  }
}{NOT NULL}
\campo{id\_quem\_fez\_acao}{INTEGER}{FK, NULL}
\campo{id\_turma}{INTEGER}{FK, NULL}
\campo{quantidade}{INTEGER}{NULL}
\campo{entidade\_alvo}{VARCHAR(40)}{NOT NULL}
\campo{id\_alvo}{VARCHAR(200)}{NOT NULL}
\campo{lida}{BOOLEAN}{DEFAULT FALSE, NOT NULL}
\campo{lida\_em}{TIMESTAMP}{NULL}
\campo{emitida\_em}{TIMESTAMP}{NOT NULL}
\campo{expira\_em}{TIMESTAMP}{ NULL}
\end{entidade}

\begin{explicacao}[Unificação das Notificações]
As quatro tabelas de notificação anteriores (\textit{Notificacao\_para\_Aluno}, \textit{Notificacao\_para\_Professor}, \textit{Notificacao\_para\_Gestor\_de\_bens\_Patrimoniais}, \textit{Notificacao\_para\_Gestor\_de\_Reagentes}) foram consolidadas em uma única entidade \textit{Notificacao}. O campo \textit{papel\_destinatario} indica em qual contexto de papel a notificação deve ser exibida, incluindo Aluno, Professor, Bolsista, Gestor\_Bens\_Patrimoniais, Gestor\_Almoxarifado e Chefe\_Geral.

No Firestore, esta entidade é implementada como sub-coleção \texttt{Usuarios/\{uid\}/Notificacoes}, eliminando a necessidade de queries em múltiplas coleções para usuários multi-role. A ação de ``Limpar tudo'' na interface marca \texttt{lida = true} e \texttt{lida\_em = NOW()} em todas as notificações ativas do usuário, preservando o histórico de auditoria. Os campos \textit{entidade\_alvo} e \textit{id\_alvo} viabilizam a navegação direta (\textit{deep linking}) para o objeto de origem ao clicar no aviso. O campo \textit{expira\_em} é nulo para notificações operacionais persistentes (como \texttt{ESCASSEZ\_ESTOQUE}), que não possuem expiração temporal automática.
\end{explicacao}

\subsection{Tabela Aluno Turma (N para N)}
\begin{entidade}{Aluno\_x\_Turma}
\campo{id\_aluno}{INTEGER}{PK, FK}
\campo{id\_turma}{INTEGER}{PK, FK}
\end{entidade}
\begin{regra}
A chave primária composta \textit{(id\_aluno, id\_turma)} impede matricular o mesmo aluno duas vezes na mesma turma.
\end{regra}

\subsection{Local}
\begin{entidade}{Local}
\campo{id}{SERIAL}{PK}
\campo{predio}{VARCHAR(30)}{NOT NULL}
\campo{andar}{VARCHAR(10)}{NOT NULL}
\campo{sala}{VARCHAR(30)}{NOT NULL}
\campo{criado\_por}{INTEGER}{FK, NOT NULL}
\campo{criado\_em}{TIMESTAMP}{NOT NULL}
\end{entidade}
\begin{regra}
A combinação de prédio, andar e sala deve ser única (\textit{UNIQUE(predio, andar, sala)}).
\end{regra}

\subsection{Post}
\begin{entidade}{Post}
\campo{id}{SERIAL}{PK}
\campo{id\_professor}{INTEGER}{FK, NOT NULL}
\campo{id\_turma}{INTEGER}{FK, NOT NULL}
\campo{id\_roteiro\_experimento}{INTEGER}{FK, NULL}
\campo{titulo}{VARCHAR(150)}{NOT NULL}
\campo{descricao}{TEXT}{NOT NULL}
\campo{criado\_em}{TIMESTAMP}{NOT NULL}
\end{entidade}

\subsection{Comentario}
\begin{entidade}{Comentario}
\campo{id}{SERIAL}{PK}
\campo{id\_post}{INTEGER}{FK, NOT NULL}
\campo{id\_usuario}{INTEGER}{FK, NOT NULL}
\campo{texto}{TEXT}{NOT NULL}
\campo{criado\_em}{TIMESTAMP}{NOT NULL}
\campo{moderado}{BOOLEAN}{DEFAULT FALSE, NOT NULL}
\campo{motivo\_moderacao}{TEXT}{NULL}
\campo{moderado\_por}{INTEGER}{FK, NULL}
\end{entidade}

\subsection{Registro de Auditoria}
\begin{entidade}{Registro\_de\_Auditoria}
\campo{id}{SERIAL}{PK}
\campo{id\_usuario}{INTEGER}{FK, NOT NULL}
\campo{acao}{VARCHAR(200)}{NOT NULL}
\campo{tipo\_entidade\_sofre\_acao}{ENUM}{
  \shortstack[l]{%
    TURMA, ALUNO, PROFESSOR, ALMOXARIFADO, RESUMO\_REAGENTE,\\
    ESPECIFICACAO\_REAGENTE, LOTE, FRASCO\_REAGENTE,\\
    EMPRESTIMO\_REAGENTE, BEM\_PATRIMONIAL, LOCAL, MATERIA,\\
    COMENTARIO, POST, ROTEIRO, REQUISICAO\_ADICAO\_BEM,\\
    REQUISICAO\_EDICAO\_BEM, USUARIO%
  }
}{NOT NULL}
\campo{id\_do\_objeto\_da\_entidade}{VARCHAR(200)}{NOT NULL}
\campo{acao\_feita\_em}{TIMESTAMP}{NOT NULL}
\campo{metadata}{JSONB}{NOT NULL}
\end{entidade}

\begin{explicacao}
\textit{acao} e \textit{id\_do\_objeto\_da\_entidade} tinham o tamanho \texttt{200} escrito como se fosse um selo de restrição em vez de parte do tipo; corrigido para \texttt{VARCHAR(200)}. O valor \texttt{REAGENTE} do enum foi substituído por valores granulares.
\textbf{Nota sobre Impressão:} A reimpressão avulsa (2ª via) de um frasco já cadastrado é auditada diretamente aqui (\texttt{tipo = FRASCO\_REAGENTE}, \texttt{acao = REIMPRESSAO\_ETIQUETA}) em vez de usar a entidade \textit{Impressao\_Etiqueta\_Frasco}.
\end{explicacao}
\newpage

\subsection{Especificação de Integridade do Modelo 3FN}
Esta subseção fecha as propriedades lógicas que devem permanecer verdadeiras antes de qualquer denormalização específica do Firestore.

\subsubsection{Cardinalidades}
\begin{longtable}{>{\raggedright\arraybackslash}p{0.36\textwidth} >{\centering\arraybackslash}p{0.14\textwidth} >{\raggedright\arraybackslash}p{0.42\textwidth}}
\toprule
\textbf{Relacionamento} & \textbf{Card.} & \textbf{Regra}\\
\midrule
\endfirsthead
\toprule
\textbf{Relacionamento} & \textbf{Card.} & \textbf{Regra} (continuação)\\
\midrule
\endhead
\bottomrule
\endfoot
Usuario -- cada papel específico & 0..1 por papel & Um usuário pode ter no máximo uma linha em cada tabela de papel; várias tabelas de papel podem estar associadas ao mesmo usuário conforme as combinações permitidas.\\
Gestor Almoxarifado -- Almoxarifado & N:N & A mesma associação não pode ocorrer duas vezes.\\
Resumo Bem -- Bem Patrimonial & 1:N & Cada bem pertence a um resumo.\\
Local -- Bem Patrimonial & 1:N & Cada bem possui um local cadastrado.\\
Resumo Reagente -- Especificação & 1:N & Um resumo agrupa especificações; cada especificação pertence a um resumo.\\
Especificação -- Substância & N:N & A composição liga múltiplas substâncias a múltiplas especificações.\\
Especificação -- Lote & 1:N & Um lote pertence a uma única especificação.\\
Lote -- Frasco & 1:N opcional & Um frasco pode ter lote NULL quando o lote não for conhecido.\\
Almoxarifado -- Frasco & 1:N & Cada frasco está em exatamente um almoxarifado.\\
Frasco -- Empréstimo & 1:N & Histórico completo; no máximo um empréstimo ativo por frasco.\\
Turma -- Aluno & N:N & Alunos podem participar de várias turmas (via sub-coleção \texttt{Turma/\{id\}/Alunos}).\\
Turma -- Post & 1:N & Cada post pertence a uma turma e ao professor autor.\\
Post -- Comentário & 1:N & Cada comentário pertence a um post.\\
Professor -- Matéria & N:N & A associação não pode ser duplicada.\\
Roteiro -- Professor & N:N & O mesmo roteiro não pode ser compartilhado duas vezes com o mesmo professor.\\
\bottomrule
\end{longtable}

\subsubsection{Nullabilidade}
Todo atributo é \textbf{NOT NULL} por padrão. Um campo só pode ser NULL quando isso representa uma ausência legítima do domínio. As exceções são:
\begin{itemize}
    \item foto de perfil;
    \item documento de baixa antes da baixa patrimonial;
    \item campos de resposta de requisição antes da resposta;
    \item campos opcionais de especificação (fabricante, código do fabricante, FDS/FISPQ e demais dados não disponíveis);
    \item CAS e fórmula química quando não cadastrados;
    \item frequência de pesagem quando o reagente não requer pesagem frequente;
    \item lote e dados de validade/pesagem quando desconhecidos ou ainda não aplicáveis;
    \item validade\_efetiva antes da abertura quando sua definição depende do prazo após abertura;
    \item dados de devolução de empréstimo antes da devolução;
    \item id\_local da linha agregada ``Geral'' no resumo patrimonial diário;
    \item id\_turma da notificação de professor quando o tipo não é COMENTARIO;
    \item atributos de composição cuja ausência é necessária para representar faixa ou notação textual do fabricante.
\end{itemize}

\subsubsection{Chaves e constraints}
\begin{itemize}
    \item Usuario.email, Bem\_Patrimonial.numero\_patrimonio, Frasco\_Reagente.codigo\_frasco, Turma.codigo\_turma, Materia.codigo\_materia e CAS não nulo são únicos.
    \item As tabelas associativas usam chave primária composta por suas duas FKs.
    \item Composicao\_Reagente usa UNIQUE(id\_especificacao\_reagente, id\_substancia\_quimica).
    \item Lote usa UNIQUE(id\_especificacao\_reagente, nome\_fornecedor, numero\_lote), evitando a identificação de dois lotes equivalentes do mesmo fornecedor e produto.
    \item Convite\_Aluno permite o mesmo e-mail em turmas diferentes, mas impede mais de um convite PENDENTE para a mesma combinação de e-mail normalizado e turma; para convite global (\textit{id\_turma} NULL), impede mais de um convite PENDENTE para o mesmo e-mail normalizado. No PostgreSQL, isso é representado por dois índices únicos parciais: \texttt{UNIQUE(email\_normalizado, id\_turma) WHERE id\_turma IS NOT NULL AND status = 'pendente'} e \texttt{UNIQUE(email\_normalizado) WHERE id\_turma IS NULL AND status = 'pendente'}. No Firestore, a Cloud Function deve realizar a mesma validação dentro de transação.
    \item Requisicao\_Edicao\_Bem\_Patrimonial permite no máximo uma requisição PENDENTE por bem.
    \item Todos os valores quantitativos de massa, volume, capacidade, contagens, mês, dia de frequência e quantidades agregadas devem obedecer aos limites não negativos ou positivos definidos pelas regras do domínio.
    \item Datas que representam intervalos devem respeitar sua ordem temporal; exemplos: fabricação <= validade e devolução >= retirada.
\end{itemize}

\subsubsection{Regras químicas fechadas}
\begin{itemize}
    \item \textbf{PURA} e \textbf{MISTURA} são tipos de substância do resumo. MISTURA exige composição na especificação; PURA não depende de múltiplas linhas de composição.
    \item A concentração, quando aplicável, é atributo da composição da especificação, e não do resumo.
    \item A densidade é atributo da especificação. É obrigatória e positiva para LIQUIDO e pode ser NULL para SOLIDO.
    \item LIQUIDO usa ml como unidade operacional; SOLIDO usa g. O par estado/unidade deve ser sempre coerente.
    \item Duas especificações com o mesmo nome podem coexistir sob o mesmo resumo quando diferirem em concentração, densidade, fabricante ou código do produto e representarem o mesmo item agrupador. Diferença de especificação não implica automaticamente novo resumo.
    \item Gases estão fora do escopo da V1. Não se deve misturar sua futura regra de pressão com os campos de pesagem de sólidos e líquidos.
\end{itemize}

\subsubsection{Máquina de estados e invariantes do frasco}
\begin{longtable}{>{\raggedright\arraybackslash}p{0.34\textwidth} >{\raggedright\arraybackslash}p{0.58\textwidth}}
\toprule
\textbf{Estado/ação} & \textbf{Invariante}\\
\midrule
\endfirsthead
\toprule
\textbf{Estado/ação} & \textbf{Invariante} (continuação)\\
\midrule
\endhead
\bottomrule
\endfoot
FECHADO + DISPONIVEL & Pode ser retirado ou aberto.\\
ABERTO + DISPONIVEL & Pode ser retirado ou ajustado por pesagem.\\
EMPRESTADO & Deve possuir exatamente um empréstimo ativo EM\_USO ou ATRASADO.\\
VAZIO/QUEBRADO & Não pode estar EMPRESTADO e deve aparecer como pendente de descarte.\\
QUARENTENA & Exige detalhe de motivo e bloqueia nova retirada.\\
VENCIDO + uso\_vencido\_autorizado = false & Exibe PENDENTE DE DESCARTE quando não estiver em quarentena; bloqueia nova retirada.\\
VENCIDO + uso\_vencido\_autorizado = true & Pode permanecer DISPONIVEL para uso excepcional após confirmação explícita no empréstimo.\\
DESCARTADO & É terminal; não retorna a DISPONIVEL.\\
Devolução & Só pode ocorrer para empréstimo EM\_USO ou ATRASADO e aplica tolerância híbrida Q06: ganho tolerado produz consumo zero e evento AJUSTE; ganho acima da tolerância bloqueia.\\
\bottomrule
\end{longtable}

\subsubsection{Fechamento dos estados automáticos}
\textit{vencido} é definido pelo servidor a partir da validade efetiva. Na abertura do frasco, \textit{validade\_efetiva} é recalculada conforme a regra de validade pós-abertura. \textit{ATRASADO} é promovido pelo job servidor quando a data de devolução prevista é ultrapassada. Nenhum desses estados pode ser inventado pelo frontend. O vencimento também é recalculado sincronicamente em operações críticas de cadastro e devolução, para que o resultado não dependa de o job diário já ter executado.

\subsection{Checklist de Fechamento do Domínio e do Modelo}
Os onze pontos que precisavam ser resolvidos antes da modelagem física foram consolidados da seguinte forma:
\begin{enumerate}
    \item \textbf{Matriz completa de permissões}: formalizada na seção de Stakeholders e referenciada como fonte normativa de autorização.
    \item \textbf{Chefe Geral e multi-role}: Chefe Geral é exclusivo por conta; eventual dupla função institucional exige contas independentes.
    \item \textbf{Obrigatoriedade de campos}: formalizada por contexto, estado, tipo de substância e operação.
    \item \textbf{Densidade}: reposicionada semanticamente como propriedade de Especificacao\_Reagente e usada nos cálculos sem duplicação no resumo.
    \item \textbf{Máquina de estados}: estados do frasco e transições de retirada, devolução, quarentena, vencimento, esvaziamento, quebra e descarte foram explicitados.
    \item \textbf{NULL/NOT NULL}: todo campo passa a ter semântica explícita, com NULL restrito a ausências legítimas do domínio.
    \item \textbf{Cardinalidades}: relacionamentos 1:N, N:N e especializações de papel foram consolidados.
    \item \textbf{Constraints}: unicidades, chaves compostas, regras temporais e limites quantitativos foram formalizados, indicando quais precisarão ser simulados por transação no Firestore.
    \item \textbf{Métricas e unidades}: limiar de escassez passa a significar quantidade de frascos; massa de balança e consumo convertido permanecem conceitos distintos.
    \item \textbf{Gases}: explicitamente fora do escopo da V1, com a extensão futura isolada para não contaminar o modelo atual.
    \item \textbf{Vencimento e destino}: cadastro, abertura e devolução distinguem explicitamente quarentena, pendência de descarte e uso vencido autorizado.
\end{enumerate}

\begin{explicacao}[Critério de passagem para a modelagem Firestore]
O modelo 3FN é considerado fechado quando cada requisito funcional consegue ser rastreado até entidade, atributo, constraint e fluxo correspondente sem depender de uma suposição não documentada. A seção seguinte pode então introduzir apenas as denormalizações necessárias às consultas do Firestore, sem alterar a semântica do modelo relacional.
\end{explicacao}

\begin{regra}[DP-C02: moderação]
Autor vê o original marcado como moderado; demais alunos e bolsistas veem aviso institucional. Professor responsável e Chefe Geral veem o original e histórico necessário à moderação/auditoria. Moderação exige motivo e moderador; proteção ocorre no servidor e nas Rules (seção 11).
\end{regra}
\begin{regra}[DP-C03: segredo de convite]
Gerar token com CSPRNG de 32 bytes e armazenar somente SHA-256 do token aleatório. Comparar em tempo constante no backend, validar expiração de sete dias, identidade verificada e consumo único transacional. Reenvio substitui hash e expiração no mesmo documento e registra o último operador; link anterior deixa de funcionar. No Firestore, docId determinístico é derivado no servidor por HMAC-SHA-256 com segredo do servidor sobre serialização inequívoca de contexto turma/global e e-mail normalizado. Não expor hash simples de e-mail previsível. HMAC é pseudonimização; documento e e-mail continuam restritos. Nunca registrar token em logs, auditoria ou banco.
\end{regra}

\begin{regra}[Q06: massa, consumo e tolerância]
Massa consumida = $\max(0, peso_{saida}-peso_{retorno})$ em g. Medida utilizada é não negativa (g para sólido, mL via densidade para líquido). Medida usada no frasco acumula apenas gramas. Tolerância normal: $\max(1\,g,0{,}005\,peso_{saida})$; higroscópico: $\max(2\,g,0{,}02\,peso_{saida})$. A fórmula de tolerância higroscópica e de variação instrumental incide estritamente sobre o \textbf{Peso Bruto de Saída ($peso_{saida}$)}. Ganho dentro da tolerância preserva peso físico real, consumo zero e evento AJUSTE/ganho\_massa\_higroscopia. Retorno abaixo da tara ($peso_{retorno} < peso_{vazio}$): diferença até 5 g exige confirmação de esgotamento para VAZIO e ajuste auditado; diferença maior com produto restante bloqueia e exige recalibração formal de tara com justificativa. Não converter ajuste em consumo negativo.
\end{regra}
\begin{regra}[Q14 e Q04: rastreabilidade na retirada]
Autoatendimento é calculado pelo servidor: só o único gestor ativo do almoxarifado pode retirar para si, com justificativa obrigatória e notificação auditável à chefia; se houver outro gestor ativo, bloquear. Para vencido ou validade desconhecida liberado com termo nas finalidades PESQUISA\_TCC\_POS ou ESTUDO\_DEGRADACAO\_RESIDUOS, exigir justificativa metodológica de 20--2000 caracteres, versão do TCR, instante servidor, UID autenticado do retirante que aceitou e referência ao Registro\_de\_Auditoria correspondente. O backend valida aceite vinculado ao frasco, finalidade e operação, nunca declaração do gestor em nome de terceiro. Sem essas condições, rejeitar. TCR registra ciência e responsabilidade, sem substituir regras institucionais de segurança química. Quarentena e pendência de descarte continuam bloqueadas.
\end{regra}
\begin{regra}[Reclassificação patrimonial]
Novo nome individual exige resolver/criar resumo alvo, preencher novo\_id\_resumo\_bem\_patrimonial e alterar só o vínculo do bem solicitado; não renomear resumo compartilhado. Versao\_bem\_origem captura a versão observada e impede aprovação sobre edição concorrente.
\end{regra}

\section{Notas de Mapeamento para Firestore}
% =================================================

\begin{explicacao}[Por que esta seção existe]
A modelagem da seção anterior é deliberadamente relacional e normalizada (3FN), pensada para ser lida e avaliada como um modelo de dados independente de tecnologia. O banco real do LCQUI, porém, é o \textbf{Firestore} (Firebase), um banco de documentos sem \texttt{JOIN} e sem \texttt{LIKE}/busca por substring. Esta seção documenta as denormalizações propositais necessárias para a implementação -- nenhuma delas deve ser lida como parte do modelo 3FN, e nenhuma delas deve alterar a seção 4.
\end{explicacao}

\subsection{Reintrodução de \textit{letra\_inicial}}
O campo \textit{letra\_inicial}, removido de \textit{Resumo\_Reagente} por ser derivável de \textit{nome} (violaria 3FN), volta a existir como campo próprio \textbf{apenas} no documento Firestore da coleção \textit{Resumo\_Reagente}. Motivo: o Firestore não tem \texttt{LIKE}/busca por substring, então a busca alfabética do catálogo precisa de um campo indexável para filtro de igualdade (ou de \textit{range query} por prefixo). Como esse filtro normalmente é combinado com outros filtros de igualdade na mesma consulta (estado físico, natureza química), um campo de igualdade simples é mais previsível dentro das limitações de índice composto do Firestore do que uma \textit{range query} de prefixo sobre \textit{nome} combinada com filtros adicionais.

\begin{regra}[Consistência obrigatória]
\textit{letra\_inicial} deve ser recalculado e reescrito pela mesma \textit{Cloud Function} que grava \textit{nome}, nunca aceito diretamente do cliente -- ver princípio ``nunca confiar no \textit{client-side}'' na seção \textit{Tecnologia e Relatórios}.
\end{regra}

\subsection{Estado físico canônico do catálogo}
DP-A01: Resumo\_Reagente.estado\_fisico é escalar SOLIDO/LIQUIDO; eh\_higroscopico também pertence ao resumo. Filtro usa igualdade: \texttt{.where("estado\_fisico", "==", "SOLIDO")}. A unidade operacional é derivada, g para sólido e mL para líquido.

\subsection{Espelhamento de \textit{nome\_equipamento} e Localização em \textit{Bem\_Patrimonial}}
No modelo 3FN, o nome do tipo de bem pertence exclusivamente a \textit{Resumo\_Bem\_Patrimonial.nome}; \textit{Bem\_Patrimonial} mantém apenas \textit{id\_resumo\_bem\_patrimonial} e \textit{id\_local}. No Firestore, porém, consultas de listagem e relatórios frequentemente partem diretamente de \textit{Bem\_Patrimonial}, sem um \textit{JOIN} para resolver o resumo ou o prédio. Por isso, o documento Firestore de \textit{Bem\_Patrimonial} deve conter os campos denormalizados \textit{nome\_equipamento}, \textit{predio}, \textit{andar} e \textit{sala}.

\begin{regra}[Fonte única de verdade para o nome do patrimônio]
\textit{nome\_equipamento} é exclusivamente um campo de leitura derivado de \textit{Resumo\_Bem\_Patrimonial.nome}. O cliente nunca pode gravá-lo diretamente. Sempre que o nome do resumo for alterado, a mesma operação transacional ou um gatilho de propagação controlado pelo servidor deve atualizar os documentos \textit{Bem\_Patrimonial} associados. Na aprovação da requisição de edição descrita na seção 10.2.5, \textit{novo\_nome} representa uma proposta de alteração do nome do resumo; o servidor resolve/cria outro resumo e altera a referência somente do bem solicitado; renomear globalmente um resumo é ação distinta do gestor.
\end{regra}

\subsection{Composição embutida na especificação}
O mapeamento físico adotado é \texttt{composicao}, array de mapas na especificação em \texttt{Resumo\_Reagente/id/Especificacoes/id}. Cada mapa referencia uma substância; o servidor impede IDs repetidos e limita o tamanho do array. Não manter simultaneamente uma coleção raiz e uma subcoleção alternativa como fontes concorrentes. A relação N:N continua normalizada na seção 4.

\subsection{Tabelas ``Materialized'' são coleções próprias, não \textit{views} do SGBD}
As tabelas descritas na seção \textit{Materialized (Views)} não são \textit{materialized views} no sentido de PostgreSQL (não existe \texttt{REFRESH MATERIALIZED VIEW} no Firestore). Cada uma delas é implementada como sua própria coleção Firestore, escrita por \textit{Cloud Functions} -- gatilhos de escrita (\texttt{onDocumentWritten}) para contadores que reagem a um evento pontual (ex.: \textit{Lote\_Materializado}), ou funções agendadas (\texttt{onSchedule}) para os resumos diários/mensais. Ver pseudo-código na seção \textit{Tecnologia e Relatórios}.

\subsection{Unicidade de \texttt{numero\_patrimonio\_proposto} no Firestore}
O PostgreSQL implementa a regra RF10b via índice único parcial
\texttt{UNIQUE(numero\_patrimonio\_proposto) WHERE status = 'pendente'},
impedindo duas requisições simultâneas para a mesma plaqueta enquanto
o bem patrimonial ainda não existe. No Firestore não há índice parcial
nem \emph{unique constraint} nativo, então a garantia é reimplementada
na Cloud Function de criação da requisição de adição.

A criação lê o documento determinístico em \texttt{Locks\_Requisicao\_Patrimonio}, verifica a plaqueta já cadastrada e grava lock e requisição na mesma transação. Aprovação e rejeição removem o lock na mesma transação da resposta. Uma consulta prévia pode melhorar a mensagem de erro, mas não substitui a chave de unicidade. Locks vinculados a requisições pendentes não expiram por idade.

\subsection{Mecanismo de Sequenciamento Atômico via Documento Singleton (\textit{Contador\_Codigo\_Frasco})}
A coleção \texttt{Contador\_Codigo\_Frasco} não existe no modelo 3FN porque bancos relacionais possuem primitivas de \texttt{SEQUENCE/SERIAL}. Para gerar o identificador atômico \texttt{LCQUI-X}, utiliza-se o documento singleton \texttt{Contador\_Codigo\_Frasco/singleton} com o campo \texttt{ultimo\_codigo\_gerado}, operado atomicamente via \texttt{runTransaction} no backend. A transação garante a unicidade; contenção e custo devem ser medidos incluindo histórico, auditoria, índices e tentativas repetidas. Não se fixa aqui uma cota universal de escritas por segundo.

\subsection{Índices Compostos Obrigatórios}
Para garantir a integridade das \textit{queries} e ordenações cruciais, os seguintes índices compostos devem ser criados explicitamente no console do Firestore:
\begin{enumerate}
    \item \texttt{Historico\_Frasco\_Reagente}: \texttt{id\_frasco\_reagente (ASC)} + \texttt{timestamp (DESC)}.
    \item \texttt{Emprestimo\_Reagente}: \texttt{id\_frasco\_reagente (ASC)} + \texttt{data\_retirada (DESC)}.
    \item \textbf{M-23} \texttt{Bem\_Patrimonial}: \texttt{predio (ASC)} + \texttt{andar (ASC)} + \texttt{sala (ASC)} -- suporta o relatório personalizado por localização (Seção~9 fluxo ALM-01).
    \item \textbf{M-23} \texttt{Bem\_Patrimonial}: \texttt{id\_resumo\_bem\_patrimonial (ASC)} + \texttt{status (ASC)} -- suporta a listagem de bens por tipo e status.
    \item \textbf{M-23 / N-08} Collection-group \texttt{Historico} (sub-coleção de \texttt{Bem\_Patrimonial}): \texttt{predio (ASC)} + \texttt{timestamp (DESC)} -- suporta relatório histórico por prédio/período.
    \item \textbf{M-23} Collection-group \texttt{Historico}: \texttt{id\_gestor (ASC)} + \texttt{timestamp (DESC)} -- suporta relatório de atividade por operador.
    \item \textbf{M-23} \texttt{Emprestimo\_Reagente}: \texttt{id\_almoxarifado (ASC)} + \texttt{data\_retirada (DESC)} -- suporta relatório mensal por almoxarifado.
    \item \texttt{Frasco\_Reagente}: \texttt{id\_almoxarifado (ASC)} + \texttt{id\_resumo\_reagente (ASC)} + \texttt{id\_especificacao\_reagente (ASC)} + \texttt{disponibilidade (ASC)} + \texttt{estado\_fisico\_frasco (ASC)} + \texttt{em\_quarentena (ASC)} + \texttt{vencido (ASC)} -- suporta job agendado de escassez.
\end{enumerate}
\subsection{Mapeamento Completo de Coleções e Sub-coleções Firestore}
\label{subsec:mapeamento-colecoes}

\begin{explicacao}[Como ler esta tabela]
Cada linha descreve uma entidade 3FN da Seção 4, a decisão de estrutura no Firestore (coleção raiz, sub-coleção ou array de maps embutido) e a justificativa. Campos de consulta necessários podem ser projetados no documento Firestore conforme o dicionário; o símbolo \textbf{Denorm:} lista esses campos. A consistência usa transação para invariantes ou propagação versionada para apresentação; snapshots históricos não são atualizados com o cadastro atual.
\end{explicacao}

\small

\small
\setlength{\tabcolsep}{4.5pt}
\small
\setlength{\tabcolsep}{4pt}
\begin{longtable}{>{\raggedright\arraybackslash}p{0.28\textwidth} >{\raggedright\arraybackslash}p{0.26\textwidth} >{\raggedright\arraybackslash}p{0.38\textwidth}}
\toprule
\textbf{Entidade 3FN} & \textbf{Estrutura Firestore} & \textbf{Justificativa / Denorm}\\
\midrule
\endfirsthead
\toprule
\textbf{Entidade 3FN} & \textbf{Estrutura Firestore} & \textbf{Justificativa / Denorm} (continuação)\\
\midrule
\endhead
\bottomrule
\endfoot
Usuario & Coleção raiz \texttt{Usuarios} (docId = Firebase Auth UID) & Acesso direto por auth; estrutura natural do SDK.\\
Chefe\_Geral & Coleção raiz \texttt{Chefe\_Geral} (docId = uid) & Lido por \texttt{atualizarCustomClaims(uid)} para compor o claim \texttt{roles} do token JWT. Claims são projeção; conta ativa e versão de permissões devem ser verificadas conforme seção 11.\\
Gestor\_Almoxarifado & Coleção raiz \texttt{Gestor\_Almoxarifado} (docId = uid) & Idem acima.\\
Gestor\_Bens\_Patrimoniais & Coleção raiz \texttt{Gestor\_Bens\_Patrimoniais} (docId = uid) & Idem.\\
Professor & Coleção raiz \texttt{Professor} (docId = uid) & Idem.\\
Aluno & Coleção raiz \texttt{Aluno} (docId = uid). Denorm: \texttt{letra\_inicial} & Idem. \texttt{letra\_inicial} permite filtro de igualdade na busca.\\
Bolsista & Coleção raiz \texttt{Bolsista} (docId = uid) & Idem.\\
Gestor\_Almoxarifado\_x\_Almoxarifado & Coleção raiz & Query bidirecional: ``gestores de X'' e ``almoxarifados de Y''.\\
Almoxarifado & Coleção raiz & Poucas dezenas de documentos.\\
Local & Coleção raiz & Referenciado por múltiplas entidades; campos \texttt{predio}, \texttt{andar}, \texttt{sala} denormalizados em \texttt{Bem\_Patrimonial} via trigger.\\
Resumo\_Bem\_Patrimonial & Coleção raiz. Denorm: \texttt{letra\_inicial\_nome} & Agrupador de bens; filtro alfabético por letra inicial.\\
Bem\_Patrimonial & Coleção raiz. Denorm: \texttt{nome\_equipamento}, \texttt{predio}, \texttt{andar}, \texttt{sala}, \texttt{letra\_inicial\_nome} & Entidade central pesquisada diretamente; evita JOINs em listagens e relatórios.\\
Historico\_Bem\_Patrimonial & Sub-coleção \nolinkurl{Bem\_Patrimonial/\{id\}/Historico}. Denorm (N-08): \texttt{predio}, \texttt{andar}, \texttt{sala} -- snapshot imutável do local no momento do evento & Sempre consultado a partir de um bem conhecido. O snapshot de localização é gravado no momento da criação do evento histórico e não é atualizado quando o bem muda de sala; isso garante rastreabilidade do bem antes e depois da mudança. Requerido pelo relatório personalizado por prédio/período.\\
Alteracao\_Bem\_Patrimonial & Campo \texttt{alteracoes} (array de maps) dentro de cada doc de \texttt{Historico} & Poucos campos por alteração; relação 1:N estreita, nunca consultada independentemente.\\
Requisicao\_Edicao\_Bem\_Patrimonial & Coleção raiz. Denorm: \texttt{nome\_equipamento}, \texttt{numero\_patrimonio} & Listada independentemente nos painéis de gestores por qualquer filtro de status.\\
Requisicao\_Adicao\_Bem\_Patrimonial & Coleção raiz & Mesma lógica; query por status e número de patrimônio proposto.\\
Impressao\_Etiqueta\_Frasco & Coleção raiz & Refere-se a impressão em lote (virgens). Reimpressões avulsas de frascos já cadastrados usam o Registro de Auditoria.\\
Contador\_Codigo\_Frasco & Documento singleton em \nolinkurl{Contador\_Codigo\_Frasco/singleton} & Sequenciador atômico operado via \texttt{runTransaction} para gerar o \texttt{LCQUI-X}.\\
Substancia\_Quimica & Coleção raiz & Catálogo global; raramente alterado.\\
Resumo\_Reagente & Coleção raiz. Denorm: \texttt{letra\_inicial} & Catálogo de reagentes; filtros de igualdade combinados exigem ambos os campos como propriedades do documento.\\
Especificacao\_Reagente & Sub-coleção \nolinkurl{Resumo\_Reagente/\{id\}/Especificacoes} & Sempre acessada a partir do resumo; nunca consultada globalmente.\\
Composicao\_Reagente & Array de maps embutido no documento de \texttt{Especificacao\_Reagente} & N pequeno e fixo (poucas substâncias por especificação); não cresce sem limite.\\
Lote & Coleção raiz. Denorm: \texttt{id\_especificacao\_reagente}, \texttt{nome\_reagente} & Relatórios consultam todos os lotes de um período sem partir de uma especificação conhecida.\\
Frasco\_Reagente & Coleção raiz. Denorm: \texttt{id\_especificacao\_reagente}, \texttt{id\_resumo\_reagente}, \texttt{nome\_reagente}, \texttt{unidade\_medida}, \texttt{id\_almoxarifado} & Entidade central; queries cruzadas constantes por almoxarifado, status e vencimento.\\
Historico\_Frasco\_Reagente & Coleção raiz (não sub-coleção). Denorm: \texttt{id\_almoxarifado} & Jobs diários consultam ``todos os históricos do dia'' sem partir de um frasco conhecido; coleção raiz simplifica esse padrão; collection-group também permitiria históricos em subcoleções.\\
Emprestimo\_Reagente & Coleção raiz. Denorm: \texttt{id\_almoxarifado}, \texttt{nome\_reagente}, \texttt{nome\_usuario\_retirou} & Relatórios por almoxarifado e por período; queries de status global (\texttt{EM\_USO}).\\
Materia & Coleção raiz & Catálogo global; raramente alterado.\\
Professor\_x\_Materia & Coleção raiz & Query bidirecional: ``matérias de X'' e ``professores de Y''.\\
Turma & Coleção raiz. Denorm: \texttt{nome\_materia} & Acesso direto por \texttt{codigo\_turma}; nome da matéria necessário na listagem sem JOIN.\\
\textbf{\texttt{Aluno\_x\_Turma}} & \textbf{Subcoleção dupla (espelhada):}\newline 1. \nolinkurl{Turma/\{id\}/Alunos}\newline 2. \nolinkurl{Usuarios/\{uid\}/Turmas} & \textbf{1. \nolinkurl{Turma/\{id\}/Alunos}:} Visão da disciplina; contagem de capacidade e listagem de chamada para o professor.\newline \textbf{2. \nolinkurl{Usuarios/\{uid\}/Turmas}:} Visão do aluno; viabiliza o feed em tempo real no menu lateral de ``Minhas Turmas'' com segurança isolada por UID e zero leituras $N+1$. Consistência garantida via transação nas Cloud Functions de matrícula/exclusão.\\
Historico\_Alunos\_Turma & Sub-coleção \nolinkurl{Turma/\{id\}/HistoricoAlunos} & Auditoria de inclusão/exclusão sempre por turma.\\
Post & Sub-coleção \nolinkurl{Turma/\{id\}/Posts} & Feed cronológico da turma; nunca consultado globalmente.\\
Comentario & Sub-coleção \nolinkurl{Turma/\{id\}/Posts/\{postId\}/Comentarios} & Thread aninhada ao post.\\
Historico\_Posts\_Turma & Sub-coleção \nolinkurl{Turma/\{id\}/Posts/\{postId\}/Historico} & Auditoria de edição sempre a partir do post.\\
Historico\_Comentario & Sub-coleção \nolinkurl{Turma/\{id\}/Posts/\{postId\}/Comentarios/\{id\}/Historico} & Auditoria de edição sempre a partir do comentário.\\
Roteiro\_Experimento & Coleção raiz & Compartilhamento cross-professor exige acesso sem partir de uma turma.\\
Convite\_Aluno & Coleção raiz & Query por e-mail normalizado e status independentemente de turma.\\
Notificacao & Sub-coleção \nolinkurl{Usuarios/\{uid\}/Notificacoes} & Query natural por usuário; funciona automaticamente para usuários multi-role sem necessidade de multiple reads.\\
Registro\_de\_Auditoria & Coleção raiz & Log global queryável por tipo de entidade e período.\\
Lote\_Materializado & Coleção raiz (docId = id do lote) & Atualizada por trigger \texttt{onFrascoCriado}/\texttt{onFrascoRemovido}.\\
Resumo\_Almoxarifado\_Diario & Coleção raiz & Job agendado diário.\\
Resumo\_Reagente\_Diario & Coleção raiz & Job agendado diário.\\
Resumo\_Bem\_Patrimonial\_Diario & Coleção raiz & Job agendado diário.\\
Atividade\_Gestor\_Almoxarifado\_Mensal & Coleção raiz & Job agendado diário (acumulado por mês).\\
Atividade\_Gestor\_Bens\_Patrimoniais\_Mensal & Coleção raiz & Job agendado diário (acumulado por mês).\\
\end{longtable}


\begin{regra}[Regra geral de denormalização]
Copiar apenas atributos necessários às consultas documentadas. Cada projeção identifica sua origem e mecanismo de atualização; snapshots de eventos são imutáveis. Vínculos e invariantes críticos são transacionais, enquanto projeções de apresentação podem ser assíncronas com reconciliação.
\end{regra}


\normalsize
\begingroup
\sloppy
\subsubsection{Dicionário de dados completo e convenções físicas}
\label{sec:dicionario}
Este dicionário especifica o alvo documental. Divergências com os caminhos do código atual são listadas na auditoria; não se pressupõe migração executada. IDs relacionais SERIAL tornam-se docIds string; FKs tornam-se IDs string, sem JOIN implícito. Datas civis usam string YYYY-MM-DD; instantes usam Timestamp UTC, exibidos em America/Sao\_Paulo. Valores quantitativos usam number finito, com precisão/escala da seção 4; rejeitar NaN, infinito, sinal e casas decimais incompatíveis. Todos os limites VARCHAR continuam válidos. O = obrigatório; N = null permitido; C = condicional. Campos calculados, autoria, timestamps, IDs e estados são gravados pelo servidor. Arrays de composição e alterações são limitados pelo esquema, nunca históricos crescentes. O docId substitui o campo id SERIAL; em associações a chave é determinística pelas FKs.

\paragraph{Usuario}\mbox{}\par
Caminho: \texttt{Usuarios/\{uid\}}.
\begin{description}
\item[\texttt{id}] docId string; O. Chave do documento, gerada pelo servidor; não duplicar como contador relacional.
\item[\texttt{nome}] string; O. Nome de apresentação; não substitui a chave do documento. Máximo 150 caracteres.
\item[\texttt{email}] string; O. Endereço normalizado para identidade/convite; não divulgar em catálogo público interno. Máximo 150 caracteres. Unicidade garantida por chave determinística/registro de unicidade no servidor.
\item[\texttt{profile\_url\_photo}] string; N. Foto opcional de apresentação do usuário; ausência usa avatar padrão, com acesso validado ao objeto.
\end{description}

\paragraph{Chefe\_Geral}\mbox{}\par
Caminho: \texttt{Chefe\_Geral/\{id\}}.
\begin{description}
\item[\texttt{id\_usuario}] string; O. Identificador da entidade usuario, cuja existência e escopo são verificados no servidor.
\end{description}

\paragraph{Gestor\_Almoxarifado}\mbox{}\par
Caminho: \texttt{Gestor\_Almoxarifado/\{id\}}.
\begin{description}
\item[\texttt{id\_usuario}] string; O. Identificador da entidade usuario, cuja existência e escopo são verificados no servidor.
\end{description}

\paragraph{Gestor\_Almoxarifado\_x\_Almoxarifado}\mbox{}\par
Caminho: \texttt{Gestor\_Almoxarifado\_x\_Almoxarifado/\{id\}}.
\begin{description}
\item[\texttt{id\_gestor\_almoxarifado}] string; O. Identificador da entidade gestor almoxarifado, cuja existência e escopo são verificados no servidor.
\item[\texttt{id\_almoxarifado}] string; O. Identificador da entidade almoxarifado, cuja existência e escopo são verificados no servidor.
\end{description}

\paragraph{Gestor\_Bens\_Patrimoniais}\mbox{}\par
Caminho: \texttt{Gestor\_Bens\_Patrimoniais/\{id\}}.
\begin{description}
\item[\texttt{id\_usuario}] string; O. Identificador da entidade usuario, cuja existência e escopo são verificados no servidor.
\end{description}

\paragraph{Bem\_Patrimonial}\mbox{}\par
Caminho: \texttt{Bem\_Patrimonial/\{id\}}.
\begin{description}
\item[\texttt{id}] docId string; O. Chave do documento, gerada pelo servidor; não duplicar como contador relacional.
\item[\texttt{id\_resumo\_bem\_patrimonial}] string; O. Identificador da entidade resumo bem patrimonial, cuja existência e escopo são verificados no servidor.
\item[\texttt{numero\_patrimonio}] string; O. Plaqueta institucional em texto, preservando zeros iniciais; identifica unicamente a unidade física. Máximo 30 caracteres.
\item[\texttt{estado\_conservacao}] string enum; O. Valor de estado conservacao, validado pelo formulário e pelo servidor conforme regras do domínio. Valores: BOM, REGULAR, RUIM.
\item[\texttt{id\_local}] string; O. Identificador da entidade local, cuja existência e escopo são verificados no servidor.
\item[\texttt{photo\_url}] string; O. Referência da foto validada; não confiar em URL arbitrária.
\item[\texttt{documento\_dado\_baixa\_pdf\_url}] string; N. Referência do comprovante PDF, nula antes da baixa e obrigatória em Ja\_dado\_baixa; preservar após vínculo.
\item[\texttt{nome\_responsavel\_sei}] string; O. Responsável institucional em texto; não é FK de usuário. Máximo 150 caracteres.
\item[\texttt{status}] string enum; O. Estado controlado pelo domínio; transições são validadas no servidor. Valores: Ativo, Inservivel, Ja\_dado\_baixa.
\item[\texttt{descricao\_complementar}] string; N. Observação opcional da unidade física, sem alterar a descrição compartilhada do modelo. Máximo 200 caracteres.
\end{description}

\paragraph{Resumo\_Bem\_Patrimonial}\mbox{}\par
Caminho: \texttt{Resumo\_Bem\_Patrimonial/\{id\}}.
\begin{description}
\item[\texttt{id}] docId string; O. Chave do documento, gerada pelo servidor; não duplicar como contador relacional.
\item[\texttt{nome}] string; O. Nome de apresentação; não substitui a chave do documento. Máximo 150 caracteres.
\item[\texttt{descricao}] string; O. Valor de descricao, validado pelo formulário e pelo servidor conforme regras do domínio.
\end{description}

\paragraph{Historico\_Bem\_Patrimonial}\mbox{}\par
Caminho: \texttt{Bem\_Patrimonial/\{id\}/Historico/\{eventoId\}}.
\begin{description}
\item[\texttt{id}] docId string; O. Chave do documento, gerada pelo servidor; não duplicar como contador relacional.
\item[\texttt{id\_bem\_patrimonial}] string; O. Identificador da entidade bem patrimonial, cuja existência e escopo são verificados no servidor.
\item[\texttt{timestamp}] Timestamp; O. Data/instante de timestamp, com ordem temporal validada e autoria do relógio do servidor para eventos.
\item[\texttt{tipo}] string enum; O. Valor de tipo, validado pelo formulário e pelo servidor conforme regras do domínio. Valores: cadastro, edicao, baixa.
\item[\texttt{id\_usuario}] string; O. Identificador da entidade usuario, cuja existência e escopo são verificados no servidor.
\end{description}

\paragraph{Alteracao\_Bem\_Patrimonial}\mbox{}\par
Caminho: \texttt{Bem\_Patrimonial/\{id\}/Historico/\{eventoId\}: alteracoes[]}.
\begin{description}
\item[\texttt{campo}] string; O. Valor de campo, validado pelo formulário e pelo servidor conforme regras do domínio. Máximo 60 caracteres.
\item[\texttt{valor\_anterior}] string; O. Valor de valor anterior, validado pelo formulário e pelo servidor conforme regras do domínio.
\item[\texttt{valor\_novo}] string; O. Valor de valor novo, validado pelo formulário e pelo servidor conforme regras do domínio.
\end{description}

\paragraph{Requisicao\_Edicao\_Bem\_Patrimonial}\mbox{}\par
Caminho: \texttt{Requisicao\_Edicao\_Bem\_Patrimonial/\{id\}}.
\begin{description}
\item[\texttt{id}] docId string; O. Chave do documento, gerada pelo servidor; não duplicar como contador relacional.
\item[\texttt{id\_bem\_patrimonial}] string; O. Identificador da entidade bem patrimonial, cuja existência e escopo são verificados no servidor.
\item[\texttt{status}] string enum; O. Estado controlado pelo domínio; transições são validadas no servidor. Valores: pendente, aprovada, rejeitada.
\item[\texttt{feita\_em}] Timestamp; O. Data/instante de feita em, com ordem temporal validada e autoria do relógio do servidor para eventos.
\item[\texttt{respondida\_em}] Timestamp; N. Data/instante de respondida em, com ordem temporal validada e autoria do relógio do servidor para eventos.
\item[\texttt{id\_usuario\_solicitante}] string; O. Identificador da entidade usuario solicitante, cuja existência e escopo são verificados no servidor.
\item[\texttt{id\_usuario\_respondente}] string; N. Identificador da entidade usuario respondente, cuja existência e escopo são verificados no servidor.
\item[\texttt{novo\_nome}] string; N. Valor novo nome, restrito ao campo correspondente da entidade de origem; preservar ausência legítima. Máximo 150 caracteres.
\item[\texttt{novo\_status}] string enum; N. Valor novo status, restrito ao campo correspondente da entidade de origem; preservar ausência legítima. Valores: Ativo, Inservivel, Ja\_dado\_baixa.
\item[\texttt{novo\_estado\_conservacao}] string enum; N. Valor novo estado conservacao, restrito ao campo correspondente da entidade de origem; preservar ausência legítima. Valores: BOM, REGULAR, RUIM.
\item[\texttt{novo\_id\_local}] string; N. Valor novo id local, restrito ao campo correspondente da entidade de origem; preservar ausência legítima.
\item[\texttt{nova\_photo\_url}] string; N. Valor nova photo url, restrito ao campo correspondente da entidade de origem; preservar ausência legítima.
\item[\texttt{motivo}] string; O. Valor de motivo, validado pelo formulário e pelo servidor conforme regras do domínio.
\item[\texttt{justificativa\_resposta}] string; N. Justificativa obrigatória ao aprovar ou rejeitar; ausente enquanto pendente.
\end{description}

\paragraph{Requisicao\_Adicao\_Bem\_Patrimonial}\mbox{}\par
Caminho: \texttt{Requisicao\_Adicao\_Bem\_Patrimonial/\{id\}}.
\begin{description}
\item[\texttt{id}] docId string; O. Chave do documento, gerada pelo servidor; não duplicar como contador relacional.
\item[\texttt{numero\_patrimonio\_proposto}] string; O. Plaqueta solicitada para o novo bem; unicidade entre pendências e bens existentes. Máximo 30 caracteres.
\item[\texttt{status}] string enum; O. Estado controlado pelo domínio; transições são validadas no servidor. Valores: pendente, aprovada, rejeitada.
\item[\texttt{feita\_em}] Timestamp; O. Data/instante de feita em, com ordem temporal validada e autoria do relógio do servidor para eventos.
\item[\texttt{id\_bem\_patrimonial\_se\_aprovado}] string; N. Identificador da entidade bem patrimonial se aprovado, cuja existência e escopo são verificados no servidor.
\item[\texttt{respondida\_em}] Timestamp; N. Data/instante de respondida em, com ordem temporal validada e autoria do relógio do servidor para eventos.
\item[\texttt{id\_usuario\_solicitante}] string; O. Identificador da entidade usuario solicitante, cuja existência e escopo são verificados no servidor.
\item[\texttt{id\_usuario\_respondente}] string; N. Identificador da entidade usuario respondente, cuja existência e escopo são verificados no servidor.
\item[\texttt{estado\_conservacao\_proposto}] string enum; O. Valor de estado conservacao proposto, validado pelo formulário e pelo servidor conforme regras do domínio. Valores: BOM, REGULAR, RUIM.
\item[\texttt{photo\_url\_proposta}] string; O. Foto obrigatória na submissão, validada pelo backend. Valor de photo url proposta, validado pelo formulário e pelo servidor conforme regras do domínio.
\item[\texttt{nome\_responsavel\_proposto}] string; O. Valor de nome responsavel proposto, validado pelo formulário e pelo servidor conforme regras do domínio. Máximo 150 caracteres.
\item[\texttt{id\_local}] string; O. Identificador da entidade local, cuja existência e escopo são verificados no servidor.
\item[\texttt{id\_resumo\_bem\_patrimonial}] string; N. Identificador da entidade resumo bem patrimonial, cuja existência e escopo são verificados no servidor.
\item[\texttt{nome\_resumo\_proposto}] string; N. Valor de nome resumo proposto, validado pelo formulário e pelo servidor conforme regras do domínio. Máximo 150 caracteres.
\item[\texttt{descricao\_resumo\_proposta}] string; N. Valor de descricao resumo proposta, validado pelo formulário e pelo servidor conforme regras do domínio.
\item[\texttt{motivo}] string; O. Valor de motivo, validado pelo formulário e pelo servidor conforme regras do domínio.
\end{description}

\paragraph{Locks\_Requisicao\_Patrimonio}\mbox{}\par
Caminho: \texttt{Locks\_Requisicao\_Patrimonio/\{id\}}.
\begin{description}
\item[\texttt{id (docId determinístico)}] string; O. Identificador da entidade id (docId determinístico), cuja existência e escopo são verificados no servidor. Máximo 100 caracteres.
\item[\texttt{criado\_em}] Timestamp; O. Data/instante de criado em, com ordem temporal validada e autoria do relógio do servidor para eventos.
\end{description}

\paragraph{Impressao\_Etiqueta\_Frasco}\mbox{}\par
Caminho: \texttt{Impressao\_Etiqueta\_Frasco/\{id\}}.
\begin{description}
\item[\texttt{id}] docId string; O. Chave do documento, gerada pelo servidor; não duplicar como contador relacional.
\item[\texttt{gerado\_em}] Timestamp; O. Data/instante de gerado em, com ordem temporal validada e autoria do relógio do servidor para eventos.
\item[\texttt{gerado\_por}] string; O. UID do operador que solicitou a geração de etiquetas, obtido da sessão autenticada.
\item[\texttt{codigo\_inicial}] number; O. Primeiro número inteiro positivo do intervalo inclusivo de etiquetas virgens.
\item[\texttt{codigo\_final}] number; O. Último número do intervalo inclusivo, maior ou igual ao inicial e limitado a 50 etiquetas.
\end{description}

\paragraph{Almoxarifado}\mbox{}\par
Caminho: \texttt{Almoxarifado/\{id\}}.
\begin{description}
\item[\texttt{id}] docId string; O. Chave do documento, gerada pelo servidor; não duplicar como contador relacional.
\item[\texttt{id\_local}] string; O. Identificador da entidade local, cuja existência e escopo são verificados no servidor.
\item[\texttt{nome}] string; O. Nome de apresentação; não substitui a chave do documento. Máximo 100 caracteres.
\item[\texttt{descricao}] string; O. Valor de descricao, validado pelo formulário e pelo servidor conforme regras do domínio. Máximo 500 caracteres.
\end{description}

\paragraph{Estoques\_Configurados}\mbox{}\par
Caminho: \texttt{Almoxarifado/\{almoxId\}/Estoques\_Configurados/\{configId\}}.
\begin{description}
\item[\texttt{configId}] docId string; O. Chave do documento, gerada de forma determinística: \texttt{\{id\_resumo\_reagente\}\_\{id\_especificacao\_reagente\}}, garantindo unicidade do par.
\item[\texttt{id\_especificacao\_reagente}] string; O. Identificador da especificação do reagente configurada.
\item[\texttt{id\_resumo\_reagente}] string; O. Identificador do resumo do reagente, mantido no documento para consultas otimizadas e manter a identidade composta.
\item[\texttt{id\_almoxarifado}] string; O. Identificador do almoxarifado (redundante à URL, porém canônico).
\item[\texttt{qtd\_limiar\_escassez}] number; O. Limiar inteiro não negativo de quantidade de frascos para este almoxarifado.
\item[\texttt{ativo}] boolean; O. Indica se o par continua configurado/intencionalmente estocado.
\item[\texttt{notificacao\_ativa}] boolean; O. Indica se a emissão automática do alerta está habilitada (permitindo silenciar sem apagar a configuração).
\end{description}

\paragraph{Substancia\_Quimica}\mbox{}\par
Caminho: \texttt{Substancia\_Quimica/\{id\}}.
\begin{description}
\item[\texttt{id}] docId string; O. Chave do documento, gerada pelo servidor; não duplicar como contador relacional.
\item[\texttt{nome}] string; O. Nome de apresentação; não substitui a chave do documento. Máximo 150 caracteres.
\item[\texttt{cas\_number}] string; N. Identificador CAS opcional da substância; quando informado, validar formato, dígito e unicidade. Máximo 15 caracteres.
\item[\texttt{formula\_quimica}] string; N. Fórmula declarada da substância; apresentação textual, sem executar marcação arbitrária. Máximo 50 caracteres.
\item[\texttt{ativo}] boolean; O. Habilitação lógica; desativação preserva identidade e histórico.
\end{description}

\paragraph{Resumo\_Reagente}\mbox{}\par
Caminho: \texttt{Resumo\_Reagente/\{id\}}.
\begin{description}
\item[\texttt{id}] docId string; O. Chave do documento, gerada pelo servidor; não duplicar como contador relacional.
\item[\texttt{nome}] string; O. Nome de apresentação; não substitui a chave do documento. Máximo 150 caracteres.
\item[\texttt{tipo\_substancia}] string enum; O. Valor de tipo substancia, validado pelo formulário e pelo servidor conforme regras do domínio. Valores: PURA, MISTURA.
\item[\texttt{natureza\_quimica}] string enum; O. Classificação curada do item do catálogo, independente de fabricante ou frasco. Valores: ORGANICO, INORGANICO, ELEMENTO, HIBRIDO.
\item[\texttt{requer\_pesagem\_frequente}] boolean; O. Indica necessidade de alertas periódicos de pesagem; true exige frequência positiva.
\item[\texttt{frequencia\_pesagem\_dias}] number; N. Intervalo inteiro positivo entre verificações; null quando não exige pesagem frequente.
\item[\texttt{estado\_fisico}] string enum; O. SOLIDO/LIQUIDO; fonte canônica do catálogo, filtro por igualdade.
\item[\texttt{eh\_higroscopico}] boolean; O, padrão false. Fonte canônica da higroscopicidade; não propagar mudanças para snapshots de frascos já cadastrados.
\end{description}

\paragraph{Composicao\_Reagente}\mbox{}\par
Caminho: \texttt{Resumo\_Reagente/\{resumoId\}/Especificacoes/\{id\}: composicao[]}.
\begin{description}
\item[\texttt{id\_substancia\_quimica}] string; O. Identificador da entidade substancia quimica, cuja existência e escopo são verificados no servidor.
\item[\texttt{valor\_min}] number; N. Concentração pontual ou limite inferior não negativo.
\item[\texttt{valor\_max}] number; N. Limite superior; quando informado exige mínimo e valor\_min <= valor\_max.
\item[\texttt{tipo\_concentracao}] string enum; O. Base da concentração; enum M\_M, V\_V, M\_V, MOL\_L, MOL\_KG, PPM ou PPB. Valores: M\_M,V\_V,M\_V,MOL\_L,MOL\_KG,PPM,PPB.
\item[\texttt{notacao\_original\_fabricante}] string; N. Rótulo original, até 50 caracteres; não substitui limites numéricos.
\end{description}

\paragraph{Especificacao\_Reagente}\mbox{}\par
Caminho: \texttt{Resumo\_Reagente/\{resumoId\}/Especificacoes/\{id\}}.
\begin{description}
\item[\texttt{id}] docId string; O. Chave do documento, gerada pelo servidor; não duplicar como contador relacional.
\item[\texttt{id\_resumo\_reagente}] string; O. Identificador da entidade resumo reagente, cuja existência e escopo são verificados no servidor.
\item[\texttt{descricao}] string; O. Valor de descricao, validado pelo formulário e pelo servidor conforme regras do domínio. Máximo 150 caracteres.
\item[\texttt{fabricante}] string; N. Nome do fabricante da especificação comercial; não é o fornecedor do lote. Máximo 150 caracteres.
\item[\texttt{codigo\_produto\_fabricante}] string; N. Código comercial do fabricante para distinguir versões de produto sob o mesmo resumo. Máximo 100 caracteres.
\item[\texttt{grau\_pureza}] string; N. Grau declarado pelo fabricante, por exemplo P.A. ou HPLC; texto opcional, não percentual obrigatório. Máximo 30 caracteres.
\item[\texttt{densidade}] number; N. Densidade em g/ml; positiva e obrigatória para líquido, imutável após confirmação.
\item[\texttt{unidade\_de\_medida}] string enum; projeção física opcional derivada pelo servidor do estado do resumo, g/ml; não é fonte canônica nem campo editável.
\item[\texttt{classe\_inflamabilidade}] string enum; O. Classe selecionada para a especificação; não derivar apenas do nome comercial. Valores: NAO\_INFLAMAVEL,CLASSE\_1,CLASSE\_2,CLASSE\_3.
\item[\texttt{eh\_controlado\_pf}] boolean; O. Indicador cadastral de controle PF declarado pelo gestor; não é inferência automática nem parecer regulatório.
\item[\texttt{eh\_controlado\_eb}] boolean; O. Indicador cadastral de controle EB declarado pelo gestor, independente do indicador PF.
\item[\texttt{link\_fds\_fispq}] string; N. Endereço HTTPS opcional para a ficha de segurança; abrir com proteção de navegação externa. Máximo 255 caracteres.
\end{description}

\paragraph{Lote}\mbox{}\par
Caminho: \texttt{Lote/\{id\}}.
\begin{description}
\item[\texttt{id}] docId string; O. Chave do documento, gerada pelo servidor; não duplicar como contador relacional.
\item[\texttt{id\_resumo\_reagente}] string; O. Projeção física imutável do documento pai da especificação, preenchida no cadastro para permitir agregações por resumo (desvio de 3FN focado em leitura). A base de dados requer backfill.
\item[\texttt{id\_especificacao\_reagente}] string; O. Identificador da entidade especificacao reagente, cuja existência e escopo são verificados no servidor.
\item[\texttt{data\_aquisicao}] Timestamp; O. Data da aquisição do lote; não pode estar no futuro no cadastro de entrada já realizada.
\item[\texttt{qtd\_frascos\_comprados}] number; O. Quantidade adquirida, distinta da contagem materializada de frascos cadastrados.
\item[\texttt{nome\_fornecedor}] string; O. Fornecedor da aquisição; participa da chave composta do lote. Máximo 80 caracteres.
\item[\texttt{numero\_lote}] string; O. Identificação do lote do fornecedor/fabricante, distinta do código LCQUI de cada frasco. Máximo 100 caracteres.
\item[\texttt{nota\_fiscal}] string; O. Referência textual do documento de aquisição; não contém o binário da nota. Máximo 100 caracteres.
\item[\texttt{data\_fabricacao}] string YYYY-MM-DD; O. Data civil declarada de fabricação; não posterior à validade.
\item[\texttt{data\_validade}] string YYYY-MM-DD; O. Data civil de validade do lote; preenche sugestão de validade fechada do frasco.
\end{description}

\paragraph{Frasco\_Reagente}\mbox{}\par
Caminho: \texttt{Frasco\_Reagente/\{id\}}.
\begin{description}
\item[\texttt{id}] docId string; O. Chave do documento, gerada pelo servidor; não duplicar como contador relacional.
\item[\texttt{id\_almoxarifado}] string; O. Identificador da entidade almoxarifado, cuja existência e escopo são verificados no servidor.
\item[\texttt{id\_lote}] string; N. Identificador da entidade lote, cuja existência e escopo são verificados no servidor.
\item[\texttt{id\_especificacao\_reagente}] string; O. Identificador da entidade especificacao reagente, cuja existência e escopo são verificados no servidor. Projeção imutável obrigatória no Firestore.
\item[\texttt{id\_resumo\_reagente}] string; O. Identificador da entidade resumo reagente, cuja existência e escopo são verificados no servidor. Projeção imutável obrigatória no Firestore.
\item[\texttt{detalhe\_local\_armazenamento}] string; N. Valor de detalhe local armazenamento, validado pelo formulário e pelo servidor conforme regras do domínio.
\item[\texttt{codigo\_frasco}] string; O. Identificador LCQUI-N único, gerado pelo servidor. Unicidade garantida por chave determinística/registro de unicidade no servidor.
\item[\texttt{data\_ultima\_pesagem}] Timestamp; N. Data/instante de data ultima pesagem, com ordem temporal validada e autoria do relógio do servidor para eventos.
\item[\texttt{conteudo\_nominal}] number; N. Conteúdo inicial declarado em ml ou g conforme especificação. Se NULL, valor original do rótulo desconhecido ou ilegível (não usar 0 para isso).
\item[\texttt{peso\_no\_cadastrado}] number; O. Peso bruto inicial em g; imutável.
\item[\texttt{peso\_atual}] number; O. Último peso bruto confirmado em g.
\item[\texttt{peso\_frasco\_vazio}] number; N. Tara em g, calculada ou informada conforme modalidade.
\item[\texttt{medida\_usada}] number; O. Acumulado gravimétrico em g; não misturar com volume.
\item[\texttt{data\_abertura}] string YYYY-MM-DD; N. Data efetiva da abertura; não inventar data histórica desconhecida (Q05).
\item[\texttt{validade\_fechado}] string YYYY-MM-DD; N. Validade declarada antes da abertura; pode ser sugerida pelo lote, com origem preservada.
\item[\texttt{validade\_efetiva}] string YYYY-MM-DD; N. Prazo aplicável calculado com abertura e validade fechada; null quando não determinável.
\item[\texttt{validade\_desconhecida}] boolean; O. Ausência explícita de informação de validade; não significa validade ilimitada.
\item[\texttt{validade\_apos\_aberto\_dias}] number; N. Prazo inteiro positivo após abertura quando declarado; compõe o mínimo da validade efetiva.
\item[\texttt{estado\_fisico\_frasco}] string enum; O. Condição do recipiente/conteúdo; quarentena e vencimento são dimensões separadas. Valores: FECHADO,ABERTO,VAZIO,QUEBRADO,DESCARTADO,EXTRAVIADO.
\item[\texttt{disponibilidade}] string enum; O. Disponível ou emprestado; exige coerência com empréstimo ativo. Valores: DISPONIVEL,EMPRESTADO.
\item[\texttt{vencido}] boolean; O. Projeção de validade comparada ao relógio do servidor; reavaliar em cada operação crítica.
\item[\texttt{em\_quarentena}] boolean; O. Bloqueio operacional independente de abertura e vencimento; true exige motivo.
\item[\texttt{uso\_vencido\_autorizado}] boolean; O. Decisão expressa de gestor para uso excepcional conforme Q04, auditada; pesquisa exige TCR e justificativa.
\item[\texttt{detalhe\_status}] string; N. Justificativa obrigatória para quarentena e decisões que a exigem.
\item[\texttt{cadastrado\_em}] Timestamp; O. Data/instante de cadastrado em, com ordem temporal validada e autoria do relógio do servidor para eventos.
\item[\texttt{cadastrado\_por}] string; O. Valor de cadastrado por, validado pelo formulário e pelo servidor conforme regras do domínio.
\item[\texttt{eh\_higroscopico}] boolean; O. Snapshot físico denormalizado, imutável, preenchido pelo backend no cadastro a partir de Resumo\_Reagente.eh\_higroscopico (fonte canônica 3FN). Alterações posteriores do resumo não reescrevem frascos existentes. Não aceitar valor do cliente nem atualização direta. Permite aplicar Q06 na devolução sem leitura adicional do resumo. Frascos legados sem snapshot exigem reconciliação explícita antes da operação.
\item[\texttt{abertura\_historica\_desconhecida}] boolean; O, padrão false. True exige data\_abertura null e estado coerente com abertura anterior, sem inventar data.
\end{description}

\paragraph{Historico\_Frasco\_Reagente}\mbox{}\par
Caminho: \texttt{Historico\_Frasco\_Reagente/\{id\}}.
\begin{description}
\item[\texttt{id}] docId string; O. Chave do documento, gerada pelo servidor; não duplicar como contador relacional.
\item[\texttt{id\_frasco\_reagente}] string; O. Identificador da entidade frasco reagente, cuja existência e escopo são verificados no servidor.
\item[\texttt{id\_almoxarifado}] string; O. Identificador da entidade almoxarifado, cuja existência e escopo são verificados no servidor.
\item[\texttt{id\_gestor}] string; O. Identificador da entidade gestor, cuja existência e escopo são verificados no servidor.
\item[\texttt{tipo}] string enum; O. Valor de tipo, validado pelo formulário e pelo servidor conforme regras do domínio.
\item[\texttt{campo\_ajustado}] string; N. Nome permitido do atributo alterado no evento; não é caminho arbitrário para atualização de documento. Máximo 30 caracteres.
\item[\texttt{id\_emprestimo\_reagente}] string; N. Identificador da entidade emprestimo reagente, cuja existência e escopo são verificados no servidor.
\item[\texttt{peso\_anterior}] number; N. Peso bruto em g imediatamente antes do evento de ajuste/pesagem, quando aplicável.
\item[\texttt{peso\_novo}] number; N. Peso bruto em g confirmado no evento, quando aplicável.
\item[\texttt{medida\_ajustada}] number; N. Diferença operacional na unidade indicada, com sinal definido pelo tipo de evento; não somar automaticamente como consumo.
\item[\texttt{unidade\_medida\_ajustada}] string enum; N. Valor de unidade medida ajustada, validado pelo formulário e pelo servidor conforme regras do domínio. Valores: ml,g. (NULL se evento não for quantitativo).
\item[\texttt{timestamp}] Timestamp; O. Data/instante de timestamp, com ordem temporal validada e autoria do relógio do servidor para eventos.
\end{description}

\paragraph{Emprestimo\_Reagente}\mbox{}\par
Caminho: \texttt{Emprestimo\_Reagente/\{id\}}.
\begin{description}
\item[\texttt{id}] docId string; O. Chave do documento, gerada pelo servidor; não duplicar como contador relacional.
\item[\texttt{id\_frasco\_reagente}] string; O. Identificador da entidade frasco reagente, cuja existência e escopo são verificados no servidor.
\item[\texttt{id\_usuario\_retirou}] string; O. Identificador da entidade usuario retirou, cuja existência e escopo são verificados no servidor.
\item[\texttt{status}] string enum; O. Estado controlado pelo domínio; transições são validadas no servidor. Valores: EM\_USO,ATRASADO,DEVOLVIDO,DEVOLVIDO\_COM\_ATRASO,ENCERRADO\_EXTRAORDINARIO.
\item[\texttt{data\_retirada}] Timestamp; O. Data/instante de data retirada, com ordem temporal validada e autoria do relógio do servidor para eventos.
\item[\texttt{data\_devolucao\_prevista}] string YYYY-MM-DD; O. Data civil cujo prazo termina no fim do dia no fuso institucional.
\item[\texttt{data\_devolucao\_efetuada}] Timestamp; N. Data/instante de data devolucao efetuada, com ordem temporal validada e autoria do relógio do servidor para eventos.
\item[\texttt{id\_local\_usado}] string; O. Identificador da entidade local usado, cuja existência e escopo são verificados no servidor.
\item[\texttt{id\_usuario\_devolveu}] string; N. Identificador da entidade usuario devolveu, cuja existência e escopo são verificados no servidor.
\item[\texttt{peso\_saida}] number; O. Peso bruto em g medido na retirada; referência imutável para o consumo deste empréstimo.
\item[\texttt{peso\_retorno}] number; N. Peso bruto em g medido na devolução; null antes do retorno e validado contra tara/saída.
\item[\texttt{medida\_utilizada}] number; N. Consumo não negativo calculado na unidade do empréstimo; massa = max(0, saída menos retorno), conversão em mL só para líquido com densidade histórica positiva.
\item[\texttt{unidade\_medida\_utilizada}] string enum; O. Valor de unidade medida utilizada, validado pelo formulário e pelo servidor conforme regras do domínio. Valores: ml,g.
\item[\texttt{peso\_perda\_evaporacao}] number; O. Perda em g registrada separadamente de consumo; não inferir arbitrariamente.
\item[\texttt{uso\_vencido\_aceito}] boolean; O. Ciência explícita do empréstimo específico; não concede nem substitui autorização do gestor.
\item[\texttt{finalidade\_uso}] string enum; O. Finalidade validada na retirada. Valores: AULA\_PRATICA, DEMONSTRACAO, PESQUISA\_TCC\_POS, ESTUDO\_DEGRADACAO\_RESIDUOS. Uso excepcional de pesquisa exige rastreabilidade TCR de Q04.
\item[\texttt{auto\_atendimento}] boolean; O, padrão false. Calculado pelo backend conforme Q14, justificativa e notificação à chefia obrigatórias quando true.
\item[\texttt{justificativa\_metodologica}] string; C no uso excepcional Q04, 20--2000 caracteres não vazios.
\item[\texttt{tcr\_versao}] string; C em Q04, até 100. Versão institucional aceita, validada pelo backend.
\item[\texttt{tcr\_aceito\_em}] Timestamp; C em Q04. Instante servidor do aceite.
\item[\texttt{tcr\_aceito\_por}] string; C em Q04. UID autenticado do solicitante/retirante, não do gestor que atende.
\item[\texttt{tcr\_auditoria\_id}] string; C em Q04. Evento imutável em Registro\_de\_Auditoria com frasco, operação, finalidade, versão, instante e aceitante; backend verifica correspondência e consumo único do aceite.
\end{description}

\paragraph{Turma}\mbox{}\par
Caminho: \texttt{Turma/\{id\}}.
\begin{description}
\item[\texttt{id}] docId string; O. Chave do documento, gerada pelo servidor; não duplicar como contador relacional.
\item[\texttt{id\_materia}] string; O. Identificador da entidade materia, cuja existência e escopo são verificados no servidor.
\item[\texttt{id\_professor}] string; O. Identificador da entidade professor, cuja existência e escopo são verificados no servidor.
\item[\texttt{status}] string enum; O. Estado controlado pelo domínio; transições são validadas no servidor. Valores: Ativo, Arquivada.
\item[\texttt{nome\_turma}] string; O. Nome de apresentação da turma, sem inferir ano/semestre a partir do texto. Máximo 100 caracteres.
\item[\texttt{ano}] number; O. Ano letivo da turma ou ano da competência do agregado mensal, conforme documento.
\item[\texttt{semestre}] number; O. Semestre letivo inteiro, exclusivamente 1 ou 2.
\item[\texttt{capacidade}] number; O. Quantidade máxima ordinária de alunos; inteiro positivo, validado na transação de ingresso.
\item[\texttt{codigo\_turma}] string; O. Código único de ingresso, não concede acesso sem validação. Máximo 20 caracteres. Unicidade garantida por chave determinística/registro de unicidade no servidor.
\item[\texttt{data\_criacao}] Timestamp; O. Data/instante de data criacao, com ordem temporal validada e autoria do relógio do servidor para eventos.
\end{description}

\paragraph{Historico\_Alunos\_Turma}\mbox{}\par
Caminho: \texttt{Turma/\{id\}/HistoricoAlunos/\{eventoId\}}.
\begin{description}
\item[\texttt{id}] docId string; O. Chave do documento, gerada pelo servidor; não duplicar como contador relacional.
\item[\texttt{id\_turma}] string; O. Identificador da entidade turma, cuja existência e escopo são verificados no servidor.
\item[\texttt{tipo}] string enum; O. Valor de tipo, validado pelo formulário e pelo servidor conforme regras do domínio. Valores: inclusao\_aluno, exclusao\_aluno.
\item[\texttt{id\_aluno}] string; O. Identificador da entidade aluno, cuja existência e escopo são verificados no servidor.
\item[\texttt{timestamp}] Timestamp; O. Data/instante de timestamp, com ordem temporal validada e autoria do relógio do servidor para eventos.
\end{description}

\paragraph{Historico\_Posts\_Turma}\mbox{}\par
Caminho: \texttt{Turma/\{id\}/Posts/\{postId\}/Historico/\{eventoId\}}.
\begin{description}
\item[\texttt{id}] docId string; O. Chave do documento, gerada pelo servidor; não duplicar como contador relacional.
\item[\texttt{id\_post}] string; O. Identificador da entidade post, cuja existência e escopo são verificados no servidor.
\item[\texttt{novo\_titulo}] string; N. Valor novo titulo, restrito ao campo correspondente da entidade de origem; preservar ausência legítima. Máximo 150 caracteres.
\item[\texttt{novo\_id\_roteiro}] string; N. Valor novo id roteiro, restrito ao campo correspondente da entidade de origem; preservar ausência legítima.
\item[\texttt{nova\_descricao}] string; N. Valor nova descricao, restrito ao campo correspondente da entidade de origem; preservar ausência legítima.
\item[\texttt{antigo\_titulo}] string; N. Valor antigo titulo, restrito ao campo correspondente da entidade de origem; preservar ausência legítima. Máximo 150 caracteres.
\item[\texttt{antigo\_id\_roteiro}] string; N. Valor antigo id roteiro, restrito ao campo correspondente da entidade de origem; preservar ausência legítima.
\item[\texttt{antiga\_descricao}] string; N. Valor antiga descricao, restrito ao campo correspondente da entidade de origem; preservar ausência legítima.
\item[\texttt{timestamp}] Timestamp; O. Data/instante de timestamp, com ordem temporal validada e autoria do relógio do servidor para eventos.
\end{description}

\paragraph{Historico\_Comentario}\mbox{}\par
Caminho: \texttt{Turma/\{id\}/Posts/\{postId\}/Comentarios/\{id\}/Historico/\{eventoId\}}.
\begin{description}
\item[\texttt{id}] docId string; O. Chave do documento, gerada pelo servidor; não duplicar como contador relacional.
\item[\texttt{id\_comentario}] string; O. Identificador da entidade comentario, cuja existência e escopo são verificados no servidor.
\item[\texttt{novo\_texto}] string; O. Valor novo texto, restrito ao campo correspondente da entidade de origem; preservar ausência legítima.
\item[\texttt{texto\_antigo}] string; O. Valor de texto antigo, validado pelo formulário e pelo servidor conforme regras do domínio.
\item[\texttt{editado\_em}] Timestamp; O. Data/instante de editado em, com ordem temporal validada e autoria do relógio do servidor para eventos.
\item[\texttt{editado\_por}] string; O. Valor de editado por, validado pelo formulário e pelo servidor conforme regras do domínio.
\end{description}

\paragraph{Roteiro\_Experimento}\mbox{}\par
Caminho: \texttt{Roteiro\_Experimento/\{id\}}.
\begin{description}
\item[\texttt{id}] docId string; O. Chave do documento, gerada pelo servidor; não duplicar como contador relacional.
\item[\texttt{id\_professor\_upload}] string; O. Identificador da entidade professor upload, cuja existência e escopo são verificados no servidor.
\item[\texttt{file\_url}] string; O. Referência do PDF; download exige revalidar acesso.
\item[\texttt{nome}] string; O. Nome de apresentação; não substitui a chave do documento. Máximo 150 caracteres.
\item[\texttt{descricao}] string; O. Valor de descricao, validado pelo formulário e pelo servidor conforme regras do domínio.
\item[\texttt{professores\_compartilhados}] array de strings; O, padrão vazio. UIDs autorizados, geridos exclusivamente pelo servidor com arrayUnion/arrayRemove; sem escrita direta cliente.
\end{description}

A associação Roteiro\_Professor\_Compartilhado existe somente no modelo relacional 3FN; no Firestore é o array ACL professores\_compartilhados em Roteiro\_Experimento. Consultar Meus por id\_professor\_upload == uid; Compartilhados comigo por array-contains no UID; Eu compartilhei filtra os roteiros próprios carregados com ACL não vazia.

\paragraph{Professor}\mbox{}\par
Caminho: \texttt{Professor/\{id\}}.
\begin{description}
\item[\texttt{id\_usuario}] string; O. Identificador da entidade usuario, cuja existência e escopo são verificados no servidor.
\item[\texttt{centro}] string enum; O. Centro institucional do professor: CCT, CCTA, CBB ou CCH. Valores: CCT, CCTA, CBB, CCH.
\item[\texttt{laboratorio}] string; O. Laboratório de lotação do professor; não fornece automaticamente prédio/andar/sala. Máximo 100 caracteres.
\end{description}

\paragraph{Materia}\mbox{}\par
Caminho: \texttt{Materia/\{id\}}.
\begin{description}
\item[\texttt{id}] docId string; O. Chave do documento, gerada pelo servidor; não duplicar como contador relacional.
\item[\texttt{nome}] string; O. Nome de apresentação; não substitui a chave do documento. Máximo 100 caracteres.
\item[\texttt{codigo\_materia}] string; O. Código acadêmico único da matéria, normalizado de modo uniforme. Máximo 10 caracteres.
\end{description}

\paragraph{Professor\_x\_Materia}\mbox{}\par
Caminho: \texttt{Professor\_x\_Materia/\{id\}}.
\begin{description}
\item[\texttt{id\_professor}] string; O. Identificador da entidade professor, cuja existência e escopo são verificados no servidor.
\item[\texttt{id\_materia}] string; O. Identificador da entidade materia, cuja existência e escopo são verificados no servidor.
\end{description}

\paragraph{Aluno}\mbox{}\par
Caminho: \texttt{Aluno/\{id\}}.
\begin{description}
\item[\texttt{id\_usuario}] string; O. Identificador da entidade usuario, cuja existência e escopo são verificados no servidor.
\item[\texttt{numero\_matricula}] string; O. Matrícula institucional textual, com zeros iniciais; obrigatória e única em Aluno, opcional e não única em Convite. Máximo 20 caracteres.
\end{description}

\paragraph{Bolsista}\mbox{}\par
Caminho: \texttt{Bolsista/\{id\}}.
\begin{description}
\item[\texttt{id\_usuario}] string; O. Identificador da entidade usuario, cuja existência e escopo são verificados no servidor.
\end{description}

\paragraph{Convite\_Aluno}\mbox{}\par
Caminho: \texttt{Convite\_Aluno/\{id\}}.
\begin{description}
\item[\texttt{id}] docId string; O. Chave do documento, gerada pelo servidor; não duplicar como contador relacional.
\item[\texttt{id\_turma}] string; N. Identificador da entidade turma, cuja existência e escopo são verificados no servidor.
\item[\texttt{email}] string; O. Endereço normalizado para identidade/convite; não divulgar em catálogo público interno. Máximo 150 caracteres.
\item[\texttt{convidado\_em}] Timestamp; O. Data/instante de convidado em, com ordem temporal validada e autoria do relógio do servidor para eventos.
\item[\texttt{status}] string enum; O. Estado controlado pelo domínio; transições são validadas no servidor. Valores: pendente, aceitado, expirado.
\item[\texttt{expira\_em}] Timestamp; O. Instante de expiração de convite/notificação; não equivale a autorização de excluir auditoria.
\item[\texttt{convidado\_por}] string; O. Valor de convidado por, validado pelo formulário e pelo servidor conforme regras do domínio.
\item[\texttt{numero\_matricula}] string; N. Matrícula institucional textual, com zeros iniciais; obrigatória e única em Aluno, opcional e não única em Convite. Máximo 20 caracteres.
\item[\texttt{token\_hash}] string; O. SHA-256 do token aleatório de 32 bytes (CSPRNG); nunca token em claro. Reenvio substitui hash e expiração atomicamente; comparação constante e aceitação única no servidor.
\item[\texttt{ultimo\_reenvio\_por}] string; N. UID servidor do operador do último reenvio. DocId derivado por HMAC-SHA-256 de contexto e e-mail normalizado com segredo servidor; não usar hash simples público de e-mail.
\end{description}

\paragraph{Notificacao}\mbox{}\par
Caminho: \texttt{Usuarios/\{uid\}/Notificacoes/\{id\}}.
\begin{description}
\item[\texttt{id}] docId string; O. Chave do documento, gerada pelo servidor; não duplicar como contador relacional.
\item[\texttt{id\_destinatario}] string; O. Identificador da entidade destinatario, cuja existência e escopo são verificados no servidor.
\item[\texttt{papel\_destinatario}] string enum; O. Contexto de exibição da notificação, sem conceder esse papel ao usuário. Inclui Aluno, Professor, Bolsista, Gestor\_Bens\_Patrimoniais e Gestor\_Almoxarifado.
\item[\texttt{tipo}] string enum; O. Valor de tipo, validado pelo formulário e pelo servidor conforme regras do domínio.
\item[\texttt{id\_quem\_fez\_acao}] string; N. Identificador da entidade quem fez acao, cuja existência e escopo são verificados no servidor.
\item[\texttt{id\_turma}] string; N. Identificador da entidade turma, cuja existência e escopo são verificados no servidor.
\item[\texttt{quantidade}] number; N. Quantidade agregada de ocorrências agrupadas na notificação, quando aplicável.
\item[\texttt{entidade\_alvo}] string; O. Tipo autorizado de recurso para navegação; mapear a rotas conhecidas, não URL arbitrária. Máximo 40 caracteres.
\item[\texttt{id\_alvo}] string; O. ID do recurso da notificação; acesso é revalidado ao abrir. Máximo 200 caracteres.
\item[\texttt{lida}] boolean; O. Estado de leitura da notificação; true preserva o registro em vez de excluí-lo.
\item[\texttt{lida\_em}] Timestamp; N. Instante servidor de marcação de leitura; null enquanto não lida.
\item[\texttt{emitida\_em}] Timestamp; O. Data/instante de emitida em, com ordem temporal validada e autoria do relógio do servidor para eventos.
\item[\texttt{expira\_em}] Timestamp; O. Instante de expiração de convite/notificação; não equivale a autorização de excluir auditoria.
\end{description}

\paragraph{Aluno\_x\_Turma}\mbox{}\par
Caminho: \texttt{Turma/\{id\}/Alunos/\{uid\}}.
\begin{description}
\item[\texttt{id\_aluno}] string; O. Identificador da entidade aluno, cuja existência e escopo são verificados no servidor.
\item[\texttt{id\_turma}] string; O. Identificador da entidade turma, cuja existência e escopo são verificados no servidor.
\end{description}

\paragraph{Local}\mbox{}\par
Caminho: \texttt{Local/\{id\}}.
\begin{description}
\item[\texttt{id}] docId string; O. Chave do documento, gerada pelo servidor; não duplicar como contador relacional.
\item[\texttt{predio}] string; O. Identificador textual do prédio no cadastro de locais. Máximo 30 caracteres.
\item[\texttt{andar}] string; O. Identificador textual do piso; permite térreo/subsolo, sem impor número inteiro. Máximo 10 caracteres.
\item[\texttt{sala}] string; O. Identificador textual da sala, único no contexto prédio/andar. Máximo 30 caracteres.
\item[\texttt{criado\_por}] string; O. Valor de criado por, validado pelo formulário e pelo servidor conforme regras do domínio.
\item[\texttt{criado\_em}] Timestamp; O. Data/instante de criado em, com ordem temporal validada e autoria do relógio do servidor para eventos.
\end{description}

\paragraph{Post}\mbox{}\par
Caminho: \texttt{Turma/\{id\}/Posts/\{postId\}}.
\begin{description}
\item[\texttt{id}] docId string; O. Chave do documento, gerada pelo servidor; não duplicar como contador relacional.
\item[\texttt{id\_professor}] string; O. Identificador da entidade professor, cuja existência e escopo são verificados no servidor.
\item[\texttt{id\_turma}] string; O. Identificador da entidade turma, cuja existência e escopo são verificados no servidor.
\item[\texttt{id\_roteiro\_experimento}] string; N. Identificador da entidade roteiro experimento, cuja existência e escopo são verificados no servidor.
\item[\texttt{titulo}] string; O. Título obrigatório de publicação, exibido no feed e na navegação de notificações. Máximo 150 caracteres.
\item[\texttt{descricao}] string; O. Valor de descricao, validado pelo formulário e pelo servidor conforme regras do domínio.
\item[\texttt{criado\_em}] Timestamp; O. Data/instante de criado em, com ordem temporal validada e autoria do relógio do servidor para eventos.
\item[\texttt{roteiro\_anexo}] map; N. Snapshot imutável da publicação: id\_roteiro (string O), nome\_arquivo (string O, até 150), tamanho\_bytes (inteiro positivo O), storage\_path (string O). Backend valida acesso ao roteiro antes de publicar. Revogação posterior da ACL não remove o anexo histórico. Aluno recebe os metadados no feed; download via obterUrlDownloadRoteiro valida vínculo atual e emite URL assinada por 15 minutos, sem leitura cliente da coleção raiz.
\end{description}

\paragraph{Comentario}\mbox{}\par
Caminho: \texttt{Turma/\{id\}/Posts/\{postId\}/Comentarios/\{id\}}.
\begin{description}
\item[\texttt{id}] docId string; O. Chave do documento, gerada pelo servidor; não duplicar como contador relacional.
\item[\texttt{id\_post}] string; O. Identificador da entidade post, cuja existência e escopo são verificados no servidor.
\item[\texttt{id\_usuario}] string; O. Identificador da entidade usuario, cuja existência e escopo são verificados no servidor.
\item[\texttt{texto}] string; O. Conteúdo textual do comentário, escapado na renderização; rejeitar texto vazio após trim.
\item[\texttt{criado\_em}] Timestamp; O. Data/instante de criado em, com ordem temporal validada e autoria do relógio do servidor para eventos.
\item[\texttt{moderado}] boolean; O, padrão false. Visibilidade conforme DP-C02.
\item[\texttt{motivo\_moderacao}] string; C quando moderado, até 2000 caracteres não vazios.
\item[\texttt{moderado\_por}] string; C quando moderado. UID autenticado do moderador, preenchido pelo servidor.
\end{description}

\paragraph{Registro\_de\_Auditoria}\mbox{}\par
Caminho: \texttt{Registro\_de\_Auditoria/\{id\}}.
\begin{description}
\item[\texttt{id}] docId string; O. Chave do documento, gerada pelo servidor; não duplicar como contador relacional.
\item[\texttt{id\_usuario}] string; O. Identificador da entidade usuario, cuja existência e escopo são verificados no servidor.
\item[\texttt{acao}] string; O. Nome da ação auditada; distinguir solicitação, conclusão e rejeição quando relevante. Máximo 200 caracteres.
\item[\texttt{tipo\_entidade\_sofre\_acao}] string enum; O. Valor de tipo entidade sofre acao, validado pelo formulário e pelo servidor conforme regras do domínio.
\item[\texttt{id\_do\_objeto\_da\_entidade}] string; O. Identificador da entidade do objeto da entidade, cuja existência e escopo são verificados no servidor. Máximo 200 caracteres.
\item[\texttt{acao\_feita\_em}] Timestamp; O. Data/instante de acao feita em, com ordem temporal validada e autoria do relógio do servidor para eventos.
\item[\texttt{metadata}] map; O. Mapa auditável e limitado ao evento, sem senha, token ou conteúdo binário.
\end{description}

\paragraph{Lote\_Materializado}\mbox{}\par
Caminho: \texttt{Lote\_Materializado/\{id\}}.
\begin{description}
\item[\texttt{id}] docId string; O. Chave do documento, gerada pelo servidor; não duplicar como contador relacional.
\item[\texttt{id\_lote}] string; O. Identificador da entidade lote, cuja existência e escopo são verificados no servidor. Unicidade garantida por chave determinística/registro de unicidade no servidor.
\item[\texttt{qtd\_frascos\_cadastrados}] number; O. Contagem de frascos cadastrados, inteira não negativa; não editável pelo cliente quando materializada.
\item[\texttt{ultimo\_reconciliador}] Timestamp; O (se já reconciliado). Marca temporal absoluta, preenchida/atualizada pelo job de reconciliação. Eventos de trigger gerados antes deste instante são ignorados, evitando que eventos assíncronos retardados somem duplamente sobre a contagem absoluta já aferida.
\end{description}

\paragraph{Auditoria\_Reconciliacao}\mbox{}\par
Caminho: \texttt{Auditoria\_Reconciliacao/\{id\}}.
\begin{description}
\item[\texttt{id}] docId string; O. Chave gerada pelo servidor para o registro da cura.
\item[\texttt{id\_lote}] string; O. Lote alvo da reconciliação onde ocorreu divergência.
\item[\texttt{qtd\_anterior}] number; O. Contagem que constava no documento \texttt{Lote\_Materializado} antes da correção (drift).
\item[\texttt{qtd\_corrigida}] number; O. Contagem absoluta efetivamente encontrada no banco de dados e aplicada (count real).
\item[\texttt{data}] Timestamp; O. Instante exato em que a correção transacional foi aplicada e a marca d'água elevada.
\end{description}

\paragraph{Resumo\_Almoxarifado\_Diario}\mbox{}\par
Caminho: \texttt{Resumo\_Almoxarifado\_Diario/\{id\}}.
\begin{description}
\item[\texttt{id}] docId string; O. Chave do documento, gerada pelo servidor; não duplicar como contador relacional.
\item[\texttt{data}] string YYYY-MM-DD; O. Data/instante de data, com ordem temporal validada e autoria do relógio do servidor para eventos.
\item[\texttt{id\_almoxarifado}] string; O. Identificador da entidade almoxarifado, cuja existência e escopo são verificados no servidor.
\item[\texttt{qtd\_frascos\_fechados\_no\_fim\_do\_dia}] number; O. Contagem de frascos fechados no fim do dia, inteira não negativa; não editável pelo cliente quando materializada.
\item[\texttt{qtd\_frascos\_abertos\_no\_fim\_do\_dia}] number; O. Contagem de frascos abertos no fim do dia, inteira não negativa; não editável pelo cliente quando materializada.
\item[\texttt{qtd\_frascos\_emprestados\_no\_fim\_do\_dia}] number; O. Contagem de frascos emprestados no fim do dia, inteira não negativa; não editável pelo cliente quando materializada.
\item[\texttt{qtd\_frascos\_cadastrados\_fechados\_durante\_o\_dia}] number; O. Contagem de frascos cadastrados fechados durante o dia, inteira não negativa; não editável pelo cliente quando materializada.
\item[\texttt{qtd\_frascos\_cadastrados\_abertos\_durante\_o\_dia}] number; O. Contagem de frascos cadastrados abertos durante o dia, inteira não negativa; não editável pelo cliente quando materializada.
\item[\texttt{qtd\_frascos\_emprestados\_durante\_o\_dia}] number; O. Contagem de frascos emprestados durante o dia, inteira não negativa; não editável pelo cliente quando materializada.
\item[\texttt{qtd\_frascos\_devolvidos\_durante\_o\_dia}] number; O. Contagem de frascos devolvidos durante o dia, inteira não negativa; não editável pelo cliente quando materializada.
\item[\texttt{volume\_total\_usado\_nos\_frascos\_devolvidos\_durante\_o\_dia}] number; O. Métrica de volume total usado nos frascos devolvidos durante o dia, não negativa e calculada a partir dos eventos/saldos do domínio.
\item[\texttt{volume\_total\_disponivel\_ml}] number; O. Métrica de volume total disponivel ml, não negativa e calculada a partir dos eventos/saldos do domínio.
\item[\texttt{volume\_utilizado\_no\_dia\_ml}] number; O. Métrica de volume utilizado no dia ml, não negativa e calculada a partir dos eventos/saldos do domínio.
\item[\texttt{massa\_total\_disponivel\_g}] number; O. Métrica de massa total disponivel g, não negativa e calculada a partir dos eventos/saldos do domínio.
\item[\texttt{massa\_utilizada\_no\_dia\_g}] number; O. Métrica de massa utilizada no dia g, não negativa e calculada a partir dos eventos/saldos do domínio.
\item[\texttt{created\_at}] Timestamp; O. Instante servidor da primeira criação da materialização.
\item[\texttt{updated\_at}] Timestamp; O. Instante servidor da última atualização da materialização; informa atualidade do resultado.
\end{description}

\paragraph{Resumo\_Reagente\_Diario}\mbox{}\par
Caminho: \texttt{Resumo\_Reagente\_Diario/\{id\}}.
\begin{description}
\item[\texttt{id}] docId string; O. Chave do documento, gerada pelo servidor; não duplicar como contador relacional.
\item[\texttt{data}] string YYYY-MM-DD; O. Data/instante de data, com ordem temporal validada e autoria do relógio do servidor para eventos.
\item[\texttt{id\_resumo\_reagente}] string; O. Identificador da entidade resumo reagente, cuja existência e escopo são verificados no servidor.
\item[\texttt{id\_almoxarifado}] string; O. Identificador da entidade almoxarifado, cuja existência e escopo são verificados no servidor.
\item[\texttt{qtd\_frascos\_fechados\_disponiveis}] number; O. Contagem de frascos fechados disponiveis, inteira não negativa; não editável pelo cliente quando materializada.
\item[\texttt{qtd\_frascos\_abertos\_disponiveis}] number; O. Contagem de frascos abertos disponiveis, inteira não negativa; não editável pelo cliente quando materializada.
\item[\texttt{qtd\_frascos\_em\_uso}] number; O. Contagem de frascos em uso, inteira não negativa; não editável pelo cliente quando materializada.
\item[\texttt{qtd\_frascos\_vazios}] number; O. Contagem de frascos vazios, inteira não negativa; não editável pelo cliente quando materializada.
\item[\texttt{qtd\_frascos\_descartados\_dia}] number; O. Contagem de frascos descartados dia, inteira não negativa; não editável pelo cliente quando materializada.
\item[\texttt{volume\_total\_disponivel\_ml}] number; O. Métrica de volume total disponivel ml, não negativa e calculada a partir dos eventos/saldos do domínio.
\item[\texttt{volume\_utilizado\_no\_dia\_ml}] number; O. Métrica de volume utilizado no dia ml, não negativa e calculada a partir dos eventos/saldos do domínio.
\item[\texttt{volume\_evaporado\_no\_dia\_ml}] number; O. Métrica de volume evaporado no dia ml, não negativa e calculada a partir dos eventos/saldos do domínio.
\item[\texttt{massa\_total\_disponivel\_g}] number; O. Métrica de massa total disponivel g, não negativa e calculada a partir dos eventos/saldos do domínio.
\item[\texttt{massa\_utilizada\_no\_dia\_g}] number; O. Métrica de massa utilizada no dia g, não negativa e calculada a partir dos eventos/saldos do domínio.
\item[\texttt{massa\_evaporada\_no\_dia\_g}] number; O. Métrica de massa evaporada no dia g, não negativa e calculada a partir dos eventos/saldos do domínio.
\item[\texttt{created\_at}] Timestamp; O. Instante servidor da primeira criação da materialização.
\item[\texttt{updated\_at}] Timestamp; O. Instante servidor da última atualização da materialização; informa atualidade do resultado.
\end{description}

\paragraph{Resumo\_Bem\_Patrimonial\_Diario}\mbox{}\par
Caminho: \texttt{Resumo\_Bem\_Patrimonial\_Diario/\{id\}}.
\begin{description}
\item[\texttt{id}] docId string; O. Chave do documento, gerada pelo servidor; não duplicar como contador relacional.
\item[\texttt{data}] string YYYY-MM-DD; O. Data/instante de data, com ordem temporal validada e autoria do relógio do servidor para eventos.
\item[\texttt{id\_local}] string; N. Identificador da entidade local, cuja existência e escopo são verificados no servidor.
\item[\texttt{qtd\_bens\_ativos\_no\_fim\_do\_dia}] number; O. Contagem de bens ativos no fim do dia, inteira não negativa; não editável pelo cliente quando materializada.
\item[\texttt{qtd\_bens\_inservivel\_no\_fim\_do\_dia}] number; O. Contagem de bens inservivel no fim do dia, inteira não negativa; não editável pelo cliente quando materializada.
\item[\texttt{qtd\_bens\_dado\_baixa\_no\_fim\_do\_dia}] number; O. Contagem de bens dado baixa no fim do dia, inteira não negativa; não editável pelo cliente quando materializada.
\item[\texttt{qtd\_bens\_cadastrados\_durante\_o\_dia}] number; O. Contagem de bens cadastrados durante o dia, inteira não negativa; não editável pelo cliente quando materializada.
\item[\texttt{qtd\_bens\_marcados\_inservivel\_durante\_o\_dia}] number; O. Contagem de bens marcados inservivel durante o dia, inteira não negativa; não editável pelo cliente quando materializada.
\item[\texttt{qtd\_bens\_dados\_baixa\_durante\_o\_dia}] number; O. Contagem de bens dados baixa durante o dia, inteira não negativa; não editável pelo cliente quando materializada.
\item[\texttt{qtd\_requisicoes\_abertas\_durante\_o\_dia}] number; O. Contagem de requisicoes abertas durante o dia, inteira não negativa; não editável pelo cliente quando materializada.
\item[\texttt{qtd\_requisicoes\_aprovadas\_durante\_o\_dia}] number; O. Contagem de requisicoes aprovadas durante o dia, inteira não negativa; não editável pelo cliente quando materializada.
\item[\texttt{qtd\_requisicoes\_rejeitadas\_durante\_o\_dia}] number; O. Contagem de requisicoes rejeitadas durante o dia, inteira não negativa; não editável pelo cliente quando materializada.
\item[\texttt{created\_at}] Timestamp; O. Instante servidor da primeira criação da materialização.
\item[\texttt{updated\_at}] Timestamp; O. Instante servidor da última atualização da materialização; informa atualidade do resultado.
\end{description}

\paragraph{Atividade\_Gestor\_Almoxarifado\_Mensal}\mbox{}\par
Caminho: \texttt{Atividade\_Gestor\_Almoxarifado\_Mensal/\{id\}}.
\begin{description}
\item[\texttt{id}] docId string; O. Chave do documento, gerada pelo servidor; não duplicar como contador relacional.
\item[\texttt{id\_gestor}] string; O. Identificador da entidade gestor, cuja existência e escopo são verificados no servidor.
\item[\texttt{ano}] number; O. Ano letivo da turma ou ano da competência do agregado mensal, conforme documento.
\item[\texttt{mes}] number; O. Mês da competência mensal, inteiro de 1 a 12.
\item[\texttt{qtd\_reagentes\_adicionados}] number; O. Contagem de reagentes adicionados, inteira não negativa; não editável pelo cliente quando materializada.
\item[\texttt{qtd\_frascos\_adicionados}] number; O. Contagem de frascos adicionados, inteira não negativa; não editável pelo cliente quando materializada.
\item[\texttt{qtd\_movimentacoes\_registradas}] number; O. Contagem de movimentacoes registradas, inteira não negativa; não editável pelo cliente quando materializada.
\item[\texttt{created\_at}] Timestamp; O. Instante servidor da primeira criação da materialização.
\item[\texttt{updated\_at}] Timestamp; O. Instante servidor da última atualização da materialização; informa atualidade do resultado.
\end{description}

\paragraph{Atividade\_Gestor\_Bens\_Patrimoniais\_Mensal}\mbox{}\par
Caminho: \texttt{Atividade\_Gestor\_Bens\_Patrimoniais\_Mensal/\{id\}}.
\begin{description}
\item[\texttt{id}] docId string; O. Chave do documento, gerada pelo servidor; não duplicar como contador relacional.
\item[\texttt{id\_gestor}] string; O. Identificador da entidade gestor, cuja existência e escopo são verificados no servidor.
\item[\texttt{ano}] number; O. Ano letivo da turma ou ano da competência do agregado mensal, conforme documento.
\item[\texttt{mes}] number; O. Mês da competência mensal, inteiro de 1 a 12.
\item[\texttt{qtd\_bens\_adicionados}] number; O. Contagem de bens adicionados, inteira não negativa; não editável pelo cliente quando materializada.
\item[\texttt{qtd\_bens\_editados}] number; O. Contagem de bens editados, inteira não negativa; não editável pelo cliente quando materializada.
\item[\texttt{qtd\_requisicoes\_aprovadas}] number; O. Contagem de requisicoes aprovadas, inteira não negativa; não editável pelo cliente quando materializada.
\item[\texttt{qtd\_requisicoes\_rejeitadas}] number; O. Contagem de requisicoes rejeitadas, inteira não negativa; não editável pelo cliente quando materializada.
\item[\texttt{created\_at}] Timestamp; O. Instante servidor da primeira criação da materialização.
\item[\texttt{updated\_at}] Timestamp; O. Instante servidor da última atualização da materialização; informa atualidade do resultado.
\end{description}

\paragraph{Complementos físicos obrigatórios e fontes de verdade.}
\begin{description}
\item[Usuarios] \texttt{ativo} (boolean, O), \texttt{versao\_permissoes} (inteiro, O), \texttt{claims\_pendentes} (boolean, O), \texttt{atualizado\_em} (Timestamp, O): controle servidor de ativação e sincronização de autorização. Não replicar papéis editáveis no navegador.
\item[Perfis] docId = UID em Chefe\_Geral, Professor, Aluno, Bolsista e gestores; \texttt{id\_usuario} é esse UID. \texttt{nome}, \texttt{email} e \texttt{letra\_inicial}, se usados em busca de perfil, são projeções de Usuarios, com leitura restrita e atualização pelo servidor.
\item[Almoxarifado] \texttt{ativo} (boolean, O) e \texttt{qtd\_gestores\_ativos} (inteiro não negativo, O) são necessários à ativação e revogação concorrente. O contador é derivado dos vínculos; ativo exige contador positivo.
\item[Turma] \texttt{qtd\_alunos} (inteiro não negativo, O) é contador transacional do vínculo ativo, não count lido fora da transação. \texttt{versao} (inteiro positivo, O) permite detectar edições concorrentes.
\item[Espelho Aluno--Turma] \texttt{Usuarios/uid/Turmas/turmaId}: \texttt{id\_turma}, \texttt{id\_professor}, \texttt{id\_materia}, \texttt{nome\_turma}, \texttt{nome\_materia} (strings O), \texttt{ano}, \texttt{semestre} (inteiros O), \texttt{status} (enum O), \texttt{ingressou\_em} (Timestamp O). Todos derivados de Turma/Matéria ou evento de ingresso; não usar espelho desatualizado como autorização. Vínculo canônico em Turma/id/Alunos/uid contém \texttt{id\_aluno} (string O), \texttt{ingressou\_em} (Timestamp O); nome/foto para colegas são projeções mínimas, sem matrícula/e-mail.
\item[Reagentes] Resumo possui \texttt{letra\_inicial} (string O, derivada do nome). O estado físico escalar e a higroscopicidade são canônicos no resumo; filtro por igualdade. Especificação embute \texttt{composicao} (array de maps com os campos da composição; IDs de substâncias únicos). O resumo pai fornece seu ID; manter ID de resumo em Lote e Frasco para resolver o caminho completo da especificação.
\item[Lote e Frasco] \texttt{id\_resumo\_reagente}, \texttt{id\_especificacao\_reagente}, \texttt{nome\_reagente} (strings O) são projeções servidor; em Frasco, especificação permanece resolvida mesmo com lote, como desnormalização deliberada. \texttt{unidade\_medida} (ml/g, O), \texttt{localizacao\_detalhada} (string até 500, N), \texttt{versao} (inteiro O), \texttt{ultima\_pesagem\_em} (Timestamp, N) e \texttt{modalidade\_cadastro} (FECHADO/CONHECE\_TARA/ESTIMA\_VOLUME/ESTIMA\_MASSA, O) sustentam UI-06. \texttt{id\_emprestimo\_ativo} (string N) é ponteiro transacional, limpo no retorno. \texttt{id\_almoxarifado} é fonte da localização operacional, nunca aceito sem escopo.
\item[Empréstimo] \texttt{id\_almoxarifado}, \texttt{id\_resumo\_reagente}, \texttt{id\_especificacao\_reagente} (strings O), \texttt{densidade\_aplicada} (number positiva C para líquido), \texttt{id\_gestor\_retirada} (string O), \texttt{id\_gestor\_devolucao} (string N até retorno). Snapshot químico da retirada preserva a memória de cálculo; não atualizar retroativamente. Retirante, devolvente e operador são papéis distintos.
\item[Patrimônio] \texttt{nome\_equipamento}, \texttt{letra\_inicial\_nome}, \texttt{predio}, \texttt{andar}, \texttt{sala} (strings O) são projeções de Resumo/Local. \texttt{versao} (inteiro O) detecta conflitos. Histórico embute \texttt{alteracoes} (array de maps campo/valor anterior/novo, admitindo null para ausência) e snapshot do local quando necessário a relatório histórico; não propagar local atual sobre fatos passados.
\item[Requisições] Adição inclui \texttt{justificativa\_resposta} (string N, obrigatória após resposta). Edição inclui \texttt{versao\_bem\_origem} (inteiro O), \texttt{novo\_id\_resumo\_bem\_patrimonial} (string N, reclassificação) e \texttt{novo\_documento\_baixa} (string N, C para proposta de baixa). Projeções nome/plaqueta são identificadores de leitura, não autorização. Resposta nunca inventa foto ausente.
\item[Convite] \texttt{exceder\_capacidade} (boolean O, default false) e \texttt{justificativa\_excecao} (string C quando true), \texttt{aceitado\_por} (string N), \texttt{aceitado\_em} (Timestamp N), \texttt{token\_hash} (string O, SHA-256 do token aleatório). Nunca guardar token em claro. Status aceitado exige UID e instante; matrícula obrigatória na identidade final.
\item[Arquivos] Roteiros, fotos e comprovantes mantêm \texttt{storage\_path} (string O quando upload), \texttt{content\_type} (string O), \texttt{tamanho\_bytes} (inteiro positivo O), \texttt{owner\_uid} (string O), \texttt{criado\_em} (Timestamp O) no registro/metadados aplicável. \texttt{file\_url}/\texttt{photo\_url} são referências legadas; URL assinada transitória não é fonte de autorização nem conteúdo persistente. Arquivo e Firestore não têm transação única: usar upload provisório e reconciliação de órfãos, preservando anexos referenciados.
\item[Etiquetas] \texttt{Contador\_Codigo\_Frasco/singleton}: \texttt{ultimo\_codigo\_gerado} (inteiro não negativo O). Não criar simultaneamente Contadores/codigo\_frasco. Impressao\_Etiqueta\_Frasco acrescenta \texttt{start\_row} (1--10), \texttt{start\_col} (1--3), \texttt{id\_operacao} (string O) e \texttt{status\_geracao} (SOLICITADA/CONCLUIDA/FALHOU). Segunda via fica na auditoria; metadata inclui operação, frasco, motivo e resultado, sem duplicar registro como lote virgem.
\item[Locks] Locks\_Requisicao\_Patrimonio contém \texttt{id\_requisicao}, \texttt{tipo} (ADICAO/EDICAO), \texttt{chave\_recurso} (string O) e \texttt{criado\_em} (Timestamp O). Lock e requisição são atômicos; limpeza só após verificar inexistência de pendência vinculada.
\item[Unicidade técnica] \texttt{Chaves\_Unicas/chaveHash}: \texttt{tipo}, \texttt{id\_recurso} (strings O), \texttt{criado\_em} (Timestamp O). Chave determinística derivada do valor normalizado e domínio (ex: \texttt{Bem\_Patrimonial\_\_\{numero\}} para plaqueta, matrícula, CAS, matéria, local, código turma, lote ou convite pendente); nunca expor e-mail na chave. Criar/liberar na mesma transação do recurso. A implementação exige \textbf{backfill} populando esta coleção para todos os registros preexistentes antes de ativar a verificação transacional.
\item[Deduplicação de Eventos] \texttt{Eventos\_Processados/eventId}: \texttt{processado\_em} (Timestamp O). Usado em transações originadas por triggers (ex: \texttt{onDocumentCreated}) para converter a entrega "at-least-once" do Cloud Functions em processamento "exactly-once".
\item[Operações técnicas] \texttt{Operacoes/id}: \texttt{uid}, \texttt{acao}, \texttt{recurso}, \texttt{hash\_entrada} (strings O), \texttt{status} (PENDENTE/CONCLUIDA/FALHOU), \texttt{resultado} (map N), \texttt{criado\_em}, \texttt{atualizado\_em} (Timestamp O). Chave é vinculada ao ator/ação; repetição com entrada diferente é rejeitada. Suporta idempotência, notificações e reconciliação Auth/Storage. Leitura restrita ao solicitante por resposta mínima; escrita servidor.
\item[Controle de papéis] \texttt{Controle\_Papeis/singleton}: \texttt{chefes\_ativos}, \texttt{gestores\_patrimoniais\_ativos} (inteiros não negativos O), \texttt{versao} (inteiro O). Atualizado na mesma transação das concessões/revogações para impedir retirada concorrente do último responsável.
\item[Materializações] Chaves por lote; data+almoxarifado; data+resumo+almoxarifado; data+local (sentinela GERAL); gestor+ano+mês. Todos os campos de métricas do dicionário são calculados no servidor. Acrescentar \texttt{calculado\_em} (Timestamp O) e \texttt{versao\_calculo} (inteiro O); reprocessar a mesma janela substitui resultado ou deduplica eventos, sem somar duas vezes. Não somar linha GERAL às linhas locais. Métrica legada de volume total usado é alias do consumo ml e não uma parcela adicional.
\end{description}

\paragraph{Consultas e integridade do dicionário.}
A entidade 3FN mantém dependências e referências canônicas; desnormalização é exceção documentada para consulta, não replicação automática de toda FK. Atualização transacional é obrigatória em vínculos, saldos e exclusões lógicas; propagação de nomes pode ser assíncrona, com versão da origem e reconciliação. Nunca aplicar trigger de nome atual a snapshot histórico. Histórico por bem admite collection-group autorizado para relatórios, portanto subcoleção não impede consulta global.

As consultas UI exigem índices versionados em firestore.indexes.json: empréstimo por almoxarifado/status/data prevista e por retirante/data; histórico de frasco por frasco/timestamp e almoxarifado/timestamp; requisição por solicitante/status/feita\_em; turma por professor/status/ano; notificações por papel/lida/emitida\_em; resumo diário por escopo/data. Acrescentar docId como desempate da paginação. Combinações opcionais de filtros precisam de inventário de consultas e validação em emulador/ambiente de integração, não promessa de qualquer combinação arbitrária.

\endgroup
\subsection{Estratégia de Busca Textual no Firestore}
\label{subsec:busca-textual}

\begin{explicacao}[Por que não existe \texttt{LIKE} no Firestore]
O Firestore não oferece busca por substring (sem \texttt{LIKE}, sem \texttt{icontains}). A estratégia adotada é: aplicar \textbf{filtros obrigatórios de igualdade} no Firestore para reduzir drasticamente o conjunto de documentos retornados, e então aplicar \textbf{substring no frontend} sobre os poucos resultados restantes.
\end{explicacao}

\subsubsection{Reagentes}
O usuário deve obrigatoriamente escolher ao menos um dos filtros: estado físico (\texttt{SOLIDO} ou \texttt{LIQUIDO}), natureza química (\texttt{ORGANICO}, \texttt{INORGANICO}, \texttt{ELEMENTO} ou \texttt{HIBRIDO}), tipo de substância (\texttt{PURA} ou \texttt{MISTURA}) e/ou letra inicial (campo \texttt{letra\_inicial} denormalizado). A query Firestore usa esses filtros de igualdade e retorna poucos documentos; o frontend aplica a busca por substring sobre o campo \texttt{nome}.

\subsubsection{Bens Patrimoniais}
Combinação obrigatória de pelo menos um: (Prédio) ou (letra inicial do nome do equipamento) ou (Status). O Firestore filtra com os campos denormalizados \texttt{predio}, \texttt{andar}, \texttt{sala}, \texttt{status} e \texttt{letra\_inicial\_nome}. O frontend aplica substring sobre \texttt{nome\_equipamento} e \texttt{numero\_patrimonio} nos resultados. Bigramas não serão implementados na V1: o filtro de (Prédio + Andar) já reduz os candidatos suficientemente.

\subsubsection{Alunos}
Busca via campo denormalizado \texttt{letra\_inicial} (filtro de igualdade). O frontend aplica substring sobre \texttt{nome} nos resultados. A \textit{range query} por prefixo (\texttt{nome >= ``A''}, \texttt{nome < ``B''}) é alternativa válida, mas não combina facilmente com outros filtros de igualdade em índice composto no Firestore; o campo \texttt{letra\_inicial} é preferível.

\subsubsection{Professores}
Poucos no sistema (dezenas). Usar listagem autorizada e paginada com campos mínimos; filtrar por substring apenas os registros carregados. UI-13 define cache e ciclo de vida.

\begin{regra}[Listeners e cache com escopo]
Aplicar UI-13: listeners limitados às telas visíveis, histórico paginado, memória por padrão e persistência somente em dispositivo confiável. Cache não garante atualização nem ausência de novas leituras faturadas ao reconectar. Encerrar listeners e limpar dados ao sair ou trocar identidade.
\end{regra}

\subsection{Espelhamento Bidirecional da Associação Aluno--Turma}
\begin{itemize}
    \item \textbf{Justificativa NoSQL:} No modelo 3FN, a tabela associativa \texttt{Aluno\_x\_Turma} responde bidirecionalmente às consultas via chave composta. No Firestore, como subcoleções não podem ser filtradas a partir da coleção pai (\texttt{Turma}), e regras de segurança não atuam como filtros de busca, a relação $N:N$ é intencionalmente denormalizada e espelhada em duas estruturas complementares:
    \begin{enumerate}
        \item \textbf{\nolinkurl{Turma/\{turmaId\}/Alunos/\{alunoId\}} (Visão da Turma):} Utilizada pelo professor para gerenciar os matriculados e validar a regra de capacidade máxima (\texttt{qtd\_alunos < capacidade}, atualizado na transação).
        \item \textbf{\nolinkurl{Usuarios/\{uid\}/Turmas/\{turmaId\}} (Visão do Aluno):} Utilizada exclusivamente pelo frontend do Aluno para alimentar o painel lateral de ``Minhas Turmas'' via listener em tempo real (\texttt{onSnapshot}), contendo os metadados essenciais (\texttt{nome\_turma}, \texttt{nome\_materia}, \texttt{ano}, \texttt{semestre}, \texttt{id\_professor}).
    \end{enumerate}
    \item \textbf{Consistência Transacional:} A criação e a remoção desse espelho são executadas de forma estritamente atômica pelo backend nas Cloud Functions. Na matrícula (\texttt{ingressarEmTurmaPorCodigo}), o espelho é criado. Na remoção (\texttt{removerAlunoTurma}), a transação deve abranger: exclusão de \nolinkurl{Turma/\{id\}/Alunos/\{uid\}}, exclusão de \nolinkurl{Usuarios/\{uid\}/Turmas/\{turmaId\}}, gravação no histórico da turma, decremento do contador de vagas e gravação da auditoria.
\end{itemize}

\newpage
% =================================================
\section{Materialized (Views)}
% =================================================
\subsection{Tabela Materialized para Lote}
\begin{entidade}{Lote\_Materializado}
\campo{id}{SERIAL}{PK}
\campo{id\_lote}{INTEGER}{FK, UNIQUE}
\campo{qtd\_frascos\_cadastrados}{INTEGER}{DEFAULT 0, NOT NULL}
\campo{ultimo\_reconciliador}{TIMESTAMP}{NULL}
\end{entidade}
\begin{explicacao}
Corrigido o typo \texttt{SERAIL} $\rightarrow$ \texttt{SERIAL}. \textit{id\_lote} recebeu \texttt{UNIQUE}: esta é uma tabela 1:1 com \textit{Lote} (uma linha por lote), então o par (\textit{id}, \textit{id\_lote}) sem essa restrição permitiria duas linhas materializadas divergentes para o mesmo lote. Mantida \textbf{exclusivamente pela transação síncrona} da \textit{Cloud Function} a cada frasco cadastrado, sem depender de triggers assíncronas. Adicionada a marca d'água \textit{ultimo\_reconciliador} que é preenchida por \textit{jobs} agendados de reconciliação para bloquear o duplo cômputo de eventos atrasados (drift silenciado, ver seção \textit{Tecnologia e Relatórios}).
\end{explicacao}

\subsection{Tabela Resumo de Almoxarifado - Diário}
\begin{entidade}{Resumo\_Almoxarifado\_Diario}
\campo{id}{SERIAL}{PK}
\campo{data}{DATE}{NOT NULL}
\campo{id\_almoxarifado}{INTEGER}{FK, NOT NULL}
\campo{qtd\_frascos\_cadastrados\_fechados\_durante\_o\_dia}{INTEGER}{DEFAULT 0, NOT NULL}
\campo{qtd\_frascos\_cadastrados\_abertos\_durante\_o\_dia}{INTEGER}{DEFAULT 0, NOT NULL}
\campo{qtd\_frascos\_emprestados\_durante\_o\_dia}{INTEGER}{DEFAULT 0, NOT NULL}
\campo{qtd\_frascos\_devolvidos\_durante\_o\_dia}{INTEGER}{DEFAULT 0, NOT NULL}
\campo{volume\_utilizado\_no\_dia\_ml}{NUMERIC(12,2)}{DEFAULT 0, NOT NULL}
\campo{volume\_evaporado\_no\_dia\_ml}{NUMERIC(12,2)}{DEFAULT 0, NOT NULL}
\campo{massa\_utilizada\_no\_dia\_g}{NUMERIC(12,2)}{DEFAULT 0, NOT NULL}
\campo{massa\_evaporada\_no\_dia\_g}{NUMERIC(12,2)}{DEFAULT 0, NOT NULL}
\campo{calculado\_em}{TIMESTAMP}{NOT NULL -- serverTimestamp da execução que produziu a projeção atual}
\campo{versao\_calculo}{TEXT}{NOT NULL -- identificador da versão da regra/algoritmo de materialização}
\campo{created\_at}{TIMESTAMP}{NOT NULL}
\campo{updated\_at}{TIMESTAMP}{NOT NULL}
\end{entidade}
\begin{regra}
\textit{UNIQUE(data, id\_almoxarifado)}. No Firestore, o documento usa ID determinístico \texttt{YYYY-MM-DD\_\_ALMOXARIFADO\_ID}. Reprocessar a mesma chave diária substitui a projeção existente; nunca cria um novo documento.
\end{regra}
\newpage
\begin{explicacao}[Separação Física Rigorosa: Massa ($g$) vs Volume ($ml$)]
Esta tabela consolida exclusivamente fatos FLOW do dia civil, sem posição acumulada de estoque.
Para eliminar distorções métricas em sólidos com densidade desconhecida ou nula, as colunas agregadas de consumo e evaporação são formalmente segregadas nas duas dimensões físicas fundamentais: reagentes líquidos acumulam em \textit{volume\_...\_ml} e reagentes sólidos acumulam em \textit{massa\_...\_g}. Nunca se tenta converter gravimetria de sólidos para volume sem densidade aferida.

O campo legado que misturava unidades foi removido por redundância; consumo usa exclusivamente as métricas dimensionadas em g e mL.
\end{explicacao}

\begin{explicacao}[Projeção FLOW e Unicidade Determinística]
\textit{Resumo\_Almoxarifado\_Diario} é uma \textbf{projeção derivada, descartável, recalculável e reconstruível} a partir das fontes históricas canônicas (\textit{Historico\_Frasco\_Reagente}, \textit{Emprestimo\_Reagente} e eventos de domínio; ver Catálogo de Métricas). \textbf{Não é fonte primária de verdade.} Não se admite correção manual incremental de totais como autoridade: qualquer divergência deve ser resolvida por reprocessamento absoluto a partir das fontes.

\textbf{Unicidade e ID determinístico.} A chave lógica diária é \textit{UNIQUE(data, id\_almoxarifado)}. No Firestore o ID do documento é determinístico no formato \texttt{YYYY-MM-DD\_\_ALMOXARIFADO\_ID}. Reprocessar a mesma data/almoxarifado deve \textbf{substituir} (sobrescrever) a projeção existente, jamais criar um novo documento por reprocessamento.

\textbf{Metadados de cálculo.} \textit{data} identifica o dia civil agregado; \textit{calculado\_em} é o \texttt{serverTimestamp} da execução que produziu a projeção atual; \textit{versao\_calculo} identifica a versão da regra/algoritmo de materialização. \textit{calculado\_em} \textbf{não é a data do fato}: não se confunde com \textit{data} nem com o timestamp físico dos eventos agregados.

\textbf{Janela FLOW.} O dia civil D em America/Sao\_Paulo usa o intervalo semiaberto do início de D até o início de D+1. Somente fatos efetivos pertencentes a D são agregados. Não há snapshot diário de estoque, leitura de estado corrente para fabricar passado, dependência de D-1 ou replay até hoje. Correções fechadas reprocessam somente as datas diretamente afetadas (origem e destino, se houver mudança de atribuição).

\textbf{Correção de fato histórico.} Uma correção administrativa que altera um fato histórico usado na agregação deve identificar a \textit{DATA\_ALVO} do fato corrigido e disparar o reprocessamento do resumo de \textbf{cada data histórica afetada}. O timestamp da correção permanece o da correção; o timestamp físico original não é reescrito.
\end{explicacao}
\newpage
\subsection{Tabela Resumo de Reagente/Solvente - Diário}
\begin{entidade}{Resumo\_Reagente\_Diario}
\campo{id}{SERIAL}{PK}
\campo{data}{DATE}{NOT NULL}
\campo{id\_resumo\_reagente}{INTEGER}{FK, NOT NULL}
\campo{id\_almoxarifado}{INTEGER}{FK, NOT NULL}
\campo{qtd\_frascos\_descartados\_dia}{INTEGER}{DEFAULT 0, NOT NULL}
\campo{qtd\_frascos\_extraviados\_dia}{INTEGER}{DEFAULT 0, NOT NULL}
\campo{volume\_utilizado\_no\_dia\_ml}{NUMERIC(12,2)}{DEFAULT 0, NOT NULL}
\campo{volume\_evaporado\_no\_dia\_ml}{NUMERIC(12,2)}{DEFAULT 0, NOT NULL}
\campo{massa\_utilizada\_no\_dia\_g}{NUMERIC(12,2)}{DEFAULT 0, NOT NULL}
\campo{massa\_evaporada\_no\_dia\_g}{NUMERIC(12,2)}{DEFAULT 0, NOT NULL}
\campo{calculado\_em}{TIMESTAMP}{NOT NULL -- serverTimestamp da execução que produziu a projeção atual}
\campo{versao\_calculo}{TEXT}{NOT NULL -- identificador da versão da regra/algoritmo de materialização}
\campo{created\_at}{TIMESTAMP}{NOT NULL}
\campo{updated\_at}{TIMESTAMP}{NOT NULL}
\end{entidade}
\begin{regra}
\textit{UNIQUE(data, id\_almoxarifado, id\_resumo\_reagente)}. No Firestore, o documento usa ID determinístico \texttt{YYYY-MM-DD\_\_ALMOXARIFADO\_ID\_\_RESUMO\_REAGENTE\_ID}. Reprocessar a mesma chave diária substitui a projeção existente; nunca cria um novo documento.
\end{regra}

\begin{explicacao}[Contratos do Estoque Atual (fora dos resumos diários)]
As regras a seguir pertencem exclusivamente às agregações atuais protegidas no backend sobre Frasco\_Reagente, não aos resumos FLOW:
\begin{itemize}
    \item \textbf{Frascos Extraviados:} Frascos com estado \texttt{EXTRAVIADO} \textbf{nunca} são contabilizados no saldo disponível, seja na quantidade de frascos ou na massa/volume total.
    \item \textbf{Conteúdo Nominal vs. Saldo:} O \textit{conteudo\_nominal} (capacidade original do fabricante) \textbf{não representa} o saldo atual do frasco e não deve ser somado no lugar de saldos medidos.
    \item \textbf{Saldos Desconhecidos:} Se o saldo de um frasco é desconhecido (ex.: frasco aberto sem \textit{peso\_atual} preciso registrado), ele não deve ser somado como zero (\texttt{0}) nos agregados totais de volume ou massa, pois isso distorceria a realidade física. Esses frascos devem ser segregados na camada de agregação, distinguindo claramente a ``quantidade conhecida'' da ``quantidade com saldo desconhecido''.
\end{itemize}
\end{explicacao}

\begin{explicacao}[Projeção FLOW e Unicidade Determinística]
\textit{Resumo\_Reagente\_Diario} é uma \textbf{projeção derivada, descartável, recalculável e reconstruível} a partir das fontes históricas canônicas (\textit{Historico\_Frasco\_Reagente}, \textit{Emprestimo\_Reagente} e eventos de domínio). \textbf{Não é fonte primária de verdade.} Não se admite correção manual incremental de totais como autoridade: divergências são resolvidas por reprocessamento absoluto a partir das fontes.

\textbf{Unicidade e ID determinístico.} A chave lógica diária é \textit{UNIQUE(data, id\_almoxarifado, id\_resumo\_reagente)}. No Firestore o ID do documento é determinístico no formato \texttt{YYYY-MM-DD\_\_ALMOXARIFADO\_ID\_\_RESUMO\_REAGENTE\_ID}. Reprocessar a mesma chave deve \textbf{substituir} (sobrescrever) a projeção existente, jamais criar um novo documento por reprocessamento.

\textbf{Metadados de cálculo.} \textit{data} identifica o dia civil agregado; \textit{calculado\_em} é o \texttt{serverTimestamp} da execução que produziu a projeção atual; \textit{versao\_calculo} identifica a versão da regra/algoritmo de materialização. \textit{calculado\_em} \textbf{não é a data do fato}.

\textbf{Janela FLOW.} O dia civil D em America/Sao\_Paulo usa o intervalo semiaberto do início de D até o início de D+1. Somente fatos efetivos pertencentes a D são agregados. Não há snapshot diário de estoque, leitura de estado corrente para fabricar passado, dependência de D-1 ou replay até hoje. Correções fechadas reprocessam somente as datas diretamente afetadas (origem e destino, se houver mudança de atribuição).

\textbf{Correção de fato histórico.} Uma correção administrativa que altera um fato histórico agregado deve identificar a \textit{DATA\_ALVO} do fato corrigido e reprocessar o resumo de \textbf{cada data histórica afetada}, sem reescrever o timestamp físico original.
\end{explicacao}

\subsection{Tabela Resumo de Bens Patrimoniais - Diário}
\begin{entidade}{Resumo\_Bem\_Patrimonial\_Diario}
\campo{id}{SERIAL}{PK}
\campo{data}{DATE}{NOT NULL}
\campo{id\_local}{INTEGER}{FK, NULL}
\campo{qtd\_bens\_ativos\_no\_fim\_do\_dia}{INTEGER}{DEFAULT 0, NOT NULL}
\campo{qtd\_bens\_inservivel\_no\_fim\_do\_dia}{INTEGER}{DEFAULT 0, NOT NULL}
\campo{qtd\_bens\_dado\_baixa\_no\_fim\_do\_dia}{INTEGER}{DEFAULT 0, NOT NULL}
\campo{qtd\_bens\_cadastrados\_durante\_o\_dia}{INTEGER}{DEFAULT 0, NOT NULL}
\campo{qtd\_bens\_marcados\_inservivel\_durante\_o\_dia}{INTEGER}{DEFAULT 0, NOT NULL}
\campo{qtd\_bens\_dados\_baixa\_durante\_o\_dia}{INTEGER}{DEFAULT 0, NOT NULL}
\campo{qtd\_requisicoes\_abertas\_durante\_o\_dia}{INTEGER}{DEFAULT 0, NOT NULL}
\campo{qtd\_requisicoes\_aprovadas\_durante\_o\_dia}{INTEGER}{DEFAULT 0, NOT NULL}
\campo{qtd\_requisicoes\_rejeitadas\_durante\_o\_dia}{INTEGER}{DEFAULT 0, NOT NULL}
\campo{created\_at}{TIMESTAMP}{NOT NULL}
\campo{updated\_at}{TIMESTAMP}{NOT NULL}
\end{entidade}

\subsection{Tabela de Atividade Mensal - Gestor de Almoxarifado}
\begin{entidade}{Atividade\_Gestor\_Almoxarifado\_Mensal}
\campo{id}{SERIAL}{PK}
\campo{id\_gestor}{INTEGER}{FK, NOT NULL}
\campo{ano}{INTEGER}{NOT NULL}
\campo{mes}{INTEGER}{NOT NULL}
\campo{qtd\_reagentes\_adicionados}{INTEGER}{DEFAULT 0, NOT NULL}
\campo{qtd\_frascos\_adicionados}{INTEGER}{DEFAULT 0, NOT NULL}
\campo{qtd\_movimentacoes\_registradas}{INTEGER}{DEFAULT 0, NOT NULL}
\campo{created\_at}{TIMESTAMP}{NOT NULL}
\campo{updated\_at}{TIMESTAMP}{NOT NULL}
\end{entidade}
\begin{regra}
\textit{UNIQUE(id\_gestor, ano, mes)}.
\end{regra}

\subsection{Tabela de Atividade Mensal - Gestor de Bens Patrimoniais}
\begin{entidade}{Atividade\_Gestor\_Bens\_Patrimoniais\_Mensal}
\campo{id}{SERIAL}{PK}
\campo{id\_gestor}{INTEGER}{FK, NOT NULL}
\campo{ano}{INTEGER}{NOT NULL}
\campo{mes}{INTEGER}{NOT NULL}
\campo{qtd\_bens\_adicionados}{INTEGER}{DEFAULT 0, NOT NULL}
\campo{qtd\_bens\_editados}{INTEGER}{DEFAULT 0, NOT NULL}
\campo{qtd\_requisicoes\_aprovadas}{INTEGER}{DEFAULT 0, NOT NULL}
\campo{qtd\_requisicoes\_rejeitadas}{INTEGER}{DEFAULT 0, NOT NULL}
\campo{created\_at}{TIMESTAMP}{NOT NULL}
\campo{updated\_at}{TIMESTAMP}{NOT NULL}
\end{entidade}
\begin{regra}
\textit{UNIQUE(id\_gestor, ano, mes)}.
\end{regra}

\subsection{Catálogo de Métricas Disponíveis por Domínio}
\begin{itemize}
    \item \textbf{\textit{Historico\_Frasco\_Reagente}} (por \textit{tipo}): frascos cadastrados, saídas, entradas, esvaziados, quebrados, descartados, vencidos, ajustes de inventário e evaporação -- por dia $\rightarrow$ alimenta \textit{Resumo\_Almoxarifado\_Diario} e \textit{Resumo\_Reagente\_Diario}; por gestor/mês $\rightarrow$ alimenta \textit{Atividade\_Gestor\_Almoxarifado\_Mensal}; contagem "agora" $\rightarrow$ \texttt{count()} server-side via backend protegido sobre \textit{Frasco\_Reagente} filtrado por \textit{vencido = true}.
    \item \textbf{\textit{Emprestimo\_Reagente}}: empréstimos abertos/devolvidos/atrasados por dia, volume e massa de consumo validado por dia $\rightarrow$ alimenta os resumos diários de reagente.
    \item \textbf{\textit{Historico\_Bem\_Patrimonial} / \textit{Alteracao\_Bem\_Patrimonial}}: bens cadastrados, editados e baixados por dia $\rightarrow$ alimenta \textit{Resumo\_Bem\_Patrimonial\_Diario}.
    \item \textbf{\textit{Requisicao\_Edicao\_Bem\_Patrimonial} / \textit{Requisicao\_Adicao\_Bem\_Patrimonial}}: requisições abertas/aprovadas/rejeitadas por dia $\rightarrow$ alimenta \textit{Resumo\_Bem\_Patrimonial\_Diario}.
    \item \textbf{\textit{Registro\_de\_Auditoria}}: log genérico entre todas as entidades.
\end{itemize}

\subsection{Estoque atual: autoridade canônica e agregação protegida}
\label{sec:estoque-atual-cache}
\textit{Frasco\_Reagente} é a autoridade do estado corrente. Não existe materialized view persistente de estoque atual. O backend executa \texttt{count()} e \texttt{sum()} sobre esse documento canônico; o navegador não executa agregações Firestore nem baixa frascos para contá-los. O contador de vínculos de lote tem finalidade cadastral distinta e não substitui estoque atual.

Para permitir \texttt{sum()}, a projeção Firestore do próprio frasco possui \texttt{saldo\_aferido\_g} e \texttt{saldo\_aferido\_ml}, derivados server-owned. O backend calcula e atualiza esses campos na mesma transação das alterações de peso, tara ou conhecimento: saldo desconhecido ou metrologicamente inconsistente produz NULL; vazio explicitamente confirmado produz zero; caso consistente usa peso menos tara e, para líquidos, divisão pela densidade válida. A origem da tara permanece explícita. Sólidos têm campo mL NULL. As somas de estoque utilizável filtram aptidão operacional (estado FECHADO/ABERTO, disponível, sem quarentena e respeitando validade/autorização); contagens totais, emprestados e desconhecidos usam filtros próprios identificados. Nenhuma quantidade desconhecida é apresentada como zero conhecido.

\paragraph{Contrato de consulta e cache.}
A callable de dashboard exige App Check, Firebase Auth, RBAC e escopo antes de aplicar limite atômico de \textbf{5 solicitações/minuto por UID}, inclusive em cache hit. Usar janela deslizante de 60 segundos, relógio do servidor e registro interno por UID; IP não é identidade primária. Rejeição ocorre antes de executar agregações. O cache compartilhado nunca dispensa autorização individual.

\texttt{Sistema\_Cache\_Dashboard} é coleção interna, efêmera, derivada e descartável. Chaves permitidas: \texttt{ESTOQUE\_\_ALMOX} e \texttt{ESTOQUE\_\_ALMOX\_\_RESUMO}, com IDs codificados deterministicamente sem colisão. O endpoint aceita somente esses escopos e filtros enumerados; qualquer filtro que altere resultado entra na chave normalizada, sem parâmetros livres que criem cardinalidade ilimitada. Campos: \texttt{escopo}, \texttt{resultado}, \texttt{calculado\_em}, \texttt{versao\_calculo}, \texttt{geracao} e \texttt{valido}. Resultado é nullable durante invalidação. A geração é token de invalidação, não métrica de estoque.

Validade exige idade estritamente menor que 30 segundos, versão compatível e ausência de invalidação posterior. A marca temporal é capturada no início do cálculo (limite conservador da idade dos dados). Remoção física por TTL do Firestore é apenas limpeza, nunca mecanismo preciso de validade. Não existe scheduler de 30 segundos: cache miss recalcula sob demanda; oito horas sem acesso geram zero recálculos.

\paragraph{Invalidação e concorrência.}
Toda mutação que altere contribuição atual atualiza, atomicamente com o frasco, as gerações das chaves do almoxarifado e do par almoxarifado/resumo afetados, marca \texttt{valido=false} e limpa o resultado. Não calcula agregações nessa transação. A invalidação cria o registro de geração se necessário, evitando corrida com primeiro preenchimento. Transferência ou alteração de associação, somente quando existir contrato suportado, invalida origem e destino. A publicação do cache relê a geração e só publica se ela não mudou desde antes das consultas; descarta cálculo obsoleto e tenta novamente de forma limitada. Agregações concorrentes ocasionais são aceitáveis; não introduzir Redis, locks distribuídos ou infraestrutura dedicada. Se a geração mudar repetidamente, responder indisponibilidade temporária em vez de publicar dado invalidado. A remoção física de um registro faz a comparação falhar; nunca tratar ausência como geração válida capturada.

Invalidam cadastro, abertura, retirada, devolução normal ou anômala, resolução de anomalia, esgotamento, pesagem que muda saldo, evaporação, ajuste metrológico, recalibração, quebra, descarte, extravio, reencontro, entrada/saída de quarentena e decisões de validade/descarte que alterem elegibilidade. Atualizações de validade pelo job também invalidam. Correção exclusivamente histórica de retirante, gestor, motivo ou auditoria e reprocessamento FLOW não invalidam estoque atual nem alteram pesos correntes. O menor escopo seguro é sempre o das chaves efetivamente afetadas.

\begin{lstlisting}[language=TypeScript, title={Consulta lazy e publicação condicionada à geração (pseudocódigo)}]
// App Check aplicado pela callable antes do handler.
async function consultarEstoque(request) {
  const escopo = await autenticarAutorizarENormalizarEscopo(request);
  await consumirLimiteUID(request.auth.uid, 5, 60); // atomico; inclui cache hit
  const ref = chaveCache(escopo);
  for (let tentativa = 0; tentativa < 2; tentativa++) {
    const base = await obterOuCriarGeracao(ref);
    if (cacheValido(base, agoraServidor(), 30)) return base;
    const inicio = agoraServidor();
    const resultado = await agregarCountSumFrascos(escopo);
    const publicado = await db.runTransaction(async tx => {
      const atual = await tx.get(ref);
      if (!atual.exists || atual.data().geracao !== base.geracao) return false;
      if (agoraServidor() - inicio >= 30000) return false;
      tx.set(ref, { escopo, resultado, calculado_em: inicio,
        versao_calculo: "estoque-v1", geracao: base.geracao, valido: true });
      return true;
    });
    if (publicado) return { resultado, calculado_em: inicio };
  }
  throw new HttpsError("unavailable", "Estoque mudou; tente novamente.");
}
\end{lstlisting}

\paragraph{Interface e relatórios.}
A UI mantém cache curto, nunca além da validade restante informada pelo servidor, cancela requisições antigas ao mudar escopo e invalida seu cache após mutação. Exibe instante de cálculo, atualização em andamento e ação Atualizar com tratamento de rate limit; não consulta automaticamente a cada 30 segundos. Dados já exibidos podem ficar desatualizados por mutação de outro usuário, portanto nunca autorizam retirada. Estado atual é fotografia corrente separada. Relatórios por período usam FLOW e fatos históricos; não prometem séries de saldo, fechados ou emprestados ao fim de cada dia. Patrimônio mantém seu contrato próprio de snapshots diários.

% =================================================
\section{Requisitos e Regras de Negócio}
% =================================================

\subsection{Requisitos funcionais}


\setlength{\tabcolsep}{5pt}
\begin{longtable}{p{0.10\textwidth} p{0.82\textwidth}}
\toprule
\textbf{ID} & \textbf{Requisito}\\
\midrule
\endfirsthead
\toprule
\textbf{ID} & \textbf{Requisito} (continuação)\\
\midrule
\endhead
\bottomrule
\endfoot
RF01 & O sistema deve permitir login por e-mail e senha.\\
RF02 & O sistema deve permitir login com Google através do provedor de autenticação configurado.\\
RF03 & O sistema deve permitir recuperação de senha.\\
RF04 & O sistema deve permitir que um usuário possua múltiplos papéis.\\
RF05 & O sistema deve permitir alternância de papel ativo no dashboard quando houver mais de um papel autorizado.\\
RF06 & O sistema deve permitir cadastro e manutenção de bens patrimoniais conforme permissões.\\
RF07 & O sistema deve manter histórico auditável de alterações patrimoniais.\\
RF08 & O sistema deve permitir requisição de adição de patrimônio.\\
RF09 & O sistema deve permitir requisição de edição de patrimônio.\\
RF10 & O sistema deve impedir mais de uma requisição de edição pendente para o mesmo bem.\\
RF11 & O sistema deve permitir aprovação ou rejeição de requisições.\\
RF12 & O sistema deve permitir registro da baixa de um bem após o procedimento institucional.\\
RF13 & O sistema deve permitir cadastro de almoxarifados.\\
RF14 & O sistema deve calcular pesos e volumes de frascos de reagentes utilizando a densidade.\\
RF15 & O sistema deve registrar empréstimo, devolução, descarte, esvaziamento e quebra de frascos.\\
RF16 & O sistema deve permitir consulta por nome, número de patrimônio, status, local e demais filtros aplicáveis.\\
RF17 & O sistema deve permitir criação de turmas com capacidade máxima estabelecida.\\
RF18 & O sistema deve permitir entrada de alunos em turmas por código ou e-mail.\\
RF19 & O sistema deve permitir criação de posts em turmas pelo professor responsável.\\
RF20 & O sistema deve permitir comentários em posts conforme regras de acesso.\\
RF21 & O sistema deve permitir upload de roteiros em PDF.\\
RF22 & O sistema deve permitir compartilhamento de roteiros entre professores.\\
RF23 & O sistema deve permitir vínculo de roteiro a turma por meio de post.\\
RF24 & O sistema deve gerar relatórios mensais e por período, incluindo consumo em mL e massa em g.\\
RF25 & O sistema deve manter dados históricos e de auditoria sem apagar fatos relevantes.\\
\bottomrule
\end{longtable}

\subsection{Regras de Negócio Gerais}
\subsubsection{Status de um frasco de reagente}
\begin{regra}
Quando um frasco de reagente é cadastrado como em quarentena, deve ser obrigatório escrever o detalhe do status explicando o motivo dele estar em quarentena. Essa justificativa humana obrigatória deve possuir no mínimo 20 caracteres após trim; mensagens automáticas do servidor não são entradas humanas.
Um frasco em quarentena não pode ser descartado diretamente. A quarentena deve ser encerrada por decisão humana explícita: \texttt{VOLTAR\_A\_DISPONIVEL} ou \texttt{PENDENTE\_DE\_DESCARTE}. No segundo caso, a decisão gera autorização operacional estruturada, mantém o frasco \texttt{INDISPONIVEL} e somente então permite \texttt{descartarFrasco}, conforme a máquina de estados da Seção~4.
\end{regra}

% M9 -- fonte normativa única de autorização
\subsection{Autorização e usuários (M9)}
\label{sec:m9-autorizacao}

Esta subseção é a fonte normativa única de autorização. Se uma tabela de
stakeholder, tela, fluxo, pseudocódigo ou Rules divergir, esta subseção
prevalece. Autenticação responde quem apresentou uma credencial Firebase
válida; autorização responde \texttt{podeExecutar(uid, operacao, recurso,
estadoAtual)}. A primeira não implica a segunda.

\paragraph{Modelo e fontes de verdade.}
O modelo é RBAC + escopo/vínculo persistido + ownership quando o domínio o
exige. Os únicos papéis são Chefe\_Geral, Gestor\_Almoxarifado,
Gestor\_Bens\_Patrimoniais, Professor, Aluno e Bolsista. Papéis atuais são
os documentos de papel cujo \texttt{id\_usuario} coincide com o UID; Bolsista
implica Aluno; Chefe\_Geral é exclusivo; Professor e Aluno não coexistem; e
Bolsista e Gestor\_Almoxarifado não coexistem. Conta, papel, vínculo,
ownership, recurso ou escopo ausente, inválido, inconsistente ou desconhecido
significa negação. Nenhum payload, cache, espelho, campo denormalizado, papel
selecionado na UI ou claim antiga cria autorização.

\begin{itemize}
  \item \texttt{Usuarios/uid.ativo=true} é pré-condição universal. A ausência
  do documento ou \texttt{ativo=false} nega.
  \item A relação \texttt{Gestor\_Almoxarifado\_x\_Almoxarifado} é o escopo
  N:N canônico do gestor; um gestor pode ter zero ou muitos vínculos.
  Chefe\_Geral tem escopo global já decidido; Gestor\_Bens\_Patrimoniais tem
  escopo global do domínio patrimonial. Não há escopo patrimonial alternativo
  nesta V1.
  \item Turma pertence a \texttt{id\_professor}; matrícula canônica é
  \texttt{Turma/id/Alunos/uid}; o espelho em
  \texttt{Usuarios/uid/Turmas/id} é somente consulta. Roteiro pertence ao
  professor que o enviou; compartilhamento persistido autoriza apenas o uso
  previsto, não redistribuição.
  \item A decisão de domínio continua separada: RBAC não torna válida
  transição M5 proibida, não corrige fato metrológico M6, não substitui
  idempotência M7 e não contorna cache/estoque/notificação M8.
\end{itemize}

\paragraph{Matriz normativa de autorização.}
Todas as linhas exigem identidade autenticada e Usuário ativo. Tx exige leitura
da autorização e do recurso na mesma transação que produz a escrita. Rules
significa leitura direta de cliente apenas se confirmar versão atual e relação
persistida; backend é Cloud Function privilegiada com validação server-side.

\scriptsize
\setlength{\tabcolsep}{2pt}
\begin{longtable}{>{\raggedright\arraybackslash}p{0.15\textwidth}>{\raggedright\arraybackslash}p{0.18\textwidth}>{\raggedright\arraybackslash}p{0.19\textwidth}>{\raggedright\arraybackslash}p{0.16\textwidth}>{\raggedright\arraybackslash}p{0.17\textwidth}}
\toprule
Ator/papel & Operação/recurso & Escopo ou ownership & Decisão/origem & Resultado \\
\midrule
Chefe\_Geral & gerir usuário, papel, vínculo e almoxarifado & global; RN-ROLE & backend, Tx & permite ou nega invariantes \\
Chefe\_Geral & operar domínios existentes & global, sem dispensar domínio & backend, Tx em mutação & permite somente matriz \\
Gestor\_Almoxarifado & cadastro, retirada, devolução, M5/M6, relatório e etiqueta & vínculo atual com almox do recurso & backend, Tx & nega outro almox \\
Gestor\_Almoxarifado & catálogo global de reagentes & papel atual; sem escopo de almox & backend & permite ações de catálogo V1 \\
Gestor\_Bens\_Patrimoniais & bem, resposta, baixa e relatório patrimonial & domínio patrimonial global & backend, Tx em mutação & permite; professor não substitui gestor \\
Professor & turma, post, roteiro e requisição próprios & professor dono; roteiro próprio/compartilhado & backend, Tx em mutação & nega recurso alheio \\
Aluno & ingresso, leitura/comentário acadêmico & matrícula atual ou convite válido & backend/Rules & nega turma, post ou roteiro alheio \\
Bolsista & direitos de Aluno e consulta/recebimento de empréstimo & matrícula e próprio empréstimo & backend/Rules & não opera balcão/administração \\
Destinatário & notificação própria & UID do caminho = destinatário & Rules/backend & lê; marcar lida só servidor \\
Sistema & trigger, job, notificação, cache & IAM mínimo; entrada server-owned & backend & sem papel humano, com dedup M7 \\
\bottomrule
\end{longtable}
\normalsize

\paragraph{Leitura e escrita.}
Catálogo de reagentes é leitura interna para papéis permitidos na Seção~3; não
expõe retirantes, auditoria ou dados de terceiros. Empréstimo completo é
legível pelo retirante, Gestor\_Almoxarifado vinculado e Chefe; bem, requisição
e histórico seguem a matriz. Consultas que poderiam retornar documento não
autorizado são negadas, pois Rules não são filtros. Escritas de domínio,
histórico, auditoria, papéis, vínculos, \texttt{Operacoes}, locks, estoques
configurados, cache, rate limit e notificações são server-owned. A única
mutação do destinatário é pedir ao backend que marque a notificação lida;
Limpar tudo não apaga documentos.

\paragraph{Claims, revogação e sessão.}
Custom Claims são projeção assinada de \texttt{roles} e
\texttt{versao\_permissoes}, gravada apenas pelo Admin SDK após o commit. O
Firebase propaga claim alterada somente para ID token novo; ID tokens têm vida
curta de cerca de uma hora. Revogar refresh token e pedir refresh melhoram a
sessão, mas não tornam claim antiga fonte autoritativa. Para leitura direta,
Rules exigem token com versão igual, Usuário ativo e relação persistida
aplicável. Para mutação crítica, backend lê documentos atuais dentro da
transação; não confia somente em claims, inclusive de Chefe\_Geral. Alteração
de papel, desativação ou remoção de vínculo incrementa versão na transação e
nega a próxima tentativa que consulte estado persistido.

\paragraph{TOCTOU, transação e idempotência.}
Para operação crítica, a transação lê Usuário, papéis, vínculo/ownership,
recurso, estado de domínio e registro M7 antes de escrever. Se revogação ou
vínculo mudar depois de uma leitura, o Firestore reexecuta a transação e a
decisão é refeita; commit posterior à revogação é negado. Verificação antes da
transação é só UX. Retry M7 concluído devolve receipt existente ao mesmo UID
sem reaplicar efeito; não permite novo efeito após perda de papel. Etapa
pendente ou efeito externo revalida antes de persistência privilegiada.

\paragraph{Security Rules, Admin SDK e UI.}
Rules protegem SDKs cliente; Admin SDK as ignora e deve aplicar esta matriz no
backend. Rules não substituem validação de operação, transação, estado de
domínio ou IAM. UI pode ocultar ação e pedir refresh, mas é UX. Collections
internas e server-owned permanecem deny-all ao cliente. Storage exige a mesma
decisão atual de papel/ownership/escopo, além de caminho, tipo, tamanho e
objeto validado.

\paragraph{Casos de regressão M9.}
Cobrir: não autenticado; sem papel; papel desconhecido; Usuário inativo;
gestor no almoxarifado correto e em outro; vínculo removido; claim antiga após
revogação; troca de papel; Professor em recurso próprio e alheio; Aluno em
turma permitida e em ação administrativa; Chefe nos limites da matriz; escrita
direta server-owned; Admin SDK sem validação; operação iniciada antes e commit
depois de revogação; retry M7 após mudança autorizativa; cache hit M8 sem
autorização; notificação de terceiro e alteração cliente-side; e tentativas de
contornar quarentena M5 ou pendência metrológica M6. Cada caso resulta em
negação ou no efeito específico da matriz, sem autorização implícita.
\subsection{Patrimônio (M10)}
\label{sec:regras-patrimonio-m10}

\begin{regra}[Autoridade, identidade e mutações patrimoniais]
\textit{Bem\_Patrimonial} é a unidade física; \textit{Resumo\_Bem\_Patrimonial}
é somente o modelo catalográfico compartilhado. O nome canônico pertence ao
resumo. \textit{nome\_equipamento} no bem e \textit{predio}/\textit{andar}/
\textit{sala} são projeções de consulta, jamais autoridades concorrentes. Uma
alteração individual de nome é reclassificação para resumo existente ou criação
de resumo novo; não renomeia o resumo de outros bens. Foto pertence apenas à
unidade física.

Toda mutação usa a ordem M9 (Auth, usuário ativo, papel/versão persistidos,
escopo global patrimonial V1 e revalidação dentro da transação) e M7
(\texttt{idOperacao}, \texttt{tipo\_operacao}, payload canônico e receipt). A
autorização é necessária, mas não dispensa unicidade, lock, versão, estado,
arquivo validado ou histórico. Gestor de Bens e Chefe podem atuar no domínio
global; Professor cria apenas requisições permitidas, nunca aprova a própria
requisição, não escreve Bem diretamente e não realiza baixa. UI e Rules são
camadas de UX/leitura; a decisão e a escrita crítica são server-owned.
\end{regra}

\begin{regra}[Plaqueta, locks e requisições]
Use \texttt{N(s)=s.trim().toUpperCase()} como única canonicalização de plaqueta.
O valor vazio ou maior que 30 caracteres após N é inválido. Bem, requisição,
lock \texttt{bem\_adicao\_\{N\}}, chave permanente
\texttt{Chaves\_Unicas/Bem\_Patrimonial\_\_\{N\}} e busca exata usam N; a
chave permanente é criada com a aprovação e jamais é liberada pela baixa.
Lock de edição é \texttt{bem\_edicao\_\{idBem\}}. Ele é criado junto da
requisição, não expira por idade e só é liberado pelo desfecho terminal que
confirma seu proprietário. Lock limita pendência; chave única protege existência
permanente; versão detecta edição concorrente; idempotência reapresenta o mesmo
resultado ao mesmo comando. Não são substitutos entre si.

Requisição de edição captura \texttt{versao\_bem\_origem}; aprovação relê o Bem
na transação e divergência termina a requisição como rejeitada por conflito,
liberando apenas seu lock. Referência de Local ou Resumo que deixou de existir é
rejeição sistêmica terminal, também com liberação vinculada. Falha transitória de
infraestrutura aborta sem transformar a requisição em rejeitada. Aprovação de
adição relê requisição, lock, chave única, Local, Resumo e foto antes de gravar,
no mesmo commit, Bem, chave, eventual resumo, resposta, histórico de cadastro e
liberação do lock.
\end{regra}

\begin{regra}[Estado, baixa, arquivo e histórico]
Conservação BOM/REGULAR/RUIM não deriva status. O ciclo V1 é
\texttt{Ativo -> Inservivel -> Ja\_dado\_baixa}; não há reversão nem salto
direto para baixa. Baixa é operação própria do Gestor/Chefe após rito SEI: exige
Bem ainda Inservivel, justificativa, PDF existente no caminho autorizado,
proprietário/origem, tamanho, assinatura \texttt{\%PDF-} e geração imutável
verificados pelo backend antes da transação. A transação grava URL/geração do
comprovante, status terminal, versão incrementada, evento \texttt{baixa} e
receipt M7. Antes da baixa, a URL é nula; depois, não pode ser reutilizada,
substituída ou apagada silenciosamente.

\texttt{Bem\_Patrimonial/\{id\}/Historico\_Patrimonio/\{eventoId\}} é a única
arquitetura normativa de histórico; coleção raiz histórica é legado divergente.
Eventos são \texttt{cadastro}, \texttt{edicao} e \texttt{baixa}; edição usa
deltas \texttt{campo, valor\_anterior, valor\_novo} do estado efetivamente
substituído. Bem e evento são atômicos; eventos usam ID determinístico da
operação/evento para retry não duplicar fatos. Snapshots de Local no evento são
imutáveis e fan-out derivado não cria histórico de negócio nem incrementa versão.
\end{regra}

\paragraph{Casos de regressão M10.}
Cobrir: criação válida, plaqueta existente e variantes de caixa/espaço, duas
adições concorrentes, resumo/local/foto ausentes ou inválidos; edição, resposta
repetida, conflito, reclassificação, local e foto; proprietário de lock errado,
rejeição/conflito liberando lock e idade sem expiração; cada transição de status;
baixa sem PDF/PDF inválido/Professor/duplicada; histórico de cadastro, delta e
baixa; retry M7, revogação M9 antes do commit, claim antiga e escrita direta de
cliente. Esses casos são critérios normativos e não atestam a implementação.

% M11 -- fonte normativa de Turmas, matrícula e convites
\subsection{Turmas, matrícula e convites (M11)}
\label{sec:regras-turmas-m11}

Esta subseção é a fonte normativa de Turma, vínculo Aluno--Turma, matrícula e
convite de ingresso (RF17, RF18, RN-TUR-01 e UI-10). Posts, comentários,
roteiros e notificações acadêmicas entram apenas nas fronteiras afetadas por
essas operações; a formalização completa desses domínios permanece fora de
M11. A autorização continua regida por M9 (Seção~\ref{sec:m9-autorizacao}), a
idempotência por M7 e a unicidade por \texttt{Chaves\_Unicas}; nada nesta
subseção cria bypass desses contratos.

\begin{regra}[Identidade e ciclo de vida da Turma]
\textit{Turma} é a entidade raiz \texttt{Turma/\{id\}}, com \texttt{id} gerado
pelo servidor. \texttt{id\_professor} é o professor dono e \texttt{id\_materia}
refere uma matéria existente, verificada no servidor, que não é recriada.
\texttt{nome\_turma} tem no máximo 100 caracteres; \texttt{ano} é inteiro;
\texttt{semestre} é exclusivamente 1 ou 2; \texttt{capacidade} é inteiro
positivo; \texttt{codigo\_turma} é único, com até 20 caracteres, gerado
exclusivamente pelo servidor; \texttt{status} pertence a \{Ativo, Arquivada\};
\texttt{data\_criacao}, \texttt{versao} (inteiro positivo) e \texttt{qtd\_alunos}
(inteiro não negativo) completam o registro.

A criação ordinária é feita pelo professor autenticado para si
(\texttt{id\_professor} = UID). Criar turma em nome de outro professor é
intervenção excepcional M9/Q13, auditada, com o dono registrado explicitamente;
a autoria não é transferida a quem não é o professor dono. Só o professor dono,
em regra, lê e altera a turma; o Chefe lê e intervém apenas nos limites de Q13.
Aluno de turma alheia não lê a turma.

Arquivar e desarquivar alteram \texttt{status} e incrementam \texttt{versao},
preservando \texttt{id}, membros, posts, comentários, histórico,
\texttt{qtd\_alunos} e \texttt{codigo\_turma}; arquivar não libera o código para
reuso. Arquivar com alunos matriculados é permitido e não é desligamento:
todos os vínculos canônicos permanecem. Turma arquivada é somente leitura
(Q08): o backend e as Rules rejeitam novo ingresso, novo convite de turma, novo
post e novo comentário, sem apagar dados. A inclusão ou remoção de membro
altera \texttt{qtd\_alunos}, vínculo e histórico, mas não é edição dos metadados
da turma e não incrementa \texttt{versao} por si; somente criação,
arquivamento/desarquivamento e edição de metadados o fazem.
\end{regra}

\begin{regra}[Matrícula canônica, espelho e contador]
\texttt{Turma/\{id\}/Alunos/\{uid\}} é o vínculo canônico e a única autoridade
de matrícula, presença e acesso acadêmico. \texttt{Usuarios/\{uid\}/Turmas/\{id\}}
é projeção de consulta que alimenta ``Minhas Turmas'' em tempo real; ela nunca
autoriza e nunca substitui o vínculo canônico (M9).

\texttt{qtd\_alunos} é o contador transacional do vínculo ativo. A leitura 3FN
\texttt{COUNT(Aluno\_x\_Turma WHERE id\_turma = X)} e \texttt{qtd\_alunos}
representam a mesma grandeza; na V1 a autoridade de vagas é
\texttt{qtd\_alunos}, lido e escrito na mesma transação que cria ou remove o
vínculo. Consultar a subcoleção \texttt{Alunos} ou recontar vínculos fora da
transação não é autoridade de vaga, pelo custo O(N) e pela concorrência.
Divergência entre \texttt{qtd\_alunos} e o número de vínculos canônicos é
inconsistência a reconciliar, nunca vaga adicional.

Ingresso, remoção e restauração de acesso atualizam, numa única transação, o
vínculo canônico, o espelho, \texttt{qtd\_alunos}, o evento em
\texttt{Historico\_Alunos\_Turma}, a auditoria e o receipt M7. Evento,
notificação ou espelho não criam matrícula por si. Nome e foto vistos por
colegas são projeções mínimas, sem matrícula nem e-mail.

Remover aluno: o professor dono desvincula (Chefe apenas por Q13). A transação
exclui o vínculo canônico e o espelho, grava \texttt{exclusao\_aluno} com
\texttt{removido\_por} e decrementa \texttt{qtd\_alunos} somente se o vínculo
existia, nunca abaixo de zero. Posts, comentários e autoria histórica são
preservados. Após remoção, o ingresso por código é bloqueado; o reingresso
exige convite explícito. Leitura de dado protegido após a remoção consulta o
vínculo canônico atual; sessão, claim ou cache antigos não restauram acesso.

Repetição (M7): reingressar quando já matriculado devolve o resultado
persistido e o acesso à turma, sem duplicar vínculo nem contagem; remover duas
vezes é idempotente e a segunda tentativa não gera decremento extra. O perdedor
de uma corrida recebe erro recuperável, nunca matrícula parcial.
\end{regra}

\begin{regra}[Capacidade, concorrência e exceção nominal]
O ingresso ordinário, por código ou por convite nominal sem exceção, exige no
commit da transação \texttt{qtd\_alunos < capacidade}. Na corrida pela última
vaga, o Firestore reexecuta a transação, um único commit cria o vínculo e o
perdedor recebe \texttt{failed-precondition} sem matrícula parcial.

Somente convite nominal válido, emitido pelo professor dono, com
\texttt{exceder\_capacidade = true} e \texttt{justificativa\_excecao}
persistida, pode ultrapassar a capacidade; a exceção é revalidada na aceitação.
Ela é específica daquele convite e daquela identidade, não é estado global da
turma nem passe livre por posse do código. Convite sem exceção continua sujeito
a vaga, e o código de turma nunca autoriza exceção.

Redução de capacidade: é proibido reduzir \texttt{capacidade} para valor menor
que \texttt{qtd\_alunos}; o servidor rejeita, fail-closed, com erro recuperável
e orienta remover membros antes ou manter a capacidade. Assim não há expulsão
implícita de alunos nem vaga negativa, e a capacidade nunca fica abaixo da
ocupação. \texttt{qtd\_alunos} nunca é negativo e pode ser reconciliado com os
vínculos canônicos.
\end{regra}

\begin{regra}[Convites: token, unicidade e fronteira externa]
\texttt{Convite\_Aluno/\{id\}} é coleção raiz; \texttt{id\_turma} \texttt{NULL}
identifica convite global. O e-mail é armazenado normalizado; o token aleatório
de 32 bytes (CSPRNG) é guardado apenas como hash SHA-256, comparado em tempo
constante, com expiração de sete dias, status \texttt{pendente}, \texttt{aceitado}
ou \texttt{expirado}, \texttt{convidado\_por}, \texttt{convidado\_em},
\texttt{numero\_matricula} opcional, \texttt{exceder\_capacidade} e
\texttt{justificativa\_excecao}, \texttt{aceitado\_por} e \texttt{aceitado\_em}.
O docId é determinístico, derivado no servidor por HMAC-SHA-256 do contexto
(turma ou global) e do e-mail normalizado com segredo do servidor; nunca se usa
hash simples e previsível de e-mail como chave nem se expõe o e-mail em claro.

Unicidade: no máximo um convite \texttt{pendente} por
(\texttt{email} normalizado, \texttt{id\_turma}); para convite global, no
máximo um pendente por e-mail normalizado com \texttt{id\_turma} \texttt{NULL}.
A chave determinística em \texttt{Chaves\_Unicas} distingue as duas situações e
é validada na transação. Convite aceito ou expirado não ocupa a pendência.
Reenvio substitui hash e expiração no mesmo documento e invalida o link
anterior. Convite de turma e convite global são distintos: aceitar convite de
turma cria vínculo, espelho e contador; aceitar convite global não cria
matrícula nem contador, apenas garante \texttt{Usuarios} e o papel Aluno, e o
aluno ingressa depois por código ou por outro convite. Matrícula é obrigatória
e única na identidade final \texttt{Aluno}, opcional e não única no convite.

Expiração operacional: \texttt{expira\_em} vencido torna o convite inaceitável
e ele deixa de contar como pendente para a unicidade; job ou reconciliação
idempotente materializa \texttt{expirado}. Toda pendência termina por aceitação,
expiração ou reenvio, nunca fica eterna.

Fronteira Auth, Firestore e e-mail: a aceitação exige sessão com e-mail
verificado correspondente. Criação condicional de \texttt{Usuario},
\texttt{Aluno}, papel e vínculo e o consumo do convite ocorrem na mesma
transação Firestore; a criação de conta Firebase Auth quando inexistente e o
envio de e-mail são etapas externas, pós-commit, recuperáveis e idempotentes.
Falha parcial não pode criar conta, papel, vínculo ou espelho divergente nem
consumir o convite sem os efeitos correspondentes. Criar o registro do convite
não é enviar o e-mail.
\end{regra}

\begin{regra}[Aceitação, reingresso e composição M7/M9]
Aceitar convite valida, na mesma transação: convite pendente não expirado,
e-mail verificado, identidade, status Ativo da turma, vaga ou exceção válida e
ausência de vínculo ativo. A aceitação é única e idempotente (M7): repetir a
aceitação devolve o resultado persistido, sem recriar vínculo nem incrementar
\texttt{qtd\_alunos}. Convite consumido, expirado, de e-mail divergente ou de
turma arquivada resulta em negação específica e recuperável. Aluno removido
anteriormente só reingressa por convite explícito, nunca pelo código.

Toda mutação M11 lê autorização e recurso na mesma transação (M9): professor
dono, conta ativa e versão de permissões, com TOCTOU coberto pela reexecução da
transação; operação iniciada antes de revogação não grava depois dela. UI-10 e
a Seção~\ref{sec:fluxos-completos} são UX: botões, código e navegação não
autorizam. O frontend não afirma envio de e-mail quando apenas o registro foi
criado e distingue carregamento, sucesso, vazio, erro e negação.
\end{regra}

\paragraph{Casos de regressão M11.}
Cobrir: corrida pela última vaga; dois convites simultâneos para o mesmo
(e-mail, turma) e para o global; replay de aceite; remoção concorrente e dupla;
arquivamento com ocupação, desarquivamento e escrita em turma arquivada; convite
global sem matrícula; ingresso de removido por código (negado) e por convite
(exceção válida); exceção de capacidade com e sem justificativa; redução de
capacidade abaixo da ocupação; contador nunca negativo; espelho ou cache
desatualizado não autoriza; retry M7 e revogação M9 antes do commit. Esses
casos são critérios normativos e não atestam a implementação.

\subsubsection{Seleção do tipo em Resumo de Reagente}
\begin{regra}
Ao cadastrar uma \textit{Especificacao\_Reagente}, o estado físico do resumo (\texttt{SOLIDO} ou \texttt{LIQUIDO}) determina quais campos operacionais de frasco serão apresentados. O tipo de substância (\texttt{PURA} ou \texttt{MISTURA}) determina as propriedades químicas obrigatórias, conforme a matriz de obrigatoriedade. DP-A01: estado físico e higroscopicidade pertencem ao resumo; unidade é derivada e densidade permanece na especificação. \\
\end{regra}
\subsubsection{Excluindo um Aluno de uma turma}
\begin{regra}
Quando o professor exclui um aluno de uma turma, ele não apga o aluno d sistema. Apenas o retira da turma. Seu histórico nos posts é mantido. Ou seja, o alunos que ainda estão na turma conseguem ver os comentários dele e também seu nome e foto de perfil nos comentários dele.
\end{regra}
\subsubsection{Chefe Geral}
\begin{regra}
O sistema pode ter mais de um chefe geral, isso não altera o sistema.
\end{regra}
\subsubsection{O responsável pelo bem patrimonial não é um Usuário do sistema LCQUI}
\begin{regra}
Por iss que "\textit{nome\_responsavel\_sei}" é um simples texto.
\end{regra}
\subsubsection{Historico de Posts na Turma}
\begin{regra}
A entidade \textit{Historico\_Posts\_Turma} registra no campo \textit{editado\_por} o identificador do usuário que realizou a edição. Normalmente é o Professor dono da turma, mas o Chefe Geral também pode editar ou remover posts em caráter excepcional de moderação institucional (Q13); nesses casos o campo \textit{editado\_por} identifica a ação da chefia e o campo \textit{acao} em \textit{Registro\_de\_Auditoria} registra a motivação. O \textit{id\_turma} não é armazenado nesta entidade 3FN porque pode ser obtido por \textit{Historico\_Posts\_Turma.id\_post} $\rightarrow$ \textit{Post.id\_turma}; se a consulta direta por turma for necessária no Firestore, esse vínculo poderá ser denormalizado na coleção de histórico.
\end{regra}
\subsubsection{Não compartilhar o mesmo roteiro 2 vezes com um professor}
\begin{regra}
Na tabela: \textit{Roteiro Professor Compartilhado} \\
UNIQUE(id\_roteiro\_experimento, id\_professor)
\end{regra}
\subsubsection{Densidade e concentração na especificação do reagente}
\begin{regra}
A densidade é uma propriedade da \textit{Especificacao\_Reagente}, e não do \textit{Resumo\_Reagente}. Ela deve ser informada quando obrigatória para a especificação e, após a confirmação cadastral, não pode ser alterada diretamente, pois participa dos cálculos de volume e consumo dos frascos. A concentração pertence à composição da especificação. Duas especificações com o mesmo nome comercial podem ter densidade, concentração, fabricante ou código de produto diferentes e, portanto, são especificações distintas sob o mesmo resumo quando ainda representam o mesmo item químico agrupador; não se deve criar um novo \textit{Resumo\_Reagente} apenas por haver diferença de concentração ou densidade.
\end{regra}
\subsubsection{Regra para código do frasco de Reagente}
\begin{regra}
O código do frasco de reagente (formato \texttt{LCQUI-N}) é gerado e garantido exclusivamente no momento do cadastro do frasco via transação atômica. A impressão de etiquetas virgens prévias fornece identificadores sugeridos para colar fisicamente, mas não cria frascos no sistema nem garante a reserva do código.
\end{regra}

\subsubsection{Segurança na Reimpressão de Frascos Cadastrados}
\begin{regra}
Para evitar a clonagem acidental de etiquetas na bancada e risco de trocas químicas, a reimpressão de segunda via é limitada a no máximo 10 frascos por sessão. O documento impresso deve conter a Ficha de Conferência e Rastreabilidade completa e 1 única etiqueta de reposição por frasco.
\end{regra}
\subsubsection{Regra em Notificação para Professor}
\begin{regra}
O campo \textit{id\_turma} é obrigatório para todas as notificações de contexto acadêmico (\texttt{COMENTARIO}, \texttt{POST}, \texttt{ADICIONADO}, \texttt{REMOVIDO}, \texttt{TURMA\_ARQUIVADA}, \texttt{TURMA\_DESARQUIVADA}) e deve ser mantido estritamente \texttt{NULL} para notificações operacionais e de almoxarifado/patrimônio. A notificação \texttt{DATA\_DEVOLUCAO\_REAGENTE} é preventiva: deve ser emitida pelo backend quando um empréstimo ainda está \texttt{EM\_USO} e sua data de devolução é hoje ou amanhã. A detecção é responsabilidade da Cloud Function; o frontend apenas apresenta e pode agrupar visualmente os registros em ``atrasados'', ``vence hoje'' e ``vence amanhã''. O frontend não decide se uma notificação deve ser criada.

O job pode usar uma única consulta por \texttt{status = EM\_USO} e \texttt{data\_devolucao\_prevista <= fim\_de\_amanha}, mas deve classificar cada empréstimo no servidor antes de criar a notificação, separando a janela de atraso da janela preventiva e aplicando idempotência. A idempotência é garantida por \textbf{docId determinístico} composto por \texttt{\{id\_emprestimo\}--\{janela\}}, onde \texttt{janela} é \texttt{VENCE\_HOJE} ou \texttt{VENCE\_AMANHA}; a Cloud Function usa \texttt{set(..., \{merge: false\})} apenas se o documento ainda não existir, evitando duplicatas em execuções repetidas do job diário. Um empréstimo que muda de janela (amanhã $\rightarrow$ hoje) cria um segundo documento; o anterior não é removido, apenas expirado.

Deve aparecer na interface:
\begin{itemize}
    \item COMENTÁRIO: x novos comentários de Fulano de Tal na turma y;
    \item REQUISICAO\_BEM: resposta de aprovação ou rejeição enviada ao Professor solicitante;
    \item REQUISICAO\_ADICAO\_BEM: nova solicitação de adição enviada aos Gestores de Bens Patrimoniais;
    \item REQUISICAO\_EDICAO\_BEM: nova solicitação de edição enviada aos Gestores de Bens Patrimoniais;
    \item DATA\_DEVOLUCAO\_REAGENTE: a data de devolução de um ou mais frasco(s) está expirando [amanhã, hoje];
    \item ROTEIRO\_COMPARTILHADO: Você tem x novo(s) roteiro(s) compartilhado(s) com você.
\end{itemize}
\end{regra}

\subsubsection{Regra para convidar um aluno que não está no sistema}
\begin{regra}
Quando o aluno é convidado sem especificar uma turma, id\_turma é NULL em \textit{Convite para Aluno}. Quando o professor vai convidar o aluno, ele pode ou não especificar uma turma, mas ao enviar um convite, tem que aparecer uma confirmação no caso de estar convidando um aluno sem uma turma para ele.
\end{regra}
\subsubsection{Regra para registrar um patrimônio como dado baixa}
\begin{regra}
É preciso fazer upload do documento pdf que comprove que o bem patrimonial foi dado baixa no sistema da UENF
\end{regra}
\subsubsection{Cálculo de Volume e Peso de Reagentes (Padrão de Laboratório)}
\begin{regra}[Controle de Volume via Pesagem]
O sistema adotará a pesagem como método principal de controle de volume dos frascos, utilizando a fórmula $Volume = \frac{Peso_{total} - Peso_{vazio}}{Densidade}$, em que $Peso_{vazio}$ é a tara real do recipiente, obtida pela medição do recipiente efetivamente vazio.

\textbf{1. Cadastro de Frasco Fechado:}\\
O usuário informa o Peso Total (frasco + reagente) medido na balança e o Conteúdo Nominal do rótulo (volume em mL para líquidos, massa em g para sólidos). O sistema deriva a tara teórica do recipiente a partir do peso total e do conteúdo nominal, quando disponível. Essa tara é apenas uma referência derivada enquanto o frasco contém produto e não é medição física definitiva do recipiente vazio:\\
\begin{itemize}
    \item Reagente Líquido ($mL$): $Peso_{vazio} = Peso_{total} - (Volume_{nominal} \times Densidade)$
    \item Reagente Sólido ($g$): $Peso_{vazio} = Peso_{total} - Massa_{nominal}$
\end{itemize}
O frasco inicia com \texttt{condicao\_inicial\_cadastro = FECHADO}, \texttt{estado\_fisico\_frasco = FECHADO} e saldo conhecido somente com dados quantitativos suficientes. Um cadastro FECHADO admite conteúdo nominal NULL quando desconhecido ou ilegível, nunca zero. Nesse caso, peso\_frasco\_vazio e origem\_tara são NULL e saldo\_desconhecido é true. Abertura ou pesagem bruta isolada não resolvem esse desconhecimento. Sem tara real, segue o mesmo ciclo de vida de JA\_ABERTO em situação metrológica equivalente: não exigir tara fabricada para quebra, descarte, extravio ou reencontro. EXTRAVIADO preserva a flag; \texttt{VAZIO} e \texttt{QUEBRADO} são estados fisicamente não utilizáveis e indisponíveis, normalmente encaminhados para descarte, mas \textbf{não} são estados terminais da máquina (\texttt{DESCARTADO} é o estado terminal); nada disso inventa peso ou quantidade química.

\textbf{2. Cadastro de Frasco Já Aberto:}\\
O frasco inicia com \texttt{condicao\_inicial\_cadastro = JA\_ABERTO}, \texttt{estado\_fisico\_frasco = ABERTO}, \texttt{peso\_frasco\_vazio = NULL} e \texttt{saldo\_desconhecido = true}. Não existem as modalidades de tara conhecida, estimativa de volume ou estimativa de massa, nem campos de peso do frasco vazio informado, volume atual estimado ou massa atual estimada. O usuário informa apenas dados reais disponíveis: o peso total atual medido na balança e, quando conhecido ou legível, o conteúdo nominal do rótulo. O conteúdo nominal é declaração original do fabricante e não representa o saldo atual; quando desconhecido ou ilegível, permanece \texttt{NULL}, nunca zero. Enquanto o recipiente não for esvaziado e pesado vazio, o saldo atual permanece desconhecido e a tara do recipiente ainda não é conhecida; não se estima saldo inicial.

\textbf{3. Empréstimo e Devolução:}\\
No empréstimo, registra-se o \textit{peso\_saida}. Na devolução, registra-se o \textit{peso\_retorno}. Quando quantitativamente validado, o consumo é documentado na tabela \textit{Emprestimo\_Reagente} como: $Volume_{utilizado} = \frac{\max(0,\ (Peso_{saida} - Peso_{retorno}^{efetivo}) - Perda_{evaporacao})}{Densidade}$ (líquidos); sólidos usam a diferença não negativa em g. A perda por evaporação é contabilizada separadamente do consumo (campo e evento \texttt{EVAPORACAO}); ganho tolerado segue Q06, com consumo zero e ajuste separado.

\textbf{4. Esgotamento e Tara Real:}\\
A determinação de que o recipiente ficou vazio é confirmada explicitamente pelo gestor na devolução, tanto para frascos originalmente \texttt{FECHADO} quanto para frascos \texttt{JA\_ABERTO}, por meio da decisão ``O frasco ficou vazio nesta devolução? [ ] Sim''. A comparação com a tara anterior não é autoridade exclusiva para declarar esgotamento.

Sem tara real anterior, ao confirmar o esgotamento, o peso aferido na devolução (\textit{peso\_retorno}) passa a ser a \textbf{tara real} do recipiente: \texttt{peso\_frasco\_vazio = peso\_retorno}, \texttt{peso\_atual = peso\_retorno}, \texttt{estado\_fisico\_frasco = VAZIO}, \texttt{disponibilidade = INDISPONIVEL} e \texttt{saldo\_desconhecido = false}. Não se altera o \texttt{conteudo\_nominal} e registra-se o evento \texttt{FICOU\_VAZIO}. Para frasco \texttt{JA\_ABERTO}, \textit{peso\_retorno} é a primeira tara real conhecida. Não existe recalibração parcial por estimativa: recalibrar tara significa recipiente efetivamente vazio com peso vazio real medido e justificativa humana conforme a lista fechada; não existe correção administrativa metrológica na V1.

\textbf{Tolerância de Devolução (Regra Q06 - PDF-025):}\\
A margem de tolerância é estritamente baseada no \textbf{peso bruto de saída} ($Peso_{saida}$). Rejeita-se explicitamente o cálculo sobre massa líquida, pois a variância incide sobre a totalidade do corpo medido na balança (incluindo o frasco).
\begin{itemize}
    \item Reagente normal: $\max(1\,g, 0{,}005 \times Peso_{saida})$
    \item Reagente higroscópico: $\max(2\,g, 0{,}02 \times Peso_{saida})$
\end{itemize}
A avaliação de ganho máximo permitido segue estritamente a fórmula: $Peso_{retorno} > Peso_{saida} + Tolerancia$. 
Um ganho de massa acima dessa tolerância \textbf{não aborta a operação via rollback}. O sistema conclui a devolução em modalidade extraordinária: o empréstimo transiciona para \texttt{DEVOLVIDO\_COM\_ANOMALIA}, e o frasco transiciona compulsoriamente para \texttt{disponibilidade = INDISPONIVEL} e \texttt{em\_quarentena = true}, com \texttt{detalhe\_status} registrando suspeita de contaminação/adulteração. É terminantemente proibido registrar evento de higroscopia (\texttt{ganho\_massa\_higroscopia}) para frascos não-higroscópicos ou com ganhos fora da tolerância legal.

Um retorno abaixo da tara real conhecida é preservado exatamente como medido; não existe limiar fixo adicional. Sem confirmação de vazio, encerrar a custódia como DEVOLVIDO\_COM\_ANOMALIA, bloquear o frasco em quarentena/INDISPONIVEL e manter consumo final pendente. A referência teórica não é limite físico absoluto. Aplicam-se os contratos de resolução da Seção 10; Q06 compara retorno com saída, não retorno com tara.
\end{regra}

\subsubsection{Exclusividade na retirada de reagentes}
\begin{regra}
O Chefe Geral atua como operador de balcão, nunca como destinatário/retirante. Apenas os usuários cadastrados como \textit{Professor} ou \textit{Bolsista} têm permissão para retirar reagentes do almoxarifado (o \textit{Emprestimo\_Reagente.id\_usuario\_retirou} deve corresponder a um \textit{Usuario} com uma dessas duas linhas associadas). Demais Alunos (mesmo de iniciação científica ou TCC, sem o papel de Bolsista) não possuem autorização no sistema para esta ação. Esta regra foi corrigida para não contradizer a seção \textit{Bolsista} (3.6), que descreve explicitamente o gestor de reagentes fazendo empréstimos para usuários Bolsista. O autoatendimento é exclusivo de usuários portadores da combinação \texttt{Professor + Gestor\_Almoxarifado}. Usuários no papel \texttt{Aluno + Gestor\_Almoxarifado} são terminantemente impedidos pelo backend de retirar reagentes para si mesmos, em respeito à matriz de segregação de funções (SOD).
\end{regra}

\subsubsection{Usuário Multi-role, exceto Chefe Geral}
\begin{regra}
Uma conta que possui o papel \textit{Chefe Geral} não pode possuir nenhum outro papel. Se a mesma pessoa exercer, na instituição, a função de Chefe Geral e também uma função representada por outro papel do sistema, deverão existir duas contas independentes, cada uma com identidade e permissões próprias.
\end{regra}
\subsubsection{Usuário Multi-role, Aluno não pode ser professor}
\begin{regra}
O usuário que é aluno não pode ser professor.
\end{regra}
\subsubsection{Usuário Multi-role, Quais multi papeis podem existir}
\begin{regra}
Pode haver:\\
Professor que é Gestor de Bens Patrimoniais.\\
Professor que é Gestor de Almoxarifado\\
Professor que é Gestor de Bens Patrimoniais e Gestor de Almoxarifado\\
Aluno que é Bolsista\\
Aluno que é Bolsista e Gestor de Bens Patrimoniais\\
Aluno que é Gestor de Almoxarifado\\
Aluno que é Gestor de Bens Patrimoniais\\
Aluno que é Gestor de Bens Patrimoniais e Gestor de Almoxarifado\\
O Usuário Aluno que também é Bolsista não pode ser Gestor de Almoxarifado. Se, na vida real, um aluno é bolsista e também precisa ser Gestor de Almoxarifado, devem ser criadas duas contas para ele no sistema.\\
No caso do Aluno que exerce algum papel de Gestor (ou ambos): o aluno pode editar bens patrimoniais e/ou reagentes, de acordo com seus papéis respectivos, porque também tem outro papel além de Aluno.
\end{regra}
\subsubsection{Aluno que é bolsista}
\begin{regra}
Apenas o chefe geral consegue colocar um aluno como bolsista.
\end{regra}
\subsubsection{Status, vencimento e descarte de frascos}
\begin{regra}
\begin{itemize}
    \item Um frasco \texttt{VAZIO} ou \texttt{QUEBRADO} deve ser exibido como ``Pendente de Descarte'', salvo se já estiver \texttt{DESCARTADO}.
    \item Um frasco \texttt{VENCIDO} não é automaticamente descartado e também não é automaticamente bloqueado para uso. Para frasco vencido, o sistema mantém a decisão explícita: \texttt{em\_quarentena = true}, ou \texttt{uso\_vencido\_autorizado = true}, ou situação de \textit{PENDENTE DE DESCARTE} com \texttt{uso\_vencido\_autorizado = false}.
    \item Em cadastro, quando a validade efetiva já tiver expirado no instante da operação, o gestor é obrigado a escolher um desses três destinos: \texttt{QUARENTENA}, \texttt{PENDENTE\_DE\_DESCARTE} ou \texttt{DISPONIVEL} com uso vencido explicitamente autorizado conforme Q04.
    \item O frasco em quarentena não pode ser retirado; o frasco pendente de descarte não pode ser retirado; o frasco vencido com uso explicitamente autorizado pode ser retirado, mas a interface deve alertar transparentemente antes da confirmação do empréstimo.
    \item Se um frasco vencer enquanto estiver emprestado, o empréstimo continua válido. A devolução \textbf{não} recalcula o vencimento com a hora atual: a passagem do tempo é responsabilidade do job de vencimento, e a devolução lê o estado canônico persistido \textit{Frasco\_Reagente.vencido}. O snapshot imutável \texttt{vencido\_na\_retirada} distingue ``venceu durante o empréstimo'' (não estava vencido na retirada) de ``já estava vencido na retirada''; validade desconhecida sem vencimento afirmativo é o terceiro caso. Em qualquer desses casos o gestor registra o destino pós-validade.
    \item A ação de descarte institucional (\texttt{descartarFrasco}) exige \texttt{em\_quarentena = false} e \texttt{disponibilidade != EMPRESTADO}, e é permitida para: (1) frascos com estado físico \texttt{VAZIO}, (2) frascos \texttt{QUEBRADO}, (3) frascos com \texttt{vencido = true} e \texttt{uso\_vencido\_autorizado = false}, e (4) frascos com autorização técnica de descarte ativa. A autorização técnica da rota (4) é produzida exclusivamente pela decisão humana \texttt{PENDENTE\_DE\_DESCARTE} de saída de quarentena e nunca por texto livre em \texttt{detalhe\_status}. \texttt{INDISPONIVEL} isolado não habilita descarte. As rotas (1)--(3) admitem conteúdo químico remanescente nos estados \texttt{ABERTO} ou \texttt{FECHADO}, formalizando a baixa do passivo químico perigoso. A quarentena deve ser resolvida antes; se \texttt{em\_quarentena = true}, nenhuma rota se aplica diretamente.
\end{itemize}
\end{regra}

\subsubsection{Extravio, reencontro e quarentena (M5)}
\begin{regra}
Consolidação normativa do ciclo de extravio, reencontro e quarentena. As operações são server-owned, exigem papel \textit{Chefe\_Geral} ou \textit{Gestor\_Almoxarifado} com vínculo ao almoxarifado do frasco e justificativa humana (mínimo de 20 caracteres após \texttt{trim}) para extravio e reencontro; a autorização completa de papéis pertence a M9.
\begin{itemize}
    \item \textbf{Extravio.} Um frasco pode ser declarado extraviado a partir de qualquer localização LOCALIZADO e estado físico, exceto \texttt{DESCARTADO} (terminal). O efeito é \texttt{situacao\_localizacao = EXTRAVIADO}, o último \texttt{estado\_fisico\_frasco} é preservado e a \texttt{disponibilidade = INDISPONIVEL}. A autorização corrente de descarte é revogada no extravio; a decisão anterior permanece no histórico auditável e não autoriza descarte automático após o reencontro. Extravio \textbf{não} é consumo, esgotamento, quebra nem peso zero: não fabrica tara, saldo, peso, consumo ou evidência metrológica. São preservados \textit{saldo}, \textit{saldo\_desconhecido}, \textit{peso\_atual}, \textit{peso\_no\_cadastrado}, tara e \textit{origem\_tara}, \textit{validade\_fechado}/\textit{validade\_efetiva}/\textit{validade\_desconhecida}, \textit{vencido}, \textit{uso\_vencido\_autorizado}, \textit{abertura\_historica\_desconhecida} e a localização. O \textit{estado físico anterior} é o último estabelecido na trilha histórica e não é apagado.
    \item \textbf{Extravio durante empréstimo ativo.} Se existir empréstimo \texttt{EM\_USO} ou \texttt{ATRASADO} para o frasco, ele é encerrado como \texttt{ENCERRADO\_EXTRAORDINARIO} com \texttt{tipo\_encerramento\_excepcional = EXTRAVIO\_SINISTRO}, \texttt{motivo\_encerramento\_excepcional} e \texttt{id\_gestor\_encerramento}. Como \textbf{não houve retorno físico}: \texttt{peso\_retorno}, \texttt{peso\_retorno\_efetivo}, \texttt{medida\_utilizada} e \texttt{id\_resolucao\_metrologica} permanecem \texttt{NULL}, \texttt{consumo\_validado = FALSE} e \texttt{data\_devolucao\_efetuada} permanece \texttt{NULL} (o instante do encerramento fica no histórico/operação; não é devolução). \texttt{peso\_saida} é preservado. \texttt{massa\_perda\_estimada\_g} só pode ser preenchida quando há tara conhecida e é uma \textbf{estimativa} gravimétrica, nunca medição nem consumo validado (fronteira M6). Empréstimo ativo que já possua \texttt{peso\_retorno} é inconsistência e exige reconciliação manual.
    \item \textbf{Reencontro.} Só um frasco cuja \texttt{situacao\_localizacao = EXTRAVIADO} pode ser reencontrado. O reencontro grava LOCALIZADO, impõe quarentena, revoga qualquer autorização corrente remanescente e preserva o último estado físico conhecido. É um \textbf{novo fato físico}: não desfaz o extravio e não reabre empréstimo encerrado. O gestor informa o estado físico constatado (\texttt{ABERTO}, \texttt{FECHADO}, \texttt{VAZIO} ou \texttt{QUEBRADO}) e o frasco passa a \texttt{em\_quarentena = TRUE} e \texttt{disponibilidade = INDISPONIVEL}. O estado constatado é validado contra o \textbf{último estado físico conhecido antes do extravio} (\texttt{estadoAntesDoExtravio}), obtido da trilha histórica ordenada sem criar coluna nova: \texttt{FECHADO} só é aceito se antes era \texttt{FECHADO} e não havia abertura conhecida; \texttt{ABERTO} só é aceito se antes era \texttt{ABERTO} ou \texttt{FECHADO} (neste caso a abertura ocorreu em período sem rastreabilidade temporal precisa); \texttt{VAZIO} só é aceito se antes já era \texttt{VAZIO} (restauração sem nova medição); \texttt{QUEBRADO} representa quebra física constatada. Transições fisicamente impossíveis são rejeitadas, como \texttt{ABERTO} $\rightarrow$ \texttt{EXTRAVIADO} $\rightarrow$ \texttt{FECHADO}, \texttt{VAZIO} $\rightarrow$ \texttt{EXTRAVIADO} $\rightarrow$ \texttt{FECHADO}/\texttt{ABERTO} e \texttt{QUEBRADO} $\rightarrow$ \texttt{EXTRAVIADO} $\rightarrow$ \texttt{ABERTO}/\texttt{FECHADO}. \texttt{saldo}, \texttt{saldo\_desconhecido}, \texttt{peso\_atual}, tara e \texttt{vencido} permanecem; uma medição válida posterior é a única forma de alterar o saldo. Se o frasco foi encontrado \texttt{ABERTO} vindo de \texttt{FECHADO} (ou sem data de abertura determinável), marca-se \texttt{abertura\_historica\_desconhecida = TRUE}, \texttt{data\_abertura = NULL}, \texttt{validade\_desconhecida = TRUE} e \texttt{validade\_efetiva = NULL}; nunca se inventa a data. Se o frasco não era previamente \texttt{VAZIO} mas parece vazio ao ser reencontrado, M5 \textbf{não} inventa esgotamento, tara ou quantidade zero: registra a observação, mantém a quarentena e a confirmação quantitativa/metrológica pertence a M6.
    \item \textbf{Validade no reencontro.} O reencontro \textbf{não} recalcula vencimento pelo relógio. A passagem do tempo é responsabilidade do job de vencimento (Seção~10.7) e o estado persistido \textit{Frasco\_Reagente.vencido} é a autoridade operacional. O reencontro apenas lê esse estado e, no caso de indeterminação exata, marca \texttt{validade\_desconhecida}; não agenda, não revalida e não antecipa as correções de validade da V2.
    \item \textbf{Quarentena.} A quarentena é uma dimensão operacional própria, distinta de \texttt{INDISPONIVEL}: um frasco em quarentena permanece \texttt{INDISPONIVEL} e nunca é liberado para nova retirada enquanto a quarentena não for resolvida. Ela é ortogonal a \textit{estado físico}, \textit{validade}, \texttt{saldo\_desconhecido} e anomalia metrológica. \texttt{vencido = TRUE} nunca é usado como mecanismo de quarentena. Enquanto durar, ficam proibidas a retirada, a devolução (não haveria empréstimo ativo) e o descarte direto.
    \item \textbf{Saídas da quarentena.} O conjunto é fechado: (i) \texttt{VOLTAR\_A\_DISPONIVEL}, que encerra a quarentena, devolve \texttt{DISPONIVEL} e exige estado físico \texttt{ABERTO}/\texttt{FECHADO}, ausência de pendência metrológica e que o frasco não esteja vencido sem uso autorizado; (ii) \texttt{PENDENTE\_DE\_DESCARTE}, que encerra a quarentena, mantém \texttt{INDISPONIVEL} e cria a autorização operacional estruturada consumida por \texttt{descartarFrasco} (HQ-M2-008/B); (iii) \textbf{permanência} em quarentena, que não é uma decisão de saída e mantém o bloqueio. \texttt{VOLTAR\_A\_DISPONIVEL} \textbf{nunca} pode liberar \texttt{VAZIO}, \texttt{QUEBRADO} ou \texttt{DESCARTADO}; para frasco reencontrado \texttt{VAZIO} ou \texttt{QUEBRADO}, a resolução só pode manter a quarentena ou produzir \texttt{PENDENTE\_DE\_DESCARTE}. Nenhuma saída revalida, renova, estende ou zera o vencimento, e nenhuma delas volta o frasco a \texttt{DISPONIVEL} por conveniência. Não existe \texttt{QUARENTENA} $\rightarrow$ \texttt{DESCARTADO} direto; o descarte técnico exige a etapa \texttt{PENDENTE\_DE\_DESCARTE}.
    \item \textbf{Rastreabilidade.} Extravio, reencontro, entrada e resolução de quarentena geram eventos imutáveis no \textit{Historico\_Frasco\_Reagente} (\texttt{EXTRAVIO}, \texttt{REENCONTRO}, \texttt{ENTROU\_EM\_QUARENTENA}, \texttt{LIBERADO\_QUARENTENA}, \texttt{PENDENTE\_DE\_DESCARTE}) com frasco, almoxarifado, ator, instante, motivo e vínculo ao empréstimo quando aplicável. O reencontro não apaga o extravio; a resolução não apaga o reencontro; nada é sobrescrito. O estado anterior é o último estabelecido na própria trilha ordenada.
\end{itemize}
As operações críticas de M5 seguem o contrato global de idempotência (Seção~\ref{sec:idempotencia-m7}): \texttt{idOperacao} obrigatório com identidade \texttt{(uid, tipo\_operacao, payload\_hash)}; retry não reaplica extravio, reencontro nem resolução de quarentena.
\end{regra}

\subsubsection{Q06, tara, medição e resolução metrológica (M6)}
\begin{regra}[Vocabulário metrológico e autoridades]
Cada campo quantitativo possui uma única autoridade. Não usar dois campos como autoridade do mesmo fato sem explicitar a diferença.
\begin{itemize}
    \item \texttt{peso\_saida}: leitura bruta da balança na retirada; imutável; base da tolerância Q06.
    \item \texttt{peso\_retorno}: observação física original da devolução; persistida e nunca reescrita por resolução posterior.
    \item \texttt{peso\_retorno\_efetivo}: valor quantitativo validado que a resolução autoriza usar no cálculo definitivo; \texttt{NULL} enquanto a pendência estiver aberta; não apaga nem substitui \texttt{peso\_retorno}.
    \item \texttt{peso\_atual}: observação corrente do frasco; não prova o peso no instante histórico da devolução.
    \item \texttt{peso\_frasco\_vazio} + \texttt{origem\_tara}: tara (\texttt{NULL}, \texttt{REFERENCIA\_TEORICA} ou \texttt{MEDIDA\_REAL}).
    \item \texttt{peso\_no\_cadastrado}: peso total inicial; imutável.
    \item \texttt{medida\_utilizada}: consumo quantitativamente validado; \texttt{NULL} \textbf{não} é zero.
    \item \texttt{peso\_perda\_evaporacao}: perda em g registrada separadamente do consumo.
    \item \texttt{massa\_perda\_estimada\_g}: estimativa de sinistro/extravio (fronteira M5); não é medição nem consumo.
    \item \texttt{saldo} (\texttt{saldo\_aferido\_g}/\texttt{saldo\_aferido\_ml}) e \texttt{saldo\_desconhecido}: saldo corrente server-owned; desconhecido \textbf{não} é zero.
    \item \texttt{consumo\_validado}: autoridade que separa resultado quantitativo encerrado de pendência aberta.
    \item \texttt{anomalia\_metrologica}: classificação textual server-owned da inconsistência.
    \item \texttt{id\_resolucao\_metrologica}: vínculo ao evento auditável que encerrou a pendência.
    \item \texttt{densidade\_aplicada}: snapshot histórico da densidade no empréstimo.
\end{itemize}
Regra geral: desconhecido $\neq$ zero, ausente $\neq$ zero, não mensurável $\neq$ zero, estimado $\neq$ medido, atual $\neq$ histórico, referência teórica $\neq$ tara real. Nunca preencher um valor apenas para completar o modelo.
\end{regra}

\begin{regra}[Q06 --- tolerância de ganho de massa]
Q06 compara \textbf{retorno com saída}, sobre o peso bruto \texttt{peso\_saida}: reagente normal $tolerancia = \max(1\,g,\ 0{,}005 \times peso_{saida})$; higroscópico $tolerancia = \max(2\,g,\ 0{,}02 \times peso_{saida})$. A condição de ganho anômalo é $peso_{retorno} > peso_{saida} + tolerancia$. O antigo limiar fixo de 5\,g está definitivamente eliminado e não pode ser recriado como definição de vazio, segunda tolerância, fallback, exceção ou justificativa. Q06 \textbf{não} é regra de tara: nunca comparar \texttt{peso\_retorno} com \texttt{peso\_frasco\_vazio}; não existe ``tolerância de tara Q06'', margem percentual para retorno abaixo da tara nem limiar fixo de vazio.
\end{regra}

\begin{regra}[Consumo, evaporação e densidade]
A perda bruta é $\Delta_{bruto} = \max(0,\ peso_{saida} - peso_{retorno}^{efetivo})$. A evaporação é registrada separadamente do consumo (campo e evento \texttt{EVAPORACAO}) e obedece à constraint $0 \le \text{Evap} \le \Delta_{bruto}$; não é autoridade de vazio nem de Q06. Se a evaporação informada exceder a perda bruta, a operação é \textbf{rejeitada} antes da conclusão quantitativa: não há clamp silencioso, conversão automática de $8\,g$ em $5\,g$, aceitação com consumo zero nem anomalia que “salve” o dado. Quando a constraint é satisfeita, $massa_{consumida} = \Delta_{bruto} - \text{Evap} \ge 0$ naturalmente; o $\max(0,\cdot)$ externo não é mecanismo para esconder evaporação impossível. Para sólidos, a medida validada é a massa em g; para líquidos, o volume é $massa_{consumida} / densidade_{aplicada}$, com \texttt{densidade\_aplicada} positiva e histórica do empréstimo (nunca uma densidade corrente arbitrária). Não misturar g e mL. Se \texttt{peso\_retorno} $\ge$ \texttt{peso\_saida} na devolução original, a perda bruta observada é zero e, com os dados daquela operação, a evaporação deve ser zero: não se afirma ganho líquido e evaporação positiva simultaneamente. Enquanto a pendência estiver aberta, \texttt{medida\_utilizada}, \texttt{peso\_retorno\_efetivo} permanecem \texttt{NULL} e \texttt{consumo\_validado = FALSE}; \texttt{NULL} nunca é interpretado como consumo zero. Quando uma rota estabelece um \texttt{peso\_retorno\_efetivo} diferente do original, a evaporação é \textbf{revalidada} contra a nova perda bruta antes de marcar \texttt{consumo\_validado = true}; se deixar de ser consistente, a pendência permanece (não se altera a medição nem a evaporação automaticamente, e não existe correção administrativa metrológica na V1).
\end{regra}

\begin{regra}[Tara --- três situações]
\begin{itemize}
    \item \textbf{Desconhecida} (\texttt{peso\_frasco\_vazio = NULL}, \texttt{origem\_tara = NULL}): não é zero; não estimar saldo atual apenas porque existe peso bruto.
    \item \textbf{Referência teórica} (\texttt{origem\_tara = REFERENCIA\_TEORICA}): valor derivado do nominal; não é medição física definitiva do recipiente vazio nem limite físico absoluto. Retorno inferior à referência teórica não é impossibilidade física e não deve ser ``corrigido'' artificialmente.
    \item \textbf{Real} (\texttt{origem\_tara = MEDIDA\_REAL}): só existe quando o recipiente foi efetivamente constatado vazio e pesado nessa condição. Nunca derivar de estimativa, massa nominal, densidade, produto ainda presente, cálculo reverso ou pesagem de recipiente não vazio.
\end{itemize}
\end{regra}

\begin{regra}[Esgotamento e tara real]
O esgotamento não é inferido automaticamente pela comparação com a tara; a autoridade é a confirmação explícita da condição física (\texttt{confirmarEsgotamento}), tanto para \texttt{FECHADO} quanto para \texttt{JA\_ABERTO}. Recipiente efetivamente vazio confirmado e pesado: \texttt{estado\_fisico\_frasco = VAZIO}, \texttt{disponibilidade = INDISPONIVEL}, \texttt{saldo\_desconhecido = false}, saldo corrente zero; se não havia tara real, \texttt{peso\_frasco\_vazio = peso vazio medido} e \texttt{origem\_tara = MEDIDA\_REAL}; \texttt{conteudo\_nominal} não é alterado; registra-se \texttt{FICOU\_VAZIO}. Tara real anterior divergente \textbf{não} é sobrescrita: registra-se \texttt{DISCREPANCIA\_TARA\_REAL} e a substituição exige \texttt{recalibrarTaraFrascoEsgotado} (recipiente vazio, justificativa humana, auditoria antes/depois, histórico da tara anterior preservado).
\end{regra}

\begin{regra}[Retorno abaixo da tara]
\begin{itemize}
    \item \textbf{Abaixo de tara real} (\texttt{origem\_tara = MEDIDA\_REAL}, \texttt{peso\_retorno < peso\_frasco\_vazio}, \texttt{confirmarEsgotamento = false}): preservar a leitura original; encerrar a custódia; empréstimo \texttt{DEVOLVIDO\_COM\_ANOMALIA}; pendência quantitativa aberta; frasco \texttt{INDISPONIVEL} e \texttt{em\_quarentena = true}; nunca criar saldo negativo nem tratar o retorno como zero.
    \item \textbf{Abaixo de referência teórica}: preservar a leitura; classificar \texttt{REFERENCIA\_TEORICA\_INCONSISTENTE}; não declarar impossibilidade física; seguir o contrato de inconsistência sem promover a referência a tara real.
\end{itemize}
\end{regra}

\begin{regra}[Ganho de massa --- três faixas]
A avaliação usa \texttt{peso\_retorno} contra \texttt{peso\_saida} e a tolerância Q06 determinada por \texttt{eh\_higroscopico} (snapshot imutável do frasco, preenchido no cadastro a partir de \textit{Resumo\_Reagente.eh\_higroscopico}; alterações posteriores do resumo não o reescrevem):
\begin{center}
\begin{tabular}{p{0.30\textwidth}p{0.28\textwidth}p{0.34\textwidth}}
\toprule
\textbf{Faixa} & \textbf{Sem higroscopia} & \textbf{Higroscópico} \\
\midrule
$peso_{retorno} \le peso_{saida}$ & consumo normal & consumo normal \\
$peso_{saida} < peso_{retorno} \le peso_{saida}+tol$ & consumo 0; ajuste/observação; \textbf{sem} evento de higroscopia & consumo 0; evento \texttt{ganho\_massa\_higroscopia} permitido \\
$peso_{retorno} > peso_{saida}+tol$ & \texttt{GANHO\_ACIMA\_Q06}: anomalia & \texttt{GANHO\_ACIMA\_Q06}: anomalia \\
\bottomrule
\end{tabular}
\end{center}
Ganho dentro da tolerância não produz consumo negativo e mantém o peso físico real; ganho acima de Q06 \textbf{não} provoca rollback da devolução física, encerra a custódia, gera \texttt{DEVOLVIDO\_COM\_ANOMALIA} e bloqueia o frasco (INDISPONIVEL/quarentena) para investigação. É proibido registrar evento de higroscopia para reagente não higroscópico ou com ganho fora da tolerância.
\end{regra}

\begin{regra}[Pendência metrológica --- \texttt{existePendenciaMetrologicaTx}]
O predicado é derivado de estado canônico, nunca de texto livre. Autoridade: \textit{Emprestimo\_Reagente}. Entrada: \texttt{idFrasco}. Retorna \texttt{true} se e somente se existe empréstimo daquele frasco com \texttt{status = DEVOLVIDO\_COM\_ANOMALIA} \textbf{e} \texttt{consumo\_validado = false}; caso contrário \texttt{false}. A pendência é um estado próprio do par canônico, não o status isolado. Dados fora do par (\texttt{consumo\_validado}=false com status que não seja \texttt{DEVOLVIDO\_COM\_ANOMALIA}) são \textbf{inconsistência de integridade}, não pendência metrológica: registram-se para correção e bloqueiam a liberação automática de quarentena até reconciliação.
\end{regra}

\begin{regra}[Resolução metrológica --- três rotas tipadas]
Não existe \texttt{resolverPendenciaMetrologica(campo, valor)}, dispatcher genérico público nem undo universal. A V1 possui exatamente \textbf{três rotas metrológicas tipadas} (HQ-M2-009/A, revisada por HQ-M6-001), cada uma com contrato próprio; todas preservam \texttt{peso\_retorno}, o status histórico e o vínculo \texttt{id\_resolucao\_metrologica}:
\begin{enumerate}
    \item \texttt{repetirPesagemDevolucaoAnomala}: relê frasco e empréstimo; registra a nova leitura; só valida o consumo histórico quando não houve movimentação física posterior que impeça provar o instante original (\texttt{houveMovimentacaoFisicaPosteriorTx}). Nesse caso, calcula $massa_{consumida}=(peso_{saida}-peso_{retorno}^{efetivo})-\text{Evap}$ com a densidade histórica, após revalidar a constraint de evaporação. Se houver movimentação física posterior, registra a observação e \textbf{mantém} a pendência.
    \item \texttt{confirmarEsgotamentoAposInspecao}: confirmação humana explícita e pesagem real do recipiente vazio; aplica o contrato de esgotamento; não infere vazio da anomalia; tara real divergente é preservada; o estado corrente vira \texttt{VAZIO}/INDISPONIVEL e o consumo histórico só é validado com evidência suficiente do retorno.
    \item \texttt{recalibrarTaraFrascoEsgotado}: somente com recipiente efetivamente vazio; justificativa humana; registra tara anterior e nova; não reprocessa automaticamente empréstimos históricos; não transforma a tara nova em evidência pretérita.
\end{enumerate}
\textbf{Não existe quarta rota metrológica.} A V1 \textbf{não} possui correção administrativa metrológica, nem genérica nem específica: após persistida, uma medição (\texttt{peso\_saida}, \texttt{peso\_retorno}, \texttt{peso\_retorno\_efetivo}, tara, evaporação ou consumo) não pode ser editada administrativamente, e não existem \texttt{CORRIGIR\_PESO\_RETORNO}, \texttt{ERRO\_TRANSCRICAO\_PESO}, \texttt{CORRIGIR\_EVAPORACAO}, \texttt{CORRIGIR\_TARA}, editor \texttt{campo+novoValor} ou \texttt{resolverPendenciaMetrologica} genérico. A operação \texttt{corrigirOperacao} permanece exclusivamente para tipos administrativos não metrológicos explicitamente suportados (hoje \texttt{RETIRADA\_RETIRANTE\_INCORRETO}) e \textbf{não} encerra pendência metrológica. Consequência: erros de transcrição de peso devem ser prevenidos \textbf{antes do commit} (UI-18); quando uma leitura errada é persistida e nenhuma das três rotas físicas tipadas puder legitimamente resolver a situação, a V1 \textbf{não} fabrica evidência para encerrar a pendência — a pendência permanece. Uma nova pesagem nunca é prova automática da medição histórica anterior.
A resolução grava evento \texttt{AJUSTE} vinculado ao empréstimo e à medição original, com autoria, instante, tara/densidade aplicadas e \texttt{id\_resolucao\_metrologica}. A data FLOW do consumo permanece a da devolução física; reprocessar somente as datas diretamente afetadas, sem replay nem snapshot de estoque. A resolução \textbf{não} libera a quarentena: a saída continua dependendo do contrato M5 (\texttt{resolverQuarentenaFrasco}). Cada rota é um comando idempotente no contrato global (Seção~\ref{sec:idempotencia-m7}): \texttt{idOperacao} obrigatório e identidade \texttt{(uid, tipo\_operacao, payload\_hash)}; retry devolve o resultado persistido sem reaplicar a resolução.
\end{regra}

\paragraph{Matriz de casos quantitativos.}
\begin{longtable}{p{0.28\textwidth}p{0.18\textwidth}p{0.16\textwidth}p{0.30\textwidth}}
\toprule
\textbf{Caso} & \textbf{Consumo} & \textbf{Pendência} & \textbf{Estado/quarentena} \\
\midrule
\endfirsthead
\toprule
\textbf{Caso} & \textbf{Consumo} & \textbf{Pendência} & \textbf{Estado/quarentena} \\
\midrule
\endhead
\bottomrule
\endfoot
1. retorno $<$ saída (normal) & validado, $>0$ & não & normal \\
2. retorno $=$ saída & 0 & não & normal \\
3. ganho $\le$ Q06, não higroscópico & 0 (ajuste) & não & normal \\
4. ganho $\le$ Q06, higroscópico & 0 (ajuste higroscópico) & não & normal \\
5. ganho $>$ Q06 & pendente & sim & INDISPONIVEL + quarentena \\
6. retorno $<$ referência teórica & pendente & sim & INDISPONIVEL + quarentena \\
7. retorno $<$ tara real s/ vazio & pendente & sim & INDISPONIVEL + quarentena \\
8. vazio confirmado sem tara real & validado & não & VAZIO/INDISP.; tara real nova \\
9. vazio confirmado c/ referência teórica & validado & não & VAZIO; tara real substitui teórica \\
10. vazio c/ tara real igual & validado & não & VAZIO; tara preservada \\
11. vazio c/ tara real divergente & validado & não & VAZIO; discrepância; recalibração \\
12. anomalia resolvida por pesagem válida & validado & não & status anômalo; quarentena mantida \\
13. pesagem não prova o instante histórico & pendente & sim & observação; quarentena \\
14. confirmação posterior de esgotamento & validado (se sem mov.) & não/sim & VAZIO; quarentena mantida \\
15. recalibração de tara real & não altera & conforme & antes/depois; sem replay \\
16. saldo desconhecido sem tara & pendente & conforme & saldo \texttt{NULL}; nunca zero \\
17. líquido c/ densidade histórica válida & validado em mL & não & normal \\
18. líquido sem densidade histórica válida & rejeitado & --- & operação recusada \\
19. resolução + permanência em quarentena & validado & não & INDISPONIVEL; M5 decide saída \\
20. extravio c/ \texttt{massa\_perda\_estimada\_g} & não & não & ENCERRADO\_EXTRAORD.; estimativa \\
21. evaporação $=$ perda bruta & 0 & não & consumo zero; válido \\
22. evaporação $<$ perda bruta & $>0$ & não & válido \\
23. evaporação $>$ perda bruta & rejeitado & --- & \textbf{bloqueia antes do commit} \\
24. ganho de massa ($peso_{retorno}\ge peso_{saida}$) c/ evaporação $>0$ & rejeitado & --- & \textbf{bloqueia antes do commit} \\
25. dupla digitação igual & conforme & conforme & prossegue \\
26. dupla digitação divergente & não persistido & --- & \textbf{bloqueia antes do commit} \\
\end{longtable}

\subsubsection{Idempotência (M7)}
\label{sec:idempotencia-m7}
\begin{regra}[Taxonomia: cinco conceitos distintos]
Idempotência é a propriedade de repetir a \textbf{mesma intenção lógica} sem produzir um segundo efeito lógico. Não confundir cinco mecanismos:
\begin{itemize}
    \item \textbf{Idempotência de comando} (cliente/usuário): repetir o mesmo comando autenticado devolve o resultado já persistido e não repete efeitos de domínio — chave de operação em \textit{Operacoes}.
    \item \textbf{Deduplicação de evento} (trigger/entrega at-least-once): o mesmo \texttt{eventId} não reaplica o efeito — \textit{Eventos\_Processados}.
    \item \textbf{Chave de unicidade}: impede criar dois recursos com o mesmo valor normalizado (plaqueta, matrícula, CAS) — \textit{Chaves\_Unicas}; não é idempotência de comando.
    \item \textbf{Lock concorrente}: serializa/limita operações concorrentes — \textit{Locks\_Requisicao\_Patrimonio}; não é idempotência.
    \item \textbf{Materialização reconstruível}: recalcular a mesma janela substitui a projeção de forma absoluta, sem acumular.
\end{itemize}
Idempotência \textbf{não} substitui pré-condições, autorização, locking nem controle de concorrência de domínio (I14).
\end{regra}

\begin{regra}[Identidade de uma operação]
A identidade semântica de um comando é o triplo:
\begin{center}
\texttt{(uid, tipo\_operacao, payload\_hash)}
\end{center}
onde \texttt{uid} é a identidade autenticada no servidor (nunca fornecida pelo payload), \texttt{tipo\_operacao} é a ação e \texttt{payload\_hash} é a impressão canônica do payload semanticamente relevante. Um \texttt{idOperacao} pertence a exatamente uma identidade. Todas as verificações de retry comparam os três campos: \textbf{repetir só é retry se uid, tipo\_operacao e payload\_hash coincidirem}. Verificar apenas \texttt{uid} ou apenas \texttt{uid+payload\_hash} é insuficiente e proibido (I4--I6). Um mesmo \texttt{idOperacao} com ator, ação ou payload diferente é rejeitado de forma fechada com \texttt{ALREADY\_EXISTS}/\texttt{already-exists} (I4--I6).
\end{regra}

\begin{regra}[\texttt{idOperacao}: formato, criação e reutilização]
\begin{itemize}
    \item \textbf{Formato}: string opaca gerada no cliente com CSPRNG (por exemplo UUID v4), 1 a 128 caracteres, conjunto seguro \texttt{[A-Za-z0-9\_-]}; não participa do próprio hash nem de nenhum payload.
    \item \textbf{Origem}: criada pelo cliente \textbf{antes da primeira tentativa mutável}, persistida junto da intenção local (sessão/rascunho) e \textbf{reutilizada} durante todo retry (timeout, perda de conexão, resposta desconhecida).
    \item \textbf{Novo id}: obrigatório apenas quando o usuário inicia deliberadamente uma \textbf{nova intenção}; o cliente nunca gera novo id apenas por ter perdido a resposta (isso derrotaria a idempotência).
    \item \textbf{Vida útil}: do primeiro envio até a obtenção de um resultado conclusivo (ou descarte explícito da intenção). Não reutilizar o mesmo id para outra intenção.
\end{itemize}
\end{regra}

\begin{regra}[Canonicalização determinística do payload]
O hash é calculado por um único algoritmo global, \texttt{hashPayload(tipoOperacao, payload)} = \texttt{SHA-256(tipoOperacao + "\char`\\n" + canonicalize(payload))}. A canonicalização \texttt{canonicalize} é definida como:
\begin{lstlisting}[language=TypeScript, title={Canonicalização canônica do payload (M7)}]
function canonicalize(value: unknown): string {
  if (value === null) return "null";
  const t = typeof value;
  if (t === "boolean" || t === "number" || t === "string")
    return JSON.stringify(value);
  if (Array.isArray(value))
    return "[" + value.map(canonicalize).join(",") + "]"; // ordem preservada
  if (t === "object") {
    const obj = value as Record<string, unknown>;
    const keys = Object.keys(obj)
      .filter(k => obj[k] !== undefined) // undefined e chave ausente são equivalentes
      .sort();                            // ordenação por code unit, recursiva
    return "{" + keys.map(k => JSON.stringify(k) + ":" + canonicalize(obj[k])).join(",") + "}";
  }
  throw new Error("canonicalize: tipo não suportado");
}
function hashPayload(tipoOperacao: string, payload: Record<string, unknown>): string {
  return sha256(tipoOperacao + "\n" + canonicalize(payload));
}
\end{lstlisting}
Consequências normativas: propriedades em ordem diferente produzem o mesmo hash; objetos aninhados são ordenados recursivamente; arrays \textbf{preservam a ordem} (quando a ordem é semanticamente relevante, reordenar é payload diferente); \texttt{null} e chave ausente são \textbf{distintos}; \texttt{undefined} é equivalente à ausência e não entra no hash; números/strings/booleanos seguem \texttt{JSON.stringify} (sem arredondamento, sem trim, sem uppercase, sem conversão de unidade); \texttt{idOperacao} é excluído do payload; campos derivados somente no servidor não entram no payload. A validação de justificativa humana (mínimo após trim) \textbf{não} altera o payload original usado no hash. Não usar o \emph{replacer} em forma de array do \texttt{JSON.stringify}, pois ele descarta propriedades de objetos aninhadas.
\end{regra}

\begin{regra}[Estados e resultado de \textit{Operacoes}]
\textit{Operacoes/\{idOperacao\}}: \texttt{uid}, \texttt{tipo\_operacao}, \texttt{payload\_hash}, \texttt{status}, \texttt{resultado} (resposta mínima), \texttt{criado\_em}, \texttt{atualizado\_em} — todos escritos pelo servidor.
\begin{itemize}
    \item \textbf{Comando Firestore atômico}: o documento é criado já como \texttt{CONCLUIDA}, na mesma transação do efeito de domínio; não há janela \texttt{PENDENTE}. É o caso da maioria das mutações.
    \item \textbf{Fluxo com etapa externa pós-commit} (PDF, upload, Auth, e-mail): cria-se \texttt{PENDENTE} antes/na transação, o processador idempotente executa o efeito externo e só então grava \texttt{CONCLUIDA} com o resultado; \texttt{FALHOU} é terminal para aquela tentativa.
    \item \textbf{Retry}: \texttt{CONCLUIDA} $\rightarrow$ devolve o \texttt{resultado} persistido sem reexecutar efeitos; \texttt{PENDENTE} $\rightarrow$ retoma/reconhece sem duplicar efeito; \texttt{FALHOU} $\rightarrow$ exige nova intenção ou retry explícito conforme a política da operação.
    \item \textbf{Dados inválidos}: documento \textit{Operacoes} incompatível/corrompido (campos ausentes) é tratado fail-closed; nunca se assume sucesso.
\end{itemize}
Retry exato \textbf{não} cria novo histórico, novo contador, novo empréstimo, novo frasco, nova notificação, novo timestamp lógico nem novo efeito externo (I1--I3, I10). O \texttt{timestamp} da operação é o da criação; o retry não o atualiza.
\end{regra}

\begin{regra}[Concorrência]
Duas requisições idênticas simultâneas são serializadas pela transação Firestore: a primeira confirma e cria \textit{Operacoes}; a segunda, relendo o documento confirmado, devolve o resultado sem segundo efeito (I9). Se a transação abortar antes do commit, nenhum efeito persiste. Se o callback transacional for repetido automaticamente pelo Firestore, nenhum efeito externo pode ocorrer dentro dele. Idempotência de comando não resolve concorrência entre \textbf{intenções diferentes} (dois gestores retirando o mesmo frasco, duas revogações removendo o último responsável, duas decisões de quarentena): isso continua exigindo pré-condições, leituras transacionais e locks.
\end{regra}

\begin{regra}[Efeitos externos]
É \textbf{proibido} executar efeito irreversível (e-mail, PDF, upload, Auth, chamada de rede) diretamente no callback transacional, pois o callback pode repetir. A transação persiste o fato e um \textbf{evento/outbox/operação}; um processador idempotente realiza o efeito após o commit. Para uploads, use caminho/objeto determinístico derivado de \texttt{idOperacao} (sobrescrita idempotente) e registre o resultado em \textit{Operacoes}. Falhas são recuperadas por reconciliação; ``Firestore commitou mas o efeito externo falhou'' mantém a operação \texttt{PENDENTE} e não anuncia conclusão prematura (I11).
\end{regra}

\begin{regra}[Triggers e entrega at-least-once]
\textit{Eventos\_Processados/\{eventId\}} deduplica eventos de trigger: a transação do trigger só aplica o efeito se ainda não processou aquele \texttt{eventId} e registra o processamento no mesmo commit. A linguagem correta é ``o \textbf{efeito observado} do mesmo \texttt{eventId} é transacionalmente único / deduplicado'', e \textbf{não} ``entrega exactly-once'': Cloud Functions/Eventarc pode reentregar o evento; o que o sistema garante é que a reentrega não reaplica o efeito (I12). Dois \texttt{eventId} distintos representam fatos distintos e não são deduplicados entre si.
\end{regra}

\begin{regra}[Jobs agendados]
Dois comportamentos distintos, nunca confundidos:
\begin{itemize}
    \item \textbf{Chave determinística + \texttt{create()}}: reexecução é \textbf{no-op} se a chave já existe (ex.: notificação com docId \texttt{\{id\_emprestimo\}--\{janela\}}; \texttt{BulkWriter.create} tratando \texttt{ALREADY\_EXISTS}/\texttt{code 6}). Reexecutar não cria novo fato.
    \item \textbf{Recomputação absoluta por \texttt{DATA\_ALVO}}: reexecução \textbf{recalcula e SUBSTITUI} a projeção da chave diária, \textbf{nunca incrementa} sobre o valor existente (I13). Dadas as mesmas fontes e \texttt{versao\_calculo}, o resultado é idêntico.
\end{itemize}
Jobs devem ser seguros para reexecução, execução parcial e \textit{backfill}; nenhum job acumula sobre materialização anterior.
\end{regra}

\begin{regra}[Notificações, etiquetas e materializações]
Notificações determinísticas usam docId determinístico e \texttt{create} (ou consumo transacional de evento), sem duplicar (I3, I13). \textbf{Etiquetas/PDF}: \texttt{idOperacao} é \textbf{obrigatório} (sem fallback \texttt{Date.now()}); a geração é etapa externa pós-commit com objeto de caminho determinístico por \texttt{idOperacao}, registrando \texttt{status\_geracao} \texttt{PENDENTE}$\rightarrow$\texttt{CONCLUIDA} e o resultado na operação; retry devolve o resultado armazenado. \textbf{Materializações}: reprocessar a mesma janela substitui a projeção de forma absoluta, sem dupla soma; dados de consumo validados não são somados novamente (I13).
\end{regra}

\begin{regra}[Relação com M5 e M6 e segurança]
O contrato global aplica-se às operações de M5 (extravio/reencontro/quarentena) e às três rotas metrológicas de M6 (\texttt{repetirPesagemDevolucaoAnomala}, \texttt{confirmarEsgotamentoAposInspecao}, \texttt{recalibrarTaraFrascoEsgotado}) \textbf{sem alterar a semântica de domínio} de M5/M6. \texttt{payload\_hash} é impressão de igualdade da operação, \textbf{não} é assinatura, MAC, prova de autoria, autorização nem controle de acesso; autoria é a identidade autenticada no servidor e autorização continua nas Rules/RBAC.
\end{regra}

\paragraph{Matriz de casos de idempotência (M7).}
\begin{longtable}{p{0.36\textwidth}p{0.22\textwidth}p{0.34\textwidth}}
\toprule
\textbf{Caso} & \textbf{Resultado} & \textbf{Observação} \\
\midrule
\endfirsthead
\toprule
\textbf{Caso} & \textbf{Resultado} & \textbf{Observação} \\
\midrule
\endhead
\bottomrule
\endfoot
1. mesma chave/ator/ação/payload após commit & resultado original & retry puro \\
2. mesma chave/ator, payload diferente & REJEITAR & \texttt{ALREADY\_EXISTS} \\
3. mesma chave, ator diferente & REJEITAR & \texttt{ALREADY\_EXISTS} \\
4. mesma chave, ação diferente & REJEITAR & \texttt{ALREADY\_EXISTS} \\
5. duas requisições idênticas simultâneas & um efeito & transação serializa \\
6. resposta HTTP perdida após commit & resultado original & retry com mesmo id \\
7. erro de validação antes da transação & nenhum efeito & nada persistido \\
8. erro dentro da transação antes do commit & nenhum efeito & rollback \\
9. retry automático do callback Firestore & nenhum efeito externo & outbox pós-commit \\
10. propriedades JSON em ordem diferente & mesmo hash & canonicalização \\
11. objeto aninhado em ordem diferente & mesmo hash & recursivo \\
12. array em ordem diferente & payload diferente & ordem preservada \\
13. \texttt{null} vs campo ausente & payload diferente & distintos \\
14. campo derivado só no servidor & fora do hash & server-owned \\
15. mesmo \texttt{eventId} duas vezes & efeito único & \textit{Eventos\_Processados} \\
16. dois \texttt{eventId} distintos & dois fatos & não deduplicar \\
17. job diário duas vezes & no-op/substituição & conforme tipo \\
18. materialização mesma \texttt{DATA\_ALVO} & substitui & sem dupla soma \\
19. notificação docId determinístico existente & no-op & \texttt{create}/\texttt{ALREADY\_EXISTS} \\
20. processador externo executado novamente & não repete efeito & idempotência do processador \\
21. \textit{Operacoes} incompatível/corrompido & fail-closed & nunca assumir sucesso \\
22. rota M6 repetida após resolução & resultado original & contrato global \\
23. etiquetas após resposta perdida & mesmo PDF/resultado & caminho determinístico por id \\
24. retry de correção administrativa & resultado original & contrato global \\
\end{longtable}

\subsubsection{Estoque atual, aptidão, escassez e notificações (M8)}
\label{sec:m8-estoque-escassez-notificacoes}

\begin{regra}[Estoque atual e resumos FLOW --- autoridades distintas]
\textit{Frasco\_Reagente} é a única autoridade do \textbf{estoque atual}. Não existe materialized view persistente de estoque atual e nenhuma regra pode reconstruí-lo por encadeamento diário. É proibido qualquer cálculo do tipo $estoqueAtual = resumoDeOntem + movimentacoesDeHoje$, bem como ``posição de D depende da posição de D-1''. O estoque atual é sempre calculado diretamente sobre o estado canônico corrente.

Os resumos \textit{Resumo\_Almoxarifado\_Diario} e \textit{Resumo\_Reagente\_Diario} são exclusivamente \textbf{projeções FLOW históricas}, derivadas, descartáveis e reconstruíveis a partir de \textit{Historico\_Frasco\_Reagente} e \textit{Emprestimo\_Reagente}. Não são autoridade do estoque atual e não voltam a receber campos STOCK (quantidades de frascos, massa ou volume ``no fim do dia''). Nenhuma leitura de estoque atual usa resumo diário como fonte e nenhuma correção exclusivamente histórica invalida cache de estoque atual se não alterar o estado corrente.
\end{regra}

\begin{regra}[Agregação protegida]
A cadeia de agregação é: navegador $\rightarrow$ backend protegido $\rightarrow$ \texttt{count()}/\texttt{sum()} $\rightarrow$ \textit{Frasco\_Reagente}. O navegador não baixa frascos para contá-los, não executa a autoridade da agregação e não trata o catálogo JSON como autoridade operacional. Nenhum contador cadastral (por exemplo, o vínculo de lote) substitui o estoque atual.
\end{regra}

\begin{regra}[Saldos quantitativos server-owned]
\texttt{saldo\_aferido\_g} e \texttt{saldo\_aferido\_ml} são campos derivados server-owned, recalculados na mesma transação das alterações de peso, tara ou conhecimento. Saldo desconhecido, metrologicamente inconsistente, ausente ou não mensurável é \texttt{NULL} e \textbf{nunca} é somado como zero conhecido; vazio explicitamente confirmado produz zero conhecido. A origem da tara permanece explícita e massa (g) e volume (mL) \textbf{nunca} são somados entre si.

A dashboard de estoque distingue conceitualmente: quantidade de frascos; quantidade com saldo conhecido; quantidade com saldo desconhecido; massa conhecida (g); volume conhecido (mL). Um frasco extraviado, ainda que fisicamente \texttt{ABERTO}, não integra o saldo disponível.
\end{regra}

\begin{regra}[Predicado único de aptidão: \texttt{frascoAptoParaUso}]
Existe uma única autoridade conceitual de aptidão, usada de forma mecanicamente coerente pela dashboard de estoque, pela escassez e pela retirada. Um frasco $f$ é \textbf{apto para uso} se, e somente se, todas as condições abaixo valem:
\begin{itemize}
    \item \texttt{situacao\_localizacao = LOCALIZADO};
    \item \texttt{disponibilidade = DISPONIVEL};
    \item \texttt{estado\_fisico\_frasco} $\in$ \{\texttt{FECHADO}, \texttt{ABERTO}\};
    \item \texttt{em\_quarentena = false};
    \item validade compatível: \texttt{vencido = false} \textbf{ou} \texttt{uso\_vencido\_autorizado = true};
    \item ausência de pendência metrológica impeditiva (\texttt{existePendenciaMetrologicaTx(f) = false}) e ausência de qualquer bloqueio de M5/M6.
\end{itemize}
Consequências: \texttt{ABERTO + EXTRAVIADO} continua fisicamente aberto, mas \textbf{não é apto}; \texttt{QUEBRADO}, \texttt{VAZIO}, \texttt{DESCARTADO}, \texttt{EMPRESTADO} e \texttt{em\_quarentena} não são aptos; frasco vencido sem autorização não é apto. A retirada pode acrescentar pré-condições próprias (escopo, finalidade, aceite de uso vencido), mas \textbf{não} pode empregar uma noção diferente de aptidão. Um frasco extraviado nunca conta como apto, mesmo que a disponibilidade tenha sido alterada por engano.
\end{regra}

\begin{regra}[\texttt{qtdAptos}]
Para o par formado por almoxarifado, resumo e especificação, \texttt{qtdAptos} é a contagem dos frascos de \textit{Frasco\_Reagente} cujo \texttt{id\_almoxarifado}, \texttt{id\_resumo\_reagente} e \texttt{id\_especificacao\_reagente} coincidem com o par e para os quais \texttt{frascoAptoParaUso(f)} é verdadeiro.
\texttt{qtdAptos} é \textbf{número de frascos aptos}, não massa ou volume, e é a única base do critério institucional de escassez. Não confundir com o estoque quantitativo da dashboard (que soma massa/volume conhecidos) apenas porque ambos consultam \textit{Frasco\_Reagente}.
\end{regra}

\begin{regra}[Configuração canônica de escassez: \textit{Estoques\_Configurados}]
A configuração corrente é \texttt{Almoxarifado/\{almoxId\}/Estoques\_Configurados/\{configId\}}, com chave determinística \texttt{\{id\_resumo\_reagente\}\_\{id\_especificacao\_reagente\}}. O documento contém \texttt{id\_resumo\_reagente}, \texttt{id\_especificacao\_reagente}, \texttt{id\_almoxarifado}, \texttt{qtd\_limiar\_escassez}, \texttt{ativo} e \texttt{notificacao\_ativa}. \texttt{qtd\_limiar\_escassez} é inteiro \textbf{não negativo} e mede quantidade de frascos; não é massa nem volume. A escolha entre limiar em frascos e limiar quantitativo em g/mL está decidida por \textbf{frascos} e não pode ser alterada silenciosamente.

Semântica das flags:
\begin{itemize}
    \item \texttt{ativo = false}: o par não participa da avaliação automática de escassez.
    \item \texttt{notificacao\_ativa = false}: a escassez pode continuar existindo conceitualmente, mas o job \textbf{não emite} alerta automático; silenciar não apaga a configuração, não altera o estoque e não deve ser exibido como inexistência de estoque.
    \item \texttt{Almoxarifado.ativo = false}: o job não executa a avaliação automática para o almoxarifado; operações de encerramento ordenado continuam regidas pela subseção ``Almoxarifado inativo''.
\end{itemize}
\end{regra}

\begin{regra}[Definição de escassez e fronteira]
Há escassez para o par (almoxarifado, resumo, especificação) se, e somente se:
\texttt{qtdAptos} $<$ \texttt{qtd\_limiar\_escassez}.
A comparação é \textbf{estritamente menor}. Portanto \texttt{qtdAptos = qtdLimiar} \textbf{não} é escassez e \texttt{qtdAptos > qtdLimiar} também não. Casos de fronteira: limiar $0$ nunca gera escassez (não há \texttt{qtdAptos} $<$ 0); limiar 1 gera escassez exatamente quando há 0 frascos aptos; com 3 frascos aptos e limiar 3, não há escassez. A unidade é frascos aptos; o critério não é satisfeito por saldo g/mL.
\end{regra}

\begin{regra}[Backfill legado como pré-condição operacional]
A escassez depende de \texttt{id\_resumo\_reagente} e \texttt{id\_especificacao\_reagente}. Frascos legados sem \texttt{id\_especificacao\_reagente} não podem ser ignorados em produção sem reconciliação, sob pena de falso alerta de escassez pela omissão de frascos aptos. Caso determinístico admite backfill legítimo (por exemplo, especificação determinada por lote conhecido); caso ambíguo é pendência manual de reconciliação e \textbf{nunca} se inventa uma especificação para satisfazer o schema. Especificar o job \textbf{não} significa que o backfill já foi executado nem que a cron está ativada: a cron de escassez só é operacional após a conclusão da reconciliação.
\end{regra}

\begin{regra}[Cache de dashboard não é autoridade da escassez]
A dashboard de estoque pode usar \texttt{Sistema\_Cache\_Dashboard} com validade semântica de até $30$ segundos. O job de escassez \textbf{não} decide escassez a partir de \texttt{Sistema\_Cache\_Dashboard.resultado}: a autoridade da escassez é o estado canônico corrente de \textit{Frasco\_Reagente} consultado diretamente pelo job. Cache e escassez são contratos distintos.
\end{regra}

\begin{regra}[Notificações de escassez e idempotência]
O tipo de notificação é \texttt{ESCASSEZ\_ESTOQUE}, persistido em \texttt{Usuarios/\{uid\}/Notificacoes/\{id\}}. Os campos pertinentes são \texttt{id\_destinatario}, \texttt{papel\_destinatario}, \texttt{tipo}, \texttt{entidade\_alvo}, \texttt{id\_alvo}, \texttt{quantidade}, \texttt{lida}, \texttt{emitida\_em} e \texttt{expira\_em}. Destinatários são \textbf{apenas} os gestores efetivamente vinculados ao almoxarifado (via \textit{Gestor\_Almoxarifado\_x\_Almoxarifado}); possuir o papel global de Gestor de Almoxarifado sem vínculo não faz o usuário receber o alerta. \texttt{expira\_em} é \texttt{NULL} para \texttt{ESCASSEZ\_ESTOQUE}, que é alerta operacional persistente.

A idempotência segue o contrato global (Seção~\ref{sec:idempotencia-m7}) por \textbf{docId determinístico}, equivalente a \texttt{escassez\_\{almox\}\_\{config\}\_\{dataCivil\}}, criado com \texttt{create()} (ou mecanismo equivalente). É proibido usar \texttt{Date.now()}, UUID novo por execução ou \texttt{doc()} aleatório para alertas periódicos deduplicados. Reexecutar o job no mesmo dia é \textbf{no-op} para a chave já existente; somente \texttt{ALREADY\_EXISTS} (``code 6'') é tratado como duplicidade esperada e \textbf{erros de outra natureza não podem ser engolidos}: devem propagar para não perder alertas. Configuração no dia civil seguinte pode gerar nova notificação (a \textit{dataCivil} muda).

A \textit{dataCivil} é derivada do fuso IANA \texttt{America/Sao\_Paulo}; não usar offset fixo \texttt{-03:00} nem \texttt{setHours} dependente do fuso do processo.
\end{regra}

\begin{regra}[Notificações: leitura, escrita e ``Limpar tudo'']
O cliente lê apenas as próprias notificações; a escrita é server-owned. Marcar como lida é operação protegida de servidor. ``Limpar tudo'' \textbf{marca como lida} as notificações ativas do usuário, preservando o histórico; não é \texttt{DELETE} e não pode ser implementado por exclusão arbitrária. A taxonomia e o mecanismo de idempotência de \texttt{ESCASSEZ\_ESTOQUE} são compatíveis com os demais tipos já compartilhados (devolução preventiva, atraso, frascos vencidos, vazios, quebrados, em quarentena, a serem pesados e autoatendimento).
\end{regra}

\paragraph{Matriz de invalidação do estoque atual.}
Legenda: \textbf{S} = sim; \textbf{N} = não; \textbf{S*} = sim somente quando a operação altera a contribuição corrente (saldo, disponibilidade, localização ou validade); \textbf{C} = condicional a contrato suportado. \texttt{Alm.} = escopo \texttt{ESTOQUE\_\_ALMOX} (mais o registro de geração do almoxarifado); \texttt{Alm+R} = escopo \texttt{ESTOQUE\_\_ALMOX\_\_RESUMO}. ``Orig.$\to$dest.'' identifica os almoxarifados cujas chaves são invalidadas quando há transferência.

\setlength{\tabcolsep}{2pt}
\footnotesize
\begin{longtable}{>{\raggedright\arraybackslash}p{0.25\textwidth} >{\raggedright\arraybackslash}p{0.05\textwidth} >{\raggedright\arraybackslash}p{0.08\textwidth} >{\raggedright\arraybackslash}p{0.09\textwidth} >{\raggedright\arraybackslash}p{0.14\textwidth} >{\raggedright\arraybackslash}p{0.12\textwidth} >{\raggedright\arraybackslash}p{0.12\textwidth}}
\toprule
\textbf{Operação} & \textbf{Est.} & \textbf{Alm.} & \textbf{Alm+R} & \textbf{Orig.$\to$dest.} & \textbf{FLOW} & \textbf{Notif.} \\
\midrule
\endfirsthead
\toprule
\textbf{Operação} & \textbf{Est.} & \textbf{Alm.} & \textbf{Alm+R} & \textbf{Orig.$\to$dest.} & \textbf{FLOW} & \textbf{Notif.} \\
\midrule
\endhead
\bottomrule
\endfoot
Cadastro de frasco & S & S & S & --- & \texttt{CADASTRO} & N \\
Abertura & S & S & S & --- & \texttt{ABERTURA} & N \\
Retirada & S & S & S & --- & \texttt{SAIU}, empr. & N \\
Devolução normal & S & S & S & --- & \texttt{ENTROU} & N \\
Devolução anômala & S & S & S & --- & \texttt{ENTROU}, anom. & N \\
Resolução metrológica & S* & S* & S* & --- & datas afetadas & N \\
Esgotamento & S & S & S & --- & \texttt{FICOU\_VAZIO} & N \\
Pesagem ordinária & S* & S* & S* & --- & \texttt{AJUSTE} & N \\
Evaporação & S* & S* & S* & --- & \texttt{EVAPORACAO} & N \\
Ajuste & S* & S* & S* & --- & \texttt{AJUSTE} & N \\
Recalibração de tara & S & S & S & --- & \texttt{AJUSTE} & N \\
Quebra & S & S & S & --- & \texttt{QUEBROU} & N \\
Descarte & S & S & S & --- & \texttt{FOI\_DESCARTADO} & N \\
Extravio & S & S & S & --- & \texttt{EXTRAVIO} & N \\
Reencontro & S & S & S & --- & \texttt{REENCONTRO} & N \\
Entrada em quarentena & S & S & S & --- & quarentena & N \\
Saída de quarentena & S & S & S & --- & liberação & N \\
Autorização/revogação de uso vencido & S & S & S & --- & evento & N \\
Autorização de descarte & S & S & S & --- & \texttt{PENDENTE} & N \\
Vencimento pelo job & S & S & S & --- & \texttt{VENCEU} & S \\
Transferência/associação & C & C & C & origem e destino & conforme & N \\
Correção só histórica & N & N & N & --- & reprocessa & N \\
Job de escassez & N & N & N & --- & N & S \\
\bottomrule
\end{longtable}
\normalsize

A invalidação é transacional com a mutação e \textbf{não} recalcula a agregação no commit. As notificações automáticas efetivamente especificadas na V1 para o escopo M8 são \texttt{FRASCOS\_VENCIDOS} (job de vencimento) e \texttt{ESCASSEZ\_ESTOQUE}; os demais tipos existem na taxonomia compartilhada, sem emissão automática redesenhada nesta rodada.

\subsubsection{Condição inicial, saldo desconhecido e validade desconhecida}
\begin{regra}
\texttt{condicao\_inicial\_cadastro} é um fato histórico imutável do \textit{Frasco\_Reagente}, com enum \texttt{FECHADO} ou \texttt{JA\_ABERTO}: \texttt{FECHADO} indica que o recipiente entrou fisicamente fechado no controle do LCQUI; \texttt{JA\_ABERTO} indica que já estava aberto quando passou ao controle do LCQUI. Não se confunde com a modalidade de saldo desconhecido: \texttt{saldo\_desconhecido} é propriedade metrológica atual, enquanto \texttt{condicao\_inicial\_cadastro} é fato histórico que não muda depois.

\texttt{saldo\_desconhecido} representa o desconhecimento \textbf{atual} da quantidade remanescente e não é flag histórica. Frasco cadastrado \texttt{JA\_ABERTO} inicia com \texttt{saldo\_desconhecido = true}; abertura de FECHADO preserva a flag enquanto não houver informação metrológica suficiente; ao esvaziar e medir a tara real, torna-se \texttt{false}. Somente \texttt{VAZIO} confirmado pela rota metrológica válida resolve o desconhecimento (saldo zero, flag falsa). \texttt{QUEBRADO}, \texttt{DESCARTADO}, \texttt{EXTRAVIADO} e quarentena preservam \texttt{saldo\_desconhecido}, seja verdadeiro ou falso. Estado físico e conhecimento metrológico são dimensões distintas: quebra, descarte e indisponibilidade não são medição, não fabricam saldo zero, tara ou pesagem. Resolver desconhecimento exige evidência metrológica válida; retirar do estoque não determina a quantidade física anteriormente existente. \texttt{DESCARTADO} continua terminal operacionalmente.

\texttt{abertura\_historica\_desconhecida} é flag histórica: quando verdadeira, \texttt{data\_abertura = NULL} e \texttt{estado\_fisico\_frasco} $\neq$ \texttt{FECHADO}. Não é apagada apenas porque o estado atual mudou para \texttt{VAZIO}, \texttt{QUEBRADO} ou \texttt{DESCARTADO}, nem porque a localização mudou para \texttt{EXTRAVIADO}; significa que já houve abertura e a data da primeira abertura é desconhecida. Nunca se inventa a data.

\texttt{validade\_desconhecida} é usada também quando não se pode determinar a validade efetiva exata por ausência da data histórica de abertura (\texttt{abertura\_historica\_desconhecida = true} e \texttt{validade\_apos\_aberto\_dias}). Salvo outra informação suficiente, \texttt{validade\_desconhecida = true} e \texttt{validade\_efetiva = NULL}. \texttt{validade\_desconhecida = true} não implica \texttt{vencido = false}: pode haver prova de vencimento sem data exata (por exemplo, validade fechada conhecida e já expirada). Devem ser tratadas separadamente ``conhecer a data exata de validade'' e ``ser possível concluir que já venceu''. Não se usa \texttt{vencido = true} como mecanismo de quarentena; \texttt{em\_quarentena}, \texttt{vencido}, \texttt{estado\_fisico\_frasco}, \texttt{disponibilidade} e \texttt{saldo\_desconhecido} são dimensões distintas.
\end{regra}

\subsubsection{Justificativas humanas obrigatórias (lista fechada)}
\begin{regra}
A exigência de justificativa humana com no mínimo 20 caracteres após \texttt{trim} aplica-se \textbf{somente} aos seguintes contratos enumerados:
\begin{enumerate}
    \item entrada manual em quarentena;
    \item saída ou liberação de quarentena;
    \item decisão de quarentena durante cadastro vencido;
    \item decisão de quarentena na primeira abertura;
    \item extravio;
    \item reencontro;
    \item descarte institucional;
    \item recalibração de tara real;
    \item autoatendimento (regra existente);
    \item justificativa metodológica Q04/Q14 (preservando o máximo específico existente);
    \item correção administrativa ou operacional de ação do gestor.
\end{enumerate}
Essa exigência não é regra genérica para toda string de \textit{motivo}, \textit{detalhe} ou \textit{observacao}; mensagens automáticas do servidor são isentas. Não se impõe mínimo global de 20 caracteres em \texttt{detalhe\_status} persistido. Pesagem de rotina não entra automaticamente nessa lista.
\end{regra}

\subsubsection{Limites de Geração de Relatórios}
\begin{regra}
Para proteger o sistema contra execuções exaustivas e travamentos em ambiente \textit{serverless}, os relatórios personalizados possuem limite de período de no máximo 31 dias corridos. Além disso, nenhuma data de início ou término pode ser superior ao dia atual, e o mês/ano de relatórios fechados não podem estar no futuro.
\end{regra}

\subsubsection{Almoxarifado inativo: operações bloqueadas e encerramento ordenado}
\label{sec:regras-almoxarifado-inativo}
\begin{regra}
Quando \texttt{Almoxarifado.ativo = false}, as seguintes operações são bloqueadas pelo servidor:
\begin{itemize}
    \item cadastro de novos frascos e lotes naquele almoxarifado;
    \item novas retiradas (abertura de \textit{Emprestimo\_Reagente}) vinculadas àquele almoxarifado;
    \item vinculação de novos gestores ao almoxarifado inativo.
\end{itemize}
As seguintes operações permanecem permitidas para encerramento ordenado:
\begin{itemize}
    \item registro de devoluções de empréstimos ativos (\texttt{EM\_USO} ou \texttt{ATRASADO});
    \item descarte e atualização de status de frascos já cadastrados;
    \item geração de relatórios históricos.
\end{itemize}
A reativação do almoxarifado (\texttt{ativo = true}) é possível e não exige recriação da entidade. A desativação é auditada em \textit{Registro\_de\_Auditoria} com ação \texttt{DESATIVACAO\_ALMOXARIFADO}.
\end{regra}

\subsection{Matriz formal de obrigatoriedade e dependência de campos}

\small

\small
\setlength{\tabcolsep}{5pt}
\begin{longtable}{>{\raggedright\arraybackslash}p{0.25\textwidth} >{\raggedright\arraybackslash}p{0.27\textwidth} >{\raggedright\arraybackslash}p{0.40\textwidth}}
\toprule
\textbf{Entidade} & \textbf{Condição} & \textbf{Obrigatoriedade}\\
\midrule
\endfirsthead
\toprule
\textbf{Entidade} & \textbf{Condição} & \textbf{Obrigatoriedade} (continuação)\\
\midrule
\endhead
\bottomrule
\endfoot
Resumo Reagente & sempre & nome, tipo e natureza.\\
Estoque Mínimo Almox. & sempre & id\_especificacao\_reagente, id\_almoxarifado e limiar de escassez positivo (quantidade de frascos).\\
Resumo Reagente & requer pesagem frequente & frequência em dias positiva.\\
Especificação & SOLIDO & unidade = g; densidade opcional.\\
Especificação & LIQUIDO & unidade = ml; densidade positiva obrigatória.\\
Especificação & MISTURA & pelo menos uma composição válida.\\
Frasco & condição inicial & condicao\_inicial\_cadastro = FECHADO ou JA\_ABERTO; fato histórico imutável, distinto da condição operacional (em quarentena ou não).\\
Frasco & em quarentena & dimensão operacional separada; detalhe do status (motivo) obrigatório.\\
Frasco & validade conhecida & validade\_fechado e validade\_efetiva devem ser consistentes; se já expirada, exige decisão de destino.\\
Frasco & validade desconhecida & validade\_desconhecida = true e validade\_efetiva = NULL; não implica vencido = false.\\
Frasco & validade expirada & exige destino QUARENTENA, PENDENTE\_DE\_DESCARTE ou DISPONIVEL com uso\_vencido\_autorizado.\\
Frasco & FECHADO & Peso total positivo obrigatório; conteúdo nominal positivo quando conhecido, NULL quando desconhecido. Tara de referência somente com dados suficientes; sem nominal, tara NULL e saldo desconhecido.\\
Frasco & abertura & data\_abertura e validade\_efetiva calculadas no servidor; aplicar mínimo entre validade fechada e prazo após abertura quando ambos existirem.\\
Frasco & JA\_ABERTO & peso total atual medido e conteudo\_nominal quando conhecido/leg\'ivel (\texttt{NULL} se desconhecido, nunca zero); saldo\_desconhecido = true e tara n\~ao conhecida at\'e o esgotamento.\\
Patrimônio & sempre & id\_resumo, numero\_patrimonio, estado\_conservacao, id\_local, photo\_url, nome\_responsavel\_sei e status.\\
Patrimônio & Ja\_dado\_baixa & documento comprobatório obrigatório.\\
Requisição edição & criação & motivo e ao menos um campo proposto.\\
Requisição & aprovada/rejeitada & respondida\_em, respondente e justificativa obrigatórios.\\
Convite & turma definida & e-mail normalizado + turma.\\
Convite & sem turma & e-mail normalizado + confirmação explícita do convite global.\\
Turma & criação & nome, capacidade positiva e professor responsável.\\
\bottomrule
\end{longtable}

\subsection{Matriz de combinações Multi-Role}
As únicas combinações permitidas para uma mesma conta são:
\begin{itemize}
    \item Professor + Gestor de Bens Patrimoniais;
    \item Professor + Gestor de Almoxarifado;
    \item Professor + Gestor de Bens Patrimoniais + Gestor de Almoxarifado;
    \item Aluno + Bolsista;
    \item Aluno + Bolsista + Gestor de Bens Patrimoniais;
    \item Aluno + Gestor de Almoxarifado;
    \item Aluno + Gestor de Bens Patrimoniais;
    \item Aluno + Gestor de Bens Patrimoniais + Gestor de Almoxarifado.
\end{itemize}
DP-C01: Aluno sem cargo de Bolsista pode exercer Gestor de Almoxarifado; Bolsista e Gestor de Almoxarifado são mutuamente exclusivos. Não são permitidos Chefe Geral combinado com qualquer outro papel, Aluno + Professor, nem Aluno + Bolsista + Gestor de Almoxarifado. A concessão ou remoção de papéis deve ser transacional e auditada.

\subsection{Fluxo obrigatório para mutações críticas}
Toda operação que altera dados relevantes segue: autenticar, resolver papéis, validar escopo, reler estado atual no servidor, validar transição e constraints, calcular derivados, persistir atomicamente quando necessário, registrar histórico/auditoria e disparar notificações/materializações. O frontend nunca é a autoridade final.


\subsection{Atribuição e revogação de papéis}
\label{sec:regras-role}
\subsubsection{RN-ROLE-01 — Chefe Geral não pode possuir múltiplos papéis}
O papel Chefe\_Geral é mutuamente exclusivo com todos os demais papéis do sistema.
Um usuário que possua Chefe\_Geral não pode possuir simultaneamente:
Professor;
Aluno;
Bolsista;
Gestor\_Almoxarifado;
Gestor\_Bens\_Patrimoniais;
qualquer outro papel do sistema.
Consequentemente, um usuário com Chefe\_Geral nunca é considerado Multi-Role.

\subsubsection{RN-ROLE-02 — Chefe Geral não pode revogar o papel de outro Chefe Geral}
Um usuário com o papel Chefe\_Geral não pode alterar ou revogar o papel Chefe\_Geral de outro usuário que também possua esse papel.
Um Chefe Geral somente pode iniciar o processo de revogação do próprio papel.

\subsubsection{RN-ROLE-03 — Auto-revogação de Chefe Geral}
Um Chefe Geral pode revogar o próprio papel Chefe\_Geral somente quando existir pelo menos outro usuário ativo com o papel Chefe\_Geral.
A operação deve ser rejeitada quando o usuário for o único Chefe Geral ativo do sistema.
A finalidade dessa regra é impedir que o sistema fique sem qualquer usuário com capacidade de administração global.

\subsubsection{RN-ROLE-04 — Revogação de Chefe Geral desativa a conta}
Como Chefe\_Geral é um papel exclusivo, sua revogação sempre resulta na ausência de outros papéis para o usuário.
Nesse caso:
o papel Chefe\_Geral é removido;
os Custom Claims são recalculados;
a conta do usuário é marcada como desativada (ativo = false);
o usuário permanece armazenado no sistema;
seus registros históricos e de auditoria são preservados;
a conta não deve ser excluída fisicamente.

\subsubsection{RN-ROLE-05 — Almoxarifado deve possuir pelo menos um gestor}
Todo almoxarifado ativo deve possuir obrigatoriamente pelo menos um Gestor\_Almoxarifado responsável.
Um Gestor de Almoxarifado somente poderá ter seu papel revogado se, após a operação, cada almoxarifado que ele gerencia continuar possuindo pelo menos um outro Gestor de Almoxarifado.
Portanto, a operação deve verificar individualmente todos os almoxarifados vinculados ao gestor.
Se ele for o único gestor de qualquer um deles, a revogação deverá ser rejeitada.

\subsubsection{RN-ROLE-06 — Gestores de Almoxarifado não possuem exclusividade sobre movimentações}
Não existe exigência de que o mesmo Gestor de Almoxarifado que registrou uma movimentação precise posteriormente registrar a movimentação complementar.
Quando um almoxarifado possuir dois ou mais gestores:
qualquer um dos gestores autorizados poderá registrar a entrada de um frasco;
qualquer um dos gestores autorizados poderá registrar a saída;
qualquer um dos gestores autorizados poderá registrar a devolução;
qualquer um dos gestores autorizados poderá registrar outras movimentações para as quais possua permissão.
O fato de um gestor ter registrado a saída de determinado frasco não impede que outro gestor autorizado registre posteriormente sua entrada ou devolução.
A autoria da operação deve permanecer registrada para fins de auditoria.

\subsubsection{RN-ROLE-07 — Revogação de Gestor de Almoxarifado em usuário Multi-Role}
Quando um usuário possuir Gestor\_Almoxarifado e pelo menos outro papel válido, a revogação de Gestor\_Almoxarifado não desativa sua conta.
O sistema deve:
validar as regras de revogação dos almoxarifados;
remover Gestor\_Almoxarifado dos papéis do usuário;
recalcular seus Custom Claims;
manter os demais papéis;
manter a conta ativa;
remover da interface as funcionalidades exclusivas de Gestor de Almoxarifado.

\subsubsection{RN-ROLE-08 — Revogação de Gestor de Bens Patrimoniais em usuário Multi-Role}
Quando um usuário possuir Gestor\_Bens\_Patrimoniais e pelo menos outro papel válido, a revogação de Gestor\_Bens\_Patrimoniais não desativa sua conta.
O sistema deve:
verificar a condição necessária para a revogação;
remover o papel;
recalcular os Custom Claims;
preservar os demais papéis;
manter a conta ativa;
retirar da interface as funcionalidades exclusivas do Gestor de Bens Patrimoniais.

\subsubsection{RN-ROLE-09 — Gestor de Bens Patrimoniais não pode ser o único}
A revogação de Gestor\_Bens\_Patrimoniais somente será permitida quando existir pelo menos outro usuário ativo que continue exercendo esse papel.
Caso o usuário seja o único Gestor de Bens Patrimoniais ativo, a revogação deverá ser rejeitada.
A finalidade é impedir que o sistema fique sem nenhum responsável pela gestão patrimonial.

\subsubsection{RN-ROLE-10 — Revogação do último papel desativa a conta}
Quando um usuário possuir apenas um papel ativo e esse papel for revogado, sua conta deverá ser desativada.
Isso se aplica a todos os papéis. A conta não deve ser excluída.
O usuário deverá permanecer no sistema para preservação de auditoria, histórico, autoria de operações, etc.

\subsubsection{RN-ROLE-11 — Revogação de papel não significa exclusão do usuário}
A revogação de qualquer papel deve preservar a identidade histórica do usuário.
Fluxo correto: Revogar papel -> Validar regras -> Remover papel -> Recalcular conjunto -> Atualizar Custom Claims -> Possui outro papel? (Sim = ativo, Não = desativado).

\subsubsection{RN-ROLE-12 — Usuário Multi-Role}
Um usuário será considerado Multi-Role quando possuir dois ou mais papéis simultaneamente, exceto Chefe\_Geral.
A remoção de um papel de um usuário Multi-Role deve afetar somente aquele papel específico.

\subsubsection{RN-ROLE-13 — Usuário com apenas um papel}
Se o único papel de um usuário for revogado, o usuário ficará com 0 papéis e sua conta será desativada, mas o usuário permanece armazenado para fins históricos e de auditoria.

\subsubsection{RN-ROLE-14 — Recalculo dos Custom Claims}
Os Custom Claims devem representar o conjunto efetivo de papéis do usuário após a operação. O backend deve recalcular o conjunto completo de papéis autorizados e então atualizar os Custom Claims. Se o conjunto resultante for vazio, a conta deverá ser marcada como desativada.

\subsubsection{RN-ROLE-15 — Auditoria das alterações de papéis}
Toda atribuição e toda revogação de papel devem gerar registro de auditoria contendo: usuário afetado, usuário que executou a operação, papel adicionado/removido, data/hora, resultado, justificativa e situação da conta. Revogações rejeitadas também devem ser registradas.

\subsubsection{RN-ROLE-16 — Invariante Bolsista}
A revogação do papel Aluno para um usuário é TERMINANTEMENTE BLOQUEADA pelo backend caso ele possua o papel Bolsista ativo. A revogação do papel de Bolsista deverá ocorrer de maneira prévia ou simultânea à perda do vínculo discente.

\paragraph{Execução concorrente e sincronização de identidade.}
A verificação do último chefe, do último gestor patrimonial e de cada almoxarifado deve compartilhar um documento de controle transacional; duas revogações concorrentes não podem ambas remover o último responsável. Revogar Aluno exige remover Bolsista na mesma operação ou rejeitar a solicitação. O modal explicita o conjunto resultante antes de confirmar. A gravação dos papéis e da auditoria ocorre no Firestore; a atualização do Firebase Auth é uma etapa externa, recuperável e idempotente. Uma falha nessa etapa mantém a autorização bloqueada até a reconciliação, sem anunciar conclusão prematura.

\paragraph{Q04 e Q14: contrato de aceite e autoatendimento.}
Para PESQUISA\_TCC\_POS ou ESTUDO\_DEGRADACAO\_RESIDUOS com frasco vencido ou validade desconhecida liberado com termo, obter aceite em sessão autenticada do retirante. Persistir tcr\_versao, tcr\_aceito\_em (servidor), tcr\_aceito\_por e tcr\_auditoria\_id no empréstimo, com evento correspondente em Registro\_de\_Auditoria; exigir justificativa\_metodologica de 20--2000 caracteres, verificando o mínimo após trim. Modal/checkbox são UX: backend valida versão vigente, identidade e vínculo do aceite ao frasco, finalidade e operação antes de confirmar a retirada. O TCR registra ciência e responsabilidade e não substitui regras institucionais de segurança química. Autoatendimento só quando não houver outro gestor ativo, calculado pelo servidor com justificativa humana obrigatória de pelo menos 20 caracteres após trim e notificação à chefia. Registrar abertura\_historica\_desconhecida quando a data anterior não for conhecida; nunca inventar a data.

\paragraph{Retenção V1 (DP-D02).}
Retenção indefinida de históricos e auditoria é política conservadora de rastreabilidade do LCQUI/UENF. A desativação lógica (\textit{soft delete}) é o padrão do sistema. A exclusão física de registros deve ocorrer somente quando houver fundamento normativo/operacional aplicável e ausência de referências impeditivas, não devendo a LGPD ser utilizada como justificativa genérica ou única base para deletar dados transacionais e históricos.

\paragraph{Proteção e Escopo do localStorage.}
O uso do \texttt{localStorage} para persistência de rascunhos de wizards deve respeitar rigorosamente os seguintes controles:
\begin{itemize}
    \item \textbf{Namespace e Segurança:} Utilizar o UID como namespace das chaves, embora isso não garanta isolamento contra XSS/same-origin. Tokens, PII, documentos fiscais completos e anexos jamais devem ser armazenados.
    \item \textbf{Ciclo de Vida:} Aplicar TTL (Time to Live) e verificá-lo tanto no \textit{boot} quanto na leitura. Ocorrerá limpeza compulsória após o sucesso do fluxo (salvamento no backend), no logout e em qualquer troca de identidade.
    \item \textbf{Robustez:} Utilizar \textit{schema versioning} e tratar cenários de múltiplas abas, limitando o escopo estritamente aos dados de catálogo (\textit{whitelist} de formulários autorizados).
\end{itemize}

\paragraph{Concessão e contadores de papéis (AUDITORIA-7, C008).}
Toda concessão de papel é transacional e compartilha \texttt{Controle\_Papeis/singleton} com a revogação. Ao criar papel ativo de Chefe Geral ou Gestor de Bens Patrimoniais, incrementar respectivamente \texttt{chefes\_ativos} ou \texttt{gestores\_patrimoniais\_ativos}, na mesma transação da identidade, papel e auditoria. Repetir concessão já efetivada não incrementa novamente. Inicialização e reconciliação dos contadores precedem o uso; Auth/Custom Claims são sincronizados depois do commit.

\begin{regra}[Decisões humanas HQ-M2-004 a HQ-M2-007 -- RESOLVED]
Cadastro FECHADO com conteúdo nominal desconhecido é permitido, com tara NULL e saldo desconhecido; abrir ou pesar bruto não resolve a ausência de informação. O ciclo de vida sem tara segue JA\_ABERTO em situação equivalente. Elimina-se integralmente o antigo limiar fixo de devolução/esgotamento; Q06 dinâmica permanece restrita à relação retorno/saída. Uma tara real anterior nunca é substituída silenciosamente: o esgotamento confirmado mantém saldo zero, e a discrepância segue recalibração auditável. Medição inconsistente é preservada com devolução física concluída, bloqueio operacional e consumo pendente até resolução.

Os dois resumos diários de reagentes são FLOW-only, independentes de posição anterior; correções fechadas reprocessam somente datas diretamente afetadas. Frasco\_Reagente é autoridade atual, acessada por agregações server-side protegidas e cache lazy de validade máxima 30 segundos ou invalidação anterior. Pesagem de rotina tem observação opcional, sem mínimo humano obrigatório; operações especiais conservam seus contratos. M2.2 não é iniciado por esta reconciliação documental.
\end{regra}

% =================================================
\section{Descrição das telas (Dashboards)}
% =================================================
\subsection{Header em comum em todos os papeis}
\begin{verbatim}
|----------------------------------------------------------------------|
|Olá, [Nome completo do usuário]                       [foto de perfil]|
|----------------------------------------------------------------------|
\end{verbatim}
Quando se clica na foto de perfil, aparece um modal que se arrasta da direita. Nele é possível deslogar e ir para "Perfil de usuário", que por sua vez abre um modal com as informações do usuário e botão para editar (foto, nome, etc.).

\subsection{Header quando o usuário é Multi-Role}
\begin{verbatim}
|----------------------------------------------------------------------|
|Olá, [Nome]  [Botão Papel 1] [Botão Papel 2]          [foto de perfil]|
|----------------------------------------------------------------------|
\end{verbatim}
Alterar os botões de papel no header modifica instantaneamente a página toda e o painel lateral para as funcionalidades daquele papel selecionado (ex: de Aluno para Gestor de Almoxarifado).

\subsection{Dashboard-Chefe Geral}
O Chefe Geral tem uma div \textit{full-width} logo abaixo do header principal, alterando conforme a aba selecionada no painel esquerdo:

\subsubsection{Aba: Almoxarifados}
Antes da área principal de busca, ele verá uma div \textit{full-width} com:
\begin{itemize}
    \item Um botão "Mais" que ao clicar terá as opções:
     \begin{itemize}
         \item Novo Almoxarifado: Abre um modal para criar um almoxarifado.
         \item Novo Gestor de Almoxarifado: Abre um modal em que ele pode escolher: transformar um aluno em gestor de almoxarifado ou criar um Usuário que é Gestor de almoxarifado.
         \item Imprimir Etiquetas: Modal bifurcado em duas abas. A Aba 1 (Novos Frascos / Virgens) oferece opções de intervalo e um Grid Visual de Offset 3×10 para escolha da célula inicial na folha A4. A Aba 2 (Frascos Cadastrados / 2ª Via) permite busca por reagente/código e seleção de até 10 frascos, gerando a respectiva Ficha de Conferência.
         \item \textbf{Nova Matéria:} Abre modal com campos Nome e Código da Matéria. O código deve ser único (validado na Cloud Function).
     \end{itemize}
    \item Um botão para "Gerenciar Gestores de Almoxarifado". Esse botão abre um novo modal que:
    \begin{itemize}
        \item Permite pesquisar e selecionar Gestores. Ao selecionar, aparece o histórico dele nós ultimos 30 dias com: (Reagentes e frascos adicionados, movimentações registradas), caso ele queira de um mês específico, terá um modal para escolher o ano e o mês.
    \end{itemize}
    \item Um botão de "Relatórios" que abre um modal com as opções para Relatórios em PDF dos almoxarifados.
\end{itemize}
Abaixo, haverá uma aba de busca onde pode-se filtrar por almoxarifado, reagente (ou ambos) e pesquisar livremente um reagente pelo nome. Ao clicar em um reagente, abre-se uma tela de detalhes:
O filtro padrão da tela de detalhes é o estado atual, mas pode ser alterado para um período específico. Ele exibe: Todas as propriedades em Resumo de reagente, inclusive as materializadas (Materialized View).

\subsubsection{Aba: Bens Patrimoniais}
Antes da aba principal de busca, ele verá uma div \textit{full-width} com:
\begin{itemize}
    \item Um botão "Mais" com as opções:
     \begin{itemize}
         \item Novo Gestor de Bens Patrimoniais.
     \end{itemize}
    \item Um botão para "Gerenciar Gestores de Bens Patrimoniais", abrindo um modal que:
    \begin{itemize}
        \item Permite pesquisar (por nome) um gestor. Ao selecionar, visualiza-se o histórico dos últimos 30 dias, ou um mês e ano esepecífico que for escolhido (requisições aprovadas/rejeitadas, bens adicionados, edições feitas).
    \end{itemize}
    \item Um botão de "Relatórios" para gerar os PDFs de bens patrimoniais (mensal, por prédio, geral).
\end{itemize}
Abaixo, a aba de busca permite filtrar por responsável do bem, estado de conservação, e pesquisa livre pelo número/nome do bem. Ao clicar em um patrimônio, abre-se o detalhamento:
Mostra o estado atual, quantos itens existem daquele tipo cadastrados, a divisão por status (Ativo, Inservível, Baixado) e a lista completa desses itens exibindo o número de patrimônio e o respectivo responsável logado no Sistema da UENF.

Ao clicar em um número de patrimônio especifico, abre um modal com as propriedades desse patrimônio, ao mudar o status dele para já dado baixa, é preciso fazer o upload do documento pdf que comprove isso.
\subsubsection{Aba: Professores}
Antes da aba de busca, uma div \textit{full-width} com:
\begin{itemize}
    \item Um botão para "Novo Professor".
\end{itemize}
Na área principal ele pode pesquisar Professores pelo nome, filtrar por Matéria lecionada, Centro e Laboratório (inclusive combinando múltiplos filtros). Ao clicar em um professor, abre-se um modal para gerenciar a conta daquele docente, ver seu histórico de ações e a lista de turmas que ele administra.

\subsubsection{Aba: Alunos}
Antes da aba de busca, uma div \textit{full-width} com:
\begin{itemize}
    \item Um botão para "Novo Aluno".
\end{itemize}
Na área principal ele pesquisa Aluno pelo nome, podendo filtrar por Matéria que cursa ou número de matrícula. Ao clicar no Aluno, abre-se um modal de gerenciamento exibindo seu histórico e as turmas em que está matriculado.

\subsection{Dashboard-Professor}
Para o papel de Professor, a dashboard foca em organização acadêmica e gestão de recursos em uso:
\begin{itemize}
    \item \textbf{Painel Lateral:} Fica do lado esquerdo, podendo ser expandido ou mantido em modo compacto. Contém dois grupos:
    \begin{itemize}
        \item \textbf{Grupo 1 (Visualizar):} Opções para acessar a área de "Reagentes" e "Bens Patrimoniais".
        \item \textbf{Grupo 2 (Turmas):} Lista navegável de todas as turmas ativas dele. Abaixo da lista, um botão "Turmas Arquivadas" (abre um modal com as turmas passadas exibindo Nome, Ano/Semestre e botão para Desarquivar). Ao desarquivar, a turma volta para a lista de turmas tanto do professor, quanto dos alunos que fazem parte dela.
    \end{itemize}
\end{itemize}

Na aba principal, abaixo do header em comum, ele verá uma div \textit{full-width} com:
\begin{itemize}
    \item \textbf{Notificações:} Botão com ícone de sino que abre um modal \textit{dropdown} alertando sobre requisições de bens respondidas e sobre reagentes emprestados cujo prazo está atrasado, vence hoje ou vence amanhã. Inclui um botão "Limpar tudo".
    \item \textbf{Botão Mais (+):} Abre \textit{dropdown} com: Nova Turma, Novo Aluno (turma opcional, com confirmação de convite global), Novo Roteiro de Experimento, Novo Pedido de Adição de Bem Patrimonial (se o usuário selecionar "Novo Local" dentro do pedido, um modal sobreposto exigirá o cadastro prévio do Prédio/Sala), e \textbf{Nova Matéria} (abre modal simples com nome e código; código deve ser único, validado pelo backend).
    \item \textbf{Gerenciar Roteiros:} Abre um modal com três abas ("Meus", "Compartilhados comigo" e "Eu compartilhei"). Permite abrir para visualização ou baixar os arquivos em PDF.
    \item \textbf{Minhas Requisições:} Modal para acompanhar os pedidos feitos de adição/edição de Bens Patrimoniais (filtrando por Pendente, Aprovada, Rejeitada).
    \item \textbf{Meus Reagentes:} Modal de acompanhamento onde o professor visualiza os reagentes atualmente sob sua posse (ordenados pela data de devolução) e o histórico de uso (Padrão: últimos 30 dias). Ao clicar num reagente histórico, ele visualiza exatamente o volume (mL) que ele consumiu nesse período.
\end{itemize}

\textbf{Área Central (Feed da Turma):}
Ao acessar a dashboard inicialmente, uma mensagem dirá: "\textit{Selecione uma turma ao lado para começarmos}".
Ao clicar em uma turma no painel lateral, a área central exibe o Feed dessa disciplina. Existe um bloco fixo superior para criar novos posts (podendo anexar PDFs novos ou selecioná-los da sua biblioteca/compartilhados). O feed lista os posts em ordem cronológica e cada post contém a área de comentários integrada para interação direta com os alunos.

\subsection{Dashboard-Gestor de Almoxarifado}
Para o perfil de Gestor de Almoxarifado, a interface é desenhada para a velocidade de atendimento de balcão e auditoria rigorosa de insumos:
\begin{itemize}
    \item \textbf{Painel Lateral:} Grupos de navegação para "Estoque" (Reagentes e Frascos), "Movimentações" (Empréstimos/Retornos), e "Descartes/Alertas".
\end{itemize}

Na aba principal, uma div \textit{full-width} com botões rápidos essenciais:
\begin{itemize}
    \item \textbf{Alertas de Gestor (Notificações):} Sino de avisos exclusivo exibindo pendências críticas: frascos declarados como vazios, vencidos ou quebrados, ou seja, aguardando mudança para o estado DESCARTADO, e reagentes que atingiram nível de escassez no estoque, frascos que precisam de pesagem de rotina, frascos que estão em quarentena.
    \item \textbf{Botão Mais (+):}
    \begin{itemize}
        \item Novo Reagente -> Abre modal exigindo dados de resumo e especificação separados conforme UI-05; CAS e concentração têm obrigatoriedade condicional, não universal. Campos: Nome, CAS, Concentração, grau de pureza (texto opcional), a quantidade em que ele é considerado escasso para alertas futuros, fórmula química , se é controlado pela pf ou pela eb, a classe de inflamabilidade, link da ficha de segurança, tipo (sólido ou líquido: pede densidade, pergunta se requer pesagem frequente, se sim, pede a frequencia em dias).
        Tem um seletor para escolher entre sólido e líquido (em temperatura ambiente):
        \begin{itemize}
            \item Para tipo sólido -> Sempre que for adicionar um frasco -> o tipo de medida é gramas (g) (setado pelo sistema)
            \item Para tipo líquido -> Sempre que for adicionar um frasco -> o tipo de medida é mililitros (ml) (setado pelo sistema)
        \end{itemize}
        Tem um seletor para a Natureza do Reagente: Orgânico, Inorgânico, Elemento ou Híbrido

        \item Adicionar Frasco: abre um modal para escolher o reagente a ser adicionado, o reagente pode ser pesquisado, por: nome, cas e formula química (deve ter um seletor para escolher, isso serve para não precisar pesquisar em todas as categorias de uma vez). Após selecionar o reagente a ser adcionado, fecha o modal de escolher reagente e passa a preencher o formulário para adicionar o frasco de reagente, o fórmulário deve conter:
         Selecionar o almoxarifao em que esse frasco vai ficar;\\
         Escrever o detalhamento do local (como: Segunda prateleira, na direita, etc) \\
         O número do lote em que esse frasco pertence. \\
         Pergunta o status dele: Fechado, Aberto, Quarentena (em caso de quarente, aparece um campo em texto para ser escrito explicando o motivo de estar em quarentena)

         \textbf{Feedback de Sucesso:} Ao salvar, exibe: "\textit{Frasco cadastrado sob o código LCQUI-XX (Reagente: YY, Data: DD/MM/AAAA). Localize a etiqueta física correspondente na folha de bancada, preencha os dados à caneta e aplique no frasco com fita adesiva.}"
         O campo estado físico no Resumo do reagente que determina os campos a serem pedidos abaixo (sólido ou líquido) \\
        líquido ->
\begin{itemize}
    \item Se é do tipo fechado, pede apenas a volume do líquido e o peso total com o frasco. (calcula o peso do frasco, mostra o peso do frasco mas não pode ser editado.)
    \item Se é do tipo Aberto, exige peso total e uma modalidade: tara conhecida ou estimativa de volume, conforme UI-06.
        \end{itemize}
        sólido -> exige peso total e massa nominal (fechado), ou tara conhecida/estimativa de massa (aberto), conforme UI-06.

        \item Imprimir Etiquetas: Modal bifurcado em duas abas. A Aba 1 (Novos Frascos / Virgens) oferece opções de intervalo (a partir do último cadastrado, após última impressão ou personalizado) e um Grid Visual de Offset 3×10 para escolha da célula inicial. A Aba 2 (Frascos Cadastrados / 2ª Via) permite busca por reagente/código e seleção de até 10 frascos, gerando a respectiva Ficha de Conferência.

    \end{itemize}
    \item \textbf{Ações de Bancada (Destacados no centro):}
    \begin{itemize}
        \item \textbf{Registrar Retirada:} Modal para escanear/buscar o frasco, selecionar o usuário retirante (Professor ou Bolsista, conforme as regras de autorização) e anotar o Peso Inicial.
        \item \textbf{Registrar Devolução:} Modal crucial. O gestor informa o frasco retornado e insere o \textbf{Peso de Retorno} da balança. O sistema faz a diferença, aplica a densidade (se for líquido), e mostra em negrito o \textit{Volume Ultilizado em mL}, ou então o \textit{Peso Ultizado em gramas}, atualizando medida no frasco e seus status automaticamente antes de fechar o modal.
    \end{itemize}
    \item \textbf{Relatórios:} Acesso à geração de PDFs mensais.
\end{itemize}
\textbf{Área Central de Busca:} Tabela geral pesquisável por nome do reagente, status do frasco ou almoxarifado físico. Ao clicar no reagente, o modal apresenta o dashboard do químico: volume total disponível, gráfico de consumo acumulado, e lista dos frascos com localização exata e status atual.

\subsection{Dashboard-Gestor de Bens Patrimoniais}
Para o Gestor de Bens Patrimoniais, o foco é em auditoria, análise comparativa e suporte processual institucional:
\begin{itemize}
    \item \textbf{Painel Lateral:} Grupos divididos em "Patrimônio Geral", "Requisições de Professores" e "Alertas importantes" (em alertas ele vai ter os patrimônios que estão em estado inservível para poder dar baixa- isso é externo ao sistema.)
\end{itemize}
Na div \textit{full-width} superior:
\begin{itemize}
    \item \textbf{Notificações:} Alerta sempre que um professor abrir uma requisição de adição de bem novo ou alteração do estado de um bem existente.
    \item \textbf{Botão Mais (+):} Adição direta de bens (ignorando requisição) e cadastro de novas descrições/resumos de itens; cadastro de gestores é exclusivo do Chefe Geral.
    \item \textbf{Analisar Requisições (Destaque):} Abre uma interface de tela dividida (Split-screen). Do lado esquerdo, os dados e foto do bem como estão atualmente no banco de dados. Do lado direito, os novos dados propostos pelo professor destacado em cores (ex: Estado de Conservação passando de Bom para Ruim). No rodapé, botões massivos de "Aprovar Alteração" ou "Rejeitar" (exigindo campo de texto com justificativa).
    \item \textbf{Gerar Relatórios:} Modal para compilação e download de PDF. Foco em permitir gerar o relatório de "Baixa de Bens Inservíveis" (contendo responsáveis) para anexar diretamente no processo interno (SEI) da Universidade.
\end{itemize}
\textbf{Área Central de Busca:} Pesquisa altamente filtrável (Prédio, Andar, Sala, Nome, Responsável, Status). O clique no bem patrimonial exibe, além das propriedades atuais do equipamento (foto, conservação, etc.), um histórico de auditoria vertical, listando quem modificou, a hora e quais campos foram alterados ao longo do tempo. Também é possível modificar manualmente um equipamento para "Ja\_dado\_baixa" caso o trâmite no SEI tenha finalizado.

\subsection{Dashboard-Aluno}
Para o perfil de Aluno, a interface remove completamente as complexidades de gestão, focando apenas no acesso aos materiais acadêmicos:
\begin{itemize}
    \item \textbf{Painel Lateral:} Um visual limpo apresentando a lista de "Minhas Turmas". No rodapé desse painel, um input text estilizado e claro dizendo "Código da Turma" com um botão "Entrar", permitindo vínculo imediato a uma nova disciplina sem modais complexos. Todo aluno tem acesso de leitura ao catálogo de reagentes conforme a seção 3. Quando também for bolsista, ele mantém o grupo Visualizar/Reagentes e ganha acompanhamento dos próprios empréstimos.

\end{itemize}
Na aba principal, abaixo do header:
\begin{itemize}
    \item \textbf{Notificações:} Sino simples indicando novos roteiros de práticas postados, ou professores que responderam às dúvidas deixadas nos comentários.
\end{itemize}
\textbf{Área Central (Visualização da Disciplina):}
\begin{itemize}
    \item Caso nenhuma turma esteja selecionada no menu lateral, a área exibe um aviso centralizado e amigável: "\textit{Selecione uma turma ao lado para ver as postagens}".
    \item Ao clicar numa turma, a tela exibe o Feed do aluno. Os posts feitos pelo professor aparecem em forma de \textit{cards} cronológicos.
    \item Se o post contém um Roteiro de Experimento anexado, um ícone de PDF chamativo aparece ao lado do nome do arquivo com um botão de ação "Baixar" (Download direto).
    \item Abaixo de cada post, a seção de comentários. O aluno vê um campo simples de \textit{input} para digitar uma dúvida ou observação, que, ao ser enviada, fica visível imediatamente na \textit{thread} daquele post para o professor e demais colegas da sala.
\end{itemize}

\subsection{Contratos detalhados das telas e modais}
\label{sec:contratos-telas}
Os contratos UI abaixo complementam os painéis anteriores, preservando sua organização. Cada abertura de modal, gaveta, página ou menu descrita acima usa o contrato correspondente. Campos não marcados como opcionais são obrigatórios; limites e enums do dicionário da seção~\ref{subsec:mapeamento-colecoes} também se aplicam ao frontend. As regras seguintes especificam o comportamento esperado e não atestam implementação.

\subsubsection{UI-01: acesso, sessão, perfil e troca de papel}
Login apresenta e-mail (tipo email, até 150 caracteres, espaços externos removidos), senha obrigatória com alternância de visibilidade, Entrar, Login com Google e Recuperar senha. Não aplicar trim à senha. Erro de credencial é genérico, sem confirmar a existência da conta. Recuperação pede somente e-mail e apresenta a mesma confirmação para endereços existentes ou não. Login não concede papel: após autenticação, carregar conta ativa e permissões antes de montar a dashboard; conta desativada vai para aviso de acesso desativado, conta sem papel para acesso não autorizado. Falha de rede oferece Tentar novamente.

A foto abre gaveta com nome, e-mail somente leitura, papéis, Perfil e Sair. Perfil contém nome obrigatório (1--150 caracteres), foto opcional com prévia, Remover foto, Cancelar e Salvar. A imagem deve ter tamanho inferior a 5 MiB; o servidor valida formato e propriedade do objeto. E-mail, matrícula e papéis não são editáveis nesse formulário. Atualizar nome propaga as projeções de busca, preservando eventos históricos. Sair encerra listeners, descarta dados da sessão, libera URLs temporárias e retorna ao login.

O seletor de papel exibe somente papéis autorizados. Ao trocar, interromper consultas do painel anterior, limpar seleções e carregar o novo escopo; um formulário modificado exige decidir entre continuar editando ou descartar. O papel preferido é apenas preferência visual. Não constitui autorização. Após alteração de permissões, renovar o token e verificar conta ativa; se o papel corrente deixou de existir, escolher outro permitido ou encerrar a sessão.

\paragraph{Padrão visual, acessibilidade e estados.}
Modal central: \texttt{fixed inset-0 z-50}, fundo escurecido, conteúdo \texttt{w-full max-w-2xl max-h-[90vh] overflow-y-auto}; formulários usam uma coluna no celular e duas em telas largas. Gaveta: largura total no celular e painel lateral no desktop. Manter título acessível, foco contido, retorno do foco ao acionador, rótulos ligados aos campos, erros textuais com \texttt{aria-describedby}, resumo de erros e foco no primeiro inválido. Escape e Cancelar fecham quando não há envio em andamento. Botões têm foco visível e texto além da cor. Salvar fica indisponível durante envio; sucesso só após resposta do servidor. Em conflito, preservar entradas e oferecer recarregar/revisar. Nunca repetir silenciosamente uma movimentação cuja resposta se perdeu: consultar seu identificador idempotente. Listas apresentam carregamento, vazio, erro, sem permissão e dados desatualizados como estados distintos.

\subsubsection{UI-02: usuários, gestores e atribuição de papéis}
Novo Professor, Novo Aluno e Novo Gestor compartilham formulário do Chefe Geral: escolher usuário existente por pesquisa autorizada ou convidar por e-mail; nome e e-mail são obrigatórios na criação. Professor exige centro (CCT, CCTA, CBB ou CCH) e laboratório até 100 caracteres; Aluno exige matrícula única até 20 caracteres no cadastro concluído. Convite pode antecipar matrícula opcional, exigida na aceitação se ainda ausente. Matérias são seleção opcional de IDs existentes, sem duplicatas. Gestor de almoxarifado aceita vínculos opcionais; sem vínculos, exibir estado vazio com orientação para procurar a chefia. Não converter Aluno em Professor implicitamente.

A lista de papéis mostra combinações permitidas, motivo de bloqueio e conjunto resultante. Conceder Bolsista exige Aluno e ausência de Gestor de Almoxarifado. Revogar abre confirmação com usuário, papel, vínculos afetados, justificativa obrigatória e consequência para a conta. Aplicar RN-ROLE-01 a RN-ROLE-15, incluindo autorrevogação e proteção do último responsável. Não oferecer revogar outro Chefe Geral. O backend valida tudo novamente e registra inclusive rejeições.

Detalhes de professor/aluno/gestor têm abas Dados, Papéis e vínculos, Turmas e Atividade. Nome, matrícula, laboratório e centro seguem os mesmos limites; somente campos autorizados são editáveis. Atividade inicia nos últimos 30 dias e consulta histórico paginado por autor e instante; o seletor mensal exige mês 1--12 e ano não futuro e usa os agregados mensais. Matrícula e e-mail não aparecem na lista de colegas. Filtros por matéria usam os vínculos Professor--Matéria e Aluno--Turma--Matéria. Q01: lotação física por prédio/gabinete/sala foi postergada à V2. Na V1, filtros de professor são Nome, Matéria, Centro e Laboratório.

\subsubsection{UI-03: almoxarifado, local e matéria}
Novo Almoxarifado exige nome (100), descrição (500), Local existente e pelo menos um gestor para ativação. Seleção múltipla lista apenas gestores ativos. Pode ser salvo inativo sem gestor; ativar exige vínculo válido. Gerenciar vínculos mostra todos os responsáveis e impede remover o último de um almoxarifado ativo, inclusive numa revogação global.

Novo Local exige prédio (30), andar (10, texto para aceitar térreo/subsolo) e sala (30). Remover espaços externos, rejeitar campos vazios e verificar unicidade normalizada da combinação no servidor. Ao salvar, retornar o ID e selecionar o local no formulário de origem, preservando o restante do pedido. Evitar dois diálogos simultâneos com foco concorrente: usar etapa interna ou suspender o diálogo pai. Autoria e data vêm da sessão e do servidor.

Nova/Editar Matéria pede nome (100) e código (10), ambos não vazios; normalizar código de forma consistente e verificar unicidade transacional. Editar mantém o ID e atualiza projeções; turmas não são recriadas. Chefe e Professor acessam essa ação; selecionar matéria existente não exige cadastrar novamente.

\subsubsection{UI-04: pesquisa, catálogo e detalhes}
Reagentes exigem pelo menos um filtro entre letra, natureza, tipo de substância e estado físico. Patrimônio exige prédio, letra ou status. Usar debounce de 300 ms como parâmetro inicial de interface, cursor estável e páginas de 25 registros. Busca textual local informa ``filtrando os itens carregados''; zero correspondências numa página não significa ausência no catálogo. Oferecer carregar próxima página e refinar filtros. Busca exata por LCQUI-N, CAS, matrícula ou plaqueta tem modo próprio, valida o formato e consulta índice/chave específica, sem varrer toda a base. Não abrir automaticamente detalhes enquanto o usuário ainda digita.

Detalhes de reagente mostram resumo, especificações separadas, composição, densidade/unidade, fabricante, controle PF/EB e link FDS; frascos são paginados por almoxarifado e estado. Mostrar saldo em g separado de mL, fechados/abertos/em uso e horário da materialização. Período histórico usa resumos diários; o estado atual usa documentos atuais. Selecionar frasco abre código, localização, lote opcional, validade, pesagem, disponibilidade e histórico. Emprestado não revela dados do retirante a quem só consulta catálogo. Ações de gestão exigem papel e vínculo; Professor/Bolsista consultam seus empréstimos, e Aluno tem leitura de catálogo conforme seção 3, sem retirada.

Patrimônio apresenta resumo catalográfico e itens físicos, foto, plaqueta, local, conservação e status separados. Responsável SEI é texto, não conta LCQUI. Detalhe individual usa ID e histórico paginado; resumo agrega contagens por status. Professor tem Solicitar edição, gestores e Chefe têm Editar e Baixa. Listas administrativas de pessoas usam respostas com campos mínimos e filtros de escopo; não baixar a coleção de identidades para simular busca.

\subsubsection{UI-05: resumo químico, substância, especificação e lote}
Novo Reagente é um assistente. Primeiro pesquisar resumo existente; criar somente se a identidade química nominal for diferente. Resumo exige nome (150), estado SOLIDO/LIQUIDO, higroscopicidade, PURA/MISTURA, natureza (ORGANICO, INORGANICO, ELEMENTO, HIBRIDO), limiar inteiro positivo de frascos e indicador de pesagem frequente; quando marcado, frequência inteira positiva em dias. Alterar esses dados exige confirmação do alcance sobre todas as especificações.

A segunda etapa seleciona ou cria especificação: descrição (150) e classe de inflamabilidade obrigatórios; estado herdado do resumo; unidade somente leitura (g/ml). Densidade positiva obrigatória para líquido; opcional para sólido. Fabricante (150), código do produto (100), grau de pureza (30) e FDS HTTPS (255) são opcionais. Grau de pureza é texto (ex.: P.A.), não percentual obrigatório. Controle PF/EB usa dois booleanos explícitos. Densidade confirmada não é editável: registrar nova especificação e preservar cálculos históricos.

Composição permite adicionar substâncias por nome/CAS, sem repetir ID. Nova substância pede nome (150), CAS opcional (15) e fórmula opcional (50); CAS informado deve ter formato e dígito verificador válidos e ser único. Mistura exige ao menos um componente. Tipo de concentração é enum da seção 4; valor\_min/valor\_max numéricos não negativos e ordenados seguem Q03, com notacao\_original\_fabricante opcional. A validação ocorre também no servidor. Não exigir CAS único de uma mistura como se fosse substância individual.

Lote é opcional no frasco. Novo Lote exige especificação, fornecedor (80), número (100), nota fiscal (100), aquisição, fabricação, validade e quantidade inteira não negativa comprada. Fabricação não excede validade; aquisição não é futura. Verificar unicidade especificação/fornecedor/número. Lotes existentes são filtrados pela especificação selecionada. Salvamento de cada entidade retorna seu ID; eventual falha da etapa seguinte permite retomar sem criar duplicatas.

\subsubsection{UI-06: cadastrar, editar e pesar frasco}
Adicionar Frasco inicia pelo seletor de resumo e especificação UI-05, com modos de busca nome/CAS/fórmula. Exige almoxarifado autorizado e ativo, localização detalhada (campo opcional até 500 caracteres, recomendado para localizar a unidade), estado FECHADO/ABERTO e modo de pesagem. Lote é opcional; selecionar Quarentena é uma decisão separada do estado físico e torna motivo obrigatório. Validade fechada é opcional se marcada desconhecida; prazo após abertura, quando informado, é inteiro positivo. Nunca tratar validade desconhecida como validade ilimitada. O servidor calcula validade efetiva; decisão sobre operação com validade desconhecida está em Q04.

Fechado: conteúdo nominal positivo (mL líquido/g sólido) e peso total positivo em g. Tara calculada é somente leitura: peso total menos conteúdo vezes densidade para líquido, ou menos conteúdo para sólido. Aberto: peso total e modo CONHECE\_TARA, ESTIMA\_VOLUME (líquido) ou ESTIMA\_MASSA (sólido). Exigir apenas a entrada correspondente; tara deve ser positiva e não maior que o peso total, conteúdo inicial positivo. Apresentar prévia identificada como estimativa; não enviar derivados como autoridade. Data de abertura anterior desconhecida não deve ser fabricada como hoje (Q05).

Se vencido, exigir destino: quarentena com motivo, pendente de descarte ou autorização didática explícita. Exibir revisão final com identidade, unidades, destino e dados de validade. Salvar cria código, frasco, histórico e auditoria de forma idempotente; apenas o servidor atribui LCQUI-N. O retorno mostra código e instrução de localizar/preencher a etiqueta física correspondente. Não gravar código sugerido por impressão anterior.

Editar propriedades permite localização, dados corrigíveis e decisões de estado autorizadas; não altera código, identidade química, peso de cadastro ou densidade. Pesagem de rotina exige novo peso em g, motivo e revisão da diferença; confirmar relê o frasco, exige disponibilidade e registra ajuste/evaporação separadamente do consumo em empréstimo. Tara incompatível bloqueia conclusão. Aplicar tolerância híbrida, esgotamento e recalibração de tara definidos em Q06; não aceitar consumo negativo silenciosamente.

\subsubsection{UI-07: retirada, devolução, quarentena e descarte}
Retirada exige código/ID do frasco, retirante ativo Professor/Bolsista, local de uso, peso de saída positivo em g, data prevista não anterior a hoje e finalidade (AULA\_PRATICA, DEMONSTRACAO, PESQUISA\_TCC\_POS, ESTUDO\_DEGRADACAO\_RESIDUOS). Nome e detalhes da especificação são somente leitura. Seleção de destinatário deve ser pesquisa autorizada com campos mínimos. Quarentena, vazio, quebrado, descartado ou pendente de descarte bloqueiam envio. Primeira abertura é confirmada e recalcula validade na mesma transação. Vencido autorizado exige ciência; pesquisa excepcional segue Q04 com TCR rastreável validado pelo backend, sem alterar a autorização do gestor. Confirmar apresenta resumo; servidor relê vínculo, frasco e destinatário, cria empréstimo e muda disponibilidade atomicamente.

Devolução seleciona empréstimo EM\_USO/ATRASADO do frasco, apresenta retirante, prazo, peso de saída e densidade fixada. Exige peso de retorno em g e identificação de quem devolve; operador é obtido da sessão. Prévia mostra diferença em g e consumo em g/mL. Peso superior ao de saída aplica tolerância híbrida Q06; ganho tolerado gera consumo zero e ajuste. Peso inferior à tara exige o fluxo de esgotamento/recalibração Q06. Ao confirmar, servidor recalcula consumo, atraso e vencimento; resultado devolve consumo final e estado atualizado. Qualquer gestor vinculado pode receber, mesmo que outro tenha registrado a saída. Frasco vencido na devolução exige destino, sem impedir o registro do retorno físico.

Alertas abrem lista filtrada de vazios, quebrados, vencidos, quarentena, pesagem ou escassez. Selecionar item abre UI-04 com ações cabíveis. Marcar vazio/quebrado exige motivo e confirmação; empréstimo ativo deve ser encerrado de modo auditado antes da transição incompatível. Quarentena e liberação exigem motivo, revalidação de validade e pesagem quando necessária. Descartar exige confirmação explícita de descarte físico concluído, motivo e registro do operador/instante; é terminal, sem excluir documento. ``Removido'' significa retirada do catálogo operacional por estado/arquivamento, nunca exclusão de identidade; política específica em Q07.

\subsubsection{UI-08: etiquetas virgens e segunda via}
Imprimir Etiquetas tem abas Novos Frascos e Frascos Cadastrados. Virgens: escolha entre próximo ao último cadastro, após última geração e intervalo personalizado; mostrar origem e instante da sugestão. Início/fim são inteiros positivos, fim maior ou igual a início e total entre 1 e 50. Grade 3 por 10 permite célula inicial por teclado ou clique, linha 1--10 e coluna 1--3, com células anteriores indisponíveis para esta impressão. Prévia informa quantidade de páginas e não reserva códigos. Última geração não prova impressão física.

Segunda via: busca autorizada por código/reagente, seleção de 1 a 10 IDs distintos e contador visível; limite bloqueia seleção adicional. Exibir identidade e localização para conferência. Servidor relê cada frasco e autoriza seu almoxarifado; falha em qualquer seleção impede gerar lote parcial silencioso. Cada página A4 contém ficha cadastral, especificação, lote/validade, último peso e tara, último empréstimo e dez últimos eventos, além de uma única etiqueta no rodapé. Ausência legítima é ``não informado'', nunca inventar saldo ou status. Auditoria registra uma segunda via por frasco.

Gerar mostra progresso indeterminado, erro recuperável e botão Baixar quando pronto; download não significa impressão concluída. Orientar impressão em tamanho real (100\%), sem ajustar à página, conferindo código antes de colar. Geometria e decisões de persistência constam na seção 10.

\subsubsection{UI-09: patrimônio e análise de requisições}
Novo Bem/Editar exige resumo existente (ou Novo Resumo com nome 150 e descrição), plaqueta textual até 30 caracteres preservando zeros iniciais, conservação BOM/REGULAR/RUIM, Local, foto, responsável SEI (150) e status. Descrição complementar é opcional (200). Foto: prévia, substituição, erro e upload inferior a 5 MiB; servidor verifica existência e autorização do objeto antes de salvar. Novo Bem não permite contornar baixa sem comprovante.

Pedido de Adição do Professor pede plaqueta, conservação, responsável, local, motivo e escolha exclusiva entre resumo existente e nome/descrição de novo resumo. DP-B02: foto é obrigatória já na submissão; o servidor rejeita ausência e valida o objeto; não usar imagem fictícia. Pedido de Edição apresenta valores atuais e campos propostos opcionais, com motivo obrigatório e ao menos uma diferença. Alterar nome individual significa reclassificação de resumo, não renomear todos os equipamentos. A seção 4 define conservação e status separadamente: RUIM não equivale automaticamente a Inservivel.

Minhas Requisições lista somente as do solicitante, por status/data, e abre detalhe com proposta, motivo, resposta, gestor e instante. Fila do gestor começa por pendentes e abre comparação em duas colunas (uma no celular); exibir campos alterados com texto/ícone, não apenas cor. Aprovar e Rejeitar exigem justificativa. Pedidos legados sem foto não podem ser aprovados até saneamento explícito; novas submissões já exigem foto e resumo completo. Confirmar relê requisição e versão do bem; se já respondida ou alterada, não sobrescrever silenciosamente. Persistir resposta, criação/edição, histórico e liberação do lock juntos, e notificar o professor de forma idempotente.

Baixa está no detalhe individual: selecionar Ja\_dado\_baixa exige PDF comprobatório inferior a 10 MiB, prévia/nome do arquivo e confirmação com plaqueta e responsável. O servidor verifica o objeto antes da transação; falha de upload mantém status anterior. Após sucesso, item sai do filtro Ativo e permanece consultável no histórico. Não excluir comprovante associado a evento de baixa.

\subsubsection{UI-10: turmas, membros e convites}
Nova Turma exige nome (100), matéria existente, ano inteiro, semestre 1/2 e capacidade inteira positiva; responsável vem do contexto autorizado. Código único é gerado no servidor e exibido com Copiar, feedback de sucesso/falha da área de transferência e alternativa de seleção manual. Detalhe de turma tem nome/período, matéria, professor, ocupação/capacidade, feed e Membros. Membros exibem nome/foto para alunos matriculados; matrícula/e-mail ficam restritos à gestão autorizada.

Arquivar abre confirmação com nome/período e explica preservação dos dados; turma sai das listas ativas. Turmas Arquivadas mostra lista paginada, ano/semestre, detalhe e Desarquivar para professor responsável. Restaurar usa mesmo ID e vínculos, sincroniza espelhos e notifica membros. Q08: turma arquivada é somente leitura; Rules e backend rejeitam escrita em posts e comentários.

Novo Aluno do Professor convida exclusivamente como Aluno, com um ou mais e-mails validados e deduplicados, matrícula opcional por pessoa e turma opcional. Sem turma, exigir confirmação de convite global; portanto o botão não fica bloqueado só por inexistirem turmas. Mostrar resultado individual por destinatário e permitir repetir apenas falhas, sem afirmar envio de e-mail se houve apenas criação de registro. Exceção de capacidade por convite nominal exige justificativa do professor (RN-TUR-01), persistida para validação na aceitação.

Aceitar convite exige sessão com e-mail verificado correspondente, convite pendente não expirado, nome e matrícula válida se faltarem. Validar conta/papéis existentes, não sobrescrever identidade. Aceitação cria vínculo e espelho atomicamente e consome convite uma vez. Convite expirado orienta solicitar novo; lotação sem exceção válida mantém erro recuperável. Ingresso por código no painel Aluno exige código não vazio até 20 caracteres, turma ativa, vaga e ausência de remoção impeditiva; aluno já matriculado recebe acesso à turma sem duplicar contagem. Removido anteriormente precisa de convite explícito. Remover membro abre confirmação com nome e turma, preserva comentários e histórico, remove vínculo/espelho e reduz contador na mesma transação.

\subsubsection{UI-11: roteiros, posts e comentários}
Novo Roteiro exige nome (150), descrição e PDF inferior a 15 MiB. Upload tem seleção/arrastar arquivo, nome/tamanho, progresso, Cancelar e Tentar novamente; não habilitar Salvar antes de validar upload e metadados no servidor. Biblioteca tem Meus, Compartilhados comigo e Eu compartilhei; cada linha mostra nome, autor, descrição e ações permitidas. Visualizar abre leitor com fallback Baixar; download revalida propriedade, compartilhamento ou matrícula numa turma cujo post vincula o roteiro. URL de arquivo não é autorização permanente.

Compartilhar abre pesquisa de professores ativos, exclui proprietário e seleciona IDs distintos. Confirmar cria relações únicas; repetir não duplica acesso. Revogar compartilhamento exige confirmação e segue Q09 quanto a posts existentes. Quem recebeu pode anexar à própria turma, mas não alterar autoria ou compartilhar em nome do dono.

Novo/Editar Post exige título (150) e descrição não vazia; roteiro é opcional e único por post. Clipe abre seletor com Meus/Compartilhados ou Novo Roteiro, preservando rascunho ao voltar. Publicar valida professor responsável, turma e acesso atual ao roteiro; guarda post, histórico de edição quando aplicável e notificações. Feed pagina por data/ID com mais recentes primeiro; navegação de notificação carrega o post alvo diretamente em vez de baixar todo o feed.

Comentários são carregados por post aberto, com texto obrigatório, autor/data e indicação de edição. Limite de texto segue Q10; rejeitar só espaços, renderizar texto escapado e não HTML arbitrário. Enviar mostra pendente e confirma após servidor validar matrícula/professor e post. Q11: autor pode editar, preservando histórico e etiqueta de edição; professor responsável modera com motivo obrigatório, sem apagar o original. Responder no feed é comentário no mesmo post, sem criar hierarquia de respostas ausente no modelo. Remoção do aluno não apaga comentários anteriores.

\subsubsection{UI-12: notificações, empréstimos próprios e relatórios}
Sino abre painel com Não lidas/Todas, tipo, resumo e instante. Consulta é pelo destinatário autenticado; Q12 inclui Bolsista e agrega alertas críticos de empréstimo independentemente do papel visual selecionado. Clicar revalida acesso ao alvo e abre o contrato correspondente; recurso indisponível apresenta mensagem sem expor dados. Limpar tudo pede confirmação de marcar notificações ativas do usuário como lidas, inclusive de outros papéis; mantém histórico e usa processamento paginado. Alertas vencidos não autorizam ações por si mesmos.

Meus Reagentes possui Em uso, ordenado por prazo, e Histórico, últimos 30 dias com filtro de período. Exibir código, reagente, peso, unidade, datas, consumo e situação; selecionar histórico abre o empréstimo e a memória de cálculo. Professor/Bolsista não registram a retirada nessa tela: procuram gestor no balcão. Não somar g e mL.

Relatórios possuem Mensal, Personalizado e, para patrimônio, Inservíveis/Bem específico. Mensal exige mês e ano não futuros; personalizado exige início/fim, início não posterior a fim, nenhuma data futura e no máximo 31 dias inclusivos. Filtros opcionais de prédio/andar/sala são dependentes: trocar prédio limpa seleções incompatíveis. Outros filtros: status, conservação, responsável, reagente e almoxarifados autorizados conforme domínio. ``Todos'' significa todos dentro do escopo do operador. Chefe tem escopo global. Relatório de um bem exige seu ID/plaqueta.

Prévia resume filtros, período, domínio e unidades; Gerar fica bloqueado se inválido. Servidor repete limites e escopo, usa agregados para totais e históricos para movimentos, distinguindo fotografia atual de eventos do período. Sem registros, informar ausência sem simular erro. Mostrar fonte/instante da materialização e período parcial quando ainda aberto. Ao concluir, oferecer Baixar/Imprimir; preservar filtros em falha e não persistir conteúdo do PDF em localStorage.

\subsubsection{UI-13: consultas, cache e armazenamento no navegador}
Chave de cache inclui UID, papel, recurso, filtros e cursor. Cache de sessão em memória é padrão; preferência de tema, menu compacto e filtros não sensíveis podem ir para localStorage com versão de formato e validação ao ler. Não armazenar senhas, tokens manualmente, e-mails, matrícula, PDFs ou autorização no localStorage. Persistência Firestore em disco só mediante opção de dispositivo confiável; bancada compartilhada usa memória. Logout/troca de identidade encerra listeners e elimina acesso aos dados anteriores; dados já baixados exigem política de limpeza, não desaparecem pela simples troca de claims.

Listeners ficam apenas nas listas visíveis e limitadas (feed ativo, notificações e estoque em atendimento); histórico e relatórios usam leitura pontual paginada. Cancelar ao fechar modal, mudar papel ou filtro. Invalidar cache afetado após mutação e revalidar dados críticos no servidor mesmo com cache recente. Offline permite consultar dados previamente autorizados com aviso de desatualização, mas não confirmar empréstimos, baixas, papéis ou matrículas. Evitar N+1: carregar referências únicas em lotes e materializações por escopo; nenhuma métrica de cache substitui validação de disponibilidade.

\paragraph{Q04 e Q14: contrato de aceite e autoatendimento.}
Para PESQUISA\_TCC\_POS ou ESTUDO\_DEGRADACAO\_RESIDUOS com frasco vencido ou validade desconhecida liberado com termo, obter aceite em sessão autenticada do retirante. Persistir tcr\_versao, tcr\_aceito\_em (servidor), tcr\_aceito\_por e tcr\_auditoria\_id no empréstimo, com evento correspondente em Registro\_de\_Auditoria; exigir justificativa\_metodologica de 20--2000 caracteres. Modal/checkbox são UX: backend valida versão vigente, identidade e vínculo do aceite ao frasco, finalidade e operação antes de confirmar a retirada. O TCR registra ciência e responsabilidade e não substitui regras institucionais de segurança química. Autoatendimento só quando não houver outro gestor ativo, calculado pelo servidor com justificativa e notificação à chefia. Registrar abertura\_historica\_desconhecida quando a data anterior não for conhecida; nunca inventar a data.

\begin{regra}[UI-11: visibilidade de moderação (DP-C02)]
Autor vê texto original com tarja ``Comentário moderado pelo docente'' e motivo. Demais alunos e bolsistas veem ``Comentário ocultado pela moderação da turma''. Professor responsável e Chefe Geral veem original e histórico necessário à moderação/auditoria. Bolsista segue visibilidade de Aluno. A resposta do backend já vem filtrada; Rules impedem obter diretamente o original restrito, inclusive histórico. Ocultação por CSS ou filtro no navegador não é proteção.
\end{regra}

\newpage
% =================================================
\section{Exemplos de fluxos}
% =================================================

\subsection{Fluxos complementares e critérios de completude}
Os fluxos são considerados completos quando descrevem ator, pré-condição, entrada, validação, mudança de estado, efeitos colaterais e resultado. As operações críticas abaixo complementam os cenários de tela já descritos.

\subsubsection{Concessão do papel Bolsista}
O Chefe Geral seleciona um usuário com papel Aluno. O backend verifica a pré-condição ($Bolsista \implies Aluno$), rejeita a operação se o usuário não for Aluno ou gerar combinação proibida (ex.: Aluno Bolsista atuando como Gestor de Almoxarifado), cria o vínculo de papel na coleção de perfis, registra auditoria e sincroniza claims e só então anuncia a permissão disponível após renovação do token.

\subsubsection{Aceitação de convite}
O aluno apresenta um convite válido. O backend verifica e-mail normalizado, status, expiração e capacidade da turma. A aceitação é idempotente: o vínculo Aluno--Turma é criado uma única vez e o convite deixa de ser pendente. Usuário e papel Aluno são criados apenas quando ainda não existirem.

\subsubsection{Arquivamento e desarquivamento de turma}
O professor responsável altera o status da turma para Arquivada. Posts, comentários, histórico e vínculos com alunos não são apagados. O evento é auditado e notificações são emitidas quando aplicável. O desarquivamento restaura o status Ativo sem recriar a turma.

\subsubsection{Retirada de reagente}
O Gestor de Almoxarifado seleciona o frasco e o Professor ou Bolsista retirante. A transação relê o frasco, valida disponibilidade e escopo do gestor, registra o peso de saída, cria o empréstimo EM\_USO e muda a disponibilidade para EMPRESTADO. Se o frasco estiver sendo aberto pela primeira vez, sua validade é recalculada antes de aprovar. Se outro processo tiver retirado o frasco antes da confirmação, a transação falha sem criar empréstimo parcial.

\subsubsection{Devolução de reagente}
O Gestor informa o frasco e o peso de retorno. A transação verifica que o empréstimo está EM\_USO ou ATRASADO, recalcula consumo e atraso com dados do servidor (prazo encerrando às 23:59:59.999 do dia previsto), atualiza o empréstimo, grava o evento no histórico e devolve o frasco à disponibilidade adequada ou ao fluxo de descarte/quarentena se vencido ou esgotado.

\begin{explicacao}[Nota sobre o Login]
Quando não houver sessão válida, a etapa inicial é COM-01: o usuário acessa a tela de Login do LCQUI, insere suas credenciais (e-mail e senha) ou escolhe o "Login com Google". O sistema exibe um \textit{spinner} (ícone de carregamento giratório) no botão e carrega conta ativa e permissões antes de abrir a dashboard correspondente. Na mesma sessão, os fluxos seguintes não exigem repetir login.
\end{explicacao}

\subsection{Fluxos do usuário Gestor de Bem Patrimonial}

\subsubsection{Adicionando um bem patrimonial do LCQUI no sistema}
Partindo da tela de login, o Gestor acessa o sistema. Caso possua múltiplos papéis, ele clica no botão "Gestor de Bens Patrimoniais" no header. Na dashboard, clica no botão "Mais (+)" e seleciona "Novo Bem Patrimonial". Um modal se sobrepõe à tela. Ele preenche o número do patrimônio, pesquisa um \textit{Resumo\_Bem} existente, escolhe o Local (prédio, andar, sala) via menus \textit{dropdown} e anexa a URL da foto. Ao clicar em "Salvar", um \textit{spinner} roda no botão e, em seguida, um \textit{toast} (notificação flutuante) verde surge na tela: "Bem cadastrado com sucesso". O modal fecha e a grid central é recarregada mostrando o novo item.

\subsubsection{Baixando a lista de bens inservíveis para dar baixa no SEI}
Partindo da tela de login, o Gestor acessa a dashboard e clica no botão "Gerar Relatórios" na parte superior. O modal de relatórios abre. Ele seleciona o tipo "Bens Inservíveis", confirma e clica em "Gerar PDF". O botão muda para o estado \textit{loading}. Após alguns segundos, o navegador exibe a notificação de download concluído ou abre o arquivo em uma nova aba, pronto para ser anexado ao Sistema Eletrônico de Informações (SEI) da UENF.

\subsubsection{Baixando o relatório mensal Geral do Mês de Agosto}
Após o login, na tela principal, o Gestor clica em "Gerar Relatórios". No modal que se abre, seleciona a aba "Relatório Geral". Nos seletores de data, escolhe "Agosto" e "2026". Ao clicar em "Gerar", um círculo de progresso é exibido no centro do modal. Finalizado o processo backend, o modal exibe a mensagem "PDF pronto" e inicia o download automaticamente.

\subsubsection{Baixando relatório de bens materiais filtrado (Prédio P5, Andar 1)}
Partindo da tela de login, o Gestor abre o modal de relatórios e vai até a aba "Personalizado". Ele preenche as datas (15/06/2026 a 15/07/2026) e abre o filtro avançado. Marca a \textit{checkbox} do "Prédio P5" e digita "1" no andar. Clica em "Gerar PDF". Após a barra de carregamento preencher, o sistema entrega o documento formatado apenas com os equipamentos deste andar específico.

\subsubsection{Verificando requisições pendentes de professores}
Após efetuar o login, a dashboard do Gestor já apresenta no painel central um card de "Painel de Requisições". O botão de notificações (sino) está com uma bolinha vermelha. O Gestor clica na notificação e depois na requisição desejada. A tela faz uma transição para o modo \textit{Split-screen}. Na metade esquerda da tela, ele vê os dados estáticos (antigos) do patrimônio em texto cinza. Na metade direita, ele vê os dados novos sugeridos pelo professor (ex: alteração de sala) destacados em fundo amarelo suave, para chamar a atenção à mudança.

\subsubsection{Aprovando uma requisição pendente}
Dando continuidade ao fluxo de verificação (após login e abertura da tela de \textit{Split-screen}), o Gestor avalia que a alteração de sala pedida pelo professor está correta. Ele clica no botão verde massivo "Aprovar Alteração" no rodapé. Um ícone de \textit{check} ($\checkmark$) animado aparece no centro da tela. A requisição desaparece da fila de pendentes e um \textit{toast} de sucesso no canto superior direito confirma que o banco de dados principal foi atualizado.

\subsubsection{Alterando um bem patrimonial para: Dado baixa}
Partindo da tela de login, o Gestor usa a busca exata pela plaqueta canônica de um equipamento \textbf{Inservivel} que já passou por todo o processo burocrático no SEI. A tela de detalhes abre o rito próprio ``Registrar baixa'', não um editor genérico de status. O Gestor anexa o PDF no caminho autorizado; o backend valida objeto, dono/origem, tamanho, bytes \texttt{\%PDF-} e geração antes da transação. A confirmação relê autorização M9, status, versão e idempotência M7, grava comprovante imutável, evento histórico \texttt{baixa} e o status terminal \texttt{Ja\_dado\_baixa}. Falha de upload ou conflito preserva o estado anterior. Após sucesso o item sai do filtro ativo, continua consultável com seu histórico e a plaqueta não volta a ficar disponível.

\subsection{Fluxos do usuário Professor}

\subsubsection{Criando uma turma}
Após fazer login, o Professor está na sua dashboard. Ele clica no painel lateral esquerdo em "Turmas" e depois no botão flutuante "+ Nova Turma". Um modal simples é exibido solicitando Nome da Turma, Matéria, Ano, Semestre e Capacidade. Ele preenche "Química Orgânica I", seleciona a matéria correspondente, informa ano 2026, semestre 2 e capacidade 30, e clica em "Salvar". Um \textit{toast} de sucesso é disparado, um novo card desliza aparecendo no painel lateral, e o \textit{codigo\_turma} recém-criado pisca rapidamente na tela com um botão de "Copiar para a área de transferência".

\subsubsection{Adicionando Alunos em uma turma}
Com login feito, o Professor clica no card da turma recém-criada no painel lateral. A área central atualiza para a visão da disciplina. No header, ele clica no botão "Novo Aluno". No modal que surge, ele copia e cola os e-mails dos alunos separados por vírgula e clica em "Enviar Convites". Uma barra de carregamento é mostrada e o sistema apresenta o resultado por destinatário. O \textit{toast} só afirma envio de e-mail quando o envio externo foi acionado; a criação do registro de convite, por si só, é reportada como convite registrado, sem anunciar e-mail enviado. Falhas parciais permitem repetir somente os destinatários que faltaram.

\subsubsection{Criando um Post em uma turma}
Após login, o Professor acessa a disciplina clicando nela. No topo do Feed de publicações (área central), ele clica na caixa de texto "Comece um novo post...". O componente se expande. Ele preenche título e orientações da aula, clica no ícone de "Clipe de Papel" (anexo), o que abre o modal da sua biblioteca de roteiros em PDF. Ele seleciona o "Roteiro\_Prática\_3.pdf" e clica em "Publicar". A tela rola suavemente para baixo e o novo card da publicação surge com um efeito de \textit{fade-in}, exibindo o texto e a prévia do PDF anexado.

\subsubsection{Excluindo um Aluno de uma turma}
O Professor loga no sistema e navega até a turma desejada. Ele clica na aba "Membros" logo acima do feed. Uma lista com nome e foto dos alunos aparece. Ao lado do nome do aluno a ser removido, ele clica no ícone vermelho de lixeira. Um modal de dupla confirmação exige que ele clique em "Sim, remover aluno". A linha do aluno desaparece da lista com uma rápida animação de encolhimento (\textit{collapse}) e um \textit{toast} avisa "Aluno removido da turma".

\subsubsection{Respondendo um comentário de um aluno}
O Professor entra no sistema e nota uma bolinha vermelha piscando no sino de Notificações. Ele clica e vê "\textit{João comentou em um Roteiro}". Ao clicar na notificação, a tela rola automaticamente até a exata postagem no feed. Ele vê a dúvida do aluno, clica na caixa de texto abaixo dizendo "Escreva uma resposta..." e aperta Enter. O comentário dele surge instantaneamente logo abaixo do aluno, com seu nome identificado com a tag "[Professor]".

\subsubsection{Pesquisando Reagente Químico para checar disponibilidade}
Partindo da tela de login, o Professor clica no atalho "Reagentes" no painel esquerdo. O catálogo JSON de Resumos e Especificações é carregado no cliente (usando o cache local quando a versão publicada não mudou) e a busca textual ocorre localmente: na grande barra de pesquisa central, ele começa a digitar "Ácido Clori..." e a tabela filtra a cada caractere, sem filtrar por tecla no Firestore e sem exigir filtro prévio. Ele encontra o item, clica sobre a linha e um modal lateral (gaveta) desliza da direita mostrando em verde "\textit{Disponível: 3 frascos fechados, 1 aberto}".

\subsubsection{Pesquisando um bem patrimonial pelo Nome}
Após login, o Professor vai à área "Bens Patrimoniais" e digita "Microscópio" na busca. A tela é atualizada mostrando os equipamentos carregados que correspondem ao texto, após filtro obrigatório, com paginação (UI-04). Ele clica em um dos cards e vê a foto e a localização exata daquele Microscópio em tela cheia.

\subsubsection{Pesquisando um bem patrimonial pelo número de patrimônio}
Professor faz login, acessa a busca de "Bens Patrimoniais" e digita exatamente "987654". O sistema detecta a correspondência exata, pula a exibição da tabela genérica e direciona o professor direto para a página exclusiva de detalhes daquele patrimônio específico.

\subsubsection{Requisitando alteração de conservação (de Bom para Ruim)}
No passo anterior, ao estar na tela de detalhes do microscópio, o Professor, após logar, percebe que a lente está trincada. Ele clica no botão amarelo "Solicitar Edição". No formulário sobreposto, ele altera o campo \textit{dropdown} "Estado de Conservação" de "Bom" para "RUIM" (conservação); Inservivel é outro campo, de status, preenche o campo de texto justificando ("Lente principal rachada") e envia. O botão vira um \textit{spinner} rapidamente e um \textit{toast} avisa: "Requisição enviada ao Gestor de Patrimônio".

\subsubsection{Upload de Roteiro e compartilhamento com outros professores}
Professor acessa a dashboard após login, clica em "Gerenciar Roteiros" (no header central). O modal abre na aba "Meus". Ele clica no ícone de Nuvem para upload, arrasta um PDF do seu computador para a área pontilhada e aguarda a barra de progresso encher (100\%). Em seguida, no próprio arquivo recém-subido, ele clica em "Compartilhar". Uma lista de professores do sistema aparece, ele marca 3 caixinhas de seleção e clica no botão "Autorizar Acesso". Um alerta flutuante informa que o roteiro agora está visível nas bibliotecas dos colegas.

\subsection{Fluxos do usuário Chefe Geral}

\subsubsection{Criando um usuário Professor e usuário Gestor de Almoxarifado}
O Chefe Geral acessa a plataforma logando com suas credenciais. Na sua dashboard, acessa a aba "Professores" e clica no botão "+ Novo Professor". No modal, insere nome, e-mail institucional, laboratório e seleciona o Centro (ex: CCT). Ao salvar, um \textit{toast} avisa do envio do convite. Da mesma forma, indo para a aba "Almoxarifados", ele clica em "Mais (+)", depois "Novo Gestor de Almoxarifado", preenche o e-mail, e um novo card de gestor é populado silenciosamente na lista do sistema, notificando o sucesso da operação.

\subsection{Fluxos do usuário Aluno}

\subsubsection{Ingressando em uma Turma}
O Aluno acessa o site do LCQUI pela primeira vez, cria a senha (ou entra via Google) e faz o login. A dashboard dele é muito mais limpa e a tela central diz: "\textit{Você ainda não tem turmas}". Ele foca sua atenção no canto inferior esquerdo onde há o input "Código da Turma". Ele digita o \textit{código} de até 20 caracteres que o professor escreveu no quadro e clica em "Entrar". O botão gira um carregamento, o sistema verifica a capacidade máxima da turma no \textit{backend}, faz o link de tabelas, e instantaneamente um Card colorido com o nome da disciplina desliza para o centro da tela com o feedback: "\textit{Matrícula realizada com sucesso!}".

\subsubsection{Baixando um roteiro e realizando um comentário}
Logado, o aluno clica na disciplina recém-adicionada. O Feed da turma é carregado na tela. Ele rola a página e acha a postagem do professor com um arquivo. Ele clica no botão vermelho e muito visível de "Baixar PDF". O \textit{download} inicia via navegador. Logo abaixo do mesmo \textit{post}, na seção de interações, ele clica na caixa de texto vazia, digita "Professor, levaremos óculos de proteção?" e pressiona Enviar. O comentário surge na tela logo após uma animação suave de transição (\textit{fade-in}).

\subsection{Fluxos do Usuário Gestor de Almoxarifado}

\subsubsection{Adicionando um frasco fechado de um reagente existente}
Partindo do login, o Gestor de Almoxarifado acessa seu painel, clica no grupo "Estoque", digita o nome de um ácido na busca e clica em "Adicionar Frasco". Ele coloca o frasco lacrado em cima da balança USB (ou lê manualmente), digita no campo "Peso Total" (ex: 610g) e no "Volume Nominal" (ex: 500mL). Ao clicar em "Salvar", o sistema deriva a tara teórica (peso total menos o conteúdo nominal, sem caráter de medição física definitiva do recipiente vazio), registra no banco e o modal se fecha com uma notificação verde "Frasco Fechado registrado no almoxarifado".

\subsubsection{Adicionando um frasco já aberto de um novo reagente}
Após o login, o Gestor nota que o reagente ainda não existe no sistema. Na dashboard central, ele clica em "Novo Reagente". Preenche as propriedades químicas obrigatórias, com destaque em vermelho para o campo "Densidade" quando líquido. Ao confirmar, o modal avança automaticamente para o passo "Cadastrar Frascos Vinculados". Como o frasco físico já está aberto, ele seleciona a condição física inicial "Já Aberto", deixa o lote não informado (sem inventar lote) e informa apenas os dados reais disponíveis: o peso total atual medido na balança e, se o rótulo estiver legível, o conteúdo nominal. O sistema exibe "Saldo atual: desconhecido" e "Tara do recipiente: ainda não conhecida" e não pede tara conhecida, volume estimado nem massa estimada. Salva a operação e a tela é atualizada sem estimar saldo inicial; o frasco permanece com \texttt{saldo\_desconhecido = true} até que seja esgotado e pesado vazio.

\subsubsection{Realizando empréstimo (Retirada) de um frasco}
Com o login concluído, um Professor ou Bolsista chega no balcão. O Gestor de Almoxarifado clica rapidamente no enorme botão verde ``Registrar Retirada'' no centro da dashboard. Um modal abre solicitando a leitura do código do frasco ou pesquisa pelo nome. Localizado o item, o Gestor seleciona o usuário autorizado, informa o local de uso e pesa o frasco.

Antes da confirmação, a interface bloqueia frascos em quarentena ou pendentes de descarte. Se o frasco estiver vencido e seu uso vencido tiver sido autorizado, deve aparecer um alerta explícito: ``Este frasco venceu em DD/MM/AAAA. O uso vencido exige autorização e ciência conforme a finalidade; pesquisa requer TCR rastreável. Deseja confirmar o empréstimo mesmo assim?''. A confirmação é enviada ao backend, que repete a validação; o frontend não pode ignorar essa regra. Se a operação também registrar a primeira abertura do frasco, a validade efetiva é recalculada no servidor dentro da mesma transação.

\subsubsection{Devolvendo um frasco (Cálculo de Consumo)}
O professor devolve o vidro à bancada. O Gestor faz login (se necessário) e clica no grande botão ``Registrar Devolução''. Busca o item (que está na lista de pendentes). O sistema solicita ``Peso de Retorno'' e apresenta a confirmação explícita ``O frasco ficou vazio nesta devolução? [ ] Sim'', disponível tanto para frascos cadastrados como FECHADO quanto para os cadastrados como JA\_ABERTO. O sistema calcula o consumo e verifica a validade efetiva com o relógio do servidor (com prazo normalizado para as 23:59:59.999 do dia previsto). A confirmação do gestor é a autoridade do esgotamento, e não a comparação com a tara anterior. Ao confirmar que o frasco ficou vazio, o peso de retorno passa a ser a tara real: o frasco vai para VAZIO/INDISPONIVEL, o saldo vira conhecido como 0\,g (sólido) ou 0\,mL (líquido), registra-se FICOU\_VAZIO e o conteúdo nominal não é alterado. Se o gestor não confirmar o esgotamento, a diferença de peso é tratada como anomalia metrológica sujeita a justificativa e o frasco mantém seu estado. Não substituir o peso medido pela tara nem concluir uma devolução fictícia para contornar a validação. Se o frasco estiver vencido, ou tiver vencido enquanto permaneceu com o professor, a devolução permite registrar a decisão do gestor entre quarentena, pendência de descarte ou continuidade de uso excepcional conforme Q04; registrar descarte físico concluído é outra operação.

\subsubsection{Corrigindo uma retirada atribuída ao usuário errado}
Após o login, o Gestor identifica no histórico do frasco que a retirada foi vinculada ao Professor errado. Ele abre ``Corrigir operação'', que exibe a operação original, o efeito que será corrigido e um campo de motivo. O gestor lê o efeito e confirma. A correção é compensatória e auditável: o lançamento incorreto é preservado, a correção é vinculada à operação original com motivo e instante do servidor, e o vínculo do retirante é corrigido sem devolução fictícia, consumo fictício ou nova pesagem artificial. O peso físico de saída, o instante físico original e o local de uso permanecem registrados; o original não é apagado.

\subsubsection{Colocando um frasco em quarentena e decidindo sua saída}
Durante uma inspeção de bancada, o Gestor decide bloquear um frasco independentemente de cadastro, abertura, devolução ou reencontro. Ele usa ``COLOCAR EM QUARENTENA'', informa o motivo humano e confirma; o sistema relê o estado atual, valida o escopo do almoxarifado e grava \texttt{em\_quarentena = true} e \texttt{disponibilidade = INDISPONIVEL}, preservando vencimento, estado físico, saldo e validade. Na decisão posterior, o gestor escolhe ``VOLTAR A DISPONIVEL'' ou ``PENDENTE DE DESCARTE'', com motivo/laudo; a saída não revalida a validade nem altera o vencimento, e uma segunda decisão concorrente é rejeitada se a pré-condição já não existir.

\subsubsection{Extravio e reencontro de um frasco}
O extravio e o reencontro são fatos distintos de liberação. Os quatro fluxos abaixo obedecem às regras da subseção M5 da Seção~7.

\textbf{Fluxo A --- disponível $\rightarrow$ extravio $\rightarrow$ reencontro $\rightarrow$ quarentena $\rightarrow$ liberação.} O Gestor registra o extravio de um frasco disponível: ele assume \texttt{situacao\_localizacao = EXTRAVIADO} e \texttt{disponibilidade = INDISPONIVEL}, preservando o estado físico, \texttt{saldo}, \texttt{saldo\_desconhecido}, peso, tara, vencimento e \texttt{abertura\_historica\_desconhecida}; a autorização corrente de descarte é revogada, mantendo a decisão anterior no histórico; nenhum peso, consumo ou tara é inventado. Dias depois, o frasco é reencontrado; o gestor informa o estado físico constatado (\texttt{ABERTO}/\texttt{FECHADO}) e o sistema aplica \texttt{em\_quarentena = true} e \texttt{INDISPONIVEL}, sem marcar vencimento por causa do reencontro e sem reabrir empréstimo algum. O saldo permanece como estava. Após inspeção, o gestor decide ``VOLTAR A DISPONIVEL'' com motivo/laudo; como não há pendência metrológica, o estado físico é ABERTO/FECHADO e o frasco não está vencido sem autorização, a quarentena é encerrada e o frasco volta a DISPONIVEL sem revalidar validade.

\textbf{Fluxo B --- emprestado $\rightarrow$ extravio durante a custódia $\rightarrow$ encerramento extraordinário $\rightarrow$ reencontro $\rightarrow$ quarentena.} O frasco está EM\_USO com o Professor quando desaparece. O gestor registra o extravio com motivo humano: a localização vira EXTRAVIADO, o estado físico conhecido é preservado e a disponibilidade vira INDISPONIVEL e o empréstimo ativo é encerrado como \texttt{ENCERRADO\_EXTRAORDINARIO} com \texttt{EXTRAVIO\_SINISTRO}. Como não houve retorno físico, \texttt{peso\_retorno}, \texttt{peso\_retorno\_efetivo}, \texttt{medida\_utilizada} e \texttt{data\_devolucao\_efetuada} permanecem nulos, \texttt{consumo\_validado = false} e \texttt{peso\_saida} é preservado; \texttt{massa\_perda\_estimada\_g} só é gravada como estimativa se houver tara conhecida. Quando o frasco é encontrado, o gestor registra o reencontro e ele entra em quarentena; o empréstimo encerrado \textbf{não} é reaberto.

\textbf{Fluxo C --- reencontrado $\rightarrow$ quarentena $\rightarrow$ condição inadequada $\rightarrow$ pendência de descarte $\rightarrow$ descarte autorizado.} Após o reencontro, a inspeção conclui que o frasco não é reaproveitável. O gestor decide ``PENDENTE DE DESCARTE'' com motivo/laudo: a quarentena é encerrada, o frasco permanece INDISPONIVEL e é criada a autorização operacional estruturada de descarte (HQ-M2-008/B). Somente depois, com a autorização ativa, ``Descartar'' é permitido e o frasco vira DESCARTADO (terminal). Não existe caminho direto de quarentena para descarte.

\textbf{Fluxo D --- reencontrado com informação quantitativa insuficiente.} O frasco é reencontrado, mas não há tara conhecida nem medição confiável do conteúdo. A quarentena permanece; o saldo continua desconhecido e \textbf{não} é fabricado. Uma pesagem de rotina posterior pode ser registrada para observar o estado, mas não resolve automaticamente a quarentena nem a pendência quantitativa. Enquanto a informação for insuficiente, o frasco não volta a DISPONIVEL e a decisão de destino aguarda subsídio; qualquer determinação quantitativa definitiva pertence à fronteira M6 (Q06/tara), sem inventar consumo, tara ou saldo. Se o gestor optar por não reaproveitar, aplica-se o Fluxo~C.

\textbf{Casos de regressão documental.} O reencontro valida o estado constatado contra o último estado conhecido antes do extravio:
\begin{itemize}
    \item \textbf{Caso 1 (válido).} \texttt{FECHADO} $\rightarrow$ \texttt{EXTRAVIADO} $\rightarrow$ \texttt{FECHADO} $\rightarrow$ quarentena, somente se nunca houve abertura conhecida.
    \item \textbf{Caso 2 (válido).} \texttt{FECHADO} $\rightarrow$ \texttt{EXTRAVIADO} $\rightarrow$ \texttt{ABERTO} $\rightarrow$ quarentena, marcando \texttt{abertura\_historica\_desconhecida = true}, \texttt{data\_abertura = null}, \texttt{validade\_desconhecida = true} e \texttt{validade\_efetiva = null}.
    \item \textbf{Caso 3 (rejeitado).} \texttt{ABERTO} $\rightarrow$ \texttt{EXTRAVIADO} $\rightarrow$ \texttt{FECHADO}: um frasco já aberto não pode ser ``ressuscitado'' como fechado.
    \item \textbf{Caso 4 (válido).} \texttt{VAZIO} $\rightarrow$ \texttt{EXTRAVIADO} $\rightarrow$ \texttt{VAZIO} $\rightarrow$ quarentena $\rightarrow$ \texttt{PENDENTE\_DE\_DESCARTE}: restauração sem nova medição, seguida de pendência de descarte.
    \item \textbf{Caso 5 (rejeitado).} \texttt{VAZIO} $\rightarrow$ \texttt{EXTRAVIADO} $\rightarrow$ \texttt{FECHADO}.
    \item \textbf{Caso 6 (válido).} \texttt{QUEBRADO} $\rightarrow$ \texttt{EXTRAVIADO} $\rightarrow$ \texttt{QUEBRADO} $\rightarrow$ quarentena $\rightarrow$ \texttt{PENDENTE\_DE\_DESCARTE}: quebra física constatada, sem recuperação operacional.
    \item \textbf{Caso 7 (rejeitado).} \texttt{QUEBRADO} $\rightarrow$ \texttt{EXTRAVIADO} $\rightarrow$ \texttt{ABERTO}.
    \item \textbf{Caso 8 (sem fabricação).} Frasco anteriormente \texttt{ABERTO} reencontrado aparentemente sem conteúdo: M5 \textbf{não} inventa \texttt{VAZIO}, \texttt{saldo = 0} nem nova tara; mantém a quarentena, registra a observação e remete a confirmação quantitativa a M6.
\end{itemize}

\textbf{Casos de regressão documental de M8 (estoque, escassez, cache e notificações).} A aptidão usa o predicado único \texttt{frascoAptoParaUso} da Seção~\ref{sec:m8-estoque-escassez-notificacoes}.

\emph{Estoque atual.}
\begin{itemize}
    \item \texttt{FECHADO + LOCALIZADO + DISPONIVEL}: apto.
    \item \texttt{ABERTO + LOCALIZADO + DISPONIVEL}: apto.
    \item \texttt{ABERTO + EXTRAVIADO}: não apto, apesar do estado físico aberto.
    \item \texttt{QUEBRADO}, \texttt{VAZIO}, \texttt{DESCARTADO}, \texttt{EMPRESTADO}: não aptos.
    \item \texttt{em\_quarentena = true}: não apto, mesmo \texttt{ABERTO}/\texttt{FECHADO}.
    \item Frasco vencido com \texttt{uso\_vencido\_autorizado = true}: apto (validade compatível); sem autorização: não apto.
    \item Saldo desconhecido: a contagem de frascos é possível conforme a métrica, mas a soma quantitativa é segregada e \textbf{nunca} tratada como zero conhecido.
    \item Massa (g) e volume (mL) nunca são somados no mesmo total.
\end{itemize}

\emph{Escassez.}
\begin{itemize}
    \item \texttt{qtdAptos < limiar}: escassez (emite alerta quando configurado).
    \item \texttt{qtdAptos = limiar}: \textbf{não} há escassez.
    \item \texttt{qtdAptos > limiar}: não há escassez.
    \item \texttt{ativo = false}: não avaliar.
    \item \texttt{notificacao\_ativa = false}: escassez pode existir conceitualmente, sem alerta automático.
    \item \texttt{Almoxarifado.ativo = false}: não executar avaliação automática.
    \item Frascos extraviados, quebrados, vazios, descartados, emprestados ou em quarentena não inflam \texttt{qtdAptos}.
\end{itemize}

\emph{Cache.}
\begin{itemize}
    \item \texttt{cache hit} válido e autorizado: responde dentro da validade restante.
    \item \texttt{cache hit} sem autorização: rejeita.
    \item \texttt{cache hit} acima do rate limit: rejeita ($5$/min/UID).
    \item \texttt{cache miss} ou inválido: recalcula sob demanda.
    \item Geração muda durante o cálculo: não publica resultado obsoleto.
    \item Duas mudanças consecutivas: responde indisponibilidade temporária, não dado obsoleto.
    \item TTL físico atrasado: não torna o cache semanticamente válido.
    \item $30$ segundos sem acesso: nenhum recálculo automático (não há scheduler).
    \item Mutação de domínio atualiza geração e marca \texttt{valido = false} na \textbf{mesma} transação, sem recalcular a agregação.
    \item Correção exclusivamente histórica não invalida estoque atual.
\end{itemize}

\emph{Notificações.}
\begin{itemize}
    \item Primeira escassez do dia: cria \texttt{ESCASSEZ\_ESTOQUE} com docId determinístico.
    \item \textit{Retry} do job no mesmo dia: não duplica (\texttt{create}/\texttt{ALREADY\_EXISTS}).
    \item Mesma configuração no dia civil seguinte: nova notificação permitida.
    \item Erro diferente de \texttt{ALREADY\_EXISTS}: \textbf{não} engolido, propaga.
    \item Gestor sem vínculo ao almoxarifado: não recebe, mesmo com o papel global.
    \item ``Limpar tudo'': marca como lida, preserva histórico, sem \texttt{DELETE}.
\end{itemize}

\emph{Interação M5/M6.}
\begin{itemize}
    \item Frasco extraviado: fora de \texttt{qtdAptos}.
    \item Reencontrado em quarentena: fora de \texttt{qtdAptos}.
    \item Quebrado: fora de \texttt{qtdAptos}.
    \item Pendência metrológica bloqueante (\texttt{DEVOLVIDO\_COM\_ANOMALIA} com \texttt{consumo\_validado = false}): fora de \texttt{qtdAptos}.
    \item Resolução metrológica que mantém quarentena: continua fora.
    \item Liberação M5 válida (\texttt{VOLTAR\_A\_DISPONIVEL}): volta a depender das demais condições de aptidão.
    \item Consumo pendente não é contado como consumo zero nem como frasco apto adicional.
\end{itemize}

O sistema renderiza na tela em negrito e com fundo destacado: \textbf{``Volume utilizado neste empréstimo: 4.5 mL''} ou \textbf{``Peso utilizado nesse empréstimo: 22 g''}. Ao clicar em ``Finalizar'', um sinal sonoro/visual confirma que a transação de estoque, consumo, histórico e auditoria foram concluídas.


\subsection{Fluxos M11 --- Turmas, matrícula e convites}
Estes exemplos são o contrato normativo da Seção~\ref{sec:regras-turmas-m11}.
Eles são ilustrativos e não atestam a implementação.

\subsubsection{Corrida pela última vaga}
Turma com \texttt{qtd\_alunos = capacidade - 1}. Dois alunos aceitam ingressar por
código no mesmo instante. Cada transação lê \texttt{qtd\_alunos} e valida
\texttt{qtd\_alunos < capacidade}; o Firestore reexecuta as transações em
conflito. A primeira commita e grava vínculo, espelho, evento e
\texttt{qtd\_alunos + 1}. A segunda, já com \texttt{qtd\_alunos = capacidade},
recebe \texttt{failed-precondition} e não cria vínculo, espelho, evento nem
contador parcial.

\subsubsection{Dois convites simultâneos}
Para o mesmo par (e-mail normalizado, turma), ou para o mesmo e-mail com
\texttt{id\_turma} \texttt{NULL}, duas criações concorrentes disputam a chave
determinística em \texttt{Chaves\_Unicas}. A primeira cria o convite
\texttt{pendente}; a segunda, na reexecução, encontra a chave ocupada e é
rejeitada sem criar convite duplicado. Turmas diferentes para o mesmo e-mail
são permitidas. Reenviar para um convite pendente substitui hash e expiração no
mesmo documento em vez de criar outro.

\subsubsection{Replay de aceite}
O aluno aceita um convite e o vínculo é criado. Ao repetir a mesma operação
M7, o backend encontra o receipt e devolve o resultado persistido: o convite
permanece aceito, o vínculo já existe, o espelho está coerente e
\texttt{qtd\_alunos} não é incrementado de novo. Convite consumido por outro
e-mail ou já expirado é negado com erro específico, sem efeito parcial.

\subsubsection{Remoção concorrente}
Professor dono e Chefe (Q13) removem o mesmo aluno simultaneamente. A primeira
transação apaga o vínculo canônico, o espelho, grava \texttt{exclusao\_aluno}
com \texttt{removido\_por} e decrementa \texttt{qtd\_alunos}. A segunda
encontra o vínculo ausente, trata o caso como idempotente e não decrementa o
contador de novo, jamais abaixo de zero. Dois retries de remoção pela mesma
origem seguem a mesma regra. Posts e comentários anteriores permanecem e nova
entrada por código é bloqueada, exigindo convite explícito.

\subsubsection{Arquivamento e desarquivamento}
O professor dono arquiva uma turma com alunos. A transação grava
\texttt{status = Arquivada} e incrementa \texttt{versao}, sem apagar membros,
posts, comentários, histórico nem alterar \texttt{codigo\_turma}. Enquanto
arquivada, o backend e as Rules rejeitam novo ingresso, novo convite, novo post
e novo comentário (Q08); os espelhos dos alunos recebem \texttt{status} por
fan-out idempotente, e espelho atrasado nunca concede escrita. O
desarquivamento restaura \texttt{Ativo} com o mesmo ID e os mesmos vínculos.

\subsubsection{Convite global}
O professor convida um e-mail sem turma, mediante confirmação explícita. O
convite tem \texttt{id\_turma} \texttt{NULL} e chave global própria. Ao aceitar,
o backend garante \texttt{Usuarios} e o papel Aluno e consome o convite, mas
\textbf{não} cria matrícula nem incrementa \texttt{qtd\_alunos}; o aluno entra
em alguma turma depois, por código ou por outro convite. Matrícula é
obrigatória e única na identidade final \texttt{Aluno}, opcional e não única no
convite.

\subsection{Fluxos M12.1 --- posts, comentários, edição e moderação}
Estes exemplos são o contrato normativo da
Seção~\ref{sec:regras-posts-comentarios-m12-1}. São ilustrativos e não atestam a
implementação.

\subsubsection{Publicar Post com e sem roteiro}
O professor dono abre a turma Ativo e publica um Post. Sem roteiro, a transação
relê o vínculo de ownership e o estado da turma, grava titulo/descricao e o
\texttt{idOperacao} M7 e emite a notificação \texttt{POST} como efeito
delimitado. Com roteiro, o backend valida o acesso atual ao roteiro, grava o
snapshot imutável \texttt{roteiro\_anexo} e só então publica; se o acesso não
puder ser validado, o Post não é criado (fronteira M12.2).

\subsubsection{Editar Post e remover da apresentação}
O professor dono edita o Post: a transação relê ownership, grava o evento em
\texttt{Historico\_Posts\_Turma} com valores antigo/novo, operador e timestamp,
marca \texttt{editado}/\texttt{editado\_em} e é idempotente M7; uma segunda
edição concorrente reexecuta a transação e preserva a versão anterior no
histórico, sem sobrescrita silenciosa. Remover da apresentação exige motivo e
não apaga o documento nem o histórico: o item some do feed de colegas e
permanece acessível ao professor dono e à auditoria.

\subsubsection{Moderar Post por Q13}
O Chefe identifica conteúdo impróprio em turma alheia e remove o Post da
apresentação com motivo. A transação exige turma Ativo, grava
\texttt{removido\_da\_apresentacao}, \texttt{motivo\_remocao},
\texttt{removido\_por}/\texttt{removido\_em}, registra o evento em
\texttt{Historico\_Posts\_Turma} (\texttt{tipo = moderacao}) e a auditoria,
sem assumir a autoria. Em turma arquivada a intervenção é negada.

\subsubsection{Comentar, editar e moderar comentário}
Aluno/Bolsista com vínculo canônico atual comenta; o texto é validado (1--2\,000,
sem HTML), escapado na renderização e gravado com \texttt{idOperacao} M7. O
autor edita o próprio comentário: grava \texttt{Historico\_Comentario},
marca \texttt{editado}/\texttt{editado\_em} e não desfaz moderação. O professor
dono (ou o Chefe por Q13) modera com motivo: \texttt{moderado = true},
\texttt{motivo\_moderacao}, \texttt{moderado\_por}/\texttt{moderado\_em}
atômicos, sem apagar o \texttt{texto}. Na listagem autorizada, o autor vê o
original marcado, o colega vê o aviso institucional e o professor dono/Chefe
veem original e histórico.

\subsubsection{Turma arquivada, remoção e replay}
Com a turma \texttt{Arquivada}, criar/editar/remover/moderar Post e
comentar/editar/moderar são negados, inclusive para o Chefe. Após a remoção de
um aluno, a leitura da turma e dos comentários pelo removido é negada, mas os
comentários anteriores permanecem para os demais. O replay da mesma operação M7
devolve o receipt e não duplica comentário, histórico ou notificação; reuso
incompatível de identidade de comando é rejeitado.

\subsubsection{Notificação sem acesso}
Uma notificação \texttt{POST}/\texttt{COMENTARIO} não carrega conteúdo
protegido. Ao clicar, o cliente revalida o acesso corrente; se o aluno foi
removido, a turma foi arquivada ou o item foi moderado, a navegação falha sem
exibir o original a terceiros; a notificação não autoriza leitura por si.

\subsection{Fluxos M12.2 --- roteiros, compartilhamento e download}
Estes exemplos são o contrato normativo da
Seção~\ref{sec:regras-roteiros-m12-2}. São ilustrativos e não atestam a
implementação.

\subsubsection{Upload e publicação de roteiro}
O professor ativo envia um PDF. O backend valida App Check/Auth/papel, tamanho
estritamente menor que 15 MiB, bytes \texttt{\%PDF-}, caminho e titularidade, e
fixa a geração validada antes de tornar o Roteiro publicável. Cancelar ou falhar
não cria registro utilizável; objeto órfão é reconciliado sem apagar objeto
referenciado. Um Roteiro não publicável não aparece em ``Meus''.

\subsubsection{Compartilhar e revogar (Q09)}
O proprietário seleciona professores ativos e autoriza o acesso; a relação
única por (roteiro, professor) é criada de forma idempotente e emite a
notificação mínima \texttt{ROTEIRO\_COMPARTILHADO}. Revogar remove o UID do
array no servidor: o destinatário deixa de ver o roteiro em ``Compartilhados
comigo'' e não pode criar novas associações, mas os Posts já publicados e seus
snapshots permanecem (Q09).

\subsubsection{Download por professor}
O proprietário baixa sempre. O destinatário baixa enquanto o compartilhamento
vigorar. Um professor terceiro não baixa. O servidor emite URL curta com a
identidade e a geração do objeto; a URL já emitida não é revogável antes de
expirar.

\subsubsection{Download por aluno, ex-aluno e turma arquivada}
O aluno com vínculo canônico atual baixa o anexo de um Post acessível, mesmo com
a turma arquivada (somente leitura). Um ex-aluno, ainda que a claim esteja
atualizada, não baixa nem lê o anexo; nenhum cache, erro ou notificação vaza o
conteúdo. Post removido da apresentação e roteiro revogado fora de Post
acessível não liberam download.

\subsubsection{Anexar, editar e desvincular Post}
O professor dono publica um Post com roteiro; o backend exige
\texttt{acessoRoteiroValidado} no commit e grava o snapshot imutável
\texttt{roteiro\_anexo} com a geração. Editar o Post mantendo ou trocando o
anexo exige acesso atual; sem ele, a edição falha fechado e o professor deve
remover o anexo. Desvincular não exige acesso e registra a versão anterior no
histórico de M12.1. O snapshot rotula a publicação, mas a leitura revalida o
acesso atual e nunca usa o snapshot como autorização permanente.

\subsection{Cobertura operacional por papel e ação}
\label{sec:fluxos-completos}
Cada fluxo parte de sessão ativa validada por COM-01 e usa os campos/estados do contrato UI referenciado na seção 8. Os nomes de coleções e os efeitos descritos seguem o dicionário da seção 5.9. Não repetir login entre ações da mesma sessão. A coluna de resultado inclui persistência e falhas relevantes.

\begin{regra}[M9 --- pré-condição comum dos fluxos]
Sessão autenticada não é autorização. Antes de cada leitura sensível ou
mutação, resolver a decisão atual da Seção~7/M9; para efeito crítico, ler
Usuário ativo, papel, escopo/ownership, recurso e estado de domínio na mesma
transação. Fluxo iniciado antes de revogação não pode gravar depois dela. Retry
M7 concluído apenas retorna o resultado já persistido ao mesmo UID.
\end{regra}

\subsubsection{COM-01 --- Todos: Login, recuperar senha, perfil, sair e alternar papel}
\textbf{Contrato:} UI-01. Login ou foto do header; preencher credenciais/perfil ou escolher papel autorizado. Auth e Usuarios determinam sessão; renovar permissões antes da dashboard. Cancelamento preserva sessão anterior; revogação limpa dados e redireciona.

\subsubsection{CHE-01 --- Chefe: Criar usuário e atribuir Bolsista ou outro papel}
\textbf{Contrato:} UI-02. Aba Pessoas > Novo > selecionar existente/convidar > preencher dados do papel > revisar > confirmar. Validar combinações e identidade, gravar perfil e auditoria, sincronizar Auth; lista atualiza após confirmação, falha de claims aparece como pendência.

\subsubsection{CHE-02 --- Chefe: Revogar papel ou autorrevogar chefia}
\textbf{Contrato:} UI-02. Pessoa > Papéis > Revogar > justificar > revisar conjunto restante. Transação verifica último chefe/gestor e vínculos; rejeição mantém papéis e audita tentativa. Último papel removido desativa conta e encerra acesso.

\subsubsection{CHE-03 --- Chefe: Criar/ativar almoxarifado e vincular gestores}
\textbf{Contrato:} UI-03. Almoxarifados > Mais > Novo > nome/descrição/local/gestores > Salvar. Criar local antes se necessário. Persistir almoxarifado e vínculos; ativo exige gestor. Gerenciar gestores permite trocar responsáveis sem deixar almoxarifado ativo órfão.

\subsubsection{CHE-04 --- Chefe: Consultar pessoas, turmas e atividade mensal}
\textbf{Contrato:} UI-02. Selecionar aba, filtro e pessoa > Atividade/Turmas > período. Ler perfis mínimos e agregados por autor/mês; abrir eventos paginados. Sem atividade exibe zero registros, sem confundir com erro de autorização.

\subsubsection{CHE-05 --- Chefe: Operar catálogo, patrimônio, matéria, etiquetas e relatórios}
\textbf{Contrato:} UI-03 a UI-09 e UI-12. Usar as mesmas entradas dos fluxos ALM/PAT/MAT abaixo com escopo global. Poder administrativo não dispensa campos, comprovante de baixa, destino de vencidos ou auditoria. Gestão de turmas/roteiros de terceiros requer Q13; não simular propriedade.

\subsubsection{MAT-01 --- Professor/Chefe: Cadastrar e editar matéria}
\textbf{Contrato:} UI-03. Mais > Nova Matéria ou detalhe > Editar > nome/código > confirmar. Reservar código único e persistir Materia; atualizar projeções e manter vínculos por ID. Código em uso retorna erro no campo.

\subsubsection{PRO-01 --- Professor: Criar turma e copiar código}
\textbf{Contrato:} UI-10. Turmas > Nova > nome/matéria/ano/semestre/capacidade > Salvar. Criar Turma e código único; selecionar turma e Copiar. Falha da área de transferência não recria turma.

\subsubsection{PRO-02 --- Professor: Arquivar e desarquivar turma}
\textbf{Contrato:} UI-10. Turma própria > Arquivar > confirmar; ou Turmas Arquivadas > selecionar > Desarquivar. Alterar status mantendo posts/membros; propagar espelhos e notificações. Q08: turmas arquivadas ficam somente em leitura.

\subsubsection{PRO-03 --- Professor: Convidar aluno com ou sem turma}
\textbf{Contrato:} UI-10. Mais/Novo Aluno > e-mails e matrícula opcional > turma opcional > confirmar ausência de turma ou justificar exceção de capacidade > Enviar. Persistir convites sem conceder papel diferente de Aluno; resultado por destinatário. Aceitação é ALU-01.

\subsubsection{PRO-04 --- Professor: Gerenciar membros e remover aluno}
\textbf{Contrato:} UI-10. Turma própria > Membros > selecionar > Remover > confirmar. O servidor remove o vínculo principal (\texttt{Turma/\{id\}/Alunos}) e também propaga a exclusão no espelho inverso (\texttt{Usuarios/\{uid\}/Turmas}) na mesma transação (N-20), atualiza contador e grava histórico/auditoria; comentários permanecem. Nova entrada por código bloqueada, retorno exige convite.

\subsubsection{PRO-05 --- Professor: Cadastrar, visualizar e baixar roteiro}
\textbf{Contrato:} UI-11. Mais > Novo Roteiro ou Gerenciar Roteiros > upload > nome/descrição/PDF > registrar. Servidor valida objeto e autoria. Biblioteca > Visualizar/Baixar revalida acesso e entrega arquivo; upload falho não cria roteiro utilizável.

\subsubsection{PRO-06 --- Professor: Compartilhar, acompanhar e revogar compartilhamento}
\textbf{Contrato:} UI-11. Biblioteca > Meus > Compartilhar > professores > confirmar. Criar vínculo único e notificação; Eu compartilhei mostra destinatários. Revogar exige confirmação e tratamento de posts existentes conforme Q09; recebido não autoriza redistribuição pelo destinatário.

\subsubsection{PRO-07 --- Professor: Publicar/editar post e vincular roteiro próprio ou recebido}
\textbf{Contrato:} UI-11. Turma própria > compositor/Editar > título/descrição > clipe > escolher roteiro ou cadastrar novo > Publicar. Revalidar acesso ao PDF e turma; persistir Post/histórico, notificar membros e atualizar feed. Falha mantém rascunho na sessão.

\subsubsection{PRO-08 --- Professor: Responder comentário e navegar por notificação}
\textbf{Contrato:} UI-11 e UI-12. Sino > aviso > post alvo > escrever comentário > Enviar. Revalidar propriedade da turma, persistir Comentario e notificar interessados; mostrar resposta no mesmo post, sem pressupor threads hierárquicas.

\subsubsection{PRO-09 --- Professor: Consultar reagentes, receber empréstimo e acompanhar consumo}
\textbf{Contrato:} UI-04, UI-07 e UI-12. Visualizar > Reagentes > filtros > detalhe > procurar gestor para retirada ALM-04. Meus Reagentes mostra empréstimos próprios e consumo histórico; somente gestor confirma movimentação. Consulta não cria reserva de estoque.

\subsubsection{PRO-10 --- Professor: Pesquisar patrimônio e requisitar adição/edição}
\textbf{Contrato:} UI-04 e UI-09. Bens > filtros/número > detalhe > Solicitar Edição, ou Mais > Pedido de Adição. Preencher proposta/motivo, criar Local se ausente, enviar. Lock e requisição são atômicos; Minhas Requisições acompanha resposta PAT-03.

\subsubsection{ALU-01 --- Aluno: Aceitar convite ou ingressar por código}
\textbf{Contrato:} UI-10. Convite > autenticar e verificar e-mail > completar nome/matrícula > Aceitar; ou painel > código > Entrar. Servidor verifica validade, combinações, turma e capacidade/exceção; cria vínculo/espelho/histórico uma vez. Convite inválido ou lotação não cria matrícula parcial.

\subsubsection{ALU-02 --- Aluno: Visualizar turmas, colegas, posts e baixar roteiro}
\textbf{Contrato:} UI-10 e UI-11. Minhas Turmas > turma > Feed/Membros > post > PDF. Ler somente turma com vínculo atual, colegas com dados mínimos e roteiro vinculado. Remoção invalida consultas e acesso ao arquivo; não confiar em espelho/cache antigo.

\subsubsection{ALU-03 --- Aluno: Comentar e ler notificações}
\textbf{Contrato:} UI-11 e UI-12. Post de turma matriculada > texto > Enviar; ou Sino > aviso > post. Validar vínculo e texto, persistir comentário e apresentar confirmação; Limpar tudo marca avisos lidos sem apagar fatos.

\subsubsection{ALU-04 --- Aluno: Consultar estoque sem retirar}
\textbf{Contrato:} UI-04. Visualizar > Reagentes > busca local (catálogo JSON) > detalhes de catálogo. Leitura conforme matriz da seção 3, sem identidade de retirantes nem botões de movimentação; concessão de Bolsista não pode ser feita pelo próprio aluno.

\subsubsection{BOL-01 --- Bolsista: Ações acadêmicas, consulta e retirada/devolução como destinatário}
\textbf{Contrato:} UI-04, UI-07, UI-10 a UI-12. Executar ALU-01 a ALU-04 pelo papel Aluno. Consultar reagente, comparecer ao balcão e informar finalidade/local/prazo ao gestor; ALM-04/05 registram saída/retorno. Acompanhar empréstimos próprios; nunca executar mutação de gestor nem acumular Gestor de Almoxarifado.

\subsubsection{ALM-01 --- Gestor Almox.: Cadastrar/editar resumo, substância, especificação e lote}
\textbf{Contrato:} UI-05. Estoque > Novo Reagente > reutilizar/criar resumo > composição/especificação > lote opcional. Validar identidade e unidades, persistir cada etapa com IDs retornados; densidade confirmada exige nova especificação. Não duplicar resumo por concentração.

\subsubsection{ALM-02 --- Gestor Almox.: Cadastrar frasco fechado ou aberto}
\textbf{Contrato:} UI-06. Estoque > Adicionar Frasco > especificação > almoxarifado/local/lote (opcional) > condição inicial e pesos > validade/destino > revisar > salvar. Transação gera LCQUI-N e grava frasco/histórico; código retornado guia etiqueta física. Repetição da requisição não gera segundo frasco.

\subsubsection{ALM-03 --- Gestor Almox.: Pesquisar, editar localização, abrir e ajustar estoque por pesagem}
\textbf{Contrato:} UI-04 e UI-06. Estoque > filtros > frasco > ação. Abrir recalcula validade; editar localização mantém identidade; pesar exige leitura e motivo, produz histórico de ajuste/evaporação. Conflito com empréstimo obriga recarregar; não sobrescrever peso inicial.

\subsubsection{ALM-04 --- Gestor Almox.: Registrar retirada e autorizar uso didático de vencido}
\textbf{Contrato:} UI-07. Registrar Retirada > código > Professor/Bolsista > local/peso/prazo/finalidade > revisão. Vencido só segue com autorização registrada e ciência explícita; pesquisa excepcional exige justificativa e TCR rastreável Q04. Transação valida disponibilidade e escopo, cria empréstimo e histórico; concorrente perdedor não retira.

\subsubsection{ALM-05 --- Gestor Almox.: Receber devolução de empréstimo próprio ou de outro gestor}
\textbf{Contrato:} UI-07. Registrar Devolução > frasco/pendência > peso/quem devolveu > destino se vencido ou de validade desconhecida > confirmar. O servidor lê o vencimento persistido do frasco (sem recalcular pelo relógio), classifica o retorno (venceu durante, já estava vencido ou validade desconhecida) e calcula consumo e atraso, encerra empréstimo e atualiza frasco/histórico. Peso incompatível é preservado, conclui devolução física com anomalia e consumo pendente para revisão; qualquer gestor vinculado pode receber.

\subsubsection{ALM-06 --- Gestor Almox.: Tratar alertas, quarentena, extravio, reencontro, vazio, quebra e descarte}
\textbf{Contrato:} UI-07, UI-14 e UI-15. Alertas > categoria > frasco > ação/motivo > confirmar. Validar transição e empréstimo ativo; gravar evento, manter identidade e invalidar saldo/alerta. Descarte é terminal; política de remoção lógica em Q07. Em caso de extravio, o frasco preserva seu estado físico, assume \texttt{situacao\_localizacao = EXTRAVIADO} e fica indisponível e preserva o valor anterior de \texttt{saldo\_desconhecido}, o vencimento e \texttt{abertura\_historica\_desconhecida}; um eventual empréstimo ativo sofre encerramento extraordinário. A autorização corrente de descarte é revogada no extravio e a decisão anterior fica no histórico. Em caso de reencontro, o gestor informa o estado físico constatado (\texttt{ABERTO}, \texttt{FECHADO}, \texttt{VAZIO} ou \texttt{QUEBRADO}), validado contra o último estado conhecido antes do extravio: \texttt{FECHADO} só se antes era \texttt{FECHADO} e sem abertura conhecida; \texttt{ABERTO} se antes era \texttt{ABERTO}/\texttt{FECHADO} (marcando abertura histórica desconhecida no segundo caso); \texttt{VAZIO} só se antes já era \texttt{VAZIO}; \texttt{QUEBRADO} representa quebra constatada. Transições impossíveis são rejeitadas. O frasco fica obrigatoriamente com \textit{em\_quarentena = true} e \texttt{INDISPONIVEL}; o saldo permanece como estava, salvo medição ou evento válido, e o empréstimo encerrado não é reaberto. Depois do reencontro, o gestor decide a saída da quarentena (``VOLTAR A DISPONIVEL'' ou ``PENDENTE DE DESCARTE'', com motivo/laudo) e a saída não renova nem recalcula validade; ``VOLTAR A DISPONIVEL'' nunca é oferecido para frasco \texttt{VAZIO} ou \texttt{QUEBRADO}. ``PENDENTE DE DESCARTE'' gera autorização operacional estruturada e mantém o frasco INDISPONIVEL; o descarte só ocorre após essa decisão, nunca diretamente da quarentena. Erros administrativos suportados, como retirada atribuída ao usuário errado, usam ``Corrigir operação'', sem apagar o lançamento original.

\subsubsection{ALM-07 --- Gestor Almox.: Imprimir etiquetas virgens ou segunda via}
\textbf{Contrato:} UI-08. Mais > Imprimir > aba > intervalo/célula ou 1--10 frascos > Gerar > Baixar. Conferir código, imprimir em 100% e colar somente após confirmar cadastro. Virgens não reservam IDs; segunda via audita cada frasco e produz uma etiqueta por ficha.

\subsubsection{ALM-08 --- Gestor Almox.: Gerar relatório de um ou todos os almoxarifados autorizados}
\textbf{Contrato:} UI-12. Relatórios > mensal/período > escopo/filtros > revisar > Gerar/Baixar. Servidor valida 31 dias e permissões, separa g/ml e usa fontes históricas adequadas; erro não anuncia PDF pronto.

\subsubsection{PAT-01 --- Gestor Patr.: Criar resumo, cadastrar e editar bem}
\textbf{Contrato:} UI-09. Patrimônio > Mais > Novo Resumo/Bem, ou detalhe > Editar > preencher/revisar > Salvar. Normalizar a plaqueta por \texttt{trim().toUpperCase()}, validar chave permanente, local, foto e responsável; alteração de nome global exige ação no resumo, reclassificação individual altera referência. Transação cria histórico de cadastro/edição e preserva antes/depois; fan-out derivado não fabrica evento nem conflito de versão.

\subsubsection{PAT-02 --- Gestor Patr.: Pesquisar, filtrar e consultar histórico}
\textbf{Contrato:} UI-04. Patrimônio > prédio/status/letra e filtros opcionais > resumo > plaqueta > detalhe/histórico. Consultar itens e eventos por cursor; mostrar conservação e status distintos, sem substituir histórico pela localização atual.

\subsubsection{PAT-03 --- Gestor Patr.: Analisar, aprovar ou rejeitar pedido}
\textbf{Contrato:} UI-09. Fila/Sino > requisição pendente > comparação > completar dados obrigatórios > justificativa > Aprovar/Rejeitar. Transação revalida M9, confere lock proprietário, versão, chave única, referências e arquivo, persiste decisão, altera/cria bem e histórico se aprovado e libera somente o lock correspondente; notifica solicitante de modo idempotente. Segundo respondente recebe resultado/receipt M7, sem dupla aprovação.

\subsubsection{PAT-04 --- Gestor Patr.: Gerar relação de inservíveis e registrar baixa após SEI}
\textbf{Contrato:} UI-09 e UI-12. Alertas/Relatórios > Inservíveis > filtros > PDF para procedimento externo. Após conclusão no SEI: bem Inservivel > Baixa > upload comprovante > confirmar. Somente confirmação server-side com PDF binariamente validado altera status; histórico e arquivo permanecem disponíveis, enquanto o Bem e sua plaqueta não são excluídos nem reutilizados.

\subsubsection{PAT-05 --- Gestor Patr.: Gerar relatório mensal, personalizado ou individual}
\textbf{Contrato:} UI-12. Relatórios ou detalhe do bem > tipo/período/filtros > Gerar > Baixar/Imprimir. Validar intervalo e datas; documento identifica filtros e instante. Mês em aberto é parcial; erro ou dados ausentes têm retorno distinto.

\paragraph{Q04 e Q14: contrato de aceite e autoatendimento.}
Para PESQUISA\_TCC\_POS ou ESTUDO\_DEGRADACAO\_RESIDUOS com frasco vencido ou validade desconhecida liberado com termo, obter aceite em sessão autenticada do retirante. Persistir tcr\_versao, tcr\_aceito\_em (servidor), tcr\_aceito\_por e tcr\_auditoria\_id no empréstimo, com evento correspondente em Registro\_de\_Auditoria; exigir justificativa\_metodologica de 20--2000 caracteres. Modal/checkbox são UX: backend valida versão vigente, identidade e vínculo do aceite ao frasco, finalidade e operação antes de confirmar a retirada. O TCR registra ciência e responsabilidade e não substitui regras institucionais de segurança química. Autoatendimento só quando não houver outro gestor ativo, calculado pelo servidor com justificativa e notificação à chefia. Registrar abertura\_historica\_desconhecida quando a data anterior não for conhecida; nunca inventar a data.


\paragraph{Exceção metrológica na devolução (HQ-M2-005).}
Se a leitura real contradiz a tara armazenada, o gestor informa exatamente a leitura da balança. O retorno físico encerra a custódia como DEVOLVIDO\_COM\_ANOMALIA, preserva a medição e bloqueia nova retirada; consumo fica pendente até resolução auditada. Referência teórica e tara real são distintas. Confirmar recipiente vazio não autoriza substituir silenciosamente tara real anterior divergente. As opções são nova pesagem, inspeção com confirmação de esgotamento ou recalibração de recipiente vazio conforme o contrato da Seção 10.

\paragraph{Emenda pré-M8: conhecimento metrológico na quebra e no descarte.}
Quebra e descarte preservam \texttt{saldo\_desconhecido} e o conhecimento quantitativo anterior, inclusive no reencontro como \texttt{QUEBRADO}. Não fabricam saldo zero, tara ou pesagem. A baixa operacional não mede a quantidade anteriormente existente. Somente vazio efetivamente confirmado pela rota metrológica válida autoriza saldo zero conhecido. Quarentena, autorização de descarte e terminalidade de \texttt{DESCARTADO} permanecem inalteradas.

\newpage
% =================================================
\section{Tecnologia e Relatórios (Vercel / Firebase)}
% =================================================

\subsection{Para relatórios}
Após avaliação das restrições do ambiente \textit{serverless} (Vercel e Firebase Functions), a biblioteca escolhida para a geração de relatórios no backend é o \textbf{\textit{pdfkit}} em conjunto com o plugin \textbf{\textit{pdfkit-table}}.

\textbf{Justificativa da escolha:}
\begin{itemize}
    \item \textbf{Baixo consumo de memória:} O \textit{pdfkit} desenha o PDF nativamente, respeitando os limites de memória da Vercel/Firebase.
    \item \textbf{Tabelas dinâmicas:} O \textit{pdfkit-table} resolve a paginação automática de tabelas longas.
    \item \textbf{Fluxo de Armazenamento:} A Cloud Function gera o PDF em um \textit{Buffer} de memória, retorna o PDF em base64 conforme o contrato atual; o cliente cria um Blob temporário para download.
\end{itemize}

\subsubsection{Assinatura Lógica de Relatórios (Hash Canônico)}
Todos os relatórios gerados pelo backend injetam um rodapé contendo um Hash Canônico. Este hash compila os parâmetros da requisição (ID, datas e registros consultados) concatenados com o Timestamp (\texttt{Date.now()}), servindo como identificador de rastreabilidade, sem equivaler a assinatura digital.

\subsubsection{Prevenção de Limites de Leitura Firestore (N+1)}
Resolver referências únicas em lotes para reduzir N+1. O tamanho é definido pela API e pelo volume medido; in e getAll têm contratos distintos, sem limite universal de 200 itens.

\subsubsection{Módulo de Geração de Etiquetas (PDF)}
Para a identificação física dos frascos, o backend também utiliza o \textbf{pdfkit} para desenhar uma matriz (Grid Visual de Offset 3×10) em tamanho A4, com etiquetas de 65,0 mm × 26,5 mm. O código de barras no formato Code 128 é gerado ativamente em memória via biblioteca \textbf{bwip-js} e embutido no buffer do PDF, garantindo alta definição na impressão e dispensa de fontes especiais no client. Para frascos já cadastrados (2ª via), o backend gera folhas individuais A4 contendo a Ficha de Conferência e Rastreabilidade com histórico recente de 10 movimentações e 1 única etiqueta de reposição centralizada no rodapé.

\begin{explicacao}[Leitura dos exemplos legados]
Os listings a seguir preservam exemplos de desenvolvimento, não uma implementação pronta para copiar. O contrato atualizado está nas seções 5.9, 8 e 11: caminhos de especificações são aninhados, datas civis são normalizadas no fuso institucional, relatórios usam base64, consultas e escritas são ordenadas em transação e claims exigem controle de revogação. Q06 utiliza tolerância híbrida, consumo não negativo e ajuste separado quando há ganho de massa tolerado. Não adotar consumo negativo, data de abertura inventada, placeholder de foto ou simples cópia de novo nome sobre o bem. Esses exemplos devem ser adaptados ao contrato antes de implementação.
\end{explicacao}

\subsection{Princípio: nunca confiar no \textit{client-side}}
Todas as \textit{Cloud Functions} desta seção seguem a mesma regra: \textbf{qualquer dado que o servidor consiga derivar ou verificar por conta própria é recalculado ou reconferido no servidor, nunca aceito como o cliente enviou}. Isso vale, em particular, para o papel do usuário (claims assinadas pelo Auth, nunca um papel declarado no payload; conta ativa verificada seletivamente conforme DP-D01), para valores calculados como peso/volume/vencimento/atraso (sempre recomputados a partir dos dados brutos e da hora do servidor), para códigos gerados pelo sistema (Regra 6.2.9), e para o estado de disponibilidade de um recurso concorrente (sempre relido dentro de uma transação antes de decidir).

\subsubsection{Funções de Apoio (Autorização)}

\begin{lstlisting}[language=TypeScript, title={Resolução de Papéis via Custom Claims (zero leituras Firestore)}]
import { onCall, HttpsError } from "firebase-functions/v2/https";
import * as admin from "firebase-admin";

// C1: Papéis lidos diretamente do JWT (Firebase Custom Claims).
// Zero leituras extras no Firestore por chamada de API.
// PREREQUISITO: toda mutação de papel deve chamar atualizarCustomClaims(uid).
function resolverPapeisDoToken(auth: {token: Record<string, unknown>} | undefined): string[] {
  if (!auth) return [];
  const roles = auth.token["roles"];
  if (!Array.isArray(roles)) return [];
  return roles as string[];
}

// Chamada após TODA concessão ou remoção de papel:
//   await atualizarCustomClaims(uid);
// Isso atualiza claims para o próximo token; não invalida automaticamente o atual; o cliente deve renovar com user.getIdToken(true).
export async function atualizarCustomClaims(uid: string): Promise<void> {
  const colecoes = ["Chefe_Geral", "Gestor_Almoxarifado", "Gestor_Bens_Patrimoniais", "Professor", "Aluno", "Bolsista"];
  const leituras = await Promise.all(
    colecoes.map((c) => admin.firestore().collection(c).doc(uid).get())
  );
  const roles = colecoes.filter((_, i) => leituras[i].exists);
  await admin.auth().setCustomUserClaims(uid, { roles });
}

async function validarPermissao(
  request: {auth?: {uid: string; token: Record<string, unknown>}},
  papeisPermitidos: string[],
  requerAtivo = false
): Promise<string[]> {
  if (!request.auth?.uid) throw new HttpsError("unauthenticated", "Usuário não autenticado.");
  if (requerAtivo) {
    const snap = await admin.firestore().collection("Usuarios").doc(request.auth.uid).get();
    if (!snap.exists || snap.data()?.ativo !== true)
      throw new HttpsError("permission-denied", "Conta desativada ou não habilitada.");
  }
  const papeis = resolverPapeisDoToken(request.auth);
  if (!papeisPermitidos.some((p) => papeis.includes(p)))
    throw new HttpsError("permission-denied", "Usuário sem papel autorizado para esta ação.");
  return papeis;
}

async function validarGestorDoAlmoxarifado(uid: string, token: Record<string, unknown>, idAlmoxarifado: string): Promise<void> {
  const papeis = resolverPapeisDoToken({ token });
  if (papeis.includes("Chefe_Geral")) return; // Chefe tem escopo global
  if (!papeis.includes("Gestor_Almoxarifado")) {
    throw new HttpsError("permission-denied", "Usuário não é Gestor de Almoxarifado.");
  }
  // 1 leitura: o vínculo de escopo (não o papel - esse já veio do token)
  const vinculo = await admin.firestore().collection("Gestor_Almoxarifado_x_Almoxarifado")
    .where("id_gestor_almoxarifado", "==", uid)
    .where("id_almoxarifado", "==", idAlmoxarifado)
    .limit(1).get();
  if (vinculo.empty) {
    throw new HttpsError("permission-denied", "Gestor não está designado para este almoxarifado.");
  }
}

// Geração atômica sequencial no Firestore via Documento Singleton (Seção 5.7)
// Executado estritamente dentro da transação de cadastro do frasco
async function obterProximoCodigoFrascoTx(
  tx: FirebaseFirestore.Transaction
): Promise<{ proximoNumero: number; novoCodigo: string }> {
  const contadorRef = admin.firestore().collection("Contador_Codigo_Frasco").doc("singleton");
  const contadorSnap = await tx.get(contadorRef);

  const ultimoCodigo = contadorSnap.data()?.ultimo_codigo_gerado || 0;
  const proximoNumero = ultimoCodigo + 1;

  tx.set(contadorRef, { ultimo_codigo_gerado: proximoNumero }, { merge: true });

  return { proximoNumero, novoCodigo: `LCQUI-${proximoNumero}` };
}
\end{lstlisting}

\begin{regra}[DP-D01: verificação seletiva e risco residual]
Usar \texttt{requerAtivo=true} em retirada, devolução, baixa patrimonial, aprovação/rejeição de requisições e concessão/revogação de papéis. As funções \texttt{concederPapel} e \texttt{revogarPapel} seguem a mesma assinatura. Leituras e mutações sem essa exigência podem continuar confiando em token válido até renovação (aproximadamente uma hora). Verificar ativo não atualiza automaticamente papéis de uma conta ainda ativa. Essa limitação é risco residual aceito da DP-D01, não revogação universal instantânea.
\end{regra}

\subsubsection{Fluxo de Reagentes: Cadastro, Retirada e Devolução}

\begin{lstlisting}[language=TypeScript, title={Cadastro de Frasco Fechado (Regra 6.2.13, item 1)}]
interface CadastroFrascoFechado {
  idResumoReagente: string;
  idEspecificacaoReagente: string;
  idAlmoxarifado: string;
  idLote?: string;
  pesoTotal: number;
  volumeNominal: number;
  validadeFechado?: string;
  validadeDesconhecida?: boolean;
  decisaoSeJaVencido?: "QUARENTENA" | "PENDENTE_DE_DESCARTE" | "DISPONIVEL";
  detalheStatus?: string;
}

export const cadastrarFrascoFechado = onCall(async (request) => {
  const dados = request.data as CadastroFrascoFechado;
  await validarPermissao(request, ["Chefe_Geral", "Gestor_Almoxarifado"]);
  await validarGestorDoAlmoxarifado(request.auth!.uid, request.auth!.token, dados.idAlmoxarifado);

  const especSnap = await admin.firestore().collection("Resumo_Reagente").doc(dados.idResumoReagente).collection("Especificacoes").doc(dados.idEspecificacaoReagente).get();
  if (!especSnap.exists) throw new HttpsError("not-found", "Especificação de reagente não encontrada.");
  const densidade = especSnap.data()!.densidade;
  const resumo = (await admin.firestore().collection("Resumo_Reagente").doc(dados.idResumoReagente).get()).data()!;
  const estadoFisico = resumo.estado_fisico;

  let pesoVazio: number;
  if (estadoFisico === "LIQUIDO") {
    if (!densidade || densidade <= 0) throw new HttpsError("failed-precondition", "Líquido exige densidade positiva.");
    pesoVazio = dados.pesoTotal - (dados.volumeNominal * densidade);
  } else {
    pesoVazio = dados.pesoTotal - dados.volumeNominal; // g
  }

  if (pesoVazio <= 0) {
    throw new HttpsError("invalid-argument", "Peso total incompatível com o volume nominal para esta densidade.");
  }

  const agora = new Date();
  const validadeFechado = dados.validadeFechado ? new Date(dados.validadeFechado) : null;
  const validadeEfetiva = dados.validadeDesconhecida ? null : validadeFechado;
  const venceuNoCadastro = Boolean(validadeEfetiva && validadeEfetiva <= agora);

  if (venceuNoCadastro && !dados.decisaoSeJaVencido) {
    throw new HttpsError("failed-precondition", "O frasco está vencido no cadastro e exige uma decisão do gestor.");
  }

  if (dados.decisaoSeJaVencido === "QUARENTENA" && !dados.detalheStatus) {
    throw new HttpsError("invalid-argument", "A entrada em quarentena exige detalhe_status.");
  }

  // Execução atômica no Firestore: Sequenciamento + Cadastro do Frasco
  return admin.firestore().runTransaction(async (tx) => {
    // 1. Gera o código sequencial estrito LCQUI-N
    const { novoCodigo } = await obterProximoCodigoFrascoTx(tx);

    // 2. Persiste o frasco com o código gerado
    const frascoRef = admin.firestore().collection("Frasco_Reagente").doc();
    tx.set(frascoRef, {
      id_almoxarifado: dados.idAlmoxarifado,
      id_lote: dados.idLote ?? null,
      id_especificacao_reagente: dados.idEspecificacaoReagente, // projeção física
      id_resumo_reagente: dados.idResumoReagente,
      eh_higroscopico: resumo.eh_higroscopico, // snapshot servidor imutável
      codigo_frasco: novoCodigo,
      conteudo_nominal: dados.volumeNominal,
      peso_no_cadastrado: dados.pesoTotal,
      peso_atual: dados.pesoTotal,
      peso_frasco_vazio: pesoVazio,
      medida_usada: 0,
      estado_fisico_frasco: "FECHADO",
      disponibilidade: "DISPONIVEL",
      validade_fechado: validadeEfetiva,
      validade_efetiva: validadeEfetiva,
      validade_desconhecida: dados.validadeDesconhecida ?? false,
      vencido: venceuNoCadastro,
      em_quarentena: venceuNoCadastro && dados.decisaoSeJaVencido === "QUARENTENA",
      uso_vencido_autorizado: venceuNoCadastro && dados.decisaoSeJaVencido === "DISPONIVEL",
      detalhe_status: venceuNoCadastro
        ? dados.detalheStatus ?? (dados.decisaoSeJaVencido === "PENDENTE_DE_DESCARTE" ? "Frasco vencido aguardando processo institucional de descarte." : null)
        : null,
      cadastrado_em: admin.firestore.FieldValue.serverTimestamp(),
      cadastrado_por: request.auth!.uid,
    });

    return { idFrasco: frascoRef.id, codigoFrasco: novoCodigo, pesoVazioCalculado: pesoVazio, venceuNoCadastro };
  });
});
\end{lstlisting}

\begin{lstlisting}[language=TypeScript, title={cadastrarFrascoAberto com bifurcação SÓLIDO/LÍQUIDO (C5)}]
// C5: modalidade para SÓLIDO usa ESTIMA_MASSA (não ESTIMA_VOLUME)
interface CadastroFrascoAberto {
  idResumoReagente: string;
  idEspecificacaoReagente: string;
  idAlmoxarifado: string;
  idLote?: string;
  // LIQUIDO: "CONHECE_TARA" | "ESTIMA_VOLUME"
  // SOLIDO:  "CONHECE_TARA" | "ESTIMA_MASSA"
  modalidade: "CONHECE_TARA" | "ESTIMA_VOLUME" | "ESTIMA_MASSA";
  pesoTotalBalanca: number;            // sempre leitura de balança em g
  pesoFrascoVazioInformado?: number;  // Opção A (ambos estados físicos)
  volumeAtualEstimado?: number;        // Opção B LIQUIDO: volume visual em mL
  massaAtualEstimada?: number;         // Opção B SOLIDO: estimativa em g
  dataAbertura?: string;
  aberturaHistoricaDesconhecida: boolean;
  validadeAberto?: string;
  validadeDesconhecida?: boolean;
  decisaoSeJaVencido?: "QUARENTENA" | "PENDENTE_DE_DESCARTE" | "DISPONIVEL";
  detalheStatus?: string;
}

export const cadastrarFrascoAberto = onCall(async (request) => {
  const dados = request.data as CadastroFrascoAberto;
  const papeis = await validarPermissao(request, ["Chefe_Geral", "Gestor_Almoxarifado"]);
  await validarGestorDoAlmoxarifado(request.auth!.uid, request.auth!.token, dados.idAlmoxarifado);

  const especSnap = await admin.firestore().collection("Resumo_Reagente").doc(dados.idResumoReagente).collection("Especificacoes").doc(dados.idEspecificacaoReagente).get();
  if (!especSnap.exists) throw new HttpsError("not-found", "Especificação não encontrada.");
  const especData = especSnap.data()!;
  const resumo = (await admin.firestore().collection("Resumo_Reagente").doc(dados.idResumoReagente).get()).data()!;
  const estadoFisico: "SOLIDO" | "LIQUIDO" = resumo.estado_fisico;
  const densidade: number | null = especData.densidade ?? null;

  // C5: bifurcação explícita por estado físico
  let pesoVazio: number;
  let conteudoNominal: number; // g para sólidos, mL para líquidos

  if (estadoFisico === "SOLIDO") {
    // Sólidos: tudo em gramas; densidade não é usada para cálculo de volume
    if (dados.modalidade === "CONHECE_TARA") {
      if (dados.pesoFrascoVazioInformado == null)
        throw new HttpsError("invalid-argument", "CONHECE_TARA exige pesoFrascoVazioInformado.");
      pesoVazio = dados.pesoFrascoVazioInformado;
      conteudoNominal = dados.pesoTotalBalanca - pesoVazio;
    } else { // ESTIMA_MASSA
      if (dados.massaAtualEstimada == null)
        throw new HttpsError("invalid-argument", "ESTIMA_MASSA exige massaAtualEstimada.");
      conteudoNominal = dados.massaAtualEstimada;
      pesoVazio = dados.pesoTotalBalanca - conteudoNominal;
    }
    if (pesoVazio <= 0)
      throw new HttpsError("invalid-argument", "Massa atual maior que o peso total medido na balança.");
  } else { // LIQUIDO
    if (!densidade || densidade <= 0)
      throw new HttpsError("failed-precondition", "Líquido exige densidade positiva cadastrada na especificação.");
    if (dados.modalidade === "CONHECE_TARA") {
      if (dados.pesoFrascoVazioInformado == null)
        throw new HttpsError("invalid-argument", "CONHECE_TARA exige pesoFrascoVazioInformado.");
      pesoVazio = dados.pesoFrascoVazioInformado;
      conteudoNominal = (dados.pesoTotalBalanca - pesoVazio) / densidade;
    } else { // ESTIMA_VOLUME
      if (dados.volumeAtualEstimado == null)
        throw new HttpsError("invalid-argument", "ESTIMA_VOLUME exige volumeAtualEstimado.");
      conteudoNominal = dados.volumeAtualEstimado;
      pesoVazio = dados.pesoTotalBalanca - (conteudoNominal * densidade);
    }
    if (conteudoNominal <= 0 || pesoVazio <= 0)
      throw new HttpsError("invalid-argument", "Dados de volume/peso inconsistentes com a densidade.");
  }

  const agora = new Date();
  const validade = dados.validadeAberto ? new Date(dados.validadeAberto) : null;
  const validadeEfetiva = dados.validadeDesconhecida ? null : validade;
  const jaVencido = Boolean(validadeEfetiva && validadeEfetiva <= agora);

  if (jaVencido && !dados.decisaoSeJaVencido)
    throw new HttpsError("failed-precondition", "O frasco está vencido no cadastro e exige uma decisão do gestor.");
  if (dados.decisaoSeJaVencido === "QUARENTENA" && !dados.detalheStatus)
    throw new HttpsError("invalid-argument", "A entrada em quarentena exige detalhe_status.");

  return admin.firestore().runTransaction(async (tx) => {
    const { novoCodigo } = await obterProximoCodigoFrascoTx(tx);
    const frascoRef = admin.firestore().collection("Frasco_Reagente").doc();

    tx.set(frascoRef, {
      id_almoxarifado: dados.idAlmoxarifado,
      id_lote: dados.idLote ?? null,
      id_especificacao_reagente: dados.idEspecificacaoReagente, // projeção física
      id_resumo_reagente: dados.idResumoReagente,
      eh_higroscopico: resumo.eh_higroscopico, // snapshot servidor imutável
      codigo_frasco: novoCodigo,
      conteudo_nominal: conteudoNominal,   // g (sólido) ou mL (líquido)
      peso_no_cadastrado: dados.pesoTotalBalanca,
      peso_atual: dados.pesoTotalBalanca,
      peso_frasco_vazio: pesoVazio,
      medida_usada: 0,
      data_abertura: dados.aberturaHistoricaDesconhecida ? null : dados.dataAbertura,
      abertura_historica_desconhecida: dados.aberturaHistoricaDesconhecida,
      estado_fisico_frasco: "ABERTO",
      disponibilidade: "DISPONIVEL",
      validade_efetiva: validadeEfetiva,
      validade_desconhecida: dados.validadeDesconhecida ?? false,
      vencido: jaVencido,
      em_quarentena: jaVencido && dados.decisaoSeJaVencido === "QUARENTENA",
      uso_vencido_autorizado: jaVencido && dados.decisaoSeJaVencido === "DISPONIVEL",
      detalhe_status: jaVencido
        ? (dados.detalheStatus ?? (dados.decisaoSeJaVencido === "PENDENTE_DE_DESCARTE"
            ? "Frasco vencido aguardando processo institucional de descarte." : null))
        : null,
      cadastrado_em: admin.firestore.FieldValue.serverTimestamp(),
      cadastrado_por: request.auth!.uid,
    });

    return { idFrasco: frascoRef.id, codigoFrasco: novoCodigo, conteudoNominal, pesoVazio, venceuNoCadastro: jaVencido };
  });
});
\end{lstlisting}

\begin{lstlisting}[language=TypeScript, title={Abertura e cálculo de validade efetiva no servidor}]
function calcularValidadeEfetivaNaAbertura(frasco: any, dataAbertura: Date): Date | null {
  if (frasco.validade_desconhecida) return null;

  const validadeFechado = frasco.validade_fechado
    ? (typeof frasco.validade_fechado.toDate === "function"
      ? frasco.validade_fechado.toDate()
      : new Date(frasco.validade_fechado))
    : null;
  const diasDepoisDeAberto = frasco.validade_apos_aberto_dias;
  const validadeDepoisDaAbertura = diasDepoisDeAberto != null
    ? new Date(dataAbertura.getTime() + diasDepoisDeAberto * 24 * 60 * 60 * 1000)
    : null;

  if (validadeFechado && validadeDepoisDaAbertura) {
    return validadeFechado <= validadeDepoisDaAbertura ? validadeFechado : validadeDepoisDaAbertura;
  }
  return validadeFechado ?? validadeDepoisDaAbertura;
}

interface AberturaFrasco {
  idFrasco: string;
  destinoSeVencerNaAbertura?: "QUARENTENA" | "PENDENTE_DE_DESCARTE" | "DISPONIVEL";
  detalheStatus?: string;
}

export const registrarAberturaFrasco = onCall(async (request) => {
  const dados = request.data as AberturaFrasco;
  await validarPermissao(request, ["Chefe_Geral", "Gestor_Almoxarifado"]);

  const frascoRef = admin.firestore().collection("Frasco_Reagente").doc(dados.idFrasco);
  return admin.firestore().runTransaction(async (tx) => {
    const snap = await tx.get(frascoRef);
    if (!snap.exists) throw new HttpsError("not-found", "Frasco não encontrado.");
    const frasco = snap.data()!;
    await validarGestorDoAlmoxarifado(request.auth!.uid, request.auth!.token, frasco.id_almoxarifado);

    if (frasco.estado_fisico_frasco !== "FECHADO") {
      throw new HttpsError("failed-precondition", "Somente um frasco FECHADO pode ser aberto por esta operação.");
    }

    const agora = new Date();
    const validadeEfetiva = calcularValidadeEfetivaNaAbertura(frasco, agora);
    const venceu = Boolean(validadeEfetiva && validadeEfetiva <= agora);

    if (venceu && !dados.destinoSeVencerNaAbertura) {
      throw new HttpsError("failed-precondition", "A abertura tornou o frasco vencido. O gestor deve escolher o destino do frasco.");
    }
    if (venceu && dados.destinoSeVencerNaAbertura === "QUARENTENA" && !dados.detalheStatus) {
      throw new HttpsError("invalid-argument", "A entrada em quarentena exige detalhe_status.");
    }

    tx.update(frascoRef, {
      estado_fisico_frasco: "ABERTO",
      data_abertura: agora,
      validade_efetiva: validadeEfetiva,
      vencido: venceu,
      em_quarentena: venceu && dados.destinoSeVencerNaAbertura === "QUARENTENA",
      uso_vencido_autorizado: venceu && dados.destinoSeVencerNaAbertura === "DISPONIVEL",
      detalhe_status: venceu
        ? (dados.detalheStatus ?? (dados.destinoSeVencerNaAbertura === "PENDENTE_DE_DESCARTE"
          ? "Frasco vencido aguardando processo institucional de descarte."
          : null))
        : null,
    });

    return { validadeEfetiva, venceu, destino: dados.destinoSeVencerNaAbertura ?? null };
  });
});
\end{lstlisting}

\begin{lstlisting}[language=TypeScript, title={Registrar Retirada (Empréstimo)}]
interface RetiradaFrasco {
  idFrasco: string;
  idUsuarioRetirou: string;
  idLocalUsado: string;
  pesoSaida: number;
  dataDevolucaoPrevista: string;
  confirmarUsoVencido?: boolean;
  abrirNoEmprestimo?: boolean;
  finalidadeUso: "AULA_PRATICA" | "DEMONSTRACAO" | "PESQUISA_TCC_POS" | "ESTUDO_DEGRADACAO_RESIDUOS";
  justificativaMetodologica?: string;
  idAceiteTcr?: string; // evento prévio autenticado pelo retirante
  justificativaAutoAtendimento?: string;
}

export const registrarRetirada = onCall(async (request) => {
  const dados = request.data as RetiradaFrasco;
  await validarPermissao(request, ["Chefe_Geral", "Gestor_Almoxarifado"], true);

  // Verificação do papel do tomador: este UID é de outro usuário
  // (não do requisitante), portanto Custom Claims não estão disponíveis.
  // Leitura Firestore é necessária aqui (1 leitura, não 6).
  const colecoesBorrower = ["Professor", "Bolsista"];
  const [profSnap, bolsSnap] = await Promise.all([
    admin.firestore().collection("Professor").doc(dados.idUsuarioRetirou).get(),
    admin.firestore().collection("Bolsista").doc(dados.idUsuarioRetirou).get(),
  ]);
  if (!profSnap.exists && !bolsSnap.exists) {
    throw new HttpsError("failed-precondition", "Usuário selecionado não é Professor nem Bolsista.");
  }

  const frascoRef = admin.firestore().collection("Frasco_Reagente").doc(dados.idFrasco);

  return admin.firestore().runTransaction(async (tx) => {
    const frascoSnap = await tx.get(frascoRef);
    if (!frascoSnap.exists) throw new HttpsError("not-found", "Frasco não encontrado.");
    const frasco = frascoSnap.data()!;

    await validarGestorDoAlmoxarifado(request.auth!.uid, request.auth!.token, frasco.id_almoxarifado);
    if (frasco.disponibilidade !== "DISPONIVEL") {
      throw new HttpsError("failed-precondition", "Frasco não está disponível para retirada.");
    }
    if (frasco.em_quarentena) {
      throw new HttpsError("failed-precondition", "Frasco em quarentena não pode ser retirado.");
    }

    const agora = new Date();
    let frascoVencido = frasco.vencido;
    let validadeEfetiva = frasco.validade_efetiva ? frasco.validade_efetiva.toDate() : null;

    if (dados.abrirNoEmprestimo && frasco.estado_fisico_frasco === "FECHADO") {
      validadeEfetiva = calcularValidadeEfetivaNaAbertura(frasco, agora);
      frascoVencido = Boolean(validadeEfetiva && validadeEfetiva <= agora);

    }

    const excepcional = frascoVencido || frasco.validade_desconhecida;
    if (excepcional && (!frasco.uso_vencido_autorizado || !dados.confirmarUsoVencido))
      throw new HttpsError("failed-precondition", "Uso excepcional exige autorização e ciência.");
    const finalidade = dados.finalidadeUso; // nunca substituir finalidade declarada
    const pesquisa = ["PESQUISA_TCC_POS", "ESTUDO_DEGRADACAO_RESIDUOS"].includes(finalidade);
    // Helpers contratuais: todas as leituras devem ocorrer antes da primeira escrita.
    // Aceite foi criado por endpoint autenticado do retirante, não pelo gestor.
    const aceite = excepcional && pesquisa
      ? await validarAceiteTcrTx(tx, dados.idAceiteTcr, {
          retirante: dados.idUsuarioRetirou, frasco: dados.idFrasco,
          finalidade, justificativa: dados.justificativaMetodologica,
        }) : null;
    const autoAtendimento = await validarAutoAtendimentoTx(tx, request.auth!.uid,
      dados.idUsuarioRetirou, frasco.id_almoxarifado, dados.justificativaAutoAtendimento);
    const quimica = await resolverSnapshotQuimicoTx(tx, frasco);

    const emprestimoRef = admin.firestore().collection("Emprestimo_Reagente").doc();
    tx.set(emprestimoRef, {
      id_frasco_reagente: dados.idFrasco,
      id_usuario_retirou: dados.idUsuarioRetirou,
      id_almoxarifado: frasco.id_almoxarifado,
      status: "EM_USO",
      data_retirada: admin.firestore.FieldValue.serverTimestamp(),
      data_devolucao_prevista: new Date(dados.dataDevolucaoPrevista),
      id_local_usado: dados.idLocalUsado,
      peso_saida: dados.pesoSaida,
      uso_vencido_aceito: Boolean(frascoVencido),
      finalidade_uso: finalidade,
      auto_atendimento: autoAtendimento,
      justificativa_metodologica: dados.justificativaMetodologica ?? null,
      tcr_versao: aceite?.versao ?? null,
      tcr_aceito_em: aceite?.aceito_em ?? null,
      tcr_aceito_por: aceite?.uid ?? null,
      tcr_auditoria_id: aceite?.auditoria_id ?? null,
      unidade_medida_utilizada: quimica.unidade,
      densidade_aplicada: quimica.densidade,

    });

    tx.update(frascoRef, {
      disponibilidade: "EMPRESTADO",
      ...(dados.abrirNoEmprestimo && frasco.estado_fisico_frasco === "FECHADO"
        ? { estado_fisico_frasco: "ABERTO", data_abertura: agora,
            validade_efetiva: validadeEfetiva, vencido: frascoVencido } : {}),
    });
    if (aceite) vincularAceiteEAuditoriaTx(tx, aceite, emprestimoRef.id);
    if (autoAtendimento) registrarNotificacaoChefiaTx(tx, emprestimoRef.id);

    const historicoRef = admin.firestore().collection("Historico_Frasco_Reagente").doc();
    tx.set(historicoRef, {
      id_frasco_reagente: dados.idFrasco,
      id_almoxarifado: frasco.id_almoxarifado,
      id_gestor: request.auth!.uid,
      tipo: "SAIU",
      id_emprestimo_reagente: emprestimoRef.id,
      timestamp: admin.firestore.FieldValue.serverTimestamp(),
    });

    return { idEmprestimo: emprestimoRef.id };
  });
});
\end{lstlisting}

\begin{lstlisting}[language=TypeScript, title={Registrar Devolução com atualização física e histórico completo}]
interface DevolucaoFrasco {
  idEmprestimo: string;
  pesoRetorno: number;
  destinoPosDevolucao?: "QUARENTENA" | "PENDENTE_DE_DESCARTE" | "DISPONIVEL";
}

export const registrarDevolucao = onCall(async (request) => {
  const dados = request.data as DevolucaoFrasco;
  await validarPermissao(request, ["Chefe_Geral", "Gestor_Almoxarifado"], true);

  const emprestimoRef = admin.firestore().collection("Emprestimo_Reagente").doc(dados.idEmprestimo);

  return admin.firestore().runTransaction(async (tx) => {
    const emprestimoSnap = await tx.get(emprestimoRef);
    if (!emprestimoSnap.exists) throw new HttpsError("not-found", "Empréstimo não encontrado.");
    const emprestimo = emprestimoSnap.data()!;
    if (emprestimo.status === "DEVOLVIDO" || emprestimo.status === "DEVOLVIDO_COM_ATRASO") {
      throw new HttpsError("failed-precondition", "Empréstimo já foi devolvido.");
    }

    const frascoRef = admin.firestore().collection("Frasco_Reagente").doc(emprestimo.id_frasco_reagente);
    const frascoSnap = await tx.get(frascoRef);
    const frasco = frascoSnap.data()!;
    await validarGestorDoAlmoxarifado(request.auth!.uid, request.auth!.token, frasco.id_almoxarifado);

    const liquido = emprestimo.unidade_medida_utilizada === "ml";
    const densidade = liquido ? emprestimo.densidade_aplicada : null;
    if (liquido && (!Number.isFinite(densidade) || densidade <= 0))
      throw new HttpsError("failed-precondition", "Densidade histórica inválida.");
    const pesoConsumido = Math.max(0, emprestimo.peso_saida - dados.pesoRetorno);

    if (typeof frasco.eh_higroscopico !== "boolean")
      throw new HttpsError("failed-precondition", "Snapshot histórico ausente: reconciliar frasco.");
    const deltaMax = frasco.eh_higroscopico
      ? Math.max(2.0, 0.02 * emprestimo.peso_saida)
      : Math.max(1.0, 0.005 * emprestimo.peso_saida);
    const limiteMaxRetorno = emprestimo.peso_saida + deltaMax;
    if (!Number.isFinite(dados.pesoRetorno) || dados.pesoRetorno <= 0 || dados.pesoRetorno > limiteMaxRetorno)
      throw new HttpsError("invalid-argument", "Peso fora da tolerância híbrida Q06.");
    const avisoHigroscopico = dados.pesoRetorno > emprestimo.peso_saida;
    const volumeUtilizado = densidade ? pesoConsumido / densidade : pesoConsumido;

    const agora = new Date();
    // Normalização: Prazo encerra-se às 23:59:59.999 do dia previsto
    const prazoLimite = new Date(emprestimo.data_devolucao_prevista.toDate());
    prazoLimite.setHours(23, 59, 59, 999);
    const atrasado = agora > prazoLimite;

    const validadeEfetiva = frasco.validade_efetiva?.toDate?.() ?? frasco.validade_efetiva ?? null;
    const venceuAgora = Boolean(validadeEfetiva && validadeEfetiva <= agora);
    const frascoVencido = Boolean(frasco.vencido || venceuAgora);

    const atualizacaoFrasco: Record<string, unknown> = {
      disponibilidade: "DISPONIVEL",
      data_ultima_pesagem: admin.firestore.FieldValue.serverTimestamp(),
      peso_atual: dados.pesoRetorno,
      medida_usada: admin.firestore.FieldValue.increment(pesoConsumido), // sempre g
      vencido: frascoVencido,
    };

    if (frascoVencido && dados.destinoPosDevolucao === "QUARENTENA") {
      atualizacaoFrasco.em_quarentena = true;
      atualizacaoFrasco.uso_vencido_autorizado = false;
      atualizacaoFrasco.detalhe_status = "Frasco vencido retido em quarentena após devolução.";
    } else if (frascoVencido && dados.destinoPosDevolucao === "DISPONIVEL") {
      atualizacaoFrasco.em_quarentena = false;
      atualizacaoFrasco.uso_vencido_autorizado = true;
    }

    tx.update(emprestimoRef, {
      status: atrasado ? "DEVOLVIDO_COM_ATRASO" : "DEVOLVIDO",
      data_devolucao_efetuada: admin.firestore.FieldValue.serverTimestamp(),
      id_usuario_devolveu: request.auth!.uid,
      peso_retorno: dados.pesoRetorno,
      medida_utilizada: volumeUtilizado,
      unidade_medida_utilizada: densidade ? "ml" : "g",
    ...(avisoHigroscopico && { aviso: "higroscopico_suspeito" }),
    });
    tx.update(frascoRef, atualizacaoFrasco);

    if (avisoHigroscopico) {
      tx.set(admin.firestore().collection("Historico_Frasco_Reagente").doc(), {
        id_frasco_reagente: frascoRef.id,
        id_almoxarifado: frasco.id_almoxarifado,
        id_gestor: request.auth!.uid,
        id_emprestimo_reagente: emprestimoRef.id,
        tipo: "AJUSTE", campo_ajustado: "ganho_massa_higroscopia",
        peso_anterior: emprestimo.peso_saida, peso_novo: dados.pesoRetorno,
        medida_ajustada: dados.pesoRetorno - emprestimo.peso_saida,
        unidade_medida_ajustada: "g",
        timestamp: admin.firestore.FieldValue.serverTimestamp(),
      });
    }

    // Registro obrigatório de devolução no histórico
    const histRef = admin.firestore().collection("Historico_Frasco_Reagente").doc();
    tx.set(histRef, {
      id_frasco_reagente: emprestimo.id_frasco_reagente,
      id_almoxarifado: frasco.id_almoxarifado,
      id_gestor: request.auth!.uid,
      tipo: "ENTROU",
      id_emprestimo_reagente: emprestimoRef.id,
      peso_anterior: emprestimo.peso_saida,
      peso_novo: dados.pesoRetorno,
      medida_ajustada: volumeUtilizado,
      unidade_medida_ajustada: densidade ? "ml" : "g",
      timestamp: admin.firestore.FieldValue.serverTimestamp(),
    });

    return { pesoConsumido, volumeUtilizado, atrasado, frascoVencido,
           avisoHigroscopico: avisoHigroscopico || false };
  });
});
\end{lstlisting}

\subsubsection{Validação síncrona de validade no cadastro e na abertura}
A validade não depende apenas do job diário. No cadastro do frasco, a Cloud Function calcula \textit{validade\_efetiva} e compara esse valor com o relógio do servidor. Se a validade já expirou, a criação só pode prosseguir com uma decisão explícita do gestor: \texttt{QUARENTENA}, \texttt{PENDENTE\_DE\_DESCARTE} ou \texttt{DISPONIVEL} com \textit{uso\_vencido\_autorizado}. A primeira opção exige \textit{detalhe\_status}. A terceira representa uma exceção operacional deliberada: reagentes vencidos podem continuar sendo utilizados em práticas menos exigentes ou demonstrações quando o gestor registra explicitamente essa autorização.

Na abertura do frasco, a função \textit{calcularValidadeEfetivaNaAbertura} aplica:
\[
\textit{validade\_efetiva} = \min(\textit{validade\_fechado},\; \textit{data\_abertura} + \textit{validade\_apos\_aberto\_dias}),
\]
quando os dois limites existem.

\subsubsection{Jobs Agendados: Vencimento, Atraso e Materialização}

\begin{lstlisting}[language=TypeScript, title={Vencimento e Atraso com busca preventiva de devoluções}]
import { onSchedule } from "firebase-functions/v2/scheduler";
import { onDocumentCreated, onDocumentDeleted, onDocumentUpdated } from "firebase-functions/v2/firestore";

export const verificarVencimentosEAtrasos = onSchedule("every day 03:00", async () => {
  const agora = new Date();
  const inicioHoje = inicioDoDiaInstitucional(agora);
  const fimDeAmanha = fimDoDiaInstitucional(adicionarDias(inicioHoje, 1));

  const emprestimosNaJanela = await admin.firestore().collection("Emprestimo_Reagente")
    .where("status", "==", "EM_USO")
    .where("data_devolucao_prevista", "<=", fimDeAmanha)
    .get();

  const batchAtrasados = admin.firestore().batch();
  const notificacoesPreventivas: Array<{ idEmprestimo: string; idUsuario: string; quando: "HOJE" | "AMANHA" }> = [];

  emprestimosNaJanela.docs.forEach((doc) => {
    const dados = doc.data();
    const prevista = dados.data_devolucao_prevista.toDate();

    if (prevista < inicioHoje) {
      batchAtrasados.update(doc.ref, { status: "ATRASADO" });
      return;
    }

    if (mesmoDia(prevista, inicioHoje)) {
      notificacoesPreventivas.push({ idEmprestimo: doc.id, idUsuario: dados.id_usuario_retirou, quando: "HOJE" });
    } else {
      notificacoesPreventivas.push({ idEmprestimo: doc.id, idUsuario: dados.id_usuario_retirou, quando: "AMANHA" });
    }
  });

  await batchAtrasados.commit();

  for (const item of notificacoesPreventivas) {
    await notificarProfessorDevolucaoUmaVez(item.idUsuario, item.idEmprestimo, item.quando);
  }

  const totalAtrasados = emprestimosNaJanela.docs.filter((doc) => {
    const prevista = doc.data().data_devolucao_prevista.toDate();
    return prevista < inicioHoje;
  }).length;
  if (totalAtrasados > 0) {
    await notificarGestoresDeReagentes("ENTREGA_ATRASADA", totalAtrasados);
  }

  const frascosVencendo = await admin.firestore().collection("Frasco_Reagente")
    .where("vencido", "==", false)
    .where("validade_efetiva", "<=", agora)
    .get();

  const batchVencidos = admin.firestore().batch();
  frascosVencendo.docs.forEach((doc) => {
    batchVencidos.update(doc.ref, { vencido: true });
    const histRef = admin.firestore().collection("Historico_Frasco_Reagente").doc();
    batchVencidos.set(histRef, {
      id_frasco_reagente: doc.id,
      id_almoxarifado: doc.data().id_almoxarifado,
      tipo: "VENCEU",
      timestamp: admin.firestore.FieldValue.serverTimestamp(),
    });
  });
  await batchVencidos.commit();
  if (!frascosVencendo.empty) {
    await notificarGestoresDeReagentes("FRASCOS_VENCIDOS", frascosVencendo.size);
  }
});

function adicionarDias(data: Date, dias: number): Date {
  const result = new Date(data);
  result.setDate(result.getDate() + dias);
  return result;
}

function mesmoDia(a: Date, b: Date): boolean {
  return a.getFullYear() === b.getFullYear() &&
         a.getMonth() === b.getMonth() &&
         a.getDate() === b.getDate();
}

function inicioDoDiaInstitucional(data: Date): Date {
  const result = new Date(data);
  result.setHours(0, 0, 0, 0);
  return result;
}

function fimDoDiaInstitucional(data: Date): Date {
  const result = new Date(data);
  result.setHours(23, 59, 59, 999);
  return result;
}

export const onFrascoCriado = onDocumentCreated("Frasco_Reagente/{frascoId}", async (event) => {
  const idLote = event.data?.data().id_lote;
  if (idLote) await atualizarContadorLote(idLote, +1);
});

export const onFrascoRemovido = onDocumentDeleted("Frasco_Reagente/{frascoId}", async (event) => {
  const idLote = event.data?.data().id_lote;
  if (idLote) await atualizarContadorLote(idLote, -1);
});

async function atualizarContadorLote(idLote: string, delta: number) {
  await admin.firestore().collection("Lote_Materializado").doc(idLote).set(
    { id_lote: idLote, qtd_frascos_cadastrados: admin.firestore.FieldValue.increment(delta) },
    { merge: true }
  );
}
\end{lstlisting}

\subsubsection{Fluxo de Bens Patrimoniais: Requisições}

\begin{lstlisting}[language=TypeScript, title={Criação e Resposta de Requisição de Edição (C2: lock determinístico)}]
export const criarRequisicaoEdicaoBem = onCall(async (request) => {
  await validarPermissao(request, ["Professor"]);
  const dados = request.data as { idBemPatrimonial: string; novoNome?: string; novoStatus?: string; novoEstadoConservacao?: string; novoIdLocal?: string; motivo: string };

  // C2: Lock determinístico - elimina phantom reads em transações Firestore.
  // Se dois professores enviarem simultaneamente, apenas um cria o lock; o outro
  // recebe ALREADY_EXISTS da transação e falha antes de criar a requisição.
  const lockId = `bem_edicao_${dados.idBemPatrimonial}`;
  const lockRef = admin.firestore().collection("Locks_Requisicao_Patrimonio").doc(lockId);
  const reqRef  = admin.firestore().collection("Requisicao_Edicao_Bem_Patrimonial").doc();

  return admin.firestore().runTransaction(async (tx) => {
    const lockSnap = await tx.get(lockRef);
    if (lockSnap.exists) {
      throw new HttpsError("failed-precondition", "Já existe uma requisição de edição pendente para este bem.");
    }
    const bemSnap = await tx.get(admin.firestore().collection("Bem_Patrimonial").doc(dados.idBemPatrimonial));
    if (!bemSnap.exists) throw new HttpsError("not-found", "Bem não encontrado.");
    // Cria lock e requisição atomicamente
    tx.set(lockRef, { criado_em: admin.firestore.FieldValue.serverTimestamp() });
    tx.set(reqRef, {
      id_bem_patrimonial: dados.idBemPatrimonial,
      novo_nome: dados.novoNome ?? null,
      versao_bem_origem: bemSnap.data()!.versao,
      novo_id_resumo_bem_patrimonial: null,
      novo_status: dados.novoStatus ?? null,
      novo_estado_conservacao: dados.novoEstadoConservacao ?? null,
      novo_id_local: dados.novoIdLocal ?? null,
      motivo: dados.motivo,
      status: "pendente",
      feita_em: admin.firestore.FieldValue.serverTimestamp(),
      id_usuario_solicitante: request.auth!.uid,
    });
    return { idRequisicao: reqRef.id };
  });
});

export const responderRequisicaoEdicaoBem = onCall(async (request) => {
  await validarPermissao(request, ["Chefe_Geral", "Gestor_Bens_Patrimoniais"], true);
  const { idRequisicao, aprovar, justificativa, idResumoAlvo } = request.data;
  return admin.firestore().runTransaction(async (tx) => {
    const reqRef = admin.firestore().collection("Requisicao_Edicao_Bem_Patrimonial").doc(idRequisicao);
    const req = (await tx.get(reqRef)).data();
    if (!req || req.status !== "pendente") throw new HttpsError("failed-precondition", "Requisição indisponível.");
    const bemRef = admin.firestore().collection("Bem_Patrimonial").doc(req.id_bem_patrimonial);
    const bem = (await tx.get(bemRef)).data();
    if (!bem || bem.versao !== req.versao_bem_origem)
      throw new HttpsError("aborted", "Bem alterado; revisar proposta.");
    // Resolve alvo existente ou prepara criação; helper só lê nesta etapa.
    const alvo = aprovar && req.novo_nome
      ? await prepararResumoAlvoTx(tx, idResumoAlvo, req.novo_nome) : null;
    const alteracoes = calcularAlteracoesPermitidas(req, bem);
    if (alvo) {
      if (alvo.novo) tx.create(alvo.ref, alvo.dados);
      alteracoes.id_resumo_bem_patrimonial = alvo.ref.id;
      alteracoes.nome_equipamento = alvo.dados.nome;
    }
    if (aprovar) {
      tx.update(bemRef, { ...alteracoes, versao: bem.versao + 1 });
      registrarHistoricoBemTx(tx, bemRef, bem, alteracoes, request.auth!.uid);
    }
    tx.update(reqRef, {
      status: aprovar ? "aprovada" : "rejeitada",
      novo_id_resumo_bem_patrimonial: alvo?.ref.id ?? null,
      respondida_em: admin.firestore.FieldValue.serverTimestamp(),
      id_usuario_respondente: request.auth!.uid, justificativa_resposta: justificativa,
    });
    tx.delete(admin.firestore().collection("Locks_Requisicao_Patrimonio").doc(`bem_edicao_${bemRef.id}`));
    return { status: aprovar ? "aprovada" : "rejeitada" };
  });
});
\end{lstlisting}

\begin{lstlisting}[language=TypeScript, title={criarRequisicaoAdicaoBem com lock determinístico (C2)}]
export const criarRequisicaoAdicaoBem = onCall(async (request) => {
  await validarPermissao(request, ["Professor"]);

  const dados = request.data as {
    numeroPatrimonioProposto: string;
    estadoConservacaoProposto: string;
    idLocal: string;
    nomeResponsavelProposto: string;
    idResumoBemPatrimonial?: string;
    nomeResumoProposto?: string;
    descricaoResumoProposta?: string;
    motivo: string;
  };

  // C2: Lock determinístico por número de patrimônio proposto
  const lockId = `bem_adicao_${dados.numeroPatrimonioProposto}`;
  const lockRef = admin.firestore().collection("Locks_Requisicao_Patrimonio").doc(lockId);
  const reqRef  = admin.firestore().collection("Requisicao_Adicao_Bem_Patrimonial").doc();

  return admin.firestore().runTransaction(async (tx) => {
    const lockSnap = await tx.get(lockRef);
    if (lockSnap.exists) {
      throw new HttpsError("failed-precondition", "Já existe requisição pendente para este número de patrimônio.");
    }
    tx.set(lockRef, { criado_em: admin.firestore.FieldValue.serverTimestamp() });
    tx.set(reqRef, {
      numero_patrimonio_proposto: dados.numeroPatrimonioProposto,
      estado_conservacao_proposto: dados.estadoConservacaoProposto,
      photo_url_proposta: null,
      nome_responsavel_proposto: dados.nomeResponsavelProposto,
      id_local: dados.idLocal,
      id_resumo_bem_patrimonial: dados.idResumoBemPatrimonial ?? null,
      nome_resumo_proposto: dados.nomeResumoProposto ?? null,
      descricao_resumo_proposta: dados.descricaoResumoProposta ?? null,
      motivo: dados.motivo,
      status: "pendente",
      feita_em: admin.firestore.FieldValue.serverTimestamp(),
      id_usuario_solicitante: request.auth!.uid,
      respondida_em: null,
      id_usuario_respondente: null,
    });
    return { idRequisicao: reqRef.id };
  });
});

export const responderRequisicaoAdicaoBem = onCall(async (request) => {
  await validarPermissao(request, ["Chefe_Geral", "Gestor_Bens_Patrimoniais"], true);
  const { idRequisicao, aprovar, justificativa } = request.data as { idRequisicao: string; aprovar: boolean; justificativa: string };
  const reqRef = admin.firestore().collection("Requisicao_Adicao_Bem_Patrimonial").doc(idRequisicao);

  return admin.firestore().runTransaction(async (tx) => {
    const reqSnap = await tx.get(reqRef);
    if (!reqSnap.exists) throw new HttpsError("not-found", "Requisição não encontrada.");
    const req = reqSnap.data()!;
    if (req.status !== "pendente") throw new HttpsError("failed-precondition", "Requisição já respondida.");

    // C2: Remove o lock ao responder (libera possível nova requisição futura)
    const lockRef = admin.firestore().collection("Locks_Requisicao_Patrimonio")
      .doc(`bem_adicao_${req.numero_patrimonio_proposto}`);
    tx.delete(lockRef);

    let idBemCriado: string | null = null;

    if (aprovar) {
      let idResumo = req.id_resumo_bem_patrimonial;

      // Se não havia resumo existente, cria atômica e transacionalmente o Resumo_Bem_Patrimonial
      if (!idResumo && req.nome_resumo_proposto) {
        const novoResumoRef = admin.firestore().collection("Resumo_Bem_Patrimonial").doc();
        tx.set(novoResumoRef, {
          nome: req.nome_resumo_proposto,
          descricao: req.descricao_resumo_proposta ?? "",
          criado_em: admin.firestore.FieldValue.serverTimestamp(),
        });
        idResumo = novoResumoRef.id;
      }

      if (!idResumo) throw new HttpsError("failed-precondition", "A aprovação exige vincular um resumo.");

      const bemRef = admin.firestore().collection("Bem_Patrimonial").doc();
      idBemCriado = bemRef.id;

      tx.set(bemRef, {
        id_resumo_bem_patrimonial: idResumo,
        nome_equipamento: req.nome_resumo_proposto ?? "Equipamento",
        numero_patrimonio: req.numero_patrimonio_proposto,
        estado_conservacao: req.estado_conservacao_proposto,
        id_local: req.id_local,
        photo_url: req.photo_url_proposta ?? "https://placeholder",
        documento_dado_baixa_pdf_url: null,
        nome_responsavel_sei: req.nome_responsavel_proposto,
        status: "Ativo",
        cadastrado_em: admin.firestore.FieldValue.serverTimestamp(),
      });

      tx.update(reqRef, {
        status: "aprovada",
        respondida_em: admin.firestore.FieldValue.serverTimestamp(),
        id_usuario_respondente: request.auth!.uid,
        id_bem_patrimonial_se_aprovado: idBemCriado,
        justificativa_resposta: justificativa,
      });
    } else {
      tx.update(reqRef, {
        status: "rejeitada",
        respondida_em: admin.firestore.FieldValue.serverTimestamp(),
        id_usuario_respondente: request.auth!.uid,
        justificativa_resposta: justificativa,
      });
    }

    return { status: aprovar ? "aprovada" : "rejeitada", idBemCriado };
  });
});

export const onResumoBemPatrimonialNomeAtualizado = onDocumentUpdated(
  "Resumo_Bem_Patrimonial/{resumoId}",
  async (event) => {
    const antes = event.data?.before.data();
    const depois = event.data?.after.data();
    if (!antes || !depois || antes.nome === depois.nome) return;

    const bens = await admin.firestore().collection("Bem_Patrimonial")
      .where("id_resumo_bem_patrimonial", "==", event.params.resumoId)
      .get();

    const batch = admin.firestore().batch();
    bens.docs.forEach((bem) => batch.update(bem.ref, {
      nome_equipamento: depois.nome,
    }));
    await batch.commit();
  }
);

export const onLocalAtualizado = onDocumentUpdated(
  "Local/{localId}",
  async (event) => {
    const antes = event.data?.before.data();
    const depois = event.data?.after.data();
    if (!antes || !depois) return;
    if (antes.predio === depois.predio && antes.andar === depois.andar && antes.sala === depois.sala) return;

    const bens = await admin.firestore().collection("Bem_Patrimonial")
      .where("id_local", "==", event.params.localId)
      .get();

    const batch = admin.firestore().batch();
    bens.docs.forEach((bem) => batch.update(bem.ref, {
      predio: depois.predio,
      andar: depois.andar,
      sala: depois.sala,
    }));
    await batch.commit();
  }
);

export const ingressarEmTurmaPorCodigo = onCall(async (request) => {
  await validarPermissao(request, ["Aluno"]);
  const { codigoTurma } = request.data as { codigoTurma: string };

  const turmaSnap = await admin.firestore().collection("Turma")
    .where("codigo_turma", "==", codigoTurma)
    .where("status", "==", "Ativo")
    .limit(1).get();
  if (turmaSnap.empty) throw new HttpsError("not-found", "Turma não encontrada ou arquivada.");
  const turmaDoc = turmaSnap.docs[0];
  const turma = turmaDoc.data();

  return admin.firestore().runTransaction(async (tx) => {
    // Checar capacidade
    const alunosSnap = await tx.get(
      turmaDoc.ref.collection("Alunos")
    );
    if (alunosSnap.size >= turma.capacidade) {
      throw new HttpsError("failed-precondition", "Turma lotada.");
    }

    // Checar se já está matriculado
    const jaMatriculado = alunosSnap.docs.some(
      (d) => d.id === request.auth!.uid
    );
    if (jaMatriculado) {
      throw new HttpsError("already-exists", "Você já está nesta turma.");
    }

    // P14: Checar se foi removido anteriormente
    const historico = await tx.get(
      turmaDoc.ref.collection("HistoricoAlunos")
        .where("id_aluno", "==", request.auth!.uid)
        .where("tipo", "==", "exclusao_aluno")
        .limit(1)
    );
    if (!historico.empty) {
      throw new HttpsError("failed-precondition",
        "Você foi removido desta turma pelo professor. O re-ingresso requer convite explícito.");
    }

    // Matricular
    tx.set(turmaDoc.ref.collection("Alunos").doc(request.auth!.uid), {
      id_aluno: request.auth!.uid,
      ingressou_em: admin.firestore.FieldValue.serverTimestamp(),
    });
    tx.set(turmaDoc.ref.collection("HistoricoAlunos").doc(), {
      id_aluno: request.auth!.uid,
      tipo: "inclusao_aluno",
      timestamp: admin.firestore.FieldValue.serverTimestamp(),
    });

    return { idTurma: turmaDoc.id, nomeTurma: turma.nome_turma };
  });
});
\end{lstlisting}

\subsubsection{Relatórios em PDF}
Pseudo-códigos das funções de relatório, revisados: os nomes de campo passam a corresponder exatamente ao modelo (\textit{data\_devolucao\_efetuada}, não \textit{data\_devolucao}; filtro por \textit{id\_almoxarifado} denormalizado em \textit{Emprestimo\_Reagente}, não por um inexistente \textit{id\_local\_saida}), e \textit{gerarRelatorioAlmoxarifado} passa a checar o escopo do gestor, não apenas o papel.

\begin{lstlisting}[language=TypeScript, title={Configuração Base e Relatório de Almoxarifado}]
import PDFDocument from "pdfkit-table";

interface FiltrosAlmoxarifado {
  idAlmoxarifado: string;
  mes: number;
  ano: number;
}

export const gerarRelatorioAlmoxarifado = onCall(async (request) => {
  const { idAlmoxarifado, mes, ano } = request.data as FiltrosAlmoxarifado;
  await validarPermissao(request, ["Chefe_Geral", "Gestor_Almoxarifado"]);
  await validarGestorDoAlmoxarifado(request.auth!.uid, request.auth!.token, idAlmoxarifado);

  const dataInicio = new Date(ano, mes - 1, 1);
  const dataFim = new Date(ano, mes, 0);

  const devolucoesSnapshot = await admin.firestore().collection("Emprestimo_Reagente")
    .where("id_almoxarifado", "==", idAlmoxarifado)
    .where("data_devolucao_efetuada", ">=", dataInicio)
    .where("data_devolucao_efetuada", "<=", dataFim)
    .get();

  const volumeTotalConsumido = devolucoesSnapshot.docs.reduce(
    (acc, doc) => acc + (doc.data().medida_utilizada || 0), 0
  );

  const historicoSnapshot = await admin.firestore().collection("Historico_Frasco_Reagente")
    .where("id_almoxarifado", "==", idAlmoxarifado)
    .where("timestamp", ">=", dataInicio)
    .where("timestamp", "<=", dataFim)
    .orderBy("timestamp", "asc").get();

  const doc = new PDFDocument({ margin: 30, size: "A4" });
  const buffer = await buildPdfBuffer(doc, historicoSnapshot.docs, volumeTotalConsumido);

  const fileName = `relatorios/almoxarifado_${idAlmoxarifado}_${mes}_${ano}_${Date.now()}.pdf`;
  const fileRef = admin.storage().bucket().file(fileName);
  await fileRef.save(buffer, { contentType: "application/pdf" });

  const [url] = await fileRef.getSignedUrl({ action: "read", expires: Date.now() + 3600000 });
  return { url };
});
\end{lstlisting}

\begin{lstlisting}[language=TypeScript, title={Relatório Mensal de um Prédio (Filtros Compostos)}]
interface FiltrosPredio {
  predio: string;
  mes: number;
  ano: number;
  andar?: string;
  sala?: string;
  estadoConservacao?: string;
  status?: string;
}

export const gerarRelatorioBensPredio = onCall(async (request) => {
  const filtros = request.data as FiltrosPredio;
  await validarPermissao(request, ["Chefe_Geral", "Gestor_Bens_Patrimoniais"]);

  let query = admin.firestore().collection("Bem_Patrimonial").where("predio", "==", filtros.predio);
  if (filtros.andar) query = query.where("andar", "==", filtros.andar);
  if (filtros.sala) query = query.where("sala", "==", filtros.sala);
  if (filtros.estadoConservacao) query = query.where("estado_conservacao", "==", filtros.estadoConservacao);
  if (filtros.status) query = query.where("status", "==", filtros.status);

  const bensSnapshot = await query.get();

  const doc = new PDFDocument({ size: "A4" });
  const buffer = await buildPredioPdfBuffer(doc, bensSnapshot.docs, filtros);

  const fileRef = admin.storage().bucket().file(`relatorios/predio_${filtros.predio}_${Date.now()}.pdf`);
  await fileRef.save(buffer, { contentType: "application/pdf" });

  const [url] = await fileRef.getSignedUrl({ action: "read", expires: Date.now() + 3600000 });
  return { url };
});
\end{lstlisting}

\begin{lstlisting}[language=TypeScript, title={Relatório Geral e Período Personalizado}]
interface FiltrosGeralEPersonalizado {
  dataInicio: string;
  dataFim: string;
  entidade: "Bens_Patrimoniais" | "Reagentes";
  prediosSelecionados?: string[];
}

export const gerarRelatorioPersonalizado = onCall(async (request) => {
  const { dataInicio, dataFim, entidade, prediosSelecionados } = request.data as FiltrosGeralEPersonalizado;
  await validarPermissao(request, ["Chefe_Geral", "Gestor_Bens_Patrimoniais", "Gestor_Almoxarifado"]);

  const dataIniObj = new Date(dataInicio);
  const dataFimObj = new Date(dataFim);
  let snapshot;

  if (entidade === "Bens_Patrimoniais") {
    let query = admin.firestore().collection("Historico_Bem_Patrimonial")
      .where("timestamp", ">=", dataIniObj)
      .where("timestamp", "<=", dataFimObj);
    if (prediosSelecionados && prediosSelecionados.length > 0) {
      query = query.where("predio", "in", prediosSelecionados);
    }
    snapshot = await query.get();
  } else {
    snapshot = await admin.firestore().collection("Emprestimo_Reagente")
      .where("data_devolucao_efetuada", ">=", dataIniObj)
      .where("data_devolucao_efetuada", "<=", dataFimObj)
      .get();
  }

  const doc = new PDFDocument({ layout: "landscape", size: "A4" });
  const buffer = await buildPersonalizadoPdfBuffer(doc, snapshot.docs, entidade);

  const fileRef = admin.storage().bucket().file(`relatorios/personalizado_${Date.now()}.pdf`);
  await fileRef.save(buffer, { contentType: "application/pdf" });
  const [url] = await fileRef.getSignedUrl({ action: "read", expires: Date.now() + 3600000 });
  return { url };
});
\end{lstlisting}

\begin{lstlisting}[language=TypeScript, title={Módulo de Etiquetas em PDF: Virgens e Reimpressão}]
// ============================================================================
// GERACAO DE ETIQUETAS VIRGENS (NOVOS FRASCOS) - ABA 1
// ============================================================================
interface DadosEtiquetasVirgens {
  codigoInicial: number;
  codigoFinal: number;
  startRow?: number; // 1 a 10 (Grid A4)
  startCol?: number; // 1 a 3
}

export const gerarPdfEtiquetasVirgens = onCall(async (request) => {
  await validarPermissao(request, ["Chefe_Geral", "Gestor_Almoxarifado"]);
  const dados = request.data as DadosEtiquetasVirgens;

  const total = dados.codigoFinal - dados.codigoInicial + 1;
  if (total <= 0 || total > 50) {
    throw new HttpsError("invalid-argument", "O lote deve conter entre 1 e 50 etiquetas por impressão.");
  }

  // Audita a sessao de geracao de etiquetas virgens (Secao 4.13)
  await admin.firestore().collection("Impressao_Etiqueta_Frasco").add({
    gerado_em: admin.firestore.FieldValue.serverTimestamp(),
    gerado_por: request.auth!.uid,
    codigo_inicial: dados.codigoInicial,
    codigo_final: dados.codigoFinal,
  });

  const codigos = Array.from({ length: total }, (_, i) => `LCQUI-${dados.codigoInicial + i}`);
  const buffer = await gerarPdfEtiquetasNovosFrascos(codigos, dados.startRow || 1, dados.startCol || 1);

  const fileName = `etiquetas/virgens_${dados.codigoInicial}_a_${dados.codigoFinal}_${Date.now()}.pdf`;
  const fileRef = admin.storage().bucket().file(fileName);
  await fileRef.save(buffer, { contentType: "application/pdf" });
  const [url] = await fileRef.getSignedUrl({ action: "read", expires: Date.now() + 3600000 });

  return { url };
});

// ============================================================================
// REIMPRESSAO E FICHA DE CONFERENCIA (FRASCOS JA CADASTRADOS) - ABA 2
// ============================================================================
interface DadosReimpressao {
  frascoIds: string[]; // Máximo 10 frascos
}

export const gerarPdfReimpressaoFrascos = onCall(async (request) => {
  await validarPermissao(request, ["Chefe_Geral", "Gestor_Almoxarifado"]);
  const { frascoIds } = request.data as DadosReimpressao;

  // Regra 7.2.10: Teto rigido de seguranca quimica
  if (!frascoIds || frascoIds.length === 0 || frascoIds.length > 10) {
    throw new HttpsError("invalid-argument", "Selecione entre 1 e 10 frascos por sessao de reimpressao.");
  }

  // 1. Busca dados completos, historico recente e ultimo emprestimo dos frascos
  const dadosFrascos = await buscarDadosCompletosParaConferencia(frascoIds);

  // 2. Registra auditoria individual de 2a via para cada frasco (Secao 4.41)
  const batch = admin.firestore().batch();
  for (const frascoId of frascoIds) {
    const auditRef = admin.firestore().collection("Registro_de_Auditoria").doc();
    batch.set(auditRef, {
      id_usuario: request.auth!.uid,
      acao: "REIMPRESSAO_ETIQUETA",
      tipo_entidade_sofre_acao: "FRASCO_REAGENTE",
      id_do_objeto_da_entidade: frascoId,
      acao_feita_em: admin.firestore.FieldValue.serverTimestamp(),
      metadata: { motivo: "segunda_via_conferencia" },
    });
  }
  await batch.commit();

  // 3. Gera PDF com Ficha de Conferencia + 1 Etiqueta Centralizada por frasco
  const buffer = await gerarPdfFrascosCadastrados(dadosFrascos);

  const fileName = `etiquetas/reimpressao_${Date.now()}.pdf`;
  const fileRef = admin.storage().bucket().file(fileName);
  await fileRef.save(buffer, { contentType: "application/pdf" });
  const [url] = await fileRef.getSignedUrl({ action: "read", expires: Date.now() + 3600000 });

  return { url };
});
\end{lstlisting}



\subsection{Consolidação do planejamento complementar e contrato de execução}
\label{sec:consolidacao}
O inventário complementar inclui Planejamento\_Etiquetas\_LCQUI.md, REGRAS\_MULTI\_ROLE.md, COMPILACAO\_NIX\_LCQUI.md, README.md e os quatro documentos de acompanhamento. Os README de scaffold do frontend não adicionam requisitos de negócio; instruções AGENTS/CLAUDE orientam desenvolvimento, não a especificação do produto. O novo arquivo DUVIDAS\_DOCUMENTACAO\_LCQUI.md registra Q01--Q14 com três alternativas cada; a opção recomendada não é aprovação.

O planejamento de etiquetas foi incorporado às seções 4 (auditoria), 5.9 (contador e campos), 7 (regras), 8 (UI-08), 9 (ALM-07), 10 (geração) e 11 (acesso). O documento de Multi-Role foi incorporado integralmente pelos IDs RN-ROLE-01 a RN-ROLE-15 na seção 7. Os Markdown de origem permanecem como histórico; esta versão explicita as decisões conflitantes em vez de criar fontes concorrentes.

\paragraph{Etiquetas e geometria.}
Papel A4 210 por 297 mm; três colunas e dez linhas; etiqueta 65 por 26,5 mm; margens laterais 7,5 mm e superior/inferior 16 mm; espaçamentos horizontal/vertical zero. Linhas cinza contínuas orientam corte. Offset é aplicado apenas na primeira página; ao completar a grade, reiniciar na célula 1/1 da seguinte. Etiqueta virgem contém código, espaço manuscrito para data/reagente e Code 128. Segunda via contém nome e data cadastrados e uma etiqueta centralizada por ficha. Gerar barras em memória com bwip-js e PDF com pdfkit; preservar proporção e área livre ao redor do código, testar leitura física e códigos longos. Casos mínimos: uma etiqueta, coluna 3, offset 4/2, última célula 10/3, 30/36 itens, nome longo, frasco sem histórico e rejeição de 11 IDs na segunda via. Os testes relatados no planejamento são evidência histórica de protótipo, não homologação da aplicação atual.

\paragraph{Contador e auditoria.}
Adota-se o caminho já usado pelo sistema, Contador\_Codigo\_Frasco/singleton e ultimo\_codigo\_gerado. Contadores/codigo\_frasco no planejamento é proposta antiga, não outro contador ativo. Cadastro incrementa contador e cria frasco atomicamente; impressão não reserva código. Lotes virgens usam Impressao\_Etiqueta\_Frasco; segunda via usa Registro\_de\_Auditoria. A sugestão contraditória de armazenar segunda via também em Impressao não é adotada. Custos reais incluem todas as escritas/leitura de domínio, índices, histórico e tentativas; não reutilizar estimativa de ``1 leitura e 2 escritas'' como custo total.

\paragraph{Relatórios, transporte e integridade.}
O backend atual retorna PDF em base64; a especificação adota esse transporte para relatórios limitados, com conversão em Blob no cliente e liberação da URL temporária. Roteiros e comprovantes permanecem no Storage. Relatórios maiores exigem decisão de infraestrutura antes de ampliar limites. Um hash simples de parâmetros/data não é assinatura digital nem prova de autoria; tratá-lo como identificador de rastreabilidade. Relatórios históricos usam eventos/snapshots da época, não somente o cadastro atual.

\paragraph{Garantias de execução.}
Exemplos de código desta seção são ilustrativos e não substituem contratos UI/dicionário. Todas as leituras transacionais precedem escritas; callbacks podem repetir e não devem enviar e-mail ou gerar PDF como efeito externo. A transação grava evento/operação e um processador idempotente realiza efeitos externos após commit. Alterações no Auth e upload no Storage exigem reconciliação, pois não participam da transação Firestore. Consultas in e getAll não têm o mesmo contrato: tamanho de lote deve ser escolhido por API e medido; 200 não é limite universal do Firestore.

\paragraph{Ambiente e documentação operacional.}
README contém instruções de aplicação e emuladores; COMPILACAO\_NIX\_LCQUI.md é guia de build documental. O flake.nix atual fornece Node/Python/Git/Docker, mas não TeX Live; não presumir que nix develop disponibiliza pdflatex. Usar TeX Live instalado ou ambiente Nix dedicado e compilar main.tex até estabilizar referências. Os quatro documentos de acompanhamento registram evidências do código, não novos requisitos independentes.

\paragraph{Referências técnicas verificadas.}
\href{https://firebase.google.com/docs/firestore/manage-data/transactions}{Firebase: transações e lotes};
\href{https://firebase.google.com/docs/auth/admin/custom-claims}{Firebase: propagação de Custom Claims};
\href{https://firebase.google.com/docs/firestore/manage-data/enable-offline}{Firebase: persistência offline}.
Consulta em 11/09/2026. Estes fundamentos apoiam a execução transacional, renovação de tokens e opção por cache em dispositivo confiável.

\paragraph{Contratos dos auxiliares Q04/Q14.}
validarAceiteTcrTx lê o evento de aceite produzido em sessão autenticada do retirante, confere versão institucional vigente, UID, frasco, finalidade, operação ainda não consumida e justificativa de 20--2000 caracteres. Não aceita timestamp ou UID arbitrário no payload do gestor. vincularAceiteEAuditoriaTx consome o aceite e registra o empréstimo correspondente na mesma transação. validarAutoAtendimentoTx conta gestores ativos, bloqueia quando existe outro e exige justificativa; registrarNotificacaoChefiaTx grava notificação/outbox idempotente. resolverSnapshotQuimicoTx resolve unidade pelo resumo e densidade pela especificação, sem gravações antes de concluir leituras. Auxiliares ilustram contratos, não atestam implementação em functions/. O TCR não substitui segurança química institucional.

\newpage
% =================================================
\section{Regras de Segurança do Firestore (Security Rules)}
% =================================================

\begin{explicacao}[Por que as Security Rules são obrigatórias]
Toda leitura via \texttt{onSnapshot} no frontend passa \textbf{diretamente} pelo Firestore SDK, \textbf{sem} passar pelas Cloud Functions. Portanto, as Cloud Functions são a camada de autorização para escrita, mas as \textit{Security Rules} são a camada de autorização para \textbf{leitura direta}. Sem rules corretas, qualquer usuário autenticado poderia ler qualquer coleção.
\end{explicacao}

\begin{regra}[Princípio: Deny-All por Padrão]
A regra raiz do arquivo \texttt{firestore.rules} deve ser:\\
\texttt{allow read, write: if false;}\\
Todas as exceções são declaradas explicitamente abaixo.
\end{regra}

\subsection{Regras por Coleção}

\small

\small
\setlength{\tabcolsep}{5pt}
\small
\setlength{\tabcolsep}{4pt}
\begin{longtable}{>{\raggedright\arraybackslash}p{0.32\textwidth} >{\raggedright\arraybackslash}p{0.62\textwidth}}
\toprule
\textbf{Coleção} & \textbf{Regra}\\
\midrule
\endfirsthead
\toprule
\textbf{Coleção} & \textbf{Regra} (continuação)\\
\midrule
\endhead
\bottomrule
\endfoot
\nolinkurl{Usuarios/\{uid\}} & Read: \texttt{request.auth.uid == uid} (usuário lê só o próprio documento). Write: negado ao cliente (só Cloud Function).\\
\nolinkurl{Usuarios/\{uid\}/Turmas/\{turmaId\}} & Read: \texttt{request.auth.uid == uid} (o aluno lê em tempo real apenas as turmas em que está matriculado). Write: negado ao cliente (gerenciada exclusivamente pelas Cloud Functions).\\
\texttt{Chefe\_Geral}, \texttt{Gestor\_Almoxarifado}, etc.~(coleções de papel) & Read: \texttt{request.auth.uid == docId} (cada usuário lê só o próprio documento de papel). Write: negado ao cliente.\\
\nolinkurl{Turma/\{turmaId\}} & Read: usuário é professor da turma (\texttt{resource.data.id\_professor == request.auth.uid}) OU está na sub-coleção \texttt{Alunos}. Write: negado ao cliente.\\
\nolinkurl{Turma/\{turmaId\}/Alunos} & Read: \texttt{request.auth.uid == resource.data.id\_aluno} OU membro ativo da mesma turma, apenas em projeção mínima sem matrícula/e-mail.\\
\nolinkurl{Turma/\{turmaId\}/Posts} & Read: membro da turma (professor ou aluno matriculado). Write: negado.\\
\nolinkurl{Turma/\{turmaId\}/Posts/\{postId\}/Comentarios} & Original moderado: somente autor, Professor responsável ou Chefe Geral com acesso ao contexto. Colegas recebem aviso por endpoint autorizado. Escrita direta negada.\\
\nolinkurl{Usuarios/\{uid\}/Notificacoes} & Read: \texttt{request.auth.uid == uid}. Write: negado ao cliente.\\
\nolinkurl{Resumo\_Reagente/\{resumoId\}} & Read: qualquer usuário com papel ativo no sistema (token de acesso válido). Write: negado ao cliente.\\
\nolinkurl{Resumo\_Reagente/\{resumoId\}/Especificacoes/\{specId\}} & Read: qualquer usuário com papel ativo no sistema (catálogo público interno). Write: negado ao cliente.\\
\texttt{Frasco\_Reagente} & Read: papel \texttt{roles} no token contém \texttt{Chefe\_Geral} OU \texttt{Gestor\_Almoxarifado} OU \texttt{Professor} OU \texttt{Bolsista} OU \texttt{Aluno}. Write: negado ao cliente.\\
\texttt{Emprestimo\_Reagente} & Read: \texttt{Chefe\_Geral} OU \texttt{Gestor\_Almoxarifado} OU o próprio usuário que retirou (\texttt{resource.data.id\_usuario\_retirou == request.auth.uid}). Write: negado.\\
\texttt{Bem\_Patrimonial} & Read: Chefe Geral, Gestor de Bens Patrimoniais ou Professor. Write: negado ao cliente.\\
\texttt{Requisicao\_Edicao\_Bem\_Patrimonial} & Read: o solicitante (\texttt{resource.data.id\_usuario\_solicitante == uid}) OU gestores de bens (papel no token). Write: negado.\\
\texttt{Requisicao\_Adicao\_Bem\_Patrimonial} & Idem acima.\\
\texttt{Locks\_Requisicao\_Patrimonio} & Read/Write: negado ao cliente (exclusivo de Cloud Function).\\
\texttt{Registro\_de\_Auditoria} & Read: \texttt{Chefe\_Geral} no token. Write: negado ao cliente.\\
\texttt{Roteiro\_Experimento} & Read: \texttt{resource.data.id\_professor\_upload == uid} OU \texttt{uid} pertence ao array \texttt{professores\_compartilhados}. Write: negado.\\
\texttt{Contador\_Codigo\_Frasco} & Read: \texttt{request.auth != null} (qualquer logado lê). Write: negado ao cliente (só Cloud Function atualiza o sequenciador).\\
\end{longtable}


\begin{regra}[Leitura de Custom Claims nas Security Rules]
O papel do usuário fica disponível nas Security Rules como \texttt{request.auth.token.roles}. Exemplo:\\
\texttt{function temPapel(papel) \{}\\
\texttt{\quad return papel in request.auth.token.roles;}\\
\texttt{\}}\\
Isso garante consistência com o mesmo Custom Claim usado pelas Cloud Functions (\texttt{resolverPapeisDoToken}).
\end{regra}


\subsection{Contrato de acesso para os fluxos detalhados}
\label{sec:acesso-detalhado}
Além do papel, validar o escopo atual. DP-D01 determina consulta de conta ativa em mutações de alto impacto conforme a seção 10; leituras e demais mutações podem confiar no JWT até renovação. Atualizar Custom Claims não invalida imediatamente o token emitido. Versão de permissões e reconciliação não constituem garantia universal já implementada. Preferência de papel no cliente só altera apresentação.

Leituras de membros, posts, comentários e PDFs de turma exigem professor responsável ou vínculo ativo; colegas recebem projeção mínima, sem e-mail/matrícula. Requisições são visíveis ao solicitante, gestor patrimonial e Chefe. Empréstimos completos são visíveis ao retirante, gestor vinculado ao almoxarifado e Chefe; demais leitores do catálogo recebem somente disponibilidade. Histórico operacional segue o mesmo escopo, e auditoria global fica com a chefia. Diretórios usados em seletores devem ter endpoints/projeções mínimas autorizadas, não leitura irrestrita de Usuarios.

Mutações de domínio e locks são exclusivas do servidor, inclusive marcar notificações lidas; não permitir delete de notificações para implementar Limpar tudo. Regras devem cobrir explicitamente Resumo\_Reagente/id/Especificacoes, Turma e suas subcoleções e Bem\_Patrimonial/id/Historico. Consultas collection-group exigem regras/índices próprios. Contador, unicidade e controles de papéis são técnicos; somente resposta mínima autorizada é exposta ao operador.

Storage valida dono, caminho, tipo e tamanho: foto inferior a 5 MiB, baixa inferior a 10 MiB e roteiro inferior a 15 MiB. Backend verifica arquivo existente e assinatura de conteúdo antes de referenciar; metadados não bastam para provar formato. Leitura de roteiro exige propriedade, compartilhamento ou post de turma acessível. Comprovante de baixa segue escopo patrimonial e não pode ser removido após vinculação. Não conceder leitura de todo Storage apenas porque o usuário está autenticado.

Estas são exigências da especificação; firestore.rules e storage.rules atuais ainda divergem, conforme AUDITORIA\_ATUALIZADA.md. A presença de deny-all genérico não anula permissões amplas de matches específicos.

\subsection{Moderação: autorização por documento e resposta filtrada}
Firestore Rules autorizam ou negam documentos completos; não fazem projeção nem substituição do campo texto. Não implementar ``texto vazio via Rules''. Para o modelo físico descrito, negar leitura direta de Comentarios e histórico a quem não pode conhecer o original; a listagem de comentários é fornecida por endpoint autenticado que valida vínculo com a turma e devolve somente aviso institucional para comentários moderados de colegas. Autor recebe original marcado; Professor responsável e Chefe Geral recebem conteúdo necessário à auditoria. Não retornar original em campos auxiliares, cache, notificações ou mensagens de erro. Consulta ampla cliente que possa incluir documentos proibidos deve ser rejeitada e substituída pelo endpoint, nunca filtrada apenas no frontend. Backend valida também autoria da edição, motivo e papel de moderação, status ativo da turma e registra histórico/auditoria.

Coleções Controle\_Papeis, Operacoes e Chaves\_Unicas são exclusivas de Admin SDK; leitura e escrita diretas de clientes são negadas. O deny-all atual já as protege, sem allow concorrente; matches explícitos constituem hardening documental. Esta seção descreve o alvo; a implementação atual de moderação permanece pendente.

\newpage
% =================================================
\section{Implementações em Estudo para Versões Futuras}
% =================================================

\begin{explicacao}[Escopo desta seção]
Nada aqui faz parte do modelo 3FN da Seção 4 ou do mapeamento físico Firestore da Seção 5, nem deve ser implementado nesta primeira versão. É o registro de uma proposta para quando reagentes gasosos entrarem no almoxarifado do LCQUI -- hoje nenhum reagente gasoso está cadastrado, por isso \texttt{GASOSO} foi removido de \textit{Resumo\_Reagente.estado\_fisico} (seção 4.17) nesta versão, em vez de deixado junto de campos pensados só para sólido/líquido.
\end{explicacao}

\subsection{Evoluções posteriores da autorização}
M9 fixa a V1 em RBAC com escopo/vínculo persistido e ownership, sem introduzir
ACL, ABAC ou capabilities. Versões futuras só poderão propor novo paradigma
depois de demonstrar requisito de domínio que a matriz M9 não expressa, com
migração, regras, backend, revogação, auditoria, idempotência e casos de
segurança próprios. Não se amplia privilégio por conveniência de interface ou
por Custom Claim. A implementação do contrato M9 nas Rules, Storage, backend,
App Check, IAM e testes é trabalho posterior à formalização executável M9.

\subsection{Por que \texttt{GASOSO} sozinho não basta}

Simplesmente devolver \texttt{GASOSO} a \textit{estado\_fisico} resolveria a classificação básica, mas um cilindro de gás carrega informação que \textbf{não} é estado físico e que não deve virar uma explosão de valores dentro do mesmo enum (\texttt{GASOSO\_COMPRIMIDO}, \texttt{GASOSO\_LIQUEFEITO}, \texttt{MISTURA\_GASOSA\_CALIBRACAO}, \ldots). São quatro dimensões independentes, que devem continuar independentes no banco:

\begin{center}
\begin{tabular}{p{0.28\textwidth}p{0.62\textwidth}}
\toprule
\textbf{Dimensão} & \textbf{Exemplos de valor} \\
\midrule
Estado físico & sólido, líquido, gasoso \\
Forma de armazenamento & gás comprimido, gás liquefeito, gás criogênico \\
Tipo de substância (já existe) & pura, mistura \\
Finalidade/aplicação & reagente comum, gás de calibração, gás de processo \\
\bottomrule
\end{tabular}
\end{center}

Um cilindro de CO\textsubscript{2} puro, por exemplo, é \textit{estado\_fisico} = gasoso, \textit{forma\_armazenamento} = liquefeito (o CO\textsubscript{2} é mantido líquido sob pressão dentro do cilindro, embora seja gasoso em condição ambiente), \textit{tipo\_substancia} = pura. Já uma mistura de calibração (ex.: 500\,ppm de CO\textsubscript{2} balanceado em N\textsubscript{2}) é \textit{estado\_fisico} = gasoso, \textit{forma\_armazenamento} = comprimido, \textit{tipo\_substancia} = mistura, \textit{finalidade} = calibração. Tratar ``mistura gasosa de calibração'' como uma categoria física própria misturaria uma propriedade química (o que a substância é) com uma propriedade comercial/de uso (para que ela serve) -- o mesmo tipo de acoplamento indevido que a decomposição do \textit{status} de \textit{Frasco\_Reagente} (seção 4.19) já evitou para outro campo.

\subsection{Proposta de campos (apenas ilustrativa)}

\begin{entidade}{Especificacao\_Reagente \textit{(campos adicionais propostos)}}
\campo{forma\_armazenamento}{ENUM}{NULL}
\subcampo{CONVENCIONAL, GAS\_COMPRIMIDO, GAS\_LIQUEFEITO, GAS\_CRIOGENICO}
\campo{finalidade}{ENUM}{NULL}
\subcampo{REAGENTE, CALIBRACAO, PROCESSO, OUTRO}
\end{entidade}

Ambos ficam em \textit{Especificacao\_Reagente}: são propriedades do produto comercial específico, não da substância química abstrata nem do item de catálogo. \texttt{NULL} para itens sólidos/líquidos, onde a dimensão não se aplica.

\subsection{Além do enum: medição de gases não é pesagem}

A Regra de Cálculo de Volume e Peso (seção 6.2.13) assume que o volume é sempre derivado de uma leitura de balança convertida por densidade. Para cilindros de gás, a prática de laboratório mais comum é ler a \textbf{pressão} no manômetro do próprio cilindro (ainda que pesar o cilindro também seja uma alternativa viável em alguns casos). Isso implica que \textit{Frasco\_Reagente} precisaria, no mínimo, de:

\begin{itemize}
    \item \textit{pressao\_nominal\_bar} -- pressão de enchimento de fábrica.
    \item \textit{pressao\_atual\_bar} -- última leitura do manômetro (substituindo/complementando \textit{peso\_no\_cadastrado}/\textit{medida\_usada} para este tipo de frasco).
    \item \textit{capacidade\_litros\_cntp} -- volume de gás equivalente em CNTP (Condições Normais de Temperatura e Pressão) quando o cilindro está cheio.
    \item \textit{tipo\_valvula} -- necessário para compatibilidade de equipamento na retirada.
\end{itemize}

O cálculo de consumo também muda: em vez de $V = \Delta m / \rho$, a estimativa correta usa a proporcionalidade entre pressão e quantidade de gás restante (lei dos gases ideais, $PV = nRT$, ou uma correlação real para gases não-ideais em alta pressão), não uma conversão por densidade. Por isso este documento não tenta encaixar gases nos campos de peso já existentes: a reintrodução de \texttt{GASOSO} deve vir acompanhada desses campos e dessa fórmula alternativa, não apenas do valor no enum.

\subsection{Horários de Turma}
A entidade \textit{Horario\_Turma} (dia da semana, hora de início/fim e local por turma) foi removida da V1 por não possuir fluxo de uso definido nas telas nem nas regras de negócio. Em versões futuras, ela pode ser reintroduzida para permitir que o sistema exiba grades horárias, detecte conflitos de sala e envie lembretes automáticos antes do início das aulas.

\subsection{Lotação física de professores (V2)}
Q01: vinculação a prédio, gabinete e sala não integra a V1 por insuficiência de dados estruturados. Integração futura de dados institucionais de lotação e alocação de salas permitirá filtragem espacial; depende de disponibilização e validação dos dados pela UENF. Na V1, busca de professor limita-se a Nome, Matéria, Centro e Laboratório.

\subsection{Edição de Resumo e Especificação de Reagente (V2)}
Na V1, \textit{Resumo\_Reagente} e \textit{Especificacao\_Reagente} podem ser cadastrados, mas \textbf{não existe edição posterior} de Resumo ou Especificação pela interface, e não há endpoint genérico de edição desses objetos. Não pertencem à V1 \textit{Resumo.version}, \textit{EditSession} nem arquitetura equivalente de travamento otimista; a concorrência de edição desses objetos fica deliberadamente para a V2.

A V2 deverá estudar edição auditável de \textit{Resumo\_Reagente} e \textit{Especificacao\_Reagente}, incluindo, quando a implementação futura for projetada: autorização; concorrência entre gestores; versionamento; preservação de histórico; auditoria; e efeitos sobre projeções/cache. Nada disso deve ser implementado agora.

\subsection{Correção auditável da validade de Frasco\_Reagente (V2)}
Na V1 não existe operação de edição ou correção posterior da validade cadastrada de um frasco. Erro excepcional desse tipo não possui correção autônoma pela interface, e não há API genérica de correção. A V2 deverá estudar três modalidades futuras:

\begin{enumerate}
    \item \texttt{DESCONHECIDA\_PARA\_CONHECIDA} -- passar a conhecer uma validade antes desconhecida;
    \item \texttt{DATA\_INCORRETA\_PARA\_DATA\_CORRETA} -- corrigir uma data cadastrada incorretamente;
    \item \texttt{CONHECIDA\_PARA\_DESCONHECIDA} -- registrar que a validade antes considerada conhecida é, na verdade, desconhecida.
\end{enumerate}

A implementação futura deverá preservar o valor anterior no histórico; preservar autoria, instante, motivo e evidência; não oferecer edição genérica \texttt{campo/valor}; recalcular campos derivados no servidor; não confundir correção documental com ``revalidar validade''; e não reescrever fatos históricos de empréstimos anteriores. Essas três modalidades não são implementadas nesta rodada.

\subsection{Captura automática de pesagens e correção metrológica auditável (versões futuras)}
Na V1, toda pesagem é \textbf{entrada manual}; a interface usa barreiras preventivas de dupla digitação independente, unidade fixa em g, contexto visual, alertas de plausibilidade e revisão final antes do commit (Seção~8, UI-18). A V1 \textbf{não} possui correção administrativa metrológica após a persistência, nem editor genérico de peso, nem tipo específico \texttt{ERRO\_TRANSCRICAO\_PESO}: uma nova pesagem não reescreve automaticamente o fato histórico anterior. A decisão humana de M6 (HQ-M6-001) mantém o conjunto metrológico fechado em três rotas físicas tipadas.

\subsubsection{Estudo futuro A — integração direta com a balança}
Estudar a integração das balanças do LCQUI diretamente ao sistema, para reduzir ou eliminar a transcrição manual. A arquitetura futura deverá avaliar, sem implementar agora: interfaces realmente disponíveis nas balanças do LCQUI (USB, serial, rede ou outro protocolo do equipamento); captura automática da leitura; identificação do instrumento; \texttt{timestamp} confiável; unidade; resolução/casas decimais; estabilidade da leitura; associação inequívoca entre frasco, operação e leitura; trilha de auditoria; autenticação/autorização; falhas de comunicação; modo offline/fallback manual; validação da interface; preservação do dado bruto; e prevenção de duplicidade/replay. Não escolher tecnologia/protocolo sem inventário real das balanças.

\subsubsection{Estudo futuro B — correção metrológica auditável}
Estudar se versões futuras devem suportar correções metrológicas administrativas excepcionalíssimas quando existir evidência independente suficiente de erro de transcrição. Não assumir que serão implementadas. Se estudadas, deverão responder: quais tipos fechados de correção existem; qual evidência é obrigatória; quem pode realizar; necessidade de segunda aprovação; preservação integral do valor original; autoria; \texttt{timestamp}; motivo; documento/evidência vinculada; efeito em FLOW; efeito em saldos; impossibilidade de alterar fisicamente o passado; e proibição de editor genérico \texttt{campo/valor}. É um estudo de versão futura, não requisito obrigatório de V1/V2.


\end{document}
